% 高一122_华师大二附中_周末作业6.tex
%   —— 高一上数学“2026-2027 年华师大二附中高一上周末作业 6”（试卷版/答案版）
% 试卷版：xelatex 高一122_华师大二附中_周末作业6.tex
% 答案版：xelatex "\def\isanswer{1}% 高一122_华师大二附中_周末作业6.tex
%   —— 高一上数学“2026-2027 年华师大二附中高一上周末作业 6”（试卷版/答案版）
% 试卷版：xelatex 高一122_华师大二附中_周末作业6.tex
% 答案版：xelatex "\def\isanswer{1}% 高一122_华师大二附中_周末作业6.tex
%   —— 高一上数学“2026-2027 年华师大二附中高一上周末作业 6”（试卷版/答案版）
% 试卷版：xelatex 高一122_华师大二附中_周末作业6.tex
% 答案版：xelatex "\def\isanswer{1}% 高一122_华师大二附中_周末作业6.tex
%   —— 高一上数学“2026-2027 年华师大二附中高一上周末作业 6”（试卷版/答案版）
% 试卷版：xelatex 高一122_华师大二附中_周末作业6.tex
% 答案版：xelatex "\def\isanswer{1}\input{高一122_华师大二附中_周末作业6.tex}"
% 紧凑版：xelatex "\def\iscompact{1}\input{高一122_华师大二附中_周末作业6.tex}"
%
% 原始材料：微信公众号“魔都数学加油站”文章，作者破晓，2026-11-09 17:52 发布；
% 连载“27学年高一122”，原文标题“27学年高一122——2026-2027年上海市华师大二附中高一上周末作业6”；
% 原文链接：https://mp.weixin.qq.com/s/pC5066JFKNkQQwjMqTK8Ow。
% 正文为 14 张 PNG 页（第 14 页 2560×3620，其余 13 页 4762×6735），
% 题目下印有【解析】，为讲评版。
% 题目、解析均据原图转录；卷面红色水印及公众号页脚未录入。全卷无插图（已逐页检索，无“如图”）。
%
% 卷首标题：原卷印作“2026-2027 年华师大二附中高一上周末作业 6”（公众号标题中的“上海市”
% 卷面标题行没有，本稿照卷面），年份区间按全稿惯例排 en dash；原卷未印副标题，本稿按题目内容
% 标为“函数”。副标题理由（以全部 14 页实际内容为准）：逐页复核 19 题，考的都是函数性质——
% 奇偶性与充要条件（Q3、Q13、Q15、Q16）、单调性与单调区间（Q4、Q7、Q14、Q17）、定义域／
% 值域／最值与复合函数（Q5、Q6、Q8、Q10、Q12、Q19、Q20）、分段函数与取整函数（Q7、Q8、
% Q11、Q15、Q19）、抽象函数与新定义（Q12、Q20、Q21）；不等式仅作求参数范围／解集的工具
% （Q9、Q14、Q16、Q18），集合与常用逻辑用语仅零星出现（Q3 的充要条件、Q11③ 的集合、
% Q20 的“非空集合 I”），均非独立章节。故主体为单一模块“函数”，与同校 高一102 周末作业 5
% 同（该函数模块卷亦按 高一83／105 成例只用模块名、不叠“集合与常用逻辑用语、不等式”）。
%
% 与原卷的排版差异：
%   ① 原文从第 3 题起，填空题编号为 3—12、选择题为 13—16、解答题为 17—21，本稿保留原题号；
%   ② 原卷每题下直接印【解析】，本稿按全稿惯例置于答案版；四道选择题括号内均未印答案、
%      答案在解析末句“故选 …”，括号内留空（全卷无只印【答案】的题、无 \ans{}）；
%   ③ 原卷未印副标题，本稿按“知识模块”口径补出；
%   ④ 原卷填空横线按全稿惯例排为 \blankline[4.4em]；选择题括号排 \mbox{（\quad）}、
%      选项段前按惯例加 \keepwithprev；解答题不留作答区（原卷为讲评版）；
%   ⑤ 原卷数集排斜体 $R$、$N_+$（下标加号），按全稿惯例改 \R、\N_+
%      （$N_+$ 下标加号与 $\N^*$ 是两种记号，照录下标、只改数集字母为正体，同 高一05 第 20 题成例）；
%      绝对值排 \lvert…\rvert、集合描述竖线排 \mid；
%      $f(x)_{min}$／$f(x)_{max}$／$y_{max}$／$(y_{max})_{min}$ 一类斜体下标按 高一96 成例
%      改正体 \min／\max；不等式一律排 \geqslant／\leqslant；
%   ⑥ 第 20 题解析标签原卷印作【解得】（全卷唯一，其余 18 题均印【解析】），本稿按全稿惯例
%      由 \solbegin 统一印【解析】，正文内容照录（行内加 % 注、并入文件头存疑清单）；
%   ⑦ 题干与解析的标点（半角／全角逗号、冒号、括号、句点、子题括号）均照原卷，
%      题目正文引号用弯引号“ ”。
%
% 子集符号：第 21 题 (2) 的 $\left[\frac13+\frac1{9b-3},1\right]\subseteq\left[\frac12,b\right]$
%   印有下划线，照录 \subseteq。全卷未出现判别式 \triangle／\Delta。
%
% 原卷存疑／笔误（均照录加注）：
%   第 11 题解析对② 的“$y=y+\Delta y$”（按上下文应为 $y=[y]+\Delta y$，原卷漏方括号）；
%   第 19 题题干“某种产品 $x$ 万件）可获利”右括号无配对（疑漏印左括号“（”）；
%   第 20 题解析标签【解得】（见排版差异⑥）；
%   第 21 题 (3) 解析“总存在唯一的 $x_2\in[0,2]$ 时，使得……”的“时”字（按上下文衍字）；
%   第 21 题 (3) 末句“综上所述，$a$ 的范围 $[-\sqrt2,2-\sqrt3]\cup[\sqrt3,2+\sqrt2]$”
%      “范围”后缺“是／为”字，照录。
%
% 行数核对：逐题数 \step 条数，合计 162 条（Q3：5；Q4：3；Q5：1；Q6：3；Q7：2；
% Q8：5；Q9：7；Q10：10；Q11：8；Q12：4；Q13：12；Q14：6；Q15：14；Q16：14；
% Q17：6；Q18：8；Q19：7；Q20：19；Q21：28），与原卷印刷行逐题相等
% （cases 块作一个印刷行计）。全卷无 \ans{}、无只印【答案】的题。
%
\documentclass[a4paper,11pt]{article}
\usepackage{common}

\begin{document}

\examtitlest[3]{2026--2027 年华师大二附中高一上周末作业 6}{函数}

\bigsec{一、填空题}

\begin{enumerate}
\setcounter{enumi}{2}

\item 已知 $p:a=0$，$q$：函数 $f(x)=x^3+x+a$ 是奇函数，则 $p$ 是 $q$ 的\blankline[4.4em]条件.

\ifanswer
\solbegin
\step{若 $a=0$，则 $f(x)=x^3+x$，}
\step{易得 $f(x)$ 的定义域为 $\R$，且 $f(-x)=-(x^3+x)=-f(x)$，则 $f(x)$ 为奇函数，}
\step{反之，若 $f(x)=x^3+x+a$ 是奇函数，}
\step{则有 $f(-x)+f(x)=-x^3-x+a+x^3+x+a=2a=0$，必有 $a=0$，}
\step{故 $p$ 是 $q$ 的充要条件.}
\solend
\fi

\item 函数 $y=\sqrt{-x^2+2x}$ 的单调递增区间是\blankline[4.4em].

\ifanswer
\solbegin
\step{由函数定义域得 $-x^2+2x\geqslant0$，解得 $0\leqslant x\leqslant2$，}
\step{二次函数 $y=-x^2+2x$ 开口向下，对称轴为 $x=1$，}
\step{结合复合函数的单调性得函数 $y=\sqrt{-x^2+2x}$ 的单调递增区间是 $[0,1]$.}
\solend
\fi

\item 设 $f(x)$ 定义域为 $[1,4]$，且 $f(x)=2x+1$，则 $f(x^2)$ 的最大值为\blankline[4.4em].

\ifanswer
\solbegin
\step{由 $x^2\in[1,4]$，得 $f(x^2)=2x^2+1\leqslant9$，当 $x=\pm2$ 时取等号，最大值为 9.}
\solend
\fi

\item 函数 $f(x)=\sqrt{x-2}-x$，$x\in(2,6]$ 的值域为\blankline[4.4em].

\ifanswer
\solbegin
\step{令 $t=\sqrt{x-2}$，则 $x=2+t^2$，由 $2<x\leqslant6$ 得 $0<t\leqslant2$，}
\step{原函数可化为 $g(t)=-t^2+t-2$，$g(0)=-2$，$g(2)=-4$，$g\left(\dfrac12\right)=-\dfrac74$，}
\step{由二次函数性质得 $g(t)\in\left[-4,-\dfrac74\right]$.}
\solend
\fi

\item 若函数 $f(x)=\lvert x-a+1\rvert$ 在区间 $[1,+\infty)$ 上是严格增函数，则实数 $a$ 的取值范围是\blankline[4.4em].

\ifanswer
\solbegin
\step{$f(x)=\begin{cases}x-a+1,x\geqslant a-1\\a-1-x,x<a-1\end{cases}$，显然 $f(x)$ 在 $[a-1,+\infty)$ 上严格增，}
\step{要使函数在区间 $[1,+\infty)$ 上是严格增函数，则 $a-1\leqslant1$，即 $a\leqslant2$.}
\solend
\fi

\item 若函数 $f(x)=\begin{cases}x^2-5,x\leqslant2\\2ax+1,x>2\end{cases}$ 的值域是 $\R$，则 $a$ 的取值范围是\blankline[4.4em].

\ifanswer
\solbegin
\step{当 $x\leqslant2$ 时，$f(x)=x^2-5$，其开口向上，对称轴为 $x=0$，值域为 $[-5,+\infty)$，}
\step{由函数 $f(x)=\begin{cases}x^2-5,x\leqslant2\\2ax+1,x>2\end{cases}$ 的值域是 $\R$，}
\step{则当 $x>2$ 时，$f(x)=2ax+1$ 的值域应包含 $(-\infty,-5)$，}
\step{所以 $f(x)$ 为严格减函数，所以 $\begin{cases}2\times2a+1\geqslant-5\\2a<0\end{cases}$，即 $\begin{cases}a\geqslant-\dfrac32\\a<0\end{cases}$，}
\step{解得 $-\dfrac32\leqslant a<0$，故 $a$ 的取值范围是 $\left[-\dfrac32,0\right)$.}
\solend
\fi

\item 已知函数 $f(x)=ax^2+\dfrac bx$ 的图象经过点 $(-1,0)$，且满足 $f(1)=2$. 若 $\forall x\in(0,+\infty)$，都有 $f(x)+f\left(\dfrac1x\right)-2m\geqslant0$ 恒成立，则实数 $m$ 的取值范围为\blankline[4.4em].

\ifanswer
\solbegin
\step{因为 $f(x)=ax^2+\dfrac bx$ 的图象经过点 $(-1,0)$，且满足 $f(1)=2$，}
\step{所以 $\begin{cases}a-b=0\\a+b=2\end{cases}$，所以 $a=1$，$b=1$，即 $f(x)=x^2+\dfrac1x$，}
\step{$\forall x\in(0,+\infty)$，$f(x)+f\left(\dfrac1x\right)-2m\geqslant0\Leftrightarrow2m\leqslant x^2+\dfrac1x+\dfrac1{x^2}+x$ 恒成立，}
\step{令函数 $g(x)=x^2+\dfrac1x+\dfrac1{x^2}+x=\left(x+\dfrac1x\right)^2+x+\dfrac1x-2$，}
\step{而当 $x>0$ 时，$x+\dfrac1x\geqslant2$，当且仅当 $x=1$ 时取等号，}
\step{因此 $g(x)\geqslant2^2+2-2=4$，当且仅当 $x=1$ 时取等号，则 $2m\leqslant4$，解得 $m\leqslant2$，}
\step{所以实数 $m\in(-\infty,2]$.}
\solend
\fi

\item 若定义在区间 $[a,b](a<b)$ 上的函数 $f(x)=k-\sqrt{x+1}$ 值域也为 $[a,b]$，则实数 $k$ 的取值范围是\blankline[4.4em].

\ifanswer
\solbegin
\step{函数 $f(x)=k-\sqrt{x+1}$ 在定义域 $[-1,+\infty)$ 严格减，}
\step{当 $f(x)$ 的定义域为 $[a,b]$ 时，$f(x)$ 的值域也为 $[a,b]$，$a<b$，}
\step{故 $f(a)=k-\sqrt{a+1}=b$，$f(b)=k-\sqrt{b+1}=a$，}
\step{两式相减得 $\sqrt{a+1}-\sqrt{b+1}=a-b=(a+1)-(b+1)$}
\step{$=(\sqrt{a+1}-\sqrt{b+1})(\sqrt{a+1}+\sqrt{b+1})$，}
\step{所以 $\sqrt{a+1}-\sqrt{b+1}<0$，即 $\sqrt{a+1}+\sqrt{b+1}=1$，}
\step{则 $k=\sqrt{b+1}+a=a+1-\sqrt{a+1}$，}
\step{令 $\lambda=\sqrt{a+1}\geqslant0$，得 $k=\lambda^2-\lambda=\left(\lambda-\dfrac12\right)^2-\dfrac14$，}
\step{又 $\sqrt{a+1}<\sqrt{b+1}$，因为 $\sqrt{a+1}+\sqrt{b+1}=1$，所以 $\sqrt{a+1}\in\left[0,\dfrac12\right)$，}
\step{所以 $0\leqslant\lambda<\dfrac12$，故实数 $k$ 的取值范围为 $\left(-\dfrac14,0\right]$.}
\solend
\fi

\item 函数 $f(x)=[x]$ 的函数值表示不超过 $x$ 的最大整数，例如：$[-3.5]=-4$，$[2.1]=2$，称 $f(x)=[x]$ 为取整函数，则下列关于取整函数的说法：①对任意的 $x\in\R$，$[-x]=-[x]$ ②对任意的 $x$、$y\in\R$，$[x]+[y]\leqslant[x+y]$ ③集合 $\{y\mid y=[nx]-n[x],x\in\R,n\in\N_+\}$ 共有 $n$ 个元素 ④$x>0$ 时，关于 $x$ 的方程 $[2x]=\left[\dfrac1x\right]$ 有无数个解 其中正确的有\blankline[4.4em].

\ifanswer
\solbegin
\step{对①，因为 $[-2.1]=-3\neq[2.1]=2$，所以①错误；}
\step{对②，设 $x=[x]+\Delta x$，$\Delta x\in[0,1)$，$y=y+\Delta y$，$\Delta y\in[0,1)$，}% 原卷印作 $y=y+\Delta y$（按上下文应为 $y=[y]+\Delta y$），照录
\step{$[x+y]=[[x]+\Delta x+[y]+\Delta y]=[x]+[y]+[\Delta x+\Delta y]$，又 $\Delta x+\Delta y\in[0,\ 2)$，}
\step{故 $\Delta x+\Delta y=0$ 或 1，故 $[x]+[y]\leqslant[x+y]$，所以②正确；}
\step{对③，$[nx]-n[x]=[n[x]+n\cdot\Delta x]-n[x]=n[x]+[n\cdot\Delta x]-n[x]=[n\cdot\Delta x]$，}
\step{因为 $\Delta x\in[0,1)$，所以 $0\leqslant n\cdot\Delta x<n$，$n\in\N_+$，所以③正确；}
\step{对④，当 $\dfrac12<x<1$ 时，$[2x]=\left[\dfrac1x\right]=1$，故方程存在无数个解，所以④正确.}
\step{故正确的是②③④.}
\solend
\fi

\item 已知函数 $f(x)$ 定义域为 $[0,1]$，且满足 $f(x)+f(1-x)=1$，$f(x)=2f\left(\dfrac x5\right)$，对任意 $x_1\leqslant x_2$，恒有 $f(x_1)\leqslant f(x_2)$，则 $f\left(\dfrac1{2026}\right)$ 的值为\blankline[4.4em].

\ifanswer
\solbegin
\step{$f(0)=2f(0)\Rightarrow f(0)=0$，从而 $f(1)=1$，$f\left(\dfrac15\right)=\dfrac12$，$f\left(\dfrac45\right)=1-f\left(\dfrac15\right)=\dfrac12$，}
\step{由单调性，$x\in\left[\dfrac15,\dfrac45\right]$ 时，$f(x)=\dfrac12$.}
\step{因为 $\dfrac15<\dfrac{625}{2026}<\dfrac45$，反复使用 $f\left(\dfrac x5\right)=\dfrac12 f(x)$，}
\step{得 $f\left(\dfrac1{2026}\right)=\dfrac1{2^4}f\left(\dfrac{625}{2026}\right)=\dfrac1{32}$.}
\solend
\fi

\end{enumerate}

\bigsec{二、选择题}

\begin{enumerate}
\setcounter{enumi}{12}

\item 已知 $f(x)$ 为 $\R$ 上的奇函数，$g(x)$ 为 $\R$ 上的偶函数，且 $g(x)\neq0$，则下列说法正确的是\mbox{（\quad）}

\keepwithprev
A. $f(x)+g(x)$ 为 $\R$ 上的奇函数\qquad\qquad B. $f(x)-g(x)$ 为 $\R$ 上的奇函数

C. $\dfrac{f(x)}{g(x)}$ 为 $\R$ 上的偶函数\qquad\qquad D. $\lvert f(x)g(x)\rvert$ 为 $\R$ 上的偶函数

\ifanswer
\solbegin
\step{因为 $f(x)$ 为 $\R$ 上的奇函数，$g(x)$ 为 $\R$ 上的偶函数，且 $g(x)\neq0$，}
\step{所以 $f(-x)=-f(x)$，$g(-x)=g(x)$，}
\step{所以 $f(-x)+g(-x)=-f(x)+g(x)\neq-[f(x)+g(x)]$，}
\step{故 $f(x)+g(x)$ 为非奇非偶函数，$A$ 错误；}
\step{同理，$f(x)-g(x)$ 为非奇非偶函数，$B$ 错误；}
\step{设 $F(x)=\dfrac{f(x)}{g(x)}$，则 $F(-x)=\dfrac{f(-x)}{g(-x)}=\dfrac{-f(x)}{g(x)}=-F(x)$，}
\step{所以 $F(x)$ 为奇函数，$C$ 错误；}
\step{设函数 $H(x)=\lvert f(x)g(x)\rvert$，}
\step{因为 $f(x)$ 为 $\R$ 上的奇函数，$g(x)$ 为 $\R$ 上的偶函数，且 $g(x)\neq0$，}
\step{所以 $f(-x)=-f(x)$，$g(-x)=g(x)$，则由函数奇偶性的定义}
\step{得 $H(-x)=\lvert f(-x)g(-x)\rvert=\lvert -f(x)g(x)\rvert=\lvert f(x)g(x)\rvert=H(x)$，$D$ 正确.}
\step{故选 $D$.}
\solend
\fi

\item 若函数 $f(x)$ 是定义域为 $(-\infty,0)\cup(0,+\infty)$ 的奇函数，在区间 $(0,+\infty)$ 上严格增，且 $f(-1)=0$，则不等式 $\dfrac{f(x-1)}{x-2}<0$ 的解集为\mbox{（\quad）}

\keepwithprev
A. $(-\infty,-2)$\qquad\qquad B. $(-2,0)$\qquad\qquad C. $(0,1)$\qquad\qquad D. $(1,2)$

\ifanswer
\solbegin
\step{若函数 $f(x)$ 是定义域为 $(-\infty,0)\cup(0,+\infty)$ 的奇函数，在区间 $(0,+\infty)$ 上严格增，}
\step{由奇函数对称性得，$f(x)$ 在 $(-\infty,0)$ 上严格增，}
\step{因为 $f(-1)=0$，所以 $f(1)=0$，}
\step{则不等式 $\dfrac{f(x-1)}{x-2}<0$ 可转化为 $\begin{cases}x>2\\f(x-1)<0\end{cases}$ 或 $\begin{cases}x<2\\f(x-1)>0\end{cases}$，}
\step{即 $\begin{cases}x>2\\x-1<-1\text{或}0<x-1<1\end{cases}$ 或 $\begin{cases}x<2\\x-1>1\text{或}-1<x-1<0\end{cases}$，解得 $0<x<1$.}
\step{故选 $C$.}
\solend
\fi

\item 若函数 $f(x)$ 同时满足：(1) 对于定义域上的任意 $x$，恒有 $f(x)+f(-x)=0$；(2) 对于定义域上的任意 $x_1$、$x_2$，当 $x_1\neq x_2$ 时，恒有 $x_1f(x_1)-x_1f(x_2)>x_2f(x_1)-x_2f(x_2)$，则称函数 $f(x)$ 为“理想函数”. 给出下列四个函数：① $f(x)=x^2$；② $f(x)=x^3$；③ $f(x)=\dfrac{2x-1}{2x+1}$；④ $f(x)=\begin{cases}x^2+4x,x\geqslant0\\-x^2+4x,x<0\end{cases}$，其中被称为“理想函数”的有\mbox{（\quad）}

\keepwithprev
A. 1 个\qquad\qquad B. 2 个\qquad\qquad C. 3 个\qquad\qquad D. 4 个

\ifanswer
\solbegin
\step{函数 $f(x)$ 同时满足：(1) 对于定义域上的任意 $x$，恒有 $f(x)+f(-x)=0$；}
\step{(2) 对于定义域上的任意 $x_1$、$x_2$，当 $x_1\neq x_2$ 时，}
\step{恒有 $x_1f(x_1)-x_1f(x_2)>x_2f(x_1)-x_2f(x_2)$，}
\step{整理得 $(x_1-x_2)[f(x_1)-f(x_2)]>0$，则称函数 $f(x)$ 为“理想函数”.}
\step{即函数满足 (1) 奇函数；(2) 函数在定义域内为严格增函数；}
\step{① $f(x)=x^2$；不满足函数为奇函数在定义域内也不严格增；}
\step{② $f(x)=x^3$；满足函数既为奇函数，也在定义域内严格增，}
\step{故该函数为理想函数；}
\step{③ $f(x)=\dfrac{2x-1}{2x+1}=\dfrac{2x+1-2}{2x+1}=-\dfrac2{2x+1}+1$，该函数不为奇函数，}
\step{但是函数在定义域内严格增，故该函数不为理想函数；}
\step{④ $f(x)=\begin{cases}x^2+4x,x\geqslant0\\-x^2+4x,x<0\end{cases}$，函数的关系式满足 $f(-x)=-f(x)$，}
\step{且在 $x\geqslant0$ 时，函数严格增，由于奇函数在对称区间上单调性相同，}
\step{故该函数为理想函数.}
\step{故选 $B$.}
\solend
\fi

\item 已知 $f(x)$ 是定义在 $\R$ 上的奇函数，当 $x_1$、$x_2\in(0,+\infty)$ 且 $x_1\neq x_2$ 时，都有\newline $\dfrac{x_2f(x_1)-x_1f(x_2)}{x_1x_2(x_1-x_2)}>0$ 成立，$f(2028)=2028$，则不等式 $f(x)-x>0$ 的解集为\mbox{（\quad）}

\keepwithprev
A. $(-\infty,-2028)\cup(2028,+\infty)$\qquad\qquad B. $\left(-\dfrac1{2028},\dfrac1{2028}\right)$

C. $(-2028,2028)$\qquad\qquad D. $(-2028,0)\cup(2028,+\infty)$

\ifanswer
\solbegin
\step{因为 $f(x)$ 是定义在 $\R$ 上的奇函数，所以 $f(-x)=-f(x)$，}
\step{令 $g(x)=\dfrac{f(x)}x$，则 $g(-x)=\dfrac{f(-x)}{-x}=\dfrac{-f(x)}{-x}=\dfrac{f(x)}x=g(x)$，}
\step{所以 $g(x)$ 为偶函数，}
\step{当 $x_1$、$x_2\in(0,+\infty)$ 且 $x_1\neq x_2$ 时，都有 $\dfrac{x_2f(x_1)-x_1f(x_2)}{x_1x_2(x_1-x_2)}>0$ 成立，}
\step{即 $\dfrac{\dfrac{f(x_1)}{x_1}-\dfrac{f(x_2)}{x_2}}{x_1-x_2}>0$，即 $g(x)$ 在 $(0,+\infty)$ 上严格增，}
\step{由偶函数对称性得 $g(x)$ 在 $(-\infty,0)$ 上严格减，}
\step{因为 $f(2028)=2028$，所以 $g(2028)=1$，}
\step{则不等式 $f(x)-x>0$ 可化为 $f(x)>x$，}
\step{当 $x=0$ 时，$f(0)=0$，显然不成立，}
\step{当 $x>0$ 时，不等式可化为 $\dfrac{f(x)}x>1$，即 $g(x)>1=g(2028)$，}
\step{所以 $\lvert x\rvert>2028$，解得 $x>2028$，}
\step{当 $x<0$ 时，不等式可化为 $\dfrac{f(x)}x<1$，即 $g(x)<1=g(2028)$，}
\step{所以 $\lvert x\rvert<2028$，解得 $-2028<x<0$，}
\step{故不等式的解集为 $(-2028,0)\cup(2028,+\infty)$. 故选 $D$.}
\solend
\fi

\end{enumerate}

\bigsec{三、解答题}

\begin{enumerate}
\setcounter{enumi}{16}

\item 已知函数 $f(x)=\dfrac{2x-3}{x+1}$.

(1) 试用函数单调性的定义证明：函数 $f(x)$ 在 $(-1,+\infty)$ 上是严格增函数；

(2) 求函数 $f(x)$ 在区间 $[1,4]$ 上的最大值和最小值.

\ifanswer
\solbegin
\step{(1) 设 $-1<x_1<x_2$，则 $f(x_1)-f(x_2)=\dfrac{2x_1-3}{x_1+1}-\dfrac{2x_2-3}{x_2+1}=\dfrac{5(x_1-x_2)}{(x_1+1)(x_2+1)}$}
\step{因为 $-1<x_1<x_2$，所以 $x_1-x_2<0$，$x_1+1>0$，$x_2+1>0$，}
\step{所以 $f(x_1)-f(x_2)<0$，所以 $f(x)$ 在 $(-1,+\infty)$ 上严格增.}
\step{(2) 由 (1) 得 $f(x)$ 在区间 $[1,4]$ 上严格增，}
\step{所以 $f(x)_{\min}=f(1)=-\dfrac12$，$f(x)_{\max}=f(4)=1$，}
\step{所以函数 $f(x)$ 在区间 $[1,4]$ 上的最大值为 1，最小值为 $-\dfrac12$.}
\solend
\fi

\item 设函数 $f(x)=x^2-mx+m$.

(1) 当 $m=2$ 时，求不等式 $f(x)<x$ 的解集；

(2) 若 $\exists x\in[1,2]$，使不等式 $f(x)<x+m-2$ 成立，求实数 $m$ 的取值范围.

\ifanswer
\solbegin
\step{(1) $m=2$ 时，函数 $f(x)=x^2-2x+2$，}
\step{不等式 $f(x)<x$ 可化为 $x^2-3x+2<0$，解得 $1<x<2$，}
\step{所以不等式的解集为 $\{x\mid1<x<2\}$；}
\step{(2) 若 $\exists x\in[1,2]$，使不等式 $f(x)<x+m-2$ 成立，即 $x^2-x+2<mx$，}
\step{所以 $m>x+\dfrac2x-1$，设 $g(x)=x+\dfrac2x-1$，$x\in[1,2]$，}
\step{则 $g(x)=x+\dfrac2x-1\geqslant2\sqrt{x\cdot\dfrac2x}-1=2\sqrt2-1$，}
\step{当且仅当 $x=\dfrac2x$，即 $x=\sqrt2$ 时取等号，}
\step{所以实数 $m$ 的取值范围是 $\{m\mid m>2\sqrt2-1\}$.}
\solend
\fi

\item 某企业原来生产某种产品 $x$ 万件）可获利 $w(x)$（万元），且满足\newline% 原卷“万件）”右括号无配对、疑漏印左括号“（”，照录
$w(x)=\begin{cases}20(x^2+17),0<x\leqslant2\\500-\dfrac{80}{x-1},2<x\leqslant5\end{cases}$.
现该企业计划引进新的生产设备和新的产品方案优化产品，优化后的产品的其他成本总投入为 $(20x+10)$ 万元. 由市场调研分析得知，当前产品供不应求. 记该企业优化后的产品的利润为 $f(x)$（单位：万元）.

(1) 求函数 $f(x)$ 的解析式；

(2) 当优化后的产品产量为多少万件时，该企业的利润 $f(x)$ 最大？最大利润是多少？请说明理由.

\ifanswer
\solbegin
\step{(1) 由题意得 $f(x)=\begin{cases}20x^2-20x+330,0<x\leqslant2\\-20x-\dfrac{80}{x-1}+490,2<x\leqslant5\end{cases}$；}
\step{(2) 当 $0<x\leqslant2$ 时，$f(x)=20(x^2-x)+330=20\left(x-\dfrac12\right)^2+325$，}
\step{故当 $x=2$ 时，$f(x)$ 取最大值，$f(2)=370$；}
\step{当 $2<x\leqslant5$ 时，$f(x)=-\left[20(x-1)+\dfrac{80}{x-1}\right]+470\leqslant-80+470=390$，}
\step{当且仅当 $20(x-1)=\dfrac{80}{x-1}$，即 $x=3$ 时，$f(3)=390$ 为最大值；}
\step{因为 $370<390$，}
\step{所以优化后的产品产量为 3 万件时，企业获最大利润 390 万元.}
\solend
\fi

\item 俄国数学家切比雪夫（1821—1894）是研究直线逼近函数的理论先驱. 对定义在非空集合 $I$ 上的函数 $f(x)$，以及函数 $g(x)=kx+b(k,b\in\R)$，切比雪夫将函数 $y=\lvert f(x)-g(x)\rvert,x\in I$ 的最大值称为 $f(x),g(x)$ 的“偏差”.

(1) 函数 $f(x)=x^2,x\in[0,1],g(x)=-x-2$，求 $f(x),g(x)$ 的“偏差”；

(2) 函数 $f(x)=\dfrac1x+1,x\in[1,2],g(x)=kx+1(k>0)$，若 $f(x),g(x)$ 的“偏差”为 2，求 $k$ 的值；

(3) 函数 $f(x)=x^2-x,x\in[0,3],g(x)=2x+b$，当 $f(x),g(x)$ 的“偏差”取最小值时，求 $b$ 的值，并求出“偏差”的最小值.

\ifanswer
% 注：原卷本题解析标签印作【解得】（全卷唯一，其余 18 题均印【解析】）；
%     本稿按全稿惯例由 \solbegin 统一印【解析】，正文内容照录。见文件头排版差异⑥。
\solbegin
\step{(1) $y=\lvert f(x)-g(x)\rvert=\lvert x^2+x+2\rvert=\left\lvert\left(x+\dfrac12\right)^2+\dfrac74\right\rvert,x\in[0,1]$，}
\step{因为 $x\in[0,1]$，所以 $\left(x+\dfrac12\right)^2+\dfrac74\in[2,4]$，}
\step{则 $y=\left\lvert\left(x+\dfrac12\right)^2+\dfrac74\right\rvert\in[2,4]$，$y_{\max}=4$，}
\step{所以函数 $f(x)$ 与 $g(x)$ 的“偏差”为 4.}
\step{(2) 令 $t=f(x)-g(x)=\dfrac1x+1-(kx+1)=\dfrac1x-kx,x\in[1,2]$，}
\step{因为 $k>0$，所以 $t=\dfrac1x-kx$ 是单调减函数，}
\step{所以 $t\in\left[\dfrac12-2k,1-k\right]$，由单调，$y=\lvert t\rvert,t\in\left[\dfrac12-2k,1-k\right]$，且 $y_{\max}=2$.}
\step{当 $\left\lvert\dfrac12-2k\right\rvert\leqslant\lvert1-k\rvert$，即 $0<k\leqslant\dfrac12$ 时，$\lvert1-k\rvert=2$，解得 $k=3$ 或 $k=-1$，}
\step{均不符合；}
\step{当 $\left\lvert\dfrac12-2k\right\rvert>\lvert1-k\rvert$，即 $k>\dfrac12$ 时，$\left\lvert\dfrac12-2k\right\rvert=2$，$\dfrac12-2k=2$ 或 $\dfrac12-2k=-2$，}
\step{解得 $k=\dfrac54$ 或 $k=-\dfrac34$（舍），所以 $k=\dfrac54$.}
\step{(3) $y=\lvert f(x)-g(x)\rvert=\lvert x^2-x-(2x+b)\rvert=\lvert x^2-3x-b\rvert$}
\step{$=\left\lvert\left(x-\dfrac32\right)^2-\dfrac94-b\right\rvert,x\in[0,3]$，}
\step{因为 $x\in[0,3]$，所以 $\left(x-\dfrac32\right)^2-\dfrac94-b\in\left[-\dfrac94-b,-b\right]$，}
\step{由 $y=\left\lvert\left(x-\dfrac32\right)^2-\dfrac94-b\right\rvert$，则 $y_{\max}=\max\left\{\lvert b\rvert,\left\lvert b+\dfrac94\right\rvert\right\}$，}
\step{令 $\lvert b\rvert\geqslant\left\lvert b+\dfrac94\right\rvert$，即 $b^2\geqslant\left(b+\dfrac94\right)^2$，解得 $b\leqslant-\dfrac98$，}
\step{$y_{\max}=\max\left\{\lvert b\rvert,\left\lvert b+\dfrac94\right\rvert\right\}=\begin{cases}-b,b\leqslant-\dfrac98\\b+\dfrac94,b>-\dfrac98\end{cases}$，}
\step{故当且仅当 $b=-\dfrac98$ 时，$(y_{\max})_{\min}=\dfrac98$，}
\step{故当 $b$ 的值为 $-\dfrac98$ 时，函数 $f(x)$ 与 $g(x)$ 的“偏差”取最小值 $\dfrac98$.}
\solend
\fi

\item 若函数 $f(x)$ 与 $g(x)$ 满足：对任意的 $x_1\in D$，总存在唯一的 $x_2\in D$，使 $f(x_1)g(x_2)=m$ 成立，则称 $f(x)$ 是 $g(x)$ 在区间 $D$ 上的“$m$ 阶伴随函数”；对任意的 $x_1\in D$，总存在唯一的 $x_2\in D$，使 $f(x_1)f(x_2)=m$ 成立，则称 $f(x)$ 是区间 $D$ 上的“$m$ 阶自伴函数”.

(1) 判断 $f(x)=x^2+1$ 是否为区间 $[0,3]$ 上的“2 阶自伴函数”？并说明理由：

(2) 若函数 $f(x)=3x-1$ 为区间 $\left[\dfrac12,b\right]$ 上的“1 阶自伴函数”，求 $b$ 的值；

(3) 若 $f(x)=\dfrac4{x+2}$ 是 $g(x)=x^2-2ax+a^2-1$ 在区间 $[0,2]$ 上的“2 阶伴随函数”，求实数 $a$ 的取值范围.

\ifanswer
\solbegin
\step{(1) 假定函数 $f(x)=x^2+1$ 是区间 $[0,3]$ 上的“2 阶自伴函数”，}
\step{取 $x_1=2,f(x_1)=f(2)=5$，由 $f(x_1)f(x_2)=2$，得 $5(x_2^2+1)=2$，}
\step{显然此方程无实数解，}
\step{所以函数 $f(x)=x^2+1$ 不是区间 $[0,3]$ 上的“2 阶自伴函数”.}
\step{(2) 函数 $f(x)=3x-1$ 为区间 $\left[\dfrac12,b\right]$ 上的“1 阶自伴函数”，}
\step{则对任意 $x_1\in\left[\dfrac12,b\right]$，总存在唯一的 $x_2\in\left[\dfrac12,b\right]$，使得 $f(x_1)f(x_2)=1$，}
\step{即 $3x_2-1=\dfrac1{3x_1-1}$，整理得 $x_2=\dfrac13+\dfrac1{3(3x_1-1)}$，}
\step{显然函数 $x_2=\dfrac13+\dfrac1{3(3x_1-1)}$ 在 $\left[\dfrac12,b\right]$ 上严格减，}
\step{且当 $x_1=\dfrac12$ 时，$x_2=1$，当 $x_1=b$ 时，$x_2=\dfrac13+\dfrac1{9b-3}$，}
\step{因此对 $\left[\dfrac12,b\right]$ 内的每一个 $x_1$，在 $\left[\dfrac13+\dfrac1{9b-3},1\right]$ 内有唯一 $x_2$ 值与之对应，}
\step{则 $x_2\in\left[\dfrac12,b\right]$，于是 $\left[\dfrac13+\dfrac1{9b-3},1\right]\subseteq\left[\dfrac12,b\right]$，应有 $\begin{cases}b\geqslant1\\\dfrac13+\dfrac1{9b-3}\geqslant\dfrac12\end{cases}$，}
\step{得 $\begin{cases}b\geqslant1\\b\leqslant1\end{cases}$，即 $b=1$，所以 $b$ 的值是 1.}
\step{(3) 由函数 $f(x)=\dfrac4{x+2}$ 在区间 $[0,2]$ 的值域为 $[1,2]$，}
\step{因为 $f(x)=\dfrac4{x+2}$ 是 $g(x)=x^2-2ax+a^2-1$ 在 $[0,2]$ 上的“2 阶伴随函数”，}
\step{则对任意的 $x_1\in[0,2]$，总存在唯一的 $x_2\in[0,2]$ 时，使得 $f(x_1)g(x_2)=2$，}% 原卷此处印“时”字（按上下文为衍字），照录
\step{所以 $g(x_2)=\dfrac2{f(x_1)}\in[1,2]$，}
\step{即 $g(x)=x^2-2ax+a^2-1$ 在 $[0,2]$ 区间上的值域必定包含区间 $[1,2]$，}
\step{且 $g(x)$ 值域在 $[1,2]$ 对应的自变量是唯一的，}
\step{又因为函数 $g(x)=x^2-2ax+a^2-1$ 开口向上，对称轴为 $x=a$，}
\step{1) 当 $a\leqslant0$ 时，$g(x)$ 在 $[0,2]$ 区间上严格增，则必有 $\begin{cases}g(x)_{\min}=g(0)\leqslant1\\g(x)_{\max}=g(2)\geqslant2\end{cases}$，}
\step{解得 $-\sqrt2\leqslant a\leqslant0$；}
\step{2) 当 $a\geqslant2$ 时，$g(x)$ 在 $[0,2]$ 区间上严格减，则必有 $\begin{cases}g(x)_{\min}=g(2)\leqslant1\\g(x)_{\max}=g(0)\geqslant2\end{cases}$，}
\step{解得 $2\leqslant a\leqslant2+\sqrt2$；}
\step{3) 当 $0<a\leqslant1$ 时，$g(x)$ 在 $[0,a]$ 上严格减，在 $[a,2]$ 上严格增，}
\step{则必有 $\begin{cases}g(0)<1\\g(x)_{\max}=g(2)\geqslant2\end{cases}$，解得 $0<a\leqslant2-\sqrt3$，}
\step{4) 当 $1<a<2$ 时，$g(x)$ 在 $[0,a]$ 上严格减，在 $[a,2]$ 上严格增，}
\step{则必有 $\begin{cases}g(2)<1\\g(x)_{\max}=g(0)\geqslant2\end{cases}$，解得 $\sqrt3\leqslant a<2$.}
\step{综上所述，$a$ 的范围 $[-\sqrt2,2-\sqrt3]\cup[\sqrt3,2+\sqrt2]$.}% 原卷“范围”后缺“是／为”字，照录
\solend
\fi

\end{enumerate}

\end{document}
"
% 紧凑版：xelatex "\def\iscompact{1}% 高一122_华师大二附中_周末作业6.tex
%   —— 高一上数学“2026-2027 年华师大二附中高一上周末作业 6”（试卷版/答案版）
% 试卷版：xelatex 高一122_华师大二附中_周末作业6.tex
% 答案版：xelatex "\def\isanswer{1}\input{高一122_华师大二附中_周末作业6.tex}"
% 紧凑版：xelatex "\def\iscompact{1}\input{高一122_华师大二附中_周末作业6.tex}"
%
% 原始材料：微信公众号“魔都数学加油站”文章，作者破晓，2026-11-09 17:52 发布；
% 连载“27学年高一122”，原文标题“27学年高一122——2026-2027年上海市华师大二附中高一上周末作业6”；
% 原文链接：https://mp.weixin.qq.com/s/pC5066JFKNkQQwjMqTK8Ow。
% 正文为 14 张 PNG 页（第 14 页 2560×3620，其余 13 页 4762×6735），
% 题目下印有【解析】，为讲评版。
% 题目、解析均据原图转录；卷面红色水印及公众号页脚未录入。全卷无插图（已逐页检索，无“如图”）。
%
% 卷首标题：原卷印作“2026-2027 年华师大二附中高一上周末作业 6”（公众号标题中的“上海市”
% 卷面标题行没有，本稿照卷面），年份区间按全稿惯例排 en dash；原卷未印副标题，本稿按题目内容
% 标为“函数”。副标题理由（以全部 14 页实际内容为准）：逐页复核 19 题，考的都是函数性质——
% 奇偶性与充要条件（Q3、Q13、Q15、Q16）、单调性与单调区间（Q4、Q7、Q14、Q17）、定义域／
% 值域／最值与复合函数（Q5、Q6、Q8、Q10、Q12、Q19、Q20）、分段函数与取整函数（Q7、Q8、
% Q11、Q15、Q19）、抽象函数与新定义（Q12、Q20、Q21）；不等式仅作求参数范围／解集的工具
% （Q9、Q14、Q16、Q18），集合与常用逻辑用语仅零星出现（Q3 的充要条件、Q11③ 的集合、
% Q20 的“非空集合 I”），均非独立章节。故主体为单一模块“函数”，与同校 高一102 周末作业 5
% 同（该函数模块卷亦按 高一83／105 成例只用模块名、不叠“集合与常用逻辑用语、不等式”）。
%
% 与原卷的排版差异：
%   ① 原文从第 3 题起，填空题编号为 3—12、选择题为 13—16、解答题为 17—21，本稿保留原题号；
%   ② 原卷每题下直接印【解析】，本稿按全稿惯例置于答案版；四道选择题括号内均未印答案、
%      答案在解析末句“故选 …”，括号内留空（全卷无只印【答案】的题、无 \ans{}）；
%   ③ 原卷未印副标题，本稿按“知识模块”口径补出；
%   ④ 原卷填空横线按全稿惯例排为 \blankline[4.4em]；选择题括号排 \mbox{（\quad）}、
%      选项段前按惯例加 \keepwithprev；解答题不留作答区（原卷为讲评版）；
%   ⑤ 原卷数集排斜体 $R$、$N_+$（下标加号），按全稿惯例改 \R、\N_+
%      （$N_+$ 下标加号与 $\N^*$ 是两种记号，照录下标、只改数集字母为正体，同 高一05 第 20 题成例）；
%      绝对值排 \lvert…\rvert、集合描述竖线排 \mid；
%      $f(x)_{min}$／$f(x)_{max}$／$y_{max}$／$(y_{max})_{min}$ 一类斜体下标按 高一96 成例
%      改正体 \min／\max；不等式一律排 \geqslant／\leqslant；
%   ⑥ 第 20 题解析标签原卷印作【解得】（全卷唯一，其余 18 题均印【解析】），本稿按全稿惯例
%      由 \solbegin 统一印【解析】，正文内容照录（行内加 % 注、并入文件头存疑清单）；
%   ⑦ 题干与解析的标点（半角／全角逗号、冒号、括号、句点、子题括号）均照原卷，
%      题目正文引号用弯引号“ ”。
%
% 子集符号：第 21 题 (2) 的 $\left[\frac13+\frac1{9b-3},1\right]\subseteq\left[\frac12,b\right]$
%   印有下划线，照录 \subseteq。全卷未出现判别式 \triangle／\Delta。
%
% 原卷存疑／笔误（均照录加注）：
%   第 11 题解析对② 的“$y=y+\Delta y$”（按上下文应为 $y=[y]+\Delta y$，原卷漏方括号）；
%   第 19 题题干“某种产品 $x$ 万件）可获利”右括号无配对（疑漏印左括号“（”）；
%   第 20 题解析标签【解得】（见排版差异⑥）；
%   第 21 题 (3) 解析“总存在唯一的 $x_2\in[0,2]$ 时，使得……”的“时”字（按上下文衍字）；
%   第 21 题 (3) 末句“综上所述，$a$ 的范围 $[-\sqrt2,2-\sqrt3]\cup[\sqrt3,2+\sqrt2]$”
%      “范围”后缺“是／为”字，照录。
%
% 行数核对：逐题数 \step 条数，合计 162 条（Q3：5；Q4：3；Q5：1；Q6：3；Q7：2；
% Q8：5；Q9：7；Q10：10；Q11：8；Q12：4；Q13：12；Q14：6；Q15：14；Q16：14；
% Q17：6；Q18：8；Q19：7；Q20：19；Q21：28），与原卷印刷行逐题相等
% （cases 块作一个印刷行计）。全卷无 \ans{}、无只印【答案】的题。
%
\documentclass[a4paper,11pt]{article}
\usepackage{common}

\begin{document}

\examtitlest[3]{2026--2027 年华师大二附中高一上周末作业 6}{函数}

\bigsec{一、填空题}

\begin{enumerate}
\setcounter{enumi}{2}

\item 已知 $p:a=0$，$q$：函数 $f(x)=x^3+x+a$ 是奇函数，则 $p$ 是 $q$ 的\blankline[4.4em]条件.

\ifanswer
\solbegin
\step{若 $a=0$，则 $f(x)=x^3+x$，}
\step{易得 $f(x)$ 的定义域为 $\R$，且 $f(-x)=-(x^3+x)=-f(x)$，则 $f(x)$ 为奇函数，}
\step{反之，若 $f(x)=x^3+x+a$ 是奇函数，}
\step{则有 $f(-x)+f(x)=-x^3-x+a+x^3+x+a=2a=0$，必有 $a=0$，}
\step{故 $p$ 是 $q$ 的充要条件.}
\solend
\fi

\item 函数 $y=\sqrt{-x^2+2x}$ 的单调递增区间是\blankline[4.4em].

\ifanswer
\solbegin
\step{由函数定义域得 $-x^2+2x\geqslant0$，解得 $0\leqslant x\leqslant2$，}
\step{二次函数 $y=-x^2+2x$ 开口向下，对称轴为 $x=1$，}
\step{结合复合函数的单调性得函数 $y=\sqrt{-x^2+2x}$ 的单调递增区间是 $[0,1]$.}
\solend
\fi

\item 设 $f(x)$ 定义域为 $[1,4]$，且 $f(x)=2x+1$，则 $f(x^2)$ 的最大值为\blankline[4.4em].

\ifanswer
\solbegin
\step{由 $x^2\in[1,4]$，得 $f(x^2)=2x^2+1\leqslant9$，当 $x=\pm2$ 时取等号，最大值为 9.}
\solend
\fi

\item 函数 $f(x)=\sqrt{x-2}-x$，$x\in(2,6]$ 的值域为\blankline[4.4em].

\ifanswer
\solbegin
\step{令 $t=\sqrt{x-2}$，则 $x=2+t^2$，由 $2<x\leqslant6$ 得 $0<t\leqslant2$，}
\step{原函数可化为 $g(t)=-t^2+t-2$，$g(0)=-2$，$g(2)=-4$，$g\left(\dfrac12\right)=-\dfrac74$，}
\step{由二次函数性质得 $g(t)\in\left[-4,-\dfrac74\right]$.}
\solend
\fi

\item 若函数 $f(x)=\lvert x-a+1\rvert$ 在区间 $[1,+\infty)$ 上是严格增函数，则实数 $a$ 的取值范围是\blankline[4.4em].

\ifanswer
\solbegin
\step{$f(x)=\begin{cases}x-a+1,x\geqslant a-1\\a-1-x,x<a-1\end{cases}$，显然 $f(x)$ 在 $[a-1,+\infty)$ 上严格增，}
\step{要使函数在区间 $[1,+\infty)$ 上是严格增函数，则 $a-1\leqslant1$，即 $a\leqslant2$.}
\solend
\fi

\item 若函数 $f(x)=\begin{cases}x^2-5,x\leqslant2\\2ax+1,x>2\end{cases}$ 的值域是 $\R$，则 $a$ 的取值范围是\blankline[4.4em].

\ifanswer
\solbegin
\step{当 $x\leqslant2$ 时，$f(x)=x^2-5$，其开口向上，对称轴为 $x=0$，值域为 $[-5,+\infty)$，}
\step{由函数 $f(x)=\begin{cases}x^2-5,x\leqslant2\\2ax+1,x>2\end{cases}$ 的值域是 $\R$，}
\step{则当 $x>2$ 时，$f(x)=2ax+1$ 的值域应包含 $(-\infty,-5)$，}
\step{所以 $f(x)$ 为严格减函数，所以 $\begin{cases}2\times2a+1\geqslant-5\\2a<0\end{cases}$，即 $\begin{cases}a\geqslant-\dfrac32\\a<0\end{cases}$，}
\step{解得 $-\dfrac32\leqslant a<0$，故 $a$ 的取值范围是 $\left[-\dfrac32,0\right)$.}
\solend
\fi

\item 已知函数 $f(x)=ax^2+\dfrac bx$ 的图象经过点 $(-1,0)$，且满足 $f(1)=2$. 若 $\forall x\in(0,+\infty)$，都有 $f(x)+f\left(\dfrac1x\right)-2m\geqslant0$ 恒成立，则实数 $m$ 的取值范围为\blankline[4.4em].

\ifanswer
\solbegin
\step{因为 $f(x)=ax^2+\dfrac bx$ 的图象经过点 $(-1,0)$，且满足 $f(1)=2$，}
\step{所以 $\begin{cases}a-b=0\\a+b=2\end{cases}$，所以 $a=1$，$b=1$，即 $f(x)=x^2+\dfrac1x$，}
\step{$\forall x\in(0,+\infty)$，$f(x)+f\left(\dfrac1x\right)-2m\geqslant0\Leftrightarrow2m\leqslant x^2+\dfrac1x+\dfrac1{x^2}+x$ 恒成立，}
\step{令函数 $g(x)=x^2+\dfrac1x+\dfrac1{x^2}+x=\left(x+\dfrac1x\right)^2+x+\dfrac1x-2$，}
\step{而当 $x>0$ 时，$x+\dfrac1x\geqslant2$，当且仅当 $x=1$ 时取等号，}
\step{因此 $g(x)\geqslant2^2+2-2=4$，当且仅当 $x=1$ 时取等号，则 $2m\leqslant4$，解得 $m\leqslant2$，}
\step{所以实数 $m\in(-\infty,2]$.}
\solend
\fi

\item 若定义在区间 $[a,b](a<b)$ 上的函数 $f(x)=k-\sqrt{x+1}$ 值域也为 $[a,b]$，则实数 $k$ 的取值范围是\blankline[4.4em].

\ifanswer
\solbegin
\step{函数 $f(x)=k-\sqrt{x+1}$ 在定义域 $[-1,+\infty)$ 严格减，}
\step{当 $f(x)$ 的定义域为 $[a,b]$ 时，$f(x)$ 的值域也为 $[a,b]$，$a<b$，}
\step{故 $f(a)=k-\sqrt{a+1}=b$，$f(b)=k-\sqrt{b+1}=a$，}
\step{两式相减得 $\sqrt{a+1}-\sqrt{b+1}=a-b=(a+1)-(b+1)$}
\step{$=(\sqrt{a+1}-\sqrt{b+1})(\sqrt{a+1}+\sqrt{b+1})$，}
\step{所以 $\sqrt{a+1}-\sqrt{b+1}<0$，即 $\sqrt{a+1}+\sqrt{b+1}=1$，}
\step{则 $k=\sqrt{b+1}+a=a+1-\sqrt{a+1}$，}
\step{令 $\lambda=\sqrt{a+1}\geqslant0$，得 $k=\lambda^2-\lambda=\left(\lambda-\dfrac12\right)^2-\dfrac14$，}
\step{又 $\sqrt{a+1}<\sqrt{b+1}$，因为 $\sqrt{a+1}+\sqrt{b+1}=1$，所以 $\sqrt{a+1}\in\left[0,\dfrac12\right)$，}
\step{所以 $0\leqslant\lambda<\dfrac12$，故实数 $k$ 的取值范围为 $\left(-\dfrac14,0\right]$.}
\solend
\fi

\item 函数 $f(x)=[x]$ 的函数值表示不超过 $x$ 的最大整数，例如：$[-3.5]=-4$，$[2.1]=2$，称 $f(x)=[x]$ 为取整函数，则下列关于取整函数的说法：①对任意的 $x\in\R$，$[-x]=-[x]$ ②对任意的 $x$、$y\in\R$，$[x]+[y]\leqslant[x+y]$ ③集合 $\{y\mid y=[nx]-n[x],x\in\R,n\in\N_+\}$ 共有 $n$ 个元素 ④$x>0$ 时，关于 $x$ 的方程 $[2x]=\left[\dfrac1x\right]$ 有无数个解 其中正确的有\blankline[4.4em].

\ifanswer
\solbegin
\step{对①，因为 $[-2.1]=-3\neq[2.1]=2$，所以①错误；}
\step{对②，设 $x=[x]+\Delta x$，$\Delta x\in[0,1)$，$y=y+\Delta y$，$\Delta y\in[0,1)$，}% 原卷印作 $y=y+\Delta y$（按上下文应为 $y=[y]+\Delta y$），照录
\step{$[x+y]=[[x]+\Delta x+[y]+\Delta y]=[x]+[y]+[\Delta x+\Delta y]$，又 $\Delta x+\Delta y\in[0,\ 2)$，}
\step{故 $\Delta x+\Delta y=0$ 或 1，故 $[x]+[y]\leqslant[x+y]$，所以②正确；}
\step{对③，$[nx]-n[x]=[n[x]+n\cdot\Delta x]-n[x]=n[x]+[n\cdot\Delta x]-n[x]=[n\cdot\Delta x]$，}
\step{因为 $\Delta x\in[0,1)$，所以 $0\leqslant n\cdot\Delta x<n$，$n\in\N_+$，所以③正确；}
\step{对④，当 $\dfrac12<x<1$ 时，$[2x]=\left[\dfrac1x\right]=1$，故方程存在无数个解，所以④正确.}
\step{故正确的是②③④.}
\solend
\fi

\item 已知函数 $f(x)$ 定义域为 $[0,1]$，且满足 $f(x)+f(1-x)=1$，$f(x)=2f\left(\dfrac x5\right)$，对任意 $x_1\leqslant x_2$，恒有 $f(x_1)\leqslant f(x_2)$，则 $f\left(\dfrac1{2026}\right)$ 的值为\blankline[4.4em].

\ifanswer
\solbegin
\step{$f(0)=2f(0)\Rightarrow f(0)=0$，从而 $f(1)=1$，$f\left(\dfrac15\right)=\dfrac12$，$f\left(\dfrac45\right)=1-f\left(\dfrac15\right)=\dfrac12$，}
\step{由单调性，$x\in\left[\dfrac15,\dfrac45\right]$ 时，$f(x)=\dfrac12$.}
\step{因为 $\dfrac15<\dfrac{625}{2026}<\dfrac45$，反复使用 $f\left(\dfrac x5\right)=\dfrac12 f(x)$，}
\step{得 $f\left(\dfrac1{2026}\right)=\dfrac1{2^4}f\left(\dfrac{625}{2026}\right)=\dfrac1{32}$.}
\solend
\fi

\end{enumerate}

\bigsec{二、选择题}

\begin{enumerate}
\setcounter{enumi}{12}

\item 已知 $f(x)$ 为 $\R$ 上的奇函数，$g(x)$ 为 $\R$ 上的偶函数，且 $g(x)\neq0$，则下列说法正确的是\mbox{（\quad）}

\keepwithprev
A. $f(x)+g(x)$ 为 $\R$ 上的奇函数\qquad\qquad B. $f(x)-g(x)$ 为 $\R$ 上的奇函数

C. $\dfrac{f(x)}{g(x)}$ 为 $\R$ 上的偶函数\qquad\qquad D. $\lvert f(x)g(x)\rvert$ 为 $\R$ 上的偶函数

\ifanswer
\solbegin
\step{因为 $f(x)$ 为 $\R$ 上的奇函数，$g(x)$ 为 $\R$ 上的偶函数，且 $g(x)\neq0$，}
\step{所以 $f(-x)=-f(x)$，$g(-x)=g(x)$，}
\step{所以 $f(-x)+g(-x)=-f(x)+g(x)\neq-[f(x)+g(x)]$，}
\step{故 $f(x)+g(x)$ 为非奇非偶函数，$A$ 错误；}
\step{同理，$f(x)-g(x)$ 为非奇非偶函数，$B$ 错误；}
\step{设 $F(x)=\dfrac{f(x)}{g(x)}$，则 $F(-x)=\dfrac{f(-x)}{g(-x)}=\dfrac{-f(x)}{g(x)}=-F(x)$，}
\step{所以 $F(x)$ 为奇函数，$C$ 错误；}
\step{设函数 $H(x)=\lvert f(x)g(x)\rvert$，}
\step{因为 $f(x)$ 为 $\R$ 上的奇函数，$g(x)$ 为 $\R$ 上的偶函数，且 $g(x)\neq0$，}
\step{所以 $f(-x)=-f(x)$，$g(-x)=g(x)$，则由函数奇偶性的定义}
\step{得 $H(-x)=\lvert f(-x)g(-x)\rvert=\lvert -f(x)g(x)\rvert=\lvert f(x)g(x)\rvert=H(x)$，$D$ 正确.}
\step{故选 $D$.}
\solend
\fi

\item 若函数 $f(x)$ 是定义域为 $(-\infty,0)\cup(0,+\infty)$ 的奇函数，在区间 $(0,+\infty)$ 上严格增，且 $f(-1)=0$，则不等式 $\dfrac{f(x-1)}{x-2}<0$ 的解集为\mbox{（\quad）}

\keepwithprev
A. $(-\infty,-2)$\qquad\qquad B. $(-2,0)$\qquad\qquad C. $(0,1)$\qquad\qquad D. $(1,2)$

\ifanswer
\solbegin
\step{若函数 $f(x)$ 是定义域为 $(-\infty,0)\cup(0,+\infty)$ 的奇函数，在区间 $(0,+\infty)$ 上严格增，}
\step{由奇函数对称性得，$f(x)$ 在 $(-\infty,0)$ 上严格增，}
\step{因为 $f(-1)=0$，所以 $f(1)=0$，}
\step{则不等式 $\dfrac{f(x-1)}{x-2}<0$ 可转化为 $\begin{cases}x>2\\f(x-1)<0\end{cases}$ 或 $\begin{cases}x<2\\f(x-1)>0\end{cases}$，}
\step{即 $\begin{cases}x>2\\x-1<-1\text{或}0<x-1<1\end{cases}$ 或 $\begin{cases}x<2\\x-1>1\text{或}-1<x-1<0\end{cases}$，解得 $0<x<1$.}
\step{故选 $C$.}
\solend
\fi

\item 若函数 $f(x)$ 同时满足：(1) 对于定义域上的任意 $x$，恒有 $f(x)+f(-x)=0$；(2) 对于定义域上的任意 $x_1$、$x_2$，当 $x_1\neq x_2$ 时，恒有 $x_1f(x_1)-x_1f(x_2)>x_2f(x_1)-x_2f(x_2)$，则称函数 $f(x)$ 为“理想函数”. 给出下列四个函数：① $f(x)=x^2$；② $f(x)=x^3$；③ $f(x)=\dfrac{2x-1}{2x+1}$；④ $f(x)=\begin{cases}x^2+4x,x\geqslant0\\-x^2+4x,x<0\end{cases}$，其中被称为“理想函数”的有\mbox{（\quad）}

\keepwithprev
A. 1 个\qquad\qquad B. 2 个\qquad\qquad C. 3 个\qquad\qquad D. 4 个

\ifanswer
\solbegin
\step{函数 $f(x)$ 同时满足：(1) 对于定义域上的任意 $x$，恒有 $f(x)+f(-x)=0$；}
\step{(2) 对于定义域上的任意 $x_1$、$x_2$，当 $x_1\neq x_2$ 时，}
\step{恒有 $x_1f(x_1)-x_1f(x_2)>x_2f(x_1)-x_2f(x_2)$，}
\step{整理得 $(x_1-x_2)[f(x_1)-f(x_2)]>0$，则称函数 $f(x)$ 为“理想函数”.}
\step{即函数满足 (1) 奇函数；(2) 函数在定义域内为严格增函数；}
\step{① $f(x)=x^2$；不满足函数为奇函数在定义域内也不严格增；}
\step{② $f(x)=x^3$；满足函数既为奇函数，也在定义域内严格增，}
\step{故该函数为理想函数；}
\step{③ $f(x)=\dfrac{2x-1}{2x+1}=\dfrac{2x+1-2}{2x+1}=-\dfrac2{2x+1}+1$，该函数不为奇函数，}
\step{但是函数在定义域内严格增，故该函数不为理想函数；}
\step{④ $f(x)=\begin{cases}x^2+4x,x\geqslant0\\-x^2+4x,x<0\end{cases}$，函数的关系式满足 $f(-x)=-f(x)$，}
\step{且在 $x\geqslant0$ 时，函数严格增，由于奇函数在对称区间上单调性相同，}
\step{故该函数为理想函数.}
\step{故选 $B$.}
\solend
\fi

\item 已知 $f(x)$ 是定义在 $\R$ 上的奇函数，当 $x_1$、$x_2\in(0,+\infty)$ 且 $x_1\neq x_2$ 时，都有\newline $\dfrac{x_2f(x_1)-x_1f(x_2)}{x_1x_2(x_1-x_2)}>0$ 成立，$f(2028)=2028$，则不等式 $f(x)-x>0$ 的解集为\mbox{（\quad）}

\keepwithprev
A. $(-\infty,-2028)\cup(2028,+\infty)$\qquad\qquad B. $\left(-\dfrac1{2028},\dfrac1{2028}\right)$

C. $(-2028,2028)$\qquad\qquad D. $(-2028,0)\cup(2028,+\infty)$

\ifanswer
\solbegin
\step{因为 $f(x)$ 是定义在 $\R$ 上的奇函数，所以 $f(-x)=-f(x)$，}
\step{令 $g(x)=\dfrac{f(x)}x$，则 $g(-x)=\dfrac{f(-x)}{-x}=\dfrac{-f(x)}{-x}=\dfrac{f(x)}x=g(x)$，}
\step{所以 $g(x)$ 为偶函数，}
\step{当 $x_1$、$x_2\in(0,+\infty)$ 且 $x_1\neq x_2$ 时，都有 $\dfrac{x_2f(x_1)-x_1f(x_2)}{x_1x_2(x_1-x_2)}>0$ 成立，}
\step{即 $\dfrac{\dfrac{f(x_1)}{x_1}-\dfrac{f(x_2)}{x_2}}{x_1-x_2}>0$，即 $g(x)$ 在 $(0,+\infty)$ 上严格增，}
\step{由偶函数对称性得 $g(x)$ 在 $(-\infty,0)$ 上严格减，}
\step{因为 $f(2028)=2028$，所以 $g(2028)=1$，}
\step{则不等式 $f(x)-x>0$ 可化为 $f(x)>x$，}
\step{当 $x=0$ 时，$f(0)=0$，显然不成立，}
\step{当 $x>0$ 时，不等式可化为 $\dfrac{f(x)}x>1$，即 $g(x)>1=g(2028)$，}
\step{所以 $\lvert x\rvert>2028$，解得 $x>2028$，}
\step{当 $x<0$ 时，不等式可化为 $\dfrac{f(x)}x<1$，即 $g(x)<1=g(2028)$，}
\step{所以 $\lvert x\rvert<2028$，解得 $-2028<x<0$，}
\step{故不等式的解集为 $(-2028,0)\cup(2028,+\infty)$. 故选 $D$.}
\solend
\fi

\end{enumerate}

\bigsec{三、解答题}

\begin{enumerate}
\setcounter{enumi}{16}

\item 已知函数 $f(x)=\dfrac{2x-3}{x+1}$.

(1) 试用函数单调性的定义证明：函数 $f(x)$ 在 $(-1,+\infty)$ 上是严格增函数；

(2) 求函数 $f(x)$ 在区间 $[1,4]$ 上的最大值和最小值.

\ifanswer
\solbegin
\step{(1) 设 $-1<x_1<x_2$，则 $f(x_1)-f(x_2)=\dfrac{2x_1-3}{x_1+1}-\dfrac{2x_2-3}{x_2+1}=\dfrac{5(x_1-x_2)}{(x_1+1)(x_2+1)}$}
\step{因为 $-1<x_1<x_2$，所以 $x_1-x_2<0$，$x_1+1>0$，$x_2+1>0$，}
\step{所以 $f(x_1)-f(x_2)<0$，所以 $f(x)$ 在 $(-1,+\infty)$ 上严格增.}
\step{(2) 由 (1) 得 $f(x)$ 在区间 $[1,4]$ 上严格增，}
\step{所以 $f(x)_{\min}=f(1)=-\dfrac12$，$f(x)_{\max}=f(4)=1$，}
\step{所以函数 $f(x)$ 在区间 $[1,4]$ 上的最大值为 1，最小值为 $-\dfrac12$.}
\solend
\fi

\item 设函数 $f(x)=x^2-mx+m$.

(1) 当 $m=2$ 时，求不等式 $f(x)<x$ 的解集；

(2) 若 $\exists x\in[1,2]$，使不等式 $f(x)<x+m-2$ 成立，求实数 $m$ 的取值范围.

\ifanswer
\solbegin
\step{(1) $m=2$ 时，函数 $f(x)=x^2-2x+2$，}
\step{不等式 $f(x)<x$ 可化为 $x^2-3x+2<0$，解得 $1<x<2$，}
\step{所以不等式的解集为 $\{x\mid1<x<2\}$；}
\step{(2) 若 $\exists x\in[1,2]$，使不等式 $f(x)<x+m-2$ 成立，即 $x^2-x+2<mx$，}
\step{所以 $m>x+\dfrac2x-1$，设 $g(x)=x+\dfrac2x-1$，$x\in[1,2]$，}
\step{则 $g(x)=x+\dfrac2x-1\geqslant2\sqrt{x\cdot\dfrac2x}-1=2\sqrt2-1$，}
\step{当且仅当 $x=\dfrac2x$，即 $x=\sqrt2$ 时取等号，}
\step{所以实数 $m$ 的取值范围是 $\{m\mid m>2\sqrt2-1\}$.}
\solend
\fi

\item 某企业原来生产某种产品 $x$ 万件）可获利 $w(x)$（万元），且满足\newline% 原卷“万件）”右括号无配对、疑漏印左括号“（”，照录
$w(x)=\begin{cases}20(x^2+17),0<x\leqslant2\\500-\dfrac{80}{x-1},2<x\leqslant5\end{cases}$.
现该企业计划引进新的生产设备和新的产品方案优化产品，优化后的产品的其他成本总投入为 $(20x+10)$ 万元. 由市场调研分析得知，当前产品供不应求. 记该企业优化后的产品的利润为 $f(x)$（单位：万元）.

(1) 求函数 $f(x)$ 的解析式；

(2) 当优化后的产品产量为多少万件时，该企业的利润 $f(x)$ 最大？最大利润是多少？请说明理由.

\ifanswer
\solbegin
\step{(1) 由题意得 $f(x)=\begin{cases}20x^2-20x+330,0<x\leqslant2\\-20x-\dfrac{80}{x-1}+490,2<x\leqslant5\end{cases}$；}
\step{(2) 当 $0<x\leqslant2$ 时，$f(x)=20(x^2-x)+330=20\left(x-\dfrac12\right)^2+325$，}
\step{故当 $x=2$ 时，$f(x)$ 取最大值，$f(2)=370$；}
\step{当 $2<x\leqslant5$ 时，$f(x)=-\left[20(x-1)+\dfrac{80}{x-1}\right]+470\leqslant-80+470=390$，}
\step{当且仅当 $20(x-1)=\dfrac{80}{x-1}$，即 $x=3$ 时，$f(3)=390$ 为最大值；}
\step{因为 $370<390$，}
\step{所以优化后的产品产量为 3 万件时，企业获最大利润 390 万元.}
\solend
\fi

\item 俄国数学家切比雪夫（1821—1894）是研究直线逼近函数的理论先驱. 对定义在非空集合 $I$ 上的函数 $f(x)$，以及函数 $g(x)=kx+b(k,b\in\R)$，切比雪夫将函数 $y=\lvert f(x)-g(x)\rvert,x\in I$ 的最大值称为 $f(x),g(x)$ 的“偏差”.

(1) 函数 $f(x)=x^2,x\in[0,1],g(x)=-x-2$，求 $f(x),g(x)$ 的“偏差”；

(2) 函数 $f(x)=\dfrac1x+1,x\in[1,2],g(x)=kx+1(k>0)$，若 $f(x),g(x)$ 的“偏差”为 2，求 $k$ 的值；

(3) 函数 $f(x)=x^2-x,x\in[0,3],g(x)=2x+b$，当 $f(x),g(x)$ 的“偏差”取最小值时，求 $b$ 的值，并求出“偏差”的最小值.

\ifanswer
% 注：原卷本题解析标签印作【解得】（全卷唯一，其余 18 题均印【解析】）；
%     本稿按全稿惯例由 \solbegin 统一印【解析】，正文内容照录。见文件头排版差异⑥。
\solbegin
\step{(1) $y=\lvert f(x)-g(x)\rvert=\lvert x^2+x+2\rvert=\left\lvert\left(x+\dfrac12\right)^2+\dfrac74\right\rvert,x\in[0,1]$，}
\step{因为 $x\in[0,1]$，所以 $\left(x+\dfrac12\right)^2+\dfrac74\in[2,4]$，}
\step{则 $y=\left\lvert\left(x+\dfrac12\right)^2+\dfrac74\right\rvert\in[2,4]$，$y_{\max}=4$，}
\step{所以函数 $f(x)$ 与 $g(x)$ 的“偏差”为 4.}
\step{(2) 令 $t=f(x)-g(x)=\dfrac1x+1-(kx+1)=\dfrac1x-kx,x\in[1,2]$，}
\step{因为 $k>0$，所以 $t=\dfrac1x-kx$ 是单调减函数，}
\step{所以 $t\in\left[\dfrac12-2k,1-k\right]$，由单调，$y=\lvert t\rvert,t\in\left[\dfrac12-2k,1-k\right]$，且 $y_{\max}=2$.}
\step{当 $\left\lvert\dfrac12-2k\right\rvert\leqslant\lvert1-k\rvert$，即 $0<k\leqslant\dfrac12$ 时，$\lvert1-k\rvert=2$，解得 $k=3$ 或 $k=-1$，}
\step{均不符合；}
\step{当 $\left\lvert\dfrac12-2k\right\rvert>\lvert1-k\rvert$，即 $k>\dfrac12$ 时，$\left\lvert\dfrac12-2k\right\rvert=2$，$\dfrac12-2k=2$ 或 $\dfrac12-2k=-2$，}
\step{解得 $k=\dfrac54$ 或 $k=-\dfrac34$（舍），所以 $k=\dfrac54$.}
\step{(3) $y=\lvert f(x)-g(x)\rvert=\lvert x^2-x-(2x+b)\rvert=\lvert x^2-3x-b\rvert$}
\step{$=\left\lvert\left(x-\dfrac32\right)^2-\dfrac94-b\right\rvert,x\in[0,3]$，}
\step{因为 $x\in[0,3]$，所以 $\left(x-\dfrac32\right)^2-\dfrac94-b\in\left[-\dfrac94-b,-b\right]$，}
\step{由 $y=\left\lvert\left(x-\dfrac32\right)^2-\dfrac94-b\right\rvert$，则 $y_{\max}=\max\left\{\lvert b\rvert,\left\lvert b+\dfrac94\right\rvert\right\}$，}
\step{令 $\lvert b\rvert\geqslant\left\lvert b+\dfrac94\right\rvert$，即 $b^2\geqslant\left(b+\dfrac94\right)^2$，解得 $b\leqslant-\dfrac98$，}
\step{$y_{\max}=\max\left\{\lvert b\rvert,\left\lvert b+\dfrac94\right\rvert\right\}=\begin{cases}-b,b\leqslant-\dfrac98\\b+\dfrac94,b>-\dfrac98\end{cases}$，}
\step{故当且仅当 $b=-\dfrac98$ 时，$(y_{\max})_{\min}=\dfrac98$，}
\step{故当 $b$ 的值为 $-\dfrac98$ 时，函数 $f(x)$ 与 $g(x)$ 的“偏差”取最小值 $\dfrac98$.}
\solend
\fi

\item 若函数 $f(x)$ 与 $g(x)$ 满足：对任意的 $x_1\in D$，总存在唯一的 $x_2\in D$，使 $f(x_1)g(x_2)=m$ 成立，则称 $f(x)$ 是 $g(x)$ 在区间 $D$ 上的“$m$ 阶伴随函数”；对任意的 $x_1\in D$，总存在唯一的 $x_2\in D$，使 $f(x_1)f(x_2)=m$ 成立，则称 $f(x)$ 是区间 $D$ 上的“$m$ 阶自伴函数”.

(1) 判断 $f(x)=x^2+1$ 是否为区间 $[0,3]$ 上的“2 阶自伴函数”？并说明理由：

(2) 若函数 $f(x)=3x-1$ 为区间 $\left[\dfrac12,b\right]$ 上的“1 阶自伴函数”，求 $b$ 的值；

(3) 若 $f(x)=\dfrac4{x+2}$ 是 $g(x)=x^2-2ax+a^2-1$ 在区间 $[0,2]$ 上的“2 阶伴随函数”，求实数 $a$ 的取值范围.

\ifanswer
\solbegin
\step{(1) 假定函数 $f(x)=x^2+1$ 是区间 $[0,3]$ 上的“2 阶自伴函数”，}
\step{取 $x_1=2,f(x_1)=f(2)=5$，由 $f(x_1)f(x_2)=2$，得 $5(x_2^2+1)=2$，}
\step{显然此方程无实数解，}
\step{所以函数 $f(x)=x^2+1$ 不是区间 $[0,3]$ 上的“2 阶自伴函数”.}
\step{(2) 函数 $f(x)=3x-1$ 为区间 $\left[\dfrac12,b\right]$ 上的“1 阶自伴函数”，}
\step{则对任意 $x_1\in\left[\dfrac12,b\right]$，总存在唯一的 $x_2\in\left[\dfrac12,b\right]$，使得 $f(x_1)f(x_2)=1$，}
\step{即 $3x_2-1=\dfrac1{3x_1-1}$，整理得 $x_2=\dfrac13+\dfrac1{3(3x_1-1)}$，}
\step{显然函数 $x_2=\dfrac13+\dfrac1{3(3x_1-1)}$ 在 $\left[\dfrac12,b\right]$ 上严格减，}
\step{且当 $x_1=\dfrac12$ 时，$x_2=1$，当 $x_1=b$ 时，$x_2=\dfrac13+\dfrac1{9b-3}$，}
\step{因此对 $\left[\dfrac12,b\right]$ 内的每一个 $x_1$，在 $\left[\dfrac13+\dfrac1{9b-3},1\right]$ 内有唯一 $x_2$ 值与之对应，}
\step{则 $x_2\in\left[\dfrac12,b\right]$，于是 $\left[\dfrac13+\dfrac1{9b-3},1\right]\subseteq\left[\dfrac12,b\right]$，应有 $\begin{cases}b\geqslant1\\\dfrac13+\dfrac1{9b-3}\geqslant\dfrac12\end{cases}$，}
\step{得 $\begin{cases}b\geqslant1\\b\leqslant1\end{cases}$，即 $b=1$，所以 $b$ 的值是 1.}
\step{(3) 由函数 $f(x)=\dfrac4{x+2}$ 在区间 $[0,2]$ 的值域为 $[1,2]$，}
\step{因为 $f(x)=\dfrac4{x+2}$ 是 $g(x)=x^2-2ax+a^2-1$ 在 $[0,2]$ 上的“2 阶伴随函数”，}
\step{则对任意的 $x_1\in[0,2]$，总存在唯一的 $x_2\in[0,2]$ 时，使得 $f(x_1)g(x_2)=2$，}% 原卷此处印“时”字（按上下文为衍字），照录
\step{所以 $g(x_2)=\dfrac2{f(x_1)}\in[1,2]$，}
\step{即 $g(x)=x^2-2ax+a^2-1$ 在 $[0,2]$ 区间上的值域必定包含区间 $[1,2]$，}
\step{且 $g(x)$ 值域在 $[1,2]$ 对应的自变量是唯一的，}
\step{又因为函数 $g(x)=x^2-2ax+a^2-1$ 开口向上，对称轴为 $x=a$，}
\step{1) 当 $a\leqslant0$ 时，$g(x)$ 在 $[0,2]$ 区间上严格增，则必有 $\begin{cases}g(x)_{\min}=g(0)\leqslant1\\g(x)_{\max}=g(2)\geqslant2\end{cases}$，}
\step{解得 $-\sqrt2\leqslant a\leqslant0$；}
\step{2) 当 $a\geqslant2$ 时，$g(x)$ 在 $[0,2]$ 区间上严格减，则必有 $\begin{cases}g(x)_{\min}=g(2)\leqslant1\\g(x)_{\max}=g(0)\geqslant2\end{cases}$，}
\step{解得 $2\leqslant a\leqslant2+\sqrt2$；}
\step{3) 当 $0<a\leqslant1$ 时，$g(x)$ 在 $[0,a]$ 上严格减，在 $[a,2]$ 上严格增，}
\step{则必有 $\begin{cases}g(0)<1\\g(x)_{\max}=g(2)\geqslant2\end{cases}$，解得 $0<a\leqslant2-\sqrt3$，}
\step{4) 当 $1<a<2$ 时，$g(x)$ 在 $[0,a]$ 上严格减，在 $[a,2]$ 上严格增，}
\step{则必有 $\begin{cases}g(2)<1\\g(x)_{\max}=g(0)\geqslant2\end{cases}$，解得 $\sqrt3\leqslant a<2$.}
\step{综上所述，$a$ 的范围 $[-\sqrt2,2-\sqrt3]\cup[\sqrt3,2+\sqrt2]$.}% 原卷“范围”后缺“是／为”字，照录
\solend
\fi

\end{enumerate}

\end{document}
"
%
% 原始材料：微信公众号“魔都数学加油站”文章，作者破晓，2026-11-09 17:52 发布；
% 连载“27学年高一122”，原文标题“27学年高一122——2026-2027年上海市华师大二附中高一上周末作业6”；
% 原文链接：https://mp.weixin.qq.com/s/pC5066JFKNkQQwjMqTK8Ow。
% 正文为 14 张 PNG 页（第 14 页 2560×3620，其余 13 页 4762×6735），
% 题目下印有【解析】，为讲评版。
% 题目、解析均据原图转录；卷面红色水印及公众号页脚未录入。全卷无插图（已逐页检索，无“如图”）。
%
% 卷首标题：原卷印作“2026-2027 年华师大二附中高一上周末作业 6”（公众号标题中的“上海市”
% 卷面标题行没有，本稿照卷面），年份区间按全稿惯例排 en dash；原卷未印副标题，本稿按题目内容
% 标为“函数”。副标题理由（以全部 14 页实际内容为准）：逐页复核 19 题，考的都是函数性质——
% 奇偶性与充要条件（Q3、Q13、Q15、Q16）、单调性与单调区间（Q4、Q7、Q14、Q17）、定义域／
% 值域／最值与复合函数（Q5、Q6、Q8、Q10、Q12、Q19、Q20）、分段函数与取整函数（Q7、Q8、
% Q11、Q15、Q19）、抽象函数与新定义（Q12、Q20、Q21）；不等式仅作求参数范围／解集的工具
% （Q9、Q14、Q16、Q18），集合与常用逻辑用语仅零星出现（Q3 的充要条件、Q11③ 的集合、
% Q20 的“非空集合 I”），均非独立章节。故主体为单一模块“函数”，与同校 高一102 周末作业 5
% 同（该函数模块卷亦按 高一83／105 成例只用模块名、不叠“集合与常用逻辑用语、不等式”）。
%
% 与原卷的排版差异：
%   ① 原文从第 3 题起，填空题编号为 3—12、选择题为 13—16、解答题为 17—21，本稿保留原题号；
%   ② 原卷每题下直接印【解析】，本稿按全稿惯例置于答案版；四道选择题括号内均未印答案、
%      答案在解析末句“故选 …”，括号内留空（全卷无只印【答案】的题、无 \ans{}）；
%   ③ 原卷未印副标题，本稿按“知识模块”口径补出；
%   ④ 原卷填空横线按全稿惯例排为 \blankline[4.4em]；选择题括号排 \mbox{（\quad）}、
%      选项段前按惯例加 \keepwithprev；解答题不留作答区（原卷为讲评版）；
%   ⑤ 原卷数集排斜体 $R$、$N_+$（下标加号），按全稿惯例改 \R、\N_+
%      （$N_+$ 下标加号与 $\N^*$ 是两种记号，照录下标、只改数集字母为正体，同 高一05 第 20 题成例）；
%      绝对值排 \lvert…\rvert、集合描述竖线排 \mid；
%      $f(x)_{min}$／$f(x)_{max}$／$y_{max}$／$(y_{max})_{min}$ 一类斜体下标按 高一96 成例
%      改正体 \min／\max；不等式一律排 \geqslant／\leqslant；
%   ⑥ 第 20 题解析标签原卷印作【解得】（全卷唯一，其余 18 题均印【解析】），本稿按全稿惯例
%      由 \solbegin 统一印【解析】，正文内容照录（行内加 % 注、并入文件头存疑清单）；
%   ⑦ 题干与解析的标点（半角／全角逗号、冒号、括号、句点、子题括号）均照原卷，
%      题目正文引号用弯引号“ ”。
%
% 子集符号：第 21 题 (2) 的 $\left[\frac13+\frac1{9b-3},1\right]\subseteq\left[\frac12,b\right]$
%   印有下划线，照录 \subseteq。全卷未出现判别式 \triangle／\Delta。
%
% 原卷存疑／笔误（均照录加注）：
%   第 11 题解析对② 的“$y=y+\Delta y$”（按上下文应为 $y=[y]+\Delta y$，原卷漏方括号）；
%   第 19 题题干“某种产品 $x$ 万件）可获利”右括号无配对（疑漏印左括号“（”）；
%   第 20 题解析标签【解得】（见排版差异⑥）；
%   第 21 题 (3) 解析“总存在唯一的 $x_2\in[0,2]$ 时，使得……”的“时”字（按上下文衍字）；
%   第 21 题 (3) 末句“综上所述，$a$ 的范围 $[-\sqrt2,2-\sqrt3]\cup[\sqrt3,2+\sqrt2]$”
%      “范围”后缺“是／为”字，照录。
%
% 行数核对：逐题数 \step 条数，合计 162 条（Q3：5；Q4：3；Q5：1；Q6：3；Q7：2；
% Q8：5；Q9：7；Q10：10；Q11：8；Q12：4；Q13：12；Q14：6；Q15：14；Q16：14；
% Q17：6；Q18：8；Q19：7；Q20：19；Q21：28），与原卷印刷行逐题相等
% （cases 块作一个印刷行计）。全卷无 \ans{}、无只印【答案】的题。
%
\documentclass[a4paper,11pt]{article}
\usepackage{common}

\begin{document}

\examtitlest[3]{2026--2027 年华师大二附中高一上周末作业 6}{函数}

\bigsec{一、填空题}

\begin{enumerate}
\setcounter{enumi}{2}

\item 已知 $p:a=0$，$q$：函数 $f(x)=x^3+x+a$ 是奇函数，则 $p$ 是 $q$ 的\blankline[4.4em]条件.

\ifanswer
\solbegin
\step{若 $a=0$，则 $f(x)=x^3+x$，}
\step{易得 $f(x)$ 的定义域为 $\R$，且 $f(-x)=-(x^3+x)=-f(x)$，则 $f(x)$ 为奇函数，}
\step{反之，若 $f(x)=x^3+x+a$ 是奇函数，}
\step{则有 $f(-x)+f(x)=-x^3-x+a+x^3+x+a=2a=0$，必有 $a=0$，}
\step{故 $p$ 是 $q$ 的充要条件.}
\solend
\fi

\item 函数 $y=\sqrt{-x^2+2x}$ 的单调递增区间是\blankline[4.4em].

\ifanswer
\solbegin
\step{由函数定义域得 $-x^2+2x\geqslant0$，解得 $0\leqslant x\leqslant2$，}
\step{二次函数 $y=-x^2+2x$ 开口向下，对称轴为 $x=1$，}
\step{结合复合函数的单调性得函数 $y=\sqrt{-x^2+2x}$ 的单调递增区间是 $[0,1]$.}
\solend
\fi

\item 设 $f(x)$ 定义域为 $[1,4]$，且 $f(x)=2x+1$，则 $f(x^2)$ 的最大值为\blankline[4.4em].

\ifanswer
\solbegin
\step{由 $x^2\in[1,4]$，得 $f(x^2)=2x^2+1\leqslant9$，当 $x=\pm2$ 时取等号，最大值为 9.}
\solend
\fi

\item 函数 $f(x)=\sqrt{x-2}-x$，$x\in(2,6]$ 的值域为\blankline[4.4em].

\ifanswer
\solbegin
\step{令 $t=\sqrt{x-2}$，则 $x=2+t^2$，由 $2<x\leqslant6$ 得 $0<t\leqslant2$，}
\step{原函数可化为 $g(t)=-t^2+t-2$，$g(0)=-2$，$g(2)=-4$，$g\left(\dfrac12\right)=-\dfrac74$，}
\step{由二次函数性质得 $g(t)\in\left[-4,-\dfrac74\right]$.}
\solend
\fi

\item 若函数 $f(x)=\lvert x-a+1\rvert$ 在区间 $[1,+\infty)$ 上是严格增函数，则实数 $a$ 的取值范围是\blankline[4.4em].

\ifanswer
\solbegin
\step{$f(x)=\begin{cases}x-a+1,x\geqslant a-1\\a-1-x,x<a-1\end{cases}$，显然 $f(x)$ 在 $[a-1,+\infty)$ 上严格增，}
\step{要使函数在区间 $[1,+\infty)$ 上是严格增函数，则 $a-1\leqslant1$，即 $a\leqslant2$.}
\solend
\fi

\item 若函数 $f(x)=\begin{cases}x^2-5,x\leqslant2\\2ax+1,x>2\end{cases}$ 的值域是 $\R$，则 $a$ 的取值范围是\blankline[4.4em].

\ifanswer
\solbegin
\step{当 $x\leqslant2$ 时，$f(x)=x^2-5$，其开口向上，对称轴为 $x=0$，值域为 $[-5,+\infty)$，}
\step{由函数 $f(x)=\begin{cases}x^2-5,x\leqslant2\\2ax+1,x>2\end{cases}$ 的值域是 $\R$，}
\step{则当 $x>2$ 时，$f(x)=2ax+1$ 的值域应包含 $(-\infty,-5)$，}
\step{所以 $f(x)$ 为严格减函数，所以 $\begin{cases}2\times2a+1\geqslant-5\\2a<0\end{cases}$，即 $\begin{cases}a\geqslant-\dfrac32\\a<0\end{cases}$，}
\step{解得 $-\dfrac32\leqslant a<0$，故 $a$ 的取值范围是 $\left[-\dfrac32,0\right)$.}
\solend
\fi

\item 已知函数 $f(x)=ax^2+\dfrac bx$ 的图象经过点 $(-1,0)$，且满足 $f(1)=2$. 若 $\forall x\in(0,+\infty)$，都有 $f(x)+f\left(\dfrac1x\right)-2m\geqslant0$ 恒成立，则实数 $m$ 的取值范围为\blankline[4.4em].

\ifanswer
\solbegin
\step{因为 $f(x)=ax^2+\dfrac bx$ 的图象经过点 $(-1,0)$，且满足 $f(1)=2$，}
\step{所以 $\begin{cases}a-b=0\\a+b=2\end{cases}$，所以 $a=1$，$b=1$，即 $f(x)=x^2+\dfrac1x$，}
\step{$\forall x\in(0,+\infty)$，$f(x)+f\left(\dfrac1x\right)-2m\geqslant0\Leftrightarrow2m\leqslant x^2+\dfrac1x+\dfrac1{x^2}+x$ 恒成立，}
\step{令函数 $g(x)=x^2+\dfrac1x+\dfrac1{x^2}+x=\left(x+\dfrac1x\right)^2+x+\dfrac1x-2$，}
\step{而当 $x>0$ 时，$x+\dfrac1x\geqslant2$，当且仅当 $x=1$ 时取等号，}
\step{因此 $g(x)\geqslant2^2+2-2=4$，当且仅当 $x=1$ 时取等号，则 $2m\leqslant4$，解得 $m\leqslant2$，}
\step{所以实数 $m\in(-\infty,2]$.}
\solend
\fi

\item 若定义在区间 $[a,b](a<b)$ 上的函数 $f(x)=k-\sqrt{x+1}$ 值域也为 $[a,b]$，则实数 $k$ 的取值范围是\blankline[4.4em].

\ifanswer
\solbegin
\step{函数 $f(x)=k-\sqrt{x+1}$ 在定义域 $[-1,+\infty)$ 严格减，}
\step{当 $f(x)$ 的定义域为 $[a,b]$ 时，$f(x)$ 的值域也为 $[a,b]$，$a<b$，}
\step{故 $f(a)=k-\sqrt{a+1}=b$，$f(b)=k-\sqrt{b+1}=a$，}
\step{两式相减得 $\sqrt{a+1}-\sqrt{b+1}=a-b=(a+1)-(b+1)$}
\step{$=(\sqrt{a+1}-\sqrt{b+1})(\sqrt{a+1}+\sqrt{b+1})$，}
\step{所以 $\sqrt{a+1}-\sqrt{b+1}<0$，即 $\sqrt{a+1}+\sqrt{b+1}=1$，}
\step{则 $k=\sqrt{b+1}+a=a+1-\sqrt{a+1}$，}
\step{令 $\lambda=\sqrt{a+1}\geqslant0$，得 $k=\lambda^2-\lambda=\left(\lambda-\dfrac12\right)^2-\dfrac14$，}
\step{又 $\sqrt{a+1}<\sqrt{b+1}$，因为 $\sqrt{a+1}+\sqrt{b+1}=1$，所以 $\sqrt{a+1}\in\left[0,\dfrac12\right)$，}
\step{所以 $0\leqslant\lambda<\dfrac12$，故实数 $k$ 的取值范围为 $\left(-\dfrac14,0\right]$.}
\solend
\fi

\item 函数 $f(x)=[x]$ 的函数值表示不超过 $x$ 的最大整数，例如：$[-3.5]=-4$，$[2.1]=2$，称 $f(x)=[x]$ 为取整函数，则下列关于取整函数的说法：①对任意的 $x\in\R$，$[-x]=-[x]$ ②对任意的 $x$、$y\in\R$，$[x]+[y]\leqslant[x+y]$ ③集合 $\{y\mid y=[nx]-n[x],x\in\R,n\in\N_+\}$ 共有 $n$ 个元素 ④$x>0$ 时，关于 $x$ 的方程 $[2x]=\left[\dfrac1x\right]$ 有无数个解 其中正确的有\blankline[4.4em].

\ifanswer
\solbegin
\step{对①，因为 $[-2.1]=-3\neq[2.1]=2$，所以①错误；}
\step{对②，设 $x=[x]+\Delta x$，$\Delta x\in[0,1)$，$y=y+\Delta y$，$\Delta y\in[0,1)$，}% 原卷印作 $y=y+\Delta y$（按上下文应为 $y=[y]+\Delta y$），照录
\step{$[x+y]=[[x]+\Delta x+[y]+\Delta y]=[x]+[y]+[\Delta x+\Delta y]$，又 $\Delta x+\Delta y\in[0,\ 2)$，}
\step{故 $\Delta x+\Delta y=0$ 或 1，故 $[x]+[y]\leqslant[x+y]$，所以②正确；}
\step{对③，$[nx]-n[x]=[n[x]+n\cdot\Delta x]-n[x]=n[x]+[n\cdot\Delta x]-n[x]=[n\cdot\Delta x]$，}
\step{因为 $\Delta x\in[0,1)$，所以 $0\leqslant n\cdot\Delta x<n$，$n\in\N_+$，所以③正确；}
\step{对④，当 $\dfrac12<x<1$ 时，$[2x]=\left[\dfrac1x\right]=1$，故方程存在无数个解，所以④正确.}
\step{故正确的是②③④.}
\solend
\fi

\item 已知函数 $f(x)$ 定义域为 $[0,1]$，且满足 $f(x)+f(1-x)=1$，$f(x)=2f\left(\dfrac x5\right)$，对任意 $x_1\leqslant x_2$，恒有 $f(x_1)\leqslant f(x_2)$，则 $f\left(\dfrac1{2026}\right)$ 的值为\blankline[4.4em].

\ifanswer
\solbegin
\step{$f(0)=2f(0)\Rightarrow f(0)=0$，从而 $f(1)=1$，$f\left(\dfrac15\right)=\dfrac12$，$f\left(\dfrac45\right)=1-f\left(\dfrac15\right)=\dfrac12$，}
\step{由单调性，$x\in\left[\dfrac15,\dfrac45\right]$ 时，$f(x)=\dfrac12$.}
\step{因为 $\dfrac15<\dfrac{625}{2026}<\dfrac45$，反复使用 $f\left(\dfrac x5\right)=\dfrac12 f(x)$，}
\step{得 $f\left(\dfrac1{2026}\right)=\dfrac1{2^4}f\left(\dfrac{625}{2026}\right)=\dfrac1{32}$.}
\solend
\fi

\end{enumerate}

\bigsec{二、选择题}

\begin{enumerate}
\setcounter{enumi}{12}

\item 已知 $f(x)$ 为 $\R$ 上的奇函数，$g(x)$ 为 $\R$ 上的偶函数，且 $g(x)\neq0$，则下列说法正确的是\mbox{（\quad）}

\keepwithprev
A. $f(x)+g(x)$ 为 $\R$ 上的奇函数\qquad\qquad B. $f(x)-g(x)$ 为 $\R$ 上的奇函数

C. $\dfrac{f(x)}{g(x)}$ 为 $\R$ 上的偶函数\qquad\qquad D. $\lvert f(x)g(x)\rvert$ 为 $\R$ 上的偶函数

\ifanswer
\solbegin
\step{因为 $f(x)$ 为 $\R$ 上的奇函数，$g(x)$ 为 $\R$ 上的偶函数，且 $g(x)\neq0$，}
\step{所以 $f(-x)=-f(x)$，$g(-x)=g(x)$，}
\step{所以 $f(-x)+g(-x)=-f(x)+g(x)\neq-[f(x)+g(x)]$，}
\step{故 $f(x)+g(x)$ 为非奇非偶函数，$A$ 错误；}
\step{同理，$f(x)-g(x)$ 为非奇非偶函数，$B$ 错误；}
\step{设 $F(x)=\dfrac{f(x)}{g(x)}$，则 $F(-x)=\dfrac{f(-x)}{g(-x)}=\dfrac{-f(x)}{g(x)}=-F(x)$，}
\step{所以 $F(x)$ 为奇函数，$C$ 错误；}
\step{设函数 $H(x)=\lvert f(x)g(x)\rvert$，}
\step{因为 $f(x)$ 为 $\R$ 上的奇函数，$g(x)$ 为 $\R$ 上的偶函数，且 $g(x)\neq0$，}
\step{所以 $f(-x)=-f(x)$，$g(-x)=g(x)$，则由函数奇偶性的定义}
\step{得 $H(-x)=\lvert f(-x)g(-x)\rvert=\lvert -f(x)g(x)\rvert=\lvert f(x)g(x)\rvert=H(x)$，$D$ 正确.}
\step{故选 $D$.}
\solend
\fi

\item 若函数 $f(x)$ 是定义域为 $(-\infty,0)\cup(0,+\infty)$ 的奇函数，在区间 $(0,+\infty)$ 上严格增，且 $f(-1)=0$，则不等式 $\dfrac{f(x-1)}{x-2}<0$ 的解集为\mbox{（\quad）}

\keepwithprev
A. $(-\infty,-2)$\qquad\qquad B. $(-2,0)$\qquad\qquad C. $(0,1)$\qquad\qquad D. $(1,2)$

\ifanswer
\solbegin
\step{若函数 $f(x)$ 是定义域为 $(-\infty,0)\cup(0,+\infty)$ 的奇函数，在区间 $(0,+\infty)$ 上严格增，}
\step{由奇函数对称性得，$f(x)$ 在 $(-\infty,0)$ 上严格增，}
\step{因为 $f(-1)=0$，所以 $f(1)=0$，}
\step{则不等式 $\dfrac{f(x-1)}{x-2}<0$ 可转化为 $\begin{cases}x>2\\f(x-1)<0\end{cases}$ 或 $\begin{cases}x<2\\f(x-1)>0\end{cases}$，}
\step{即 $\begin{cases}x>2\\x-1<-1\text{或}0<x-1<1\end{cases}$ 或 $\begin{cases}x<2\\x-1>1\text{或}-1<x-1<0\end{cases}$，解得 $0<x<1$.}
\step{故选 $C$.}
\solend
\fi

\item 若函数 $f(x)$ 同时满足：(1) 对于定义域上的任意 $x$，恒有 $f(x)+f(-x)=0$；(2) 对于定义域上的任意 $x_1$、$x_2$，当 $x_1\neq x_2$ 时，恒有 $x_1f(x_1)-x_1f(x_2)>x_2f(x_1)-x_2f(x_2)$，则称函数 $f(x)$ 为“理想函数”. 给出下列四个函数：① $f(x)=x^2$；② $f(x)=x^3$；③ $f(x)=\dfrac{2x-1}{2x+1}$；④ $f(x)=\begin{cases}x^2+4x,x\geqslant0\\-x^2+4x,x<0\end{cases}$，其中被称为“理想函数”的有\mbox{（\quad）}

\keepwithprev
A. 1 个\qquad\qquad B. 2 个\qquad\qquad C. 3 个\qquad\qquad D. 4 个

\ifanswer
\solbegin
\step{函数 $f(x)$ 同时满足：(1) 对于定义域上的任意 $x$，恒有 $f(x)+f(-x)=0$；}
\step{(2) 对于定义域上的任意 $x_1$、$x_2$，当 $x_1\neq x_2$ 时，}
\step{恒有 $x_1f(x_1)-x_1f(x_2)>x_2f(x_1)-x_2f(x_2)$，}
\step{整理得 $(x_1-x_2)[f(x_1)-f(x_2)]>0$，则称函数 $f(x)$ 为“理想函数”.}
\step{即函数满足 (1) 奇函数；(2) 函数在定义域内为严格增函数；}
\step{① $f(x)=x^2$；不满足函数为奇函数在定义域内也不严格增；}
\step{② $f(x)=x^3$；满足函数既为奇函数，也在定义域内严格增，}
\step{故该函数为理想函数；}
\step{③ $f(x)=\dfrac{2x-1}{2x+1}=\dfrac{2x+1-2}{2x+1}=-\dfrac2{2x+1}+1$，该函数不为奇函数，}
\step{但是函数在定义域内严格增，故该函数不为理想函数；}
\step{④ $f(x)=\begin{cases}x^2+4x,x\geqslant0\\-x^2+4x,x<0\end{cases}$，函数的关系式满足 $f(-x)=-f(x)$，}
\step{且在 $x\geqslant0$ 时，函数严格增，由于奇函数在对称区间上单调性相同，}
\step{故该函数为理想函数.}
\step{故选 $B$.}
\solend
\fi

\item 已知 $f(x)$ 是定义在 $\R$ 上的奇函数，当 $x_1$、$x_2\in(0,+\infty)$ 且 $x_1\neq x_2$ 时，都有\newline $\dfrac{x_2f(x_1)-x_1f(x_2)}{x_1x_2(x_1-x_2)}>0$ 成立，$f(2028)=2028$，则不等式 $f(x)-x>0$ 的解集为\mbox{（\quad）}

\keepwithprev
A. $(-\infty,-2028)\cup(2028,+\infty)$\qquad\qquad B. $\left(-\dfrac1{2028},\dfrac1{2028}\right)$

C. $(-2028,2028)$\qquad\qquad D. $(-2028,0)\cup(2028,+\infty)$

\ifanswer
\solbegin
\step{因为 $f(x)$ 是定义在 $\R$ 上的奇函数，所以 $f(-x)=-f(x)$，}
\step{令 $g(x)=\dfrac{f(x)}x$，则 $g(-x)=\dfrac{f(-x)}{-x}=\dfrac{-f(x)}{-x}=\dfrac{f(x)}x=g(x)$，}
\step{所以 $g(x)$ 为偶函数，}
\step{当 $x_1$、$x_2\in(0,+\infty)$ 且 $x_1\neq x_2$ 时，都有 $\dfrac{x_2f(x_1)-x_1f(x_2)}{x_1x_2(x_1-x_2)}>0$ 成立，}
\step{即 $\dfrac{\dfrac{f(x_1)}{x_1}-\dfrac{f(x_2)}{x_2}}{x_1-x_2}>0$，即 $g(x)$ 在 $(0,+\infty)$ 上严格增，}
\step{由偶函数对称性得 $g(x)$ 在 $(-\infty,0)$ 上严格减，}
\step{因为 $f(2028)=2028$，所以 $g(2028)=1$，}
\step{则不等式 $f(x)-x>0$ 可化为 $f(x)>x$，}
\step{当 $x=0$ 时，$f(0)=0$，显然不成立，}
\step{当 $x>0$ 时，不等式可化为 $\dfrac{f(x)}x>1$，即 $g(x)>1=g(2028)$，}
\step{所以 $\lvert x\rvert>2028$，解得 $x>2028$，}
\step{当 $x<0$ 时，不等式可化为 $\dfrac{f(x)}x<1$，即 $g(x)<1=g(2028)$，}
\step{所以 $\lvert x\rvert<2028$，解得 $-2028<x<0$，}
\step{故不等式的解集为 $(-2028,0)\cup(2028,+\infty)$. 故选 $D$.}
\solend
\fi

\end{enumerate}

\bigsec{三、解答题}

\begin{enumerate}
\setcounter{enumi}{16}

\item 已知函数 $f(x)=\dfrac{2x-3}{x+1}$.

(1) 试用函数单调性的定义证明：函数 $f(x)$ 在 $(-1,+\infty)$ 上是严格增函数；

(2) 求函数 $f(x)$ 在区间 $[1,4]$ 上的最大值和最小值.

\ifanswer
\solbegin
\step{(1) 设 $-1<x_1<x_2$，则 $f(x_1)-f(x_2)=\dfrac{2x_1-3}{x_1+1}-\dfrac{2x_2-3}{x_2+1}=\dfrac{5(x_1-x_2)}{(x_1+1)(x_2+1)}$}
\step{因为 $-1<x_1<x_2$，所以 $x_1-x_2<0$，$x_1+1>0$，$x_2+1>0$，}
\step{所以 $f(x_1)-f(x_2)<0$，所以 $f(x)$ 在 $(-1,+\infty)$ 上严格增.}
\step{(2) 由 (1) 得 $f(x)$ 在区间 $[1,4]$ 上严格增，}
\step{所以 $f(x)_{\min}=f(1)=-\dfrac12$，$f(x)_{\max}=f(4)=1$，}
\step{所以函数 $f(x)$ 在区间 $[1,4]$ 上的最大值为 1，最小值为 $-\dfrac12$.}
\solend
\fi

\item 设函数 $f(x)=x^2-mx+m$.

(1) 当 $m=2$ 时，求不等式 $f(x)<x$ 的解集；

(2) 若 $\exists x\in[1,2]$，使不等式 $f(x)<x+m-2$ 成立，求实数 $m$ 的取值范围.

\ifanswer
\solbegin
\step{(1) $m=2$ 时，函数 $f(x)=x^2-2x+2$，}
\step{不等式 $f(x)<x$ 可化为 $x^2-3x+2<0$，解得 $1<x<2$，}
\step{所以不等式的解集为 $\{x\mid1<x<2\}$；}
\step{(2) 若 $\exists x\in[1,2]$，使不等式 $f(x)<x+m-2$ 成立，即 $x^2-x+2<mx$，}
\step{所以 $m>x+\dfrac2x-1$，设 $g(x)=x+\dfrac2x-1$，$x\in[1,2]$，}
\step{则 $g(x)=x+\dfrac2x-1\geqslant2\sqrt{x\cdot\dfrac2x}-1=2\sqrt2-1$，}
\step{当且仅当 $x=\dfrac2x$，即 $x=\sqrt2$ 时取等号，}
\step{所以实数 $m$ 的取值范围是 $\{m\mid m>2\sqrt2-1\}$.}
\solend
\fi

\item 某企业原来生产某种产品 $x$ 万件）可获利 $w(x)$（万元），且满足\newline% 原卷“万件）”右括号无配对、疑漏印左括号“（”，照录
$w(x)=\begin{cases}20(x^2+17),0<x\leqslant2\\500-\dfrac{80}{x-1},2<x\leqslant5\end{cases}$.
现该企业计划引进新的生产设备和新的产品方案优化产品，优化后的产品的其他成本总投入为 $(20x+10)$ 万元. 由市场调研分析得知，当前产品供不应求. 记该企业优化后的产品的利润为 $f(x)$（单位：万元）.

(1) 求函数 $f(x)$ 的解析式；

(2) 当优化后的产品产量为多少万件时，该企业的利润 $f(x)$ 最大？最大利润是多少？请说明理由.

\ifanswer
\solbegin
\step{(1) 由题意得 $f(x)=\begin{cases}20x^2-20x+330,0<x\leqslant2\\-20x-\dfrac{80}{x-1}+490,2<x\leqslant5\end{cases}$；}
\step{(2) 当 $0<x\leqslant2$ 时，$f(x)=20(x^2-x)+330=20\left(x-\dfrac12\right)^2+325$，}
\step{故当 $x=2$ 时，$f(x)$ 取最大值，$f(2)=370$；}
\step{当 $2<x\leqslant5$ 时，$f(x)=-\left[20(x-1)+\dfrac{80}{x-1}\right]+470\leqslant-80+470=390$，}
\step{当且仅当 $20(x-1)=\dfrac{80}{x-1}$，即 $x=3$ 时，$f(3)=390$ 为最大值；}
\step{因为 $370<390$，}
\step{所以优化后的产品产量为 3 万件时，企业获最大利润 390 万元.}
\solend
\fi

\item 俄国数学家切比雪夫（1821—1894）是研究直线逼近函数的理论先驱. 对定义在非空集合 $I$ 上的函数 $f(x)$，以及函数 $g(x)=kx+b(k,b\in\R)$，切比雪夫将函数 $y=\lvert f(x)-g(x)\rvert,x\in I$ 的最大值称为 $f(x),g(x)$ 的“偏差”.

(1) 函数 $f(x)=x^2,x\in[0,1],g(x)=-x-2$，求 $f(x),g(x)$ 的“偏差”；

(2) 函数 $f(x)=\dfrac1x+1,x\in[1,2],g(x)=kx+1(k>0)$，若 $f(x),g(x)$ 的“偏差”为 2，求 $k$ 的值；

(3) 函数 $f(x)=x^2-x,x\in[0,3],g(x)=2x+b$，当 $f(x),g(x)$ 的“偏差”取最小值时，求 $b$ 的值，并求出“偏差”的最小值.

\ifanswer
% 注：原卷本题解析标签印作【解得】（全卷唯一，其余 18 题均印【解析】）；
%     本稿按全稿惯例由 \solbegin 统一印【解析】，正文内容照录。见文件头排版差异⑥。
\solbegin
\step{(1) $y=\lvert f(x)-g(x)\rvert=\lvert x^2+x+2\rvert=\left\lvert\left(x+\dfrac12\right)^2+\dfrac74\right\rvert,x\in[0,1]$，}
\step{因为 $x\in[0,1]$，所以 $\left(x+\dfrac12\right)^2+\dfrac74\in[2,4]$，}
\step{则 $y=\left\lvert\left(x+\dfrac12\right)^2+\dfrac74\right\rvert\in[2,4]$，$y_{\max}=4$，}
\step{所以函数 $f(x)$ 与 $g(x)$ 的“偏差”为 4.}
\step{(2) 令 $t=f(x)-g(x)=\dfrac1x+1-(kx+1)=\dfrac1x-kx,x\in[1,2]$，}
\step{因为 $k>0$，所以 $t=\dfrac1x-kx$ 是单调减函数，}
\step{所以 $t\in\left[\dfrac12-2k,1-k\right]$，由单调，$y=\lvert t\rvert,t\in\left[\dfrac12-2k,1-k\right]$，且 $y_{\max}=2$.}
\step{当 $\left\lvert\dfrac12-2k\right\rvert\leqslant\lvert1-k\rvert$，即 $0<k\leqslant\dfrac12$ 时，$\lvert1-k\rvert=2$，解得 $k=3$ 或 $k=-1$，}
\step{均不符合；}
\step{当 $\left\lvert\dfrac12-2k\right\rvert>\lvert1-k\rvert$，即 $k>\dfrac12$ 时，$\left\lvert\dfrac12-2k\right\rvert=2$，$\dfrac12-2k=2$ 或 $\dfrac12-2k=-2$，}
\step{解得 $k=\dfrac54$ 或 $k=-\dfrac34$（舍），所以 $k=\dfrac54$.}
\step{(3) $y=\lvert f(x)-g(x)\rvert=\lvert x^2-x-(2x+b)\rvert=\lvert x^2-3x-b\rvert$}
\step{$=\left\lvert\left(x-\dfrac32\right)^2-\dfrac94-b\right\rvert,x\in[0,3]$，}
\step{因为 $x\in[0,3]$，所以 $\left(x-\dfrac32\right)^2-\dfrac94-b\in\left[-\dfrac94-b,-b\right]$，}
\step{由 $y=\left\lvert\left(x-\dfrac32\right)^2-\dfrac94-b\right\rvert$，则 $y_{\max}=\max\left\{\lvert b\rvert,\left\lvert b+\dfrac94\right\rvert\right\}$，}
\step{令 $\lvert b\rvert\geqslant\left\lvert b+\dfrac94\right\rvert$，即 $b^2\geqslant\left(b+\dfrac94\right)^2$，解得 $b\leqslant-\dfrac98$，}
\step{$y_{\max}=\max\left\{\lvert b\rvert,\left\lvert b+\dfrac94\right\rvert\right\}=\begin{cases}-b,b\leqslant-\dfrac98\\b+\dfrac94,b>-\dfrac98\end{cases}$，}
\step{故当且仅当 $b=-\dfrac98$ 时，$(y_{\max})_{\min}=\dfrac98$，}
\step{故当 $b$ 的值为 $-\dfrac98$ 时，函数 $f(x)$ 与 $g(x)$ 的“偏差”取最小值 $\dfrac98$.}
\solend
\fi

\item 若函数 $f(x)$ 与 $g(x)$ 满足：对任意的 $x_1\in D$，总存在唯一的 $x_2\in D$，使 $f(x_1)g(x_2)=m$ 成立，则称 $f(x)$ 是 $g(x)$ 在区间 $D$ 上的“$m$ 阶伴随函数”；对任意的 $x_1\in D$，总存在唯一的 $x_2\in D$，使 $f(x_1)f(x_2)=m$ 成立，则称 $f(x)$ 是区间 $D$ 上的“$m$ 阶自伴函数”.

(1) 判断 $f(x)=x^2+1$ 是否为区间 $[0,3]$ 上的“2 阶自伴函数”？并说明理由：

(2) 若函数 $f(x)=3x-1$ 为区间 $\left[\dfrac12,b\right]$ 上的“1 阶自伴函数”，求 $b$ 的值；

(3) 若 $f(x)=\dfrac4{x+2}$ 是 $g(x)=x^2-2ax+a^2-1$ 在区间 $[0,2]$ 上的“2 阶伴随函数”，求实数 $a$ 的取值范围.

\ifanswer
\solbegin
\step{(1) 假定函数 $f(x)=x^2+1$ 是区间 $[0,3]$ 上的“2 阶自伴函数”，}
\step{取 $x_1=2,f(x_1)=f(2)=5$，由 $f(x_1)f(x_2)=2$，得 $5(x_2^2+1)=2$，}
\step{显然此方程无实数解，}
\step{所以函数 $f(x)=x^2+1$ 不是区间 $[0,3]$ 上的“2 阶自伴函数”.}
\step{(2) 函数 $f(x)=3x-1$ 为区间 $\left[\dfrac12,b\right]$ 上的“1 阶自伴函数”，}
\step{则对任意 $x_1\in\left[\dfrac12,b\right]$，总存在唯一的 $x_2\in\left[\dfrac12,b\right]$，使得 $f(x_1)f(x_2)=1$，}
\step{即 $3x_2-1=\dfrac1{3x_1-1}$，整理得 $x_2=\dfrac13+\dfrac1{3(3x_1-1)}$，}
\step{显然函数 $x_2=\dfrac13+\dfrac1{3(3x_1-1)}$ 在 $\left[\dfrac12,b\right]$ 上严格减，}
\step{且当 $x_1=\dfrac12$ 时，$x_2=1$，当 $x_1=b$ 时，$x_2=\dfrac13+\dfrac1{9b-3}$，}
\step{因此对 $\left[\dfrac12,b\right]$ 内的每一个 $x_1$，在 $\left[\dfrac13+\dfrac1{9b-3},1\right]$ 内有唯一 $x_2$ 值与之对应，}
\step{则 $x_2\in\left[\dfrac12,b\right]$，于是 $\left[\dfrac13+\dfrac1{9b-3},1\right]\subseteq\left[\dfrac12,b\right]$，应有 $\begin{cases}b\geqslant1\\\dfrac13+\dfrac1{9b-3}\geqslant\dfrac12\end{cases}$，}
\step{得 $\begin{cases}b\geqslant1\\b\leqslant1\end{cases}$，即 $b=1$，所以 $b$ 的值是 1.}
\step{(3) 由函数 $f(x)=\dfrac4{x+2}$ 在区间 $[0,2]$ 的值域为 $[1,2]$，}
\step{因为 $f(x)=\dfrac4{x+2}$ 是 $g(x)=x^2-2ax+a^2-1$ 在 $[0,2]$ 上的“2 阶伴随函数”，}
\step{则对任意的 $x_1\in[0,2]$，总存在唯一的 $x_2\in[0,2]$ 时，使得 $f(x_1)g(x_2)=2$，}% 原卷此处印“时”字（按上下文为衍字），照录
\step{所以 $g(x_2)=\dfrac2{f(x_1)}\in[1,2]$，}
\step{即 $g(x)=x^2-2ax+a^2-1$ 在 $[0,2]$ 区间上的值域必定包含区间 $[1,2]$，}
\step{且 $g(x)$ 值域在 $[1,2]$ 对应的自变量是唯一的，}
\step{又因为函数 $g(x)=x^2-2ax+a^2-1$ 开口向上，对称轴为 $x=a$，}
\step{1) 当 $a\leqslant0$ 时，$g(x)$ 在 $[0,2]$ 区间上严格增，则必有 $\begin{cases}g(x)_{\min}=g(0)\leqslant1\\g(x)_{\max}=g(2)\geqslant2\end{cases}$，}
\step{解得 $-\sqrt2\leqslant a\leqslant0$；}
\step{2) 当 $a\geqslant2$ 时，$g(x)$ 在 $[0,2]$ 区间上严格减，则必有 $\begin{cases}g(x)_{\min}=g(2)\leqslant1\\g(x)_{\max}=g(0)\geqslant2\end{cases}$，}
\step{解得 $2\leqslant a\leqslant2+\sqrt2$；}
\step{3) 当 $0<a\leqslant1$ 时，$g(x)$ 在 $[0,a]$ 上严格减，在 $[a,2]$ 上严格增，}
\step{则必有 $\begin{cases}g(0)<1\\g(x)_{\max}=g(2)\geqslant2\end{cases}$，解得 $0<a\leqslant2-\sqrt3$，}
\step{4) 当 $1<a<2$ 时，$g(x)$ 在 $[0,a]$ 上严格减，在 $[a,2]$ 上严格增，}
\step{则必有 $\begin{cases}g(2)<1\\g(x)_{\max}=g(0)\geqslant2\end{cases}$，解得 $\sqrt3\leqslant a<2$.}
\step{综上所述，$a$ 的范围 $[-\sqrt2,2-\sqrt3]\cup[\sqrt3,2+\sqrt2]$.}% 原卷“范围”后缺“是／为”字，照录
\solend
\fi

\end{enumerate}

\end{document}
"
% 紧凑版：xelatex "\def\iscompact{1}% 高一122_华师大二附中_周末作业6.tex
%   —— 高一上数学“2026-2027 年华师大二附中高一上周末作业 6”（试卷版/答案版）
% 试卷版：xelatex 高一122_华师大二附中_周末作业6.tex
% 答案版：xelatex "\def\isanswer{1}% 高一122_华师大二附中_周末作业6.tex
%   —— 高一上数学“2026-2027 年华师大二附中高一上周末作业 6”（试卷版/答案版）
% 试卷版：xelatex 高一122_华师大二附中_周末作业6.tex
% 答案版：xelatex "\def\isanswer{1}\input{高一122_华师大二附中_周末作业6.tex}"
% 紧凑版：xelatex "\def\iscompact{1}\input{高一122_华师大二附中_周末作业6.tex}"
%
% 原始材料：微信公众号“魔都数学加油站”文章，作者破晓，2026-11-09 17:52 发布；
% 连载“27学年高一122”，原文标题“27学年高一122——2026-2027年上海市华师大二附中高一上周末作业6”；
% 原文链接：https://mp.weixin.qq.com/s/pC5066JFKNkQQwjMqTK8Ow。
% 正文为 14 张 PNG 页（第 14 页 2560×3620，其余 13 页 4762×6735），
% 题目下印有【解析】，为讲评版。
% 题目、解析均据原图转录；卷面红色水印及公众号页脚未录入。全卷无插图（已逐页检索，无“如图”）。
%
% 卷首标题：原卷印作“2026-2027 年华师大二附中高一上周末作业 6”（公众号标题中的“上海市”
% 卷面标题行没有，本稿照卷面），年份区间按全稿惯例排 en dash；原卷未印副标题，本稿按题目内容
% 标为“函数”。副标题理由（以全部 14 页实际内容为准）：逐页复核 19 题，考的都是函数性质——
% 奇偶性与充要条件（Q3、Q13、Q15、Q16）、单调性与单调区间（Q4、Q7、Q14、Q17）、定义域／
% 值域／最值与复合函数（Q5、Q6、Q8、Q10、Q12、Q19、Q20）、分段函数与取整函数（Q7、Q8、
% Q11、Q15、Q19）、抽象函数与新定义（Q12、Q20、Q21）；不等式仅作求参数范围／解集的工具
% （Q9、Q14、Q16、Q18），集合与常用逻辑用语仅零星出现（Q3 的充要条件、Q11③ 的集合、
% Q20 的“非空集合 I”），均非独立章节。故主体为单一模块“函数”，与同校 高一102 周末作业 5
% 同（该函数模块卷亦按 高一83／105 成例只用模块名、不叠“集合与常用逻辑用语、不等式”）。
%
% 与原卷的排版差异：
%   ① 原文从第 3 题起，填空题编号为 3—12、选择题为 13—16、解答题为 17—21，本稿保留原题号；
%   ② 原卷每题下直接印【解析】，本稿按全稿惯例置于答案版；四道选择题括号内均未印答案、
%      答案在解析末句“故选 …”，括号内留空（全卷无只印【答案】的题、无 \ans{}）；
%   ③ 原卷未印副标题，本稿按“知识模块”口径补出；
%   ④ 原卷填空横线按全稿惯例排为 \blankline[4.4em]；选择题括号排 \mbox{（\quad）}、
%      选项段前按惯例加 \keepwithprev；解答题不留作答区（原卷为讲评版）；
%   ⑤ 原卷数集排斜体 $R$、$N_+$（下标加号），按全稿惯例改 \R、\N_+
%      （$N_+$ 下标加号与 $\N^*$ 是两种记号，照录下标、只改数集字母为正体，同 高一05 第 20 题成例）；
%      绝对值排 \lvert…\rvert、集合描述竖线排 \mid；
%      $f(x)_{min}$／$f(x)_{max}$／$y_{max}$／$(y_{max})_{min}$ 一类斜体下标按 高一96 成例
%      改正体 \min／\max；不等式一律排 \geqslant／\leqslant；
%   ⑥ 第 20 题解析标签原卷印作【解得】（全卷唯一，其余 18 题均印【解析】），本稿按全稿惯例
%      由 \solbegin 统一印【解析】，正文内容照录（行内加 % 注、并入文件头存疑清单）；
%   ⑦ 题干与解析的标点（半角／全角逗号、冒号、括号、句点、子题括号）均照原卷，
%      题目正文引号用弯引号“ ”。
%
% 子集符号：第 21 题 (2) 的 $\left[\frac13+\frac1{9b-3},1\right]\subseteq\left[\frac12,b\right]$
%   印有下划线，照录 \subseteq。全卷未出现判别式 \triangle／\Delta。
%
% 原卷存疑／笔误（均照录加注）：
%   第 11 题解析对② 的“$y=y+\Delta y$”（按上下文应为 $y=[y]+\Delta y$，原卷漏方括号）；
%   第 19 题题干“某种产品 $x$ 万件）可获利”右括号无配对（疑漏印左括号“（”）；
%   第 20 题解析标签【解得】（见排版差异⑥）；
%   第 21 题 (3) 解析“总存在唯一的 $x_2\in[0,2]$ 时，使得……”的“时”字（按上下文衍字）；
%   第 21 题 (3) 末句“综上所述，$a$ 的范围 $[-\sqrt2,2-\sqrt3]\cup[\sqrt3,2+\sqrt2]$”
%      “范围”后缺“是／为”字，照录。
%
% 行数核对：逐题数 \step 条数，合计 162 条（Q3：5；Q4：3；Q5：1；Q6：3；Q7：2；
% Q8：5；Q9：7；Q10：10；Q11：8；Q12：4；Q13：12；Q14：6；Q15：14；Q16：14；
% Q17：6；Q18：8；Q19：7；Q20：19；Q21：28），与原卷印刷行逐题相等
% （cases 块作一个印刷行计）。全卷无 \ans{}、无只印【答案】的题。
%
\documentclass[a4paper,11pt]{article}
\usepackage{common}

\begin{document}

\examtitlest[3]{2026--2027 年华师大二附中高一上周末作业 6}{函数}

\bigsec{一、填空题}

\begin{enumerate}
\setcounter{enumi}{2}

\item 已知 $p:a=0$，$q$：函数 $f(x)=x^3+x+a$ 是奇函数，则 $p$ 是 $q$ 的\blankline[4.4em]条件.

\ifanswer
\solbegin
\step{若 $a=0$，则 $f(x)=x^3+x$，}
\step{易得 $f(x)$ 的定义域为 $\R$，且 $f(-x)=-(x^3+x)=-f(x)$，则 $f(x)$ 为奇函数，}
\step{反之，若 $f(x)=x^3+x+a$ 是奇函数，}
\step{则有 $f(-x)+f(x)=-x^3-x+a+x^3+x+a=2a=0$，必有 $a=0$，}
\step{故 $p$ 是 $q$ 的充要条件.}
\solend
\fi

\item 函数 $y=\sqrt{-x^2+2x}$ 的单调递增区间是\blankline[4.4em].

\ifanswer
\solbegin
\step{由函数定义域得 $-x^2+2x\geqslant0$，解得 $0\leqslant x\leqslant2$，}
\step{二次函数 $y=-x^2+2x$ 开口向下，对称轴为 $x=1$，}
\step{结合复合函数的单调性得函数 $y=\sqrt{-x^2+2x}$ 的单调递增区间是 $[0,1]$.}
\solend
\fi

\item 设 $f(x)$ 定义域为 $[1,4]$，且 $f(x)=2x+1$，则 $f(x^2)$ 的最大值为\blankline[4.4em].

\ifanswer
\solbegin
\step{由 $x^2\in[1,4]$，得 $f(x^2)=2x^2+1\leqslant9$，当 $x=\pm2$ 时取等号，最大值为 9.}
\solend
\fi

\item 函数 $f(x)=\sqrt{x-2}-x$，$x\in(2,6]$ 的值域为\blankline[4.4em].

\ifanswer
\solbegin
\step{令 $t=\sqrt{x-2}$，则 $x=2+t^2$，由 $2<x\leqslant6$ 得 $0<t\leqslant2$，}
\step{原函数可化为 $g(t)=-t^2+t-2$，$g(0)=-2$，$g(2)=-4$，$g\left(\dfrac12\right)=-\dfrac74$，}
\step{由二次函数性质得 $g(t)\in\left[-4,-\dfrac74\right]$.}
\solend
\fi

\item 若函数 $f(x)=\lvert x-a+1\rvert$ 在区间 $[1,+\infty)$ 上是严格增函数，则实数 $a$ 的取值范围是\blankline[4.4em].

\ifanswer
\solbegin
\step{$f(x)=\begin{cases}x-a+1,x\geqslant a-1\\a-1-x,x<a-1\end{cases}$，显然 $f(x)$ 在 $[a-1,+\infty)$ 上严格增，}
\step{要使函数在区间 $[1,+\infty)$ 上是严格增函数，则 $a-1\leqslant1$，即 $a\leqslant2$.}
\solend
\fi

\item 若函数 $f(x)=\begin{cases}x^2-5,x\leqslant2\\2ax+1,x>2\end{cases}$ 的值域是 $\R$，则 $a$ 的取值范围是\blankline[4.4em].

\ifanswer
\solbegin
\step{当 $x\leqslant2$ 时，$f(x)=x^2-5$，其开口向上，对称轴为 $x=0$，值域为 $[-5,+\infty)$，}
\step{由函数 $f(x)=\begin{cases}x^2-5,x\leqslant2\\2ax+1,x>2\end{cases}$ 的值域是 $\R$，}
\step{则当 $x>2$ 时，$f(x)=2ax+1$ 的值域应包含 $(-\infty,-5)$，}
\step{所以 $f(x)$ 为严格减函数，所以 $\begin{cases}2\times2a+1\geqslant-5\\2a<0\end{cases}$，即 $\begin{cases}a\geqslant-\dfrac32\\a<0\end{cases}$，}
\step{解得 $-\dfrac32\leqslant a<0$，故 $a$ 的取值范围是 $\left[-\dfrac32,0\right)$.}
\solend
\fi

\item 已知函数 $f(x)=ax^2+\dfrac bx$ 的图象经过点 $(-1,0)$，且满足 $f(1)=2$. 若 $\forall x\in(0,+\infty)$，都有 $f(x)+f\left(\dfrac1x\right)-2m\geqslant0$ 恒成立，则实数 $m$ 的取值范围为\blankline[4.4em].

\ifanswer
\solbegin
\step{因为 $f(x)=ax^2+\dfrac bx$ 的图象经过点 $(-1,0)$，且满足 $f(1)=2$，}
\step{所以 $\begin{cases}a-b=0\\a+b=2\end{cases}$，所以 $a=1$，$b=1$，即 $f(x)=x^2+\dfrac1x$，}
\step{$\forall x\in(0,+\infty)$，$f(x)+f\left(\dfrac1x\right)-2m\geqslant0\Leftrightarrow2m\leqslant x^2+\dfrac1x+\dfrac1{x^2}+x$ 恒成立，}
\step{令函数 $g(x)=x^2+\dfrac1x+\dfrac1{x^2}+x=\left(x+\dfrac1x\right)^2+x+\dfrac1x-2$，}
\step{而当 $x>0$ 时，$x+\dfrac1x\geqslant2$，当且仅当 $x=1$ 时取等号，}
\step{因此 $g(x)\geqslant2^2+2-2=4$，当且仅当 $x=1$ 时取等号，则 $2m\leqslant4$，解得 $m\leqslant2$，}
\step{所以实数 $m\in(-\infty,2]$.}
\solend
\fi

\item 若定义在区间 $[a,b](a<b)$ 上的函数 $f(x)=k-\sqrt{x+1}$ 值域也为 $[a,b]$，则实数 $k$ 的取值范围是\blankline[4.4em].

\ifanswer
\solbegin
\step{函数 $f(x)=k-\sqrt{x+1}$ 在定义域 $[-1,+\infty)$ 严格减，}
\step{当 $f(x)$ 的定义域为 $[a,b]$ 时，$f(x)$ 的值域也为 $[a,b]$，$a<b$，}
\step{故 $f(a)=k-\sqrt{a+1}=b$，$f(b)=k-\sqrt{b+1}=a$，}
\step{两式相减得 $\sqrt{a+1}-\sqrt{b+1}=a-b=(a+1)-(b+1)$}
\step{$=(\sqrt{a+1}-\sqrt{b+1})(\sqrt{a+1}+\sqrt{b+1})$，}
\step{所以 $\sqrt{a+1}-\sqrt{b+1}<0$，即 $\sqrt{a+1}+\sqrt{b+1}=1$，}
\step{则 $k=\sqrt{b+1}+a=a+1-\sqrt{a+1}$，}
\step{令 $\lambda=\sqrt{a+1}\geqslant0$，得 $k=\lambda^2-\lambda=\left(\lambda-\dfrac12\right)^2-\dfrac14$，}
\step{又 $\sqrt{a+1}<\sqrt{b+1}$，因为 $\sqrt{a+1}+\sqrt{b+1}=1$，所以 $\sqrt{a+1}\in\left[0,\dfrac12\right)$，}
\step{所以 $0\leqslant\lambda<\dfrac12$，故实数 $k$ 的取值范围为 $\left(-\dfrac14,0\right]$.}
\solend
\fi

\item 函数 $f(x)=[x]$ 的函数值表示不超过 $x$ 的最大整数，例如：$[-3.5]=-4$，$[2.1]=2$，称 $f(x)=[x]$ 为取整函数，则下列关于取整函数的说法：①对任意的 $x\in\R$，$[-x]=-[x]$ ②对任意的 $x$、$y\in\R$，$[x]+[y]\leqslant[x+y]$ ③集合 $\{y\mid y=[nx]-n[x],x\in\R,n\in\N_+\}$ 共有 $n$ 个元素 ④$x>0$ 时，关于 $x$ 的方程 $[2x]=\left[\dfrac1x\right]$ 有无数个解 其中正确的有\blankline[4.4em].

\ifanswer
\solbegin
\step{对①，因为 $[-2.1]=-3\neq[2.1]=2$，所以①错误；}
\step{对②，设 $x=[x]+\Delta x$，$\Delta x\in[0,1)$，$y=y+\Delta y$，$\Delta y\in[0,1)$，}% 原卷印作 $y=y+\Delta y$（按上下文应为 $y=[y]+\Delta y$），照录
\step{$[x+y]=[[x]+\Delta x+[y]+\Delta y]=[x]+[y]+[\Delta x+\Delta y]$，又 $\Delta x+\Delta y\in[0,\ 2)$，}
\step{故 $\Delta x+\Delta y=0$ 或 1，故 $[x]+[y]\leqslant[x+y]$，所以②正确；}
\step{对③，$[nx]-n[x]=[n[x]+n\cdot\Delta x]-n[x]=n[x]+[n\cdot\Delta x]-n[x]=[n\cdot\Delta x]$，}
\step{因为 $\Delta x\in[0,1)$，所以 $0\leqslant n\cdot\Delta x<n$，$n\in\N_+$，所以③正确；}
\step{对④，当 $\dfrac12<x<1$ 时，$[2x]=\left[\dfrac1x\right]=1$，故方程存在无数个解，所以④正确.}
\step{故正确的是②③④.}
\solend
\fi

\item 已知函数 $f(x)$ 定义域为 $[0,1]$，且满足 $f(x)+f(1-x)=1$，$f(x)=2f\left(\dfrac x5\right)$，对任意 $x_1\leqslant x_2$，恒有 $f(x_1)\leqslant f(x_2)$，则 $f\left(\dfrac1{2026}\right)$ 的值为\blankline[4.4em].

\ifanswer
\solbegin
\step{$f(0)=2f(0)\Rightarrow f(0)=0$，从而 $f(1)=1$，$f\left(\dfrac15\right)=\dfrac12$，$f\left(\dfrac45\right)=1-f\left(\dfrac15\right)=\dfrac12$，}
\step{由单调性，$x\in\left[\dfrac15,\dfrac45\right]$ 时，$f(x)=\dfrac12$.}
\step{因为 $\dfrac15<\dfrac{625}{2026}<\dfrac45$，反复使用 $f\left(\dfrac x5\right)=\dfrac12 f(x)$，}
\step{得 $f\left(\dfrac1{2026}\right)=\dfrac1{2^4}f\left(\dfrac{625}{2026}\right)=\dfrac1{32}$.}
\solend
\fi

\end{enumerate}

\bigsec{二、选择题}

\begin{enumerate}
\setcounter{enumi}{12}

\item 已知 $f(x)$ 为 $\R$ 上的奇函数，$g(x)$ 为 $\R$ 上的偶函数，且 $g(x)\neq0$，则下列说法正确的是\mbox{（\quad）}

\keepwithprev
A. $f(x)+g(x)$ 为 $\R$ 上的奇函数\qquad\qquad B. $f(x)-g(x)$ 为 $\R$ 上的奇函数

C. $\dfrac{f(x)}{g(x)}$ 为 $\R$ 上的偶函数\qquad\qquad D. $\lvert f(x)g(x)\rvert$ 为 $\R$ 上的偶函数

\ifanswer
\solbegin
\step{因为 $f(x)$ 为 $\R$ 上的奇函数，$g(x)$ 为 $\R$ 上的偶函数，且 $g(x)\neq0$，}
\step{所以 $f(-x)=-f(x)$，$g(-x)=g(x)$，}
\step{所以 $f(-x)+g(-x)=-f(x)+g(x)\neq-[f(x)+g(x)]$，}
\step{故 $f(x)+g(x)$ 为非奇非偶函数，$A$ 错误；}
\step{同理，$f(x)-g(x)$ 为非奇非偶函数，$B$ 错误；}
\step{设 $F(x)=\dfrac{f(x)}{g(x)}$，则 $F(-x)=\dfrac{f(-x)}{g(-x)}=\dfrac{-f(x)}{g(x)}=-F(x)$，}
\step{所以 $F(x)$ 为奇函数，$C$ 错误；}
\step{设函数 $H(x)=\lvert f(x)g(x)\rvert$，}
\step{因为 $f(x)$ 为 $\R$ 上的奇函数，$g(x)$ 为 $\R$ 上的偶函数，且 $g(x)\neq0$，}
\step{所以 $f(-x)=-f(x)$，$g(-x)=g(x)$，则由函数奇偶性的定义}
\step{得 $H(-x)=\lvert f(-x)g(-x)\rvert=\lvert -f(x)g(x)\rvert=\lvert f(x)g(x)\rvert=H(x)$，$D$ 正确.}
\step{故选 $D$.}
\solend
\fi

\item 若函数 $f(x)$ 是定义域为 $(-\infty,0)\cup(0,+\infty)$ 的奇函数，在区间 $(0,+\infty)$ 上严格增，且 $f(-1)=0$，则不等式 $\dfrac{f(x-1)}{x-2}<0$ 的解集为\mbox{（\quad）}

\keepwithprev
A. $(-\infty,-2)$\qquad\qquad B. $(-2,0)$\qquad\qquad C. $(0,1)$\qquad\qquad D. $(1,2)$

\ifanswer
\solbegin
\step{若函数 $f(x)$ 是定义域为 $(-\infty,0)\cup(0,+\infty)$ 的奇函数，在区间 $(0,+\infty)$ 上严格增，}
\step{由奇函数对称性得，$f(x)$ 在 $(-\infty,0)$ 上严格增，}
\step{因为 $f(-1)=0$，所以 $f(1)=0$，}
\step{则不等式 $\dfrac{f(x-1)}{x-2}<0$ 可转化为 $\begin{cases}x>2\\f(x-1)<0\end{cases}$ 或 $\begin{cases}x<2\\f(x-1)>0\end{cases}$，}
\step{即 $\begin{cases}x>2\\x-1<-1\text{或}0<x-1<1\end{cases}$ 或 $\begin{cases}x<2\\x-1>1\text{或}-1<x-1<0\end{cases}$，解得 $0<x<1$.}
\step{故选 $C$.}
\solend
\fi

\item 若函数 $f(x)$ 同时满足：(1) 对于定义域上的任意 $x$，恒有 $f(x)+f(-x)=0$；(2) 对于定义域上的任意 $x_1$、$x_2$，当 $x_1\neq x_2$ 时，恒有 $x_1f(x_1)-x_1f(x_2)>x_2f(x_1)-x_2f(x_2)$，则称函数 $f(x)$ 为“理想函数”. 给出下列四个函数：① $f(x)=x^2$；② $f(x)=x^3$；③ $f(x)=\dfrac{2x-1}{2x+1}$；④ $f(x)=\begin{cases}x^2+4x,x\geqslant0\\-x^2+4x,x<0\end{cases}$，其中被称为“理想函数”的有\mbox{（\quad）}

\keepwithprev
A. 1 个\qquad\qquad B. 2 个\qquad\qquad C. 3 个\qquad\qquad D. 4 个

\ifanswer
\solbegin
\step{函数 $f(x)$ 同时满足：(1) 对于定义域上的任意 $x$，恒有 $f(x)+f(-x)=0$；}
\step{(2) 对于定义域上的任意 $x_1$、$x_2$，当 $x_1\neq x_2$ 时，}
\step{恒有 $x_1f(x_1)-x_1f(x_2)>x_2f(x_1)-x_2f(x_2)$，}
\step{整理得 $(x_1-x_2)[f(x_1)-f(x_2)]>0$，则称函数 $f(x)$ 为“理想函数”.}
\step{即函数满足 (1) 奇函数；(2) 函数在定义域内为严格增函数；}
\step{① $f(x)=x^2$；不满足函数为奇函数在定义域内也不严格增；}
\step{② $f(x)=x^3$；满足函数既为奇函数，也在定义域内严格增，}
\step{故该函数为理想函数；}
\step{③ $f(x)=\dfrac{2x-1}{2x+1}=\dfrac{2x+1-2}{2x+1}=-\dfrac2{2x+1}+1$，该函数不为奇函数，}
\step{但是函数在定义域内严格增，故该函数不为理想函数；}
\step{④ $f(x)=\begin{cases}x^2+4x,x\geqslant0\\-x^2+4x,x<0\end{cases}$，函数的关系式满足 $f(-x)=-f(x)$，}
\step{且在 $x\geqslant0$ 时，函数严格增，由于奇函数在对称区间上单调性相同，}
\step{故该函数为理想函数.}
\step{故选 $B$.}
\solend
\fi

\item 已知 $f(x)$ 是定义在 $\R$ 上的奇函数，当 $x_1$、$x_2\in(0,+\infty)$ 且 $x_1\neq x_2$ 时，都有\newline $\dfrac{x_2f(x_1)-x_1f(x_2)}{x_1x_2(x_1-x_2)}>0$ 成立，$f(2028)=2028$，则不等式 $f(x)-x>0$ 的解集为\mbox{（\quad）}

\keepwithprev
A. $(-\infty,-2028)\cup(2028,+\infty)$\qquad\qquad B. $\left(-\dfrac1{2028},\dfrac1{2028}\right)$

C. $(-2028,2028)$\qquad\qquad D. $(-2028,0)\cup(2028,+\infty)$

\ifanswer
\solbegin
\step{因为 $f(x)$ 是定义在 $\R$ 上的奇函数，所以 $f(-x)=-f(x)$，}
\step{令 $g(x)=\dfrac{f(x)}x$，则 $g(-x)=\dfrac{f(-x)}{-x}=\dfrac{-f(x)}{-x}=\dfrac{f(x)}x=g(x)$，}
\step{所以 $g(x)$ 为偶函数，}
\step{当 $x_1$、$x_2\in(0,+\infty)$ 且 $x_1\neq x_2$ 时，都有 $\dfrac{x_2f(x_1)-x_1f(x_2)}{x_1x_2(x_1-x_2)}>0$ 成立，}
\step{即 $\dfrac{\dfrac{f(x_1)}{x_1}-\dfrac{f(x_2)}{x_2}}{x_1-x_2}>0$，即 $g(x)$ 在 $(0,+\infty)$ 上严格增，}
\step{由偶函数对称性得 $g(x)$ 在 $(-\infty,0)$ 上严格减，}
\step{因为 $f(2028)=2028$，所以 $g(2028)=1$，}
\step{则不等式 $f(x)-x>0$ 可化为 $f(x)>x$，}
\step{当 $x=0$ 时，$f(0)=0$，显然不成立，}
\step{当 $x>0$ 时，不等式可化为 $\dfrac{f(x)}x>1$，即 $g(x)>1=g(2028)$，}
\step{所以 $\lvert x\rvert>2028$，解得 $x>2028$，}
\step{当 $x<0$ 时，不等式可化为 $\dfrac{f(x)}x<1$，即 $g(x)<1=g(2028)$，}
\step{所以 $\lvert x\rvert<2028$，解得 $-2028<x<0$，}
\step{故不等式的解集为 $(-2028,0)\cup(2028,+\infty)$. 故选 $D$.}
\solend
\fi

\end{enumerate}

\bigsec{三、解答题}

\begin{enumerate}
\setcounter{enumi}{16}

\item 已知函数 $f(x)=\dfrac{2x-3}{x+1}$.

(1) 试用函数单调性的定义证明：函数 $f(x)$ 在 $(-1,+\infty)$ 上是严格增函数；

(2) 求函数 $f(x)$ 在区间 $[1,4]$ 上的最大值和最小值.

\ifanswer
\solbegin
\step{(1) 设 $-1<x_1<x_2$，则 $f(x_1)-f(x_2)=\dfrac{2x_1-3}{x_1+1}-\dfrac{2x_2-3}{x_2+1}=\dfrac{5(x_1-x_2)}{(x_1+1)(x_2+1)}$}
\step{因为 $-1<x_1<x_2$，所以 $x_1-x_2<0$，$x_1+1>0$，$x_2+1>0$，}
\step{所以 $f(x_1)-f(x_2)<0$，所以 $f(x)$ 在 $(-1,+\infty)$ 上严格增.}
\step{(2) 由 (1) 得 $f(x)$ 在区间 $[1,4]$ 上严格增，}
\step{所以 $f(x)_{\min}=f(1)=-\dfrac12$，$f(x)_{\max}=f(4)=1$，}
\step{所以函数 $f(x)$ 在区间 $[1,4]$ 上的最大值为 1，最小值为 $-\dfrac12$.}
\solend
\fi

\item 设函数 $f(x)=x^2-mx+m$.

(1) 当 $m=2$ 时，求不等式 $f(x)<x$ 的解集；

(2) 若 $\exists x\in[1,2]$，使不等式 $f(x)<x+m-2$ 成立，求实数 $m$ 的取值范围.

\ifanswer
\solbegin
\step{(1) $m=2$ 时，函数 $f(x)=x^2-2x+2$，}
\step{不等式 $f(x)<x$ 可化为 $x^2-3x+2<0$，解得 $1<x<2$，}
\step{所以不等式的解集为 $\{x\mid1<x<2\}$；}
\step{(2) 若 $\exists x\in[1,2]$，使不等式 $f(x)<x+m-2$ 成立，即 $x^2-x+2<mx$，}
\step{所以 $m>x+\dfrac2x-1$，设 $g(x)=x+\dfrac2x-1$，$x\in[1,2]$，}
\step{则 $g(x)=x+\dfrac2x-1\geqslant2\sqrt{x\cdot\dfrac2x}-1=2\sqrt2-1$，}
\step{当且仅当 $x=\dfrac2x$，即 $x=\sqrt2$ 时取等号，}
\step{所以实数 $m$ 的取值范围是 $\{m\mid m>2\sqrt2-1\}$.}
\solend
\fi

\item 某企业原来生产某种产品 $x$ 万件）可获利 $w(x)$（万元），且满足\newline% 原卷“万件）”右括号无配对、疑漏印左括号“（”，照录
$w(x)=\begin{cases}20(x^2+17),0<x\leqslant2\\500-\dfrac{80}{x-1},2<x\leqslant5\end{cases}$.
现该企业计划引进新的生产设备和新的产品方案优化产品，优化后的产品的其他成本总投入为 $(20x+10)$ 万元. 由市场调研分析得知，当前产品供不应求. 记该企业优化后的产品的利润为 $f(x)$（单位：万元）.

(1) 求函数 $f(x)$ 的解析式；

(2) 当优化后的产品产量为多少万件时，该企业的利润 $f(x)$ 最大？最大利润是多少？请说明理由.

\ifanswer
\solbegin
\step{(1) 由题意得 $f(x)=\begin{cases}20x^2-20x+330,0<x\leqslant2\\-20x-\dfrac{80}{x-1}+490,2<x\leqslant5\end{cases}$；}
\step{(2) 当 $0<x\leqslant2$ 时，$f(x)=20(x^2-x)+330=20\left(x-\dfrac12\right)^2+325$，}
\step{故当 $x=2$ 时，$f(x)$ 取最大值，$f(2)=370$；}
\step{当 $2<x\leqslant5$ 时，$f(x)=-\left[20(x-1)+\dfrac{80}{x-1}\right]+470\leqslant-80+470=390$，}
\step{当且仅当 $20(x-1)=\dfrac{80}{x-1}$，即 $x=3$ 时，$f(3)=390$ 为最大值；}
\step{因为 $370<390$，}
\step{所以优化后的产品产量为 3 万件时，企业获最大利润 390 万元.}
\solend
\fi

\item 俄国数学家切比雪夫（1821—1894）是研究直线逼近函数的理论先驱. 对定义在非空集合 $I$ 上的函数 $f(x)$，以及函数 $g(x)=kx+b(k,b\in\R)$，切比雪夫将函数 $y=\lvert f(x)-g(x)\rvert,x\in I$ 的最大值称为 $f(x),g(x)$ 的“偏差”.

(1) 函数 $f(x)=x^2,x\in[0,1],g(x)=-x-2$，求 $f(x),g(x)$ 的“偏差”；

(2) 函数 $f(x)=\dfrac1x+1,x\in[1,2],g(x)=kx+1(k>0)$，若 $f(x),g(x)$ 的“偏差”为 2，求 $k$ 的值；

(3) 函数 $f(x)=x^2-x,x\in[0,3],g(x)=2x+b$，当 $f(x),g(x)$ 的“偏差”取最小值时，求 $b$ 的值，并求出“偏差”的最小值.

\ifanswer
% 注：原卷本题解析标签印作【解得】（全卷唯一，其余 18 题均印【解析】）；
%     本稿按全稿惯例由 \solbegin 统一印【解析】，正文内容照录。见文件头排版差异⑥。
\solbegin
\step{(1) $y=\lvert f(x)-g(x)\rvert=\lvert x^2+x+2\rvert=\left\lvert\left(x+\dfrac12\right)^2+\dfrac74\right\rvert,x\in[0,1]$，}
\step{因为 $x\in[0,1]$，所以 $\left(x+\dfrac12\right)^2+\dfrac74\in[2,4]$，}
\step{则 $y=\left\lvert\left(x+\dfrac12\right)^2+\dfrac74\right\rvert\in[2,4]$，$y_{\max}=4$，}
\step{所以函数 $f(x)$ 与 $g(x)$ 的“偏差”为 4.}
\step{(2) 令 $t=f(x)-g(x)=\dfrac1x+1-(kx+1)=\dfrac1x-kx,x\in[1,2]$，}
\step{因为 $k>0$，所以 $t=\dfrac1x-kx$ 是单调减函数，}
\step{所以 $t\in\left[\dfrac12-2k,1-k\right]$，由单调，$y=\lvert t\rvert,t\in\left[\dfrac12-2k,1-k\right]$，且 $y_{\max}=2$.}
\step{当 $\left\lvert\dfrac12-2k\right\rvert\leqslant\lvert1-k\rvert$，即 $0<k\leqslant\dfrac12$ 时，$\lvert1-k\rvert=2$，解得 $k=3$ 或 $k=-1$，}
\step{均不符合；}
\step{当 $\left\lvert\dfrac12-2k\right\rvert>\lvert1-k\rvert$，即 $k>\dfrac12$ 时，$\left\lvert\dfrac12-2k\right\rvert=2$，$\dfrac12-2k=2$ 或 $\dfrac12-2k=-2$，}
\step{解得 $k=\dfrac54$ 或 $k=-\dfrac34$（舍），所以 $k=\dfrac54$.}
\step{(3) $y=\lvert f(x)-g(x)\rvert=\lvert x^2-x-(2x+b)\rvert=\lvert x^2-3x-b\rvert$}
\step{$=\left\lvert\left(x-\dfrac32\right)^2-\dfrac94-b\right\rvert,x\in[0,3]$，}
\step{因为 $x\in[0,3]$，所以 $\left(x-\dfrac32\right)^2-\dfrac94-b\in\left[-\dfrac94-b,-b\right]$，}
\step{由 $y=\left\lvert\left(x-\dfrac32\right)^2-\dfrac94-b\right\rvert$，则 $y_{\max}=\max\left\{\lvert b\rvert,\left\lvert b+\dfrac94\right\rvert\right\}$，}
\step{令 $\lvert b\rvert\geqslant\left\lvert b+\dfrac94\right\rvert$，即 $b^2\geqslant\left(b+\dfrac94\right)^2$，解得 $b\leqslant-\dfrac98$，}
\step{$y_{\max}=\max\left\{\lvert b\rvert,\left\lvert b+\dfrac94\right\rvert\right\}=\begin{cases}-b,b\leqslant-\dfrac98\\b+\dfrac94,b>-\dfrac98\end{cases}$，}
\step{故当且仅当 $b=-\dfrac98$ 时，$(y_{\max})_{\min}=\dfrac98$，}
\step{故当 $b$ 的值为 $-\dfrac98$ 时，函数 $f(x)$ 与 $g(x)$ 的“偏差”取最小值 $\dfrac98$.}
\solend
\fi

\item 若函数 $f(x)$ 与 $g(x)$ 满足：对任意的 $x_1\in D$，总存在唯一的 $x_2\in D$，使 $f(x_1)g(x_2)=m$ 成立，则称 $f(x)$ 是 $g(x)$ 在区间 $D$ 上的“$m$ 阶伴随函数”；对任意的 $x_1\in D$，总存在唯一的 $x_2\in D$，使 $f(x_1)f(x_2)=m$ 成立，则称 $f(x)$ 是区间 $D$ 上的“$m$ 阶自伴函数”.

(1) 判断 $f(x)=x^2+1$ 是否为区间 $[0,3]$ 上的“2 阶自伴函数”？并说明理由：

(2) 若函数 $f(x)=3x-1$ 为区间 $\left[\dfrac12,b\right]$ 上的“1 阶自伴函数”，求 $b$ 的值；

(3) 若 $f(x)=\dfrac4{x+2}$ 是 $g(x)=x^2-2ax+a^2-1$ 在区间 $[0,2]$ 上的“2 阶伴随函数”，求实数 $a$ 的取值范围.

\ifanswer
\solbegin
\step{(1) 假定函数 $f(x)=x^2+1$ 是区间 $[0,3]$ 上的“2 阶自伴函数”，}
\step{取 $x_1=2,f(x_1)=f(2)=5$，由 $f(x_1)f(x_2)=2$，得 $5(x_2^2+1)=2$，}
\step{显然此方程无实数解，}
\step{所以函数 $f(x)=x^2+1$ 不是区间 $[0,3]$ 上的“2 阶自伴函数”.}
\step{(2) 函数 $f(x)=3x-1$ 为区间 $\left[\dfrac12,b\right]$ 上的“1 阶自伴函数”，}
\step{则对任意 $x_1\in\left[\dfrac12,b\right]$，总存在唯一的 $x_2\in\left[\dfrac12,b\right]$，使得 $f(x_1)f(x_2)=1$，}
\step{即 $3x_2-1=\dfrac1{3x_1-1}$，整理得 $x_2=\dfrac13+\dfrac1{3(3x_1-1)}$，}
\step{显然函数 $x_2=\dfrac13+\dfrac1{3(3x_1-1)}$ 在 $\left[\dfrac12,b\right]$ 上严格减，}
\step{且当 $x_1=\dfrac12$ 时，$x_2=1$，当 $x_1=b$ 时，$x_2=\dfrac13+\dfrac1{9b-3}$，}
\step{因此对 $\left[\dfrac12,b\right]$ 内的每一个 $x_1$，在 $\left[\dfrac13+\dfrac1{9b-3},1\right]$ 内有唯一 $x_2$ 值与之对应，}
\step{则 $x_2\in\left[\dfrac12,b\right]$，于是 $\left[\dfrac13+\dfrac1{9b-3},1\right]\subseteq\left[\dfrac12,b\right]$，应有 $\begin{cases}b\geqslant1\\\dfrac13+\dfrac1{9b-3}\geqslant\dfrac12\end{cases}$，}
\step{得 $\begin{cases}b\geqslant1\\b\leqslant1\end{cases}$，即 $b=1$，所以 $b$ 的值是 1.}
\step{(3) 由函数 $f(x)=\dfrac4{x+2}$ 在区间 $[0,2]$ 的值域为 $[1,2]$，}
\step{因为 $f(x)=\dfrac4{x+2}$ 是 $g(x)=x^2-2ax+a^2-1$ 在 $[0,2]$ 上的“2 阶伴随函数”，}
\step{则对任意的 $x_1\in[0,2]$，总存在唯一的 $x_2\in[0,2]$ 时，使得 $f(x_1)g(x_2)=2$，}% 原卷此处印“时”字（按上下文为衍字），照录
\step{所以 $g(x_2)=\dfrac2{f(x_1)}\in[1,2]$，}
\step{即 $g(x)=x^2-2ax+a^2-1$ 在 $[0,2]$ 区间上的值域必定包含区间 $[1,2]$，}
\step{且 $g(x)$ 值域在 $[1,2]$ 对应的自变量是唯一的，}
\step{又因为函数 $g(x)=x^2-2ax+a^2-1$ 开口向上，对称轴为 $x=a$，}
\step{1) 当 $a\leqslant0$ 时，$g(x)$ 在 $[0,2]$ 区间上严格增，则必有 $\begin{cases}g(x)_{\min}=g(0)\leqslant1\\g(x)_{\max}=g(2)\geqslant2\end{cases}$，}
\step{解得 $-\sqrt2\leqslant a\leqslant0$；}
\step{2) 当 $a\geqslant2$ 时，$g(x)$ 在 $[0,2]$ 区间上严格减，则必有 $\begin{cases}g(x)_{\min}=g(2)\leqslant1\\g(x)_{\max}=g(0)\geqslant2\end{cases}$，}
\step{解得 $2\leqslant a\leqslant2+\sqrt2$；}
\step{3) 当 $0<a\leqslant1$ 时，$g(x)$ 在 $[0,a]$ 上严格减，在 $[a,2]$ 上严格增，}
\step{则必有 $\begin{cases}g(0)<1\\g(x)_{\max}=g(2)\geqslant2\end{cases}$，解得 $0<a\leqslant2-\sqrt3$，}
\step{4) 当 $1<a<2$ 时，$g(x)$ 在 $[0,a]$ 上严格减，在 $[a,2]$ 上严格增，}
\step{则必有 $\begin{cases}g(2)<1\\g(x)_{\max}=g(0)\geqslant2\end{cases}$，解得 $\sqrt3\leqslant a<2$.}
\step{综上所述，$a$ 的范围 $[-\sqrt2,2-\sqrt3]\cup[\sqrt3,2+\sqrt2]$.}% 原卷“范围”后缺“是／为”字，照录
\solend
\fi

\end{enumerate}

\end{document}
"
% 紧凑版：xelatex "\def\iscompact{1}% 高一122_华师大二附中_周末作业6.tex
%   —— 高一上数学“2026-2027 年华师大二附中高一上周末作业 6”（试卷版/答案版）
% 试卷版：xelatex 高一122_华师大二附中_周末作业6.tex
% 答案版：xelatex "\def\isanswer{1}\input{高一122_华师大二附中_周末作业6.tex}"
% 紧凑版：xelatex "\def\iscompact{1}\input{高一122_华师大二附中_周末作业6.tex}"
%
% 原始材料：微信公众号“魔都数学加油站”文章，作者破晓，2026-11-09 17:52 发布；
% 连载“27学年高一122”，原文标题“27学年高一122——2026-2027年上海市华师大二附中高一上周末作业6”；
% 原文链接：https://mp.weixin.qq.com/s/pC5066JFKNkQQwjMqTK8Ow。
% 正文为 14 张 PNG 页（第 14 页 2560×3620，其余 13 页 4762×6735），
% 题目下印有【解析】，为讲评版。
% 题目、解析均据原图转录；卷面红色水印及公众号页脚未录入。全卷无插图（已逐页检索，无“如图”）。
%
% 卷首标题：原卷印作“2026-2027 年华师大二附中高一上周末作业 6”（公众号标题中的“上海市”
% 卷面标题行没有，本稿照卷面），年份区间按全稿惯例排 en dash；原卷未印副标题，本稿按题目内容
% 标为“函数”。副标题理由（以全部 14 页实际内容为准）：逐页复核 19 题，考的都是函数性质——
% 奇偶性与充要条件（Q3、Q13、Q15、Q16）、单调性与单调区间（Q4、Q7、Q14、Q17）、定义域／
% 值域／最值与复合函数（Q5、Q6、Q8、Q10、Q12、Q19、Q20）、分段函数与取整函数（Q7、Q8、
% Q11、Q15、Q19）、抽象函数与新定义（Q12、Q20、Q21）；不等式仅作求参数范围／解集的工具
% （Q9、Q14、Q16、Q18），集合与常用逻辑用语仅零星出现（Q3 的充要条件、Q11③ 的集合、
% Q20 的“非空集合 I”），均非独立章节。故主体为单一模块“函数”，与同校 高一102 周末作业 5
% 同（该函数模块卷亦按 高一83／105 成例只用模块名、不叠“集合与常用逻辑用语、不等式”）。
%
% 与原卷的排版差异：
%   ① 原文从第 3 题起，填空题编号为 3—12、选择题为 13—16、解答题为 17—21，本稿保留原题号；
%   ② 原卷每题下直接印【解析】，本稿按全稿惯例置于答案版；四道选择题括号内均未印答案、
%      答案在解析末句“故选 …”，括号内留空（全卷无只印【答案】的题、无 \ans{}）；
%   ③ 原卷未印副标题，本稿按“知识模块”口径补出；
%   ④ 原卷填空横线按全稿惯例排为 \blankline[4.4em]；选择题括号排 \mbox{（\quad）}、
%      选项段前按惯例加 \keepwithprev；解答题不留作答区（原卷为讲评版）；
%   ⑤ 原卷数集排斜体 $R$、$N_+$（下标加号），按全稿惯例改 \R、\N_+
%      （$N_+$ 下标加号与 $\N^*$ 是两种记号，照录下标、只改数集字母为正体，同 高一05 第 20 题成例）；
%      绝对值排 \lvert…\rvert、集合描述竖线排 \mid；
%      $f(x)_{min}$／$f(x)_{max}$／$y_{max}$／$(y_{max})_{min}$ 一类斜体下标按 高一96 成例
%      改正体 \min／\max；不等式一律排 \geqslant／\leqslant；
%   ⑥ 第 20 题解析标签原卷印作【解得】（全卷唯一，其余 18 题均印【解析】），本稿按全稿惯例
%      由 \solbegin 统一印【解析】，正文内容照录（行内加 % 注、并入文件头存疑清单）；
%   ⑦ 题干与解析的标点（半角／全角逗号、冒号、括号、句点、子题括号）均照原卷，
%      题目正文引号用弯引号“ ”。
%
% 子集符号：第 21 题 (2) 的 $\left[\frac13+\frac1{9b-3},1\right]\subseteq\left[\frac12,b\right]$
%   印有下划线，照录 \subseteq。全卷未出现判别式 \triangle／\Delta。
%
% 原卷存疑／笔误（均照录加注）：
%   第 11 题解析对② 的“$y=y+\Delta y$”（按上下文应为 $y=[y]+\Delta y$，原卷漏方括号）；
%   第 19 题题干“某种产品 $x$ 万件）可获利”右括号无配对（疑漏印左括号“（”）；
%   第 20 题解析标签【解得】（见排版差异⑥）；
%   第 21 题 (3) 解析“总存在唯一的 $x_2\in[0,2]$ 时，使得……”的“时”字（按上下文衍字）；
%   第 21 题 (3) 末句“综上所述，$a$ 的范围 $[-\sqrt2,2-\sqrt3]\cup[\sqrt3,2+\sqrt2]$”
%      “范围”后缺“是／为”字，照录。
%
% 行数核对：逐题数 \step 条数，合计 162 条（Q3：5；Q4：3；Q5：1；Q6：3；Q7：2；
% Q8：5；Q9：7；Q10：10；Q11：8；Q12：4；Q13：12；Q14：6；Q15：14；Q16：14；
% Q17：6；Q18：8；Q19：7；Q20：19；Q21：28），与原卷印刷行逐题相等
% （cases 块作一个印刷行计）。全卷无 \ans{}、无只印【答案】的题。
%
\documentclass[a4paper,11pt]{article}
\usepackage{common}

\begin{document}

\examtitlest[3]{2026--2027 年华师大二附中高一上周末作业 6}{函数}

\bigsec{一、填空题}

\begin{enumerate}
\setcounter{enumi}{2}

\item 已知 $p:a=0$，$q$：函数 $f(x)=x^3+x+a$ 是奇函数，则 $p$ 是 $q$ 的\blankline[4.4em]条件.

\ifanswer
\solbegin
\step{若 $a=0$，则 $f(x)=x^3+x$，}
\step{易得 $f(x)$ 的定义域为 $\R$，且 $f(-x)=-(x^3+x)=-f(x)$，则 $f(x)$ 为奇函数，}
\step{反之，若 $f(x)=x^3+x+a$ 是奇函数，}
\step{则有 $f(-x)+f(x)=-x^3-x+a+x^3+x+a=2a=0$，必有 $a=0$，}
\step{故 $p$ 是 $q$ 的充要条件.}
\solend
\fi

\item 函数 $y=\sqrt{-x^2+2x}$ 的单调递增区间是\blankline[4.4em].

\ifanswer
\solbegin
\step{由函数定义域得 $-x^2+2x\geqslant0$，解得 $0\leqslant x\leqslant2$，}
\step{二次函数 $y=-x^2+2x$ 开口向下，对称轴为 $x=1$，}
\step{结合复合函数的单调性得函数 $y=\sqrt{-x^2+2x}$ 的单调递增区间是 $[0,1]$.}
\solend
\fi

\item 设 $f(x)$ 定义域为 $[1,4]$，且 $f(x)=2x+1$，则 $f(x^2)$ 的最大值为\blankline[4.4em].

\ifanswer
\solbegin
\step{由 $x^2\in[1,4]$，得 $f(x^2)=2x^2+1\leqslant9$，当 $x=\pm2$ 时取等号，最大值为 9.}
\solend
\fi

\item 函数 $f(x)=\sqrt{x-2}-x$，$x\in(2,6]$ 的值域为\blankline[4.4em].

\ifanswer
\solbegin
\step{令 $t=\sqrt{x-2}$，则 $x=2+t^2$，由 $2<x\leqslant6$ 得 $0<t\leqslant2$，}
\step{原函数可化为 $g(t)=-t^2+t-2$，$g(0)=-2$，$g(2)=-4$，$g\left(\dfrac12\right)=-\dfrac74$，}
\step{由二次函数性质得 $g(t)\in\left[-4,-\dfrac74\right]$.}
\solend
\fi

\item 若函数 $f(x)=\lvert x-a+1\rvert$ 在区间 $[1,+\infty)$ 上是严格增函数，则实数 $a$ 的取值范围是\blankline[4.4em].

\ifanswer
\solbegin
\step{$f(x)=\begin{cases}x-a+1,x\geqslant a-1\\a-1-x,x<a-1\end{cases}$，显然 $f(x)$ 在 $[a-1,+\infty)$ 上严格增，}
\step{要使函数在区间 $[1,+\infty)$ 上是严格增函数，则 $a-1\leqslant1$，即 $a\leqslant2$.}
\solend
\fi

\item 若函数 $f(x)=\begin{cases}x^2-5,x\leqslant2\\2ax+1,x>2\end{cases}$ 的值域是 $\R$，则 $a$ 的取值范围是\blankline[4.4em].

\ifanswer
\solbegin
\step{当 $x\leqslant2$ 时，$f(x)=x^2-5$，其开口向上，对称轴为 $x=0$，值域为 $[-5,+\infty)$，}
\step{由函数 $f(x)=\begin{cases}x^2-5,x\leqslant2\\2ax+1,x>2\end{cases}$ 的值域是 $\R$，}
\step{则当 $x>2$ 时，$f(x)=2ax+1$ 的值域应包含 $(-\infty,-5)$，}
\step{所以 $f(x)$ 为严格减函数，所以 $\begin{cases}2\times2a+1\geqslant-5\\2a<0\end{cases}$，即 $\begin{cases}a\geqslant-\dfrac32\\a<0\end{cases}$，}
\step{解得 $-\dfrac32\leqslant a<0$，故 $a$ 的取值范围是 $\left[-\dfrac32,0\right)$.}
\solend
\fi

\item 已知函数 $f(x)=ax^2+\dfrac bx$ 的图象经过点 $(-1,0)$，且满足 $f(1)=2$. 若 $\forall x\in(0,+\infty)$，都有 $f(x)+f\left(\dfrac1x\right)-2m\geqslant0$ 恒成立，则实数 $m$ 的取值范围为\blankline[4.4em].

\ifanswer
\solbegin
\step{因为 $f(x)=ax^2+\dfrac bx$ 的图象经过点 $(-1,0)$，且满足 $f(1)=2$，}
\step{所以 $\begin{cases}a-b=0\\a+b=2\end{cases}$，所以 $a=1$，$b=1$，即 $f(x)=x^2+\dfrac1x$，}
\step{$\forall x\in(0,+\infty)$，$f(x)+f\left(\dfrac1x\right)-2m\geqslant0\Leftrightarrow2m\leqslant x^2+\dfrac1x+\dfrac1{x^2}+x$ 恒成立，}
\step{令函数 $g(x)=x^2+\dfrac1x+\dfrac1{x^2}+x=\left(x+\dfrac1x\right)^2+x+\dfrac1x-2$，}
\step{而当 $x>0$ 时，$x+\dfrac1x\geqslant2$，当且仅当 $x=1$ 时取等号，}
\step{因此 $g(x)\geqslant2^2+2-2=4$，当且仅当 $x=1$ 时取等号，则 $2m\leqslant4$，解得 $m\leqslant2$，}
\step{所以实数 $m\in(-\infty,2]$.}
\solend
\fi

\item 若定义在区间 $[a,b](a<b)$ 上的函数 $f(x)=k-\sqrt{x+1}$ 值域也为 $[a,b]$，则实数 $k$ 的取值范围是\blankline[4.4em].

\ifanswer
\solbegin
\step{函数 $f(x)=k-\sqrt{x+1}$ 在定义域 $[-1,+\infty)$ 严格减，}
\step{当 $f(x)$ 的定义域为 $[a,b]$ 时，$f(x)$ 的值域也为 $[a,b]$，$a<b$，}
\step{故 $f(a)=k-\sqrt{a+1}=b$，$f(b)=k-\sqrt{b+1}=a$，}
\step{两式相减得 $\sqrt{a+1}-\sqrt{b+1}=a-b=(a+1)-(b+1)$}
\step{$=(\sqrt{a+1}-\sqrt{b+1})(\sqrt{a+1}+\sqrt{b+1})$，}
\step{所以 $\sqrt{a+1}-\sqrt{b+1}<0$，即 $\sqrt{a+1}+\sqrt{b+1}=1$，}
\step{则 $k=\sqrt{b+1}+a=a+1-\sqrt{a+1}$，}
\step{令 $\lambda=\sqrt{a+1}\geqslant0$，得 $k=\lambda^2-\lambda=\left(\lambda-\dfrac12\right)^2-\dfrac14$，}
\step{又 $\sqrt{a+1}<\sqrt{b+1}$，因为 $\sqrt{a+1}+\sqrt{b+1}=1$，所以 $\sqrt{a+1}\in\left[0,\dfrac12\right)$，}
\step{所以 $0\leqslant\lambda<\dfrac12$，故实数 $k$ 的取值范围为 $\left(-\dfrac14,0\right]$.}
\solend
\fi

\item 函数 $f(x)=[x]$ 的函数值表示不超过 $x$ 的最大整数，例如：$[-3.5]=-4$，$[2.1]=2$，称 $f(x)=[x]$ 为取整函数，则下列关于取整函数的说法：①对任意的 $x\in\R$，$[-x]=-[x]$ ②对任意的 $x$、$y\in\R$，$[x]+[y]\leqslant[x+y]$ ③集合 $\{y\mid y=[nx]-n[x],x\in\R,n\in\N_+\}$ 共有 $n$ 个元素 ④$x>0$ 时，关于 $x$ 的方程 $[2x]=\left[\dfrac1x\right]$ 有无数个解 其中正确的有\blankline[4.4em].

\ifanswer
\solbegin
\step{对①，因为 $[-2.1]=-3\neq[2.1]=2$，所以①错误；}
\step{对②，设 $x=[x]+\Delta x$，$\Delta x\in[0,1)$，$y=y+\Delta y$，$\Delta y\in[0,1)$，}% 原卷印作 $y=y+\Delta y$（按上下文应为 $y=[y]+\Delta y$），照录
\step{$[x+y]=[[x]+\Delta x+[y]+\Delta y]=[x]+[y]+[\Delta x+\Delta y]$，又 $\Delta x+\Delta y\in[0,\ 2)$，}
\step{故 $\Delta x+\Delta y=0$ 或 1，故 $[x]+[y]\leqslant[x+y]$，所以②正确；}
\step{对③，$[nx]-n[x]=[n[x]+n\cdot\Delta x]-n[x]=n[x]+[n\cdot\Delta x]-n[x]=[n\cdot\Delta x]$，}
\step{因为 $\Delta x\in[0,1)$，所以 $0\leqslant n\cdot\Delta x<n$，$n\in\N_+$，所以③正确；}
\step{对④，当 $\dfrac12<x<1$ 时，$[2x]=\left[\dfrac1x\right]=1$，故方程存在无数个解，所以④正确.}
\step{故正确的是②③④.}
\solend
\fi

\item 已知函数 $f(x)$ 定义域为 $[0,1]$，且满足 $f(x)+f(1-x)=1$，$f(x)=2f\left(\dfrac x5\right)$，对任意 $x_1\leqslant x_2$，恒有 $f(x_1)\leqslant f(x_2)$，则 $f\left(\dfrac1{2026}\right)$ 的值为\blankline[4.4em].

\ifanswer
\solbegin
\step{$f(0)=2f(0)\Rightarrow f(0)=0$，从而 $f(1)=1$，$f\left(\dfrac15\right)=\dfrac12$，$f\left(\dfrac45\right)=1-f\left(\dfrac15\right)=\dfrac12$，}
\step{由单调性，$x\in\left[\dfrac15,\dfrac45\right]$ 时，$f(x)=\dfrac12$.}
\step{因为 $\dfrac15<\dfrac{625}{2026}<\dfrac45$，反复使用 $f\left(\dfrac x5\right)=\dfrac12 f(x)$，}
\step{得 $f\left(\dfrac1{2026}\right)=\dfrac1{2^4}f\left(\dfrac{625}{2026}\right)=\dfrac1{32}$.}
\solend
\fi

\end{enumerate}

\bigsec{二、选择题}

\begin{enumerate}
\setcounter{enumi}{12}

\item 已知 $f(x)$ 为 $\R$ 上的奇函数，$g(x)$ 为 $\R$ 上的偶函数，且 $g(x)\neq0$，则下列说法正确的是\mbox{（\quad）}

\keepwithprev
A. $f(x)+g(x)$ 为 $\R$ 上的奇函数\qquad\qquad B. $f(x)-g(x)$ 为 $\R$ 上的奇函数

C. $\dfrac{f(x)}{g(x)}$ 为 $\R$ 上的偶函数\qquad\qquad D. $\lvert f(x)g(x)\rvert$ 为 $\R$ 上的偶函数

\ifanswer
\solbegin
\step{因为 $f(x)$ 为 $\R$ 上的奇函数，$g(x)$ 为 $\R$ 上的偶函数，且 $g(x)\neq0$，}
\step{所以 $f(-x)=-f(x)$，$g(-x)=g(x)$，}
\step{所以 $f(-x)+g(-x)=-f(x)+g(x)\neq-[f(x)+g(x)]$，}
\step{故 $f(x)+g(x)$ 为非奇非偶函数，$A$ 错误；}
\step{同理，$f(x)-g(x)$ 为非奇非偶函数，$B$ 错误；}
\step{设 $F(x)=\dfrac{f(x)}{g(x)}$，则 $F(-x)=\dfrac{f(-x)}{g(-x)}=\dfrac{-f(x)}{g(x)}=-F(x)$，}
\step{所以 $F(x)$ 为奇函数，$C$ 错误；}
\step{设函数 $H(x)=\lvert f(x)g(x)\rvert$，}
\step{因为 $f(x)$ 为 $\R$ 上的奇函数，$g(x)$ 为 $\R$ 上的偶函数，且 $g(x)\neq0$，}
\step{所以 $f(-x)=-f(x)$，$g(-x)=g(x)$，则由函数奇偶性的定义}
\step{得 $H(-x)=\lvert f(-x)g(-x)\rvert=\lvert -f(x)g(x)\rvert=\lvert f(x)g(x)\rvert=H(x)$，$D$ 正确.}
\step{故选 $D$.}
\solend
\fi

\item 若函数 $f(x)$ 是定义域为 $(-\infty,0)\cup(0,+\infty)$ 的奇函数，在区间 $(0,+\infty)$ 上严格增，且 $f(-1)=0$，则不等式 $\dfrac{f(x-1)}{x-2}<0$ 的解集为\mbox{（\quad）}

\keepwithprev
A. $(-\infty,-2)$\qquad\qquad B. $(-2,0)$\qquad\qquad C. $(0,1)$\qquad\qquad D. $(1,2)$

\ifanswer
\solbegin
\step{若函数 $f(x)$ 是定义域为 $(-\infty,0)\cup(0,+\infty)$ 的奇函数，在区间 $(0,+\infty)$ 上严格增，}
\step{由奇函数对称性得，$f(x)$ 在 $(-\infty,0)$ 上严格增，}
\step{因为 $f(-1)=0$，所以 $f(1)=0$，}
\step{则不等式 $\dfrac{f(x-1)}{x-2}<0$ 可转化为 $\begin{cases}x>2\\f(x-1)<0\end{cases}$ 或 $\begin{cases}x<2\\f(x-1)>0\end{cases}$，}
\step{即 $\begin{cases}x>2\\x-1<-1\text{或}0<x-1<1\end{cases}$ 或 $\begin{cases}x<2\\x-1>1\text{或}-1<x-1<0\end{cases}$，解得 $0<x<1$.}
\step{故选 $C$.}
\solend
\fi

\item 若函数 $f(x)$ 同时满足：(1) 对于定义域上的任意 $x$，恒有 $f(x)+f(-x)=0$；(2) 对于定义域上的任意 $x_1$、$x_2$，当 $x_1\neq x_2$ 时，恒有 $x_1f(x_1)-x_1f(x_2)>x_2f(x_1)-x_2f(x_2)$，则称函数 $f(x)$ 为“理想函数”. 给出下列四个函数：① $f(x)=x^2$；② $f(x)=x^3$；③ $f(x)=\dfrac{2x-1}{2x+1}$；④ $f(x)=\begin{cases}x^2+4x,x\geqslant0\\-x^2+4x,x<0\end{cases}$，其中被称为“理想函数”的有\mbox{（\quad）}

\keepwithprev
A. 1 个\qquad\qquad B. 2 个\qquad\qquad C. 3 个\qquad\qquad D. 4 个

\ifanswer
\solbegin
\step{函数 $f(x)$ 同时满足：(1) 对于定义域上的任意 $x$，恒有 $f(x)+f(-x)=0$；}
\step{(2) 对于定义域上的任意 $x_1$、$x_2$，当 $x_1\neq x_2$ 时，}
\step{恒有 $x_1f(x_1)-x_1f(x_2)>x_2f(x_1)-x_2f(x_2)$，}
\step{整理得 $(x_1-x_2)[f(x_1)-f(x_2)]>0$，则称函数 $f(x)$ 为“理想函数”.}
\step{即函数满足 (1) 奇函数；(2) 函数在定义域内为严格增函数；}
\step{① $f(x)=x^2$；不满足函数为奇函数在定义域内也不严格增；}
\step{② $f(x)=x^3$；满足函数既为奇函数，也在定义域内严格增，}
\step{故该函数为理想函数；}
\step{③ $f(x)=\dfrac{2x-1}{2x+1}=\dfrac{2x+1-2}{2x+1}=-\dfrac2{2x+1}+1$，该函数不为奇函数，}
\step{但是函数在定义域内严格增，故该函数不为理想函数；}
\step{④ $f(x)=\begin{cases}x^2+4x,x\geqslant0\\-x^2+4x,x<0\end{cases}$，函数的关系式满足 $f(-x)=-f(x)$，}
\step{且在 $x\geqslant0$ 时，函数严格增，由于奇函数在对称区间上单调性相同，}
\step{故该函数为理想函数.}
\step{故选 $B$.}
\solend
\fi

\item 已知 $f(x)$ 是定义在 $\R$ 上的奇函数，当 $x_1$、$x_2\in(0,+\infty)$ 且 $x_1\neq x_2$ 时，都有\newline $\dfrac{x_2f(x_1)-x_1f(x_2)}{x_1x_2(x_1-x_2)}>0$ 成立，$f(2028)=2028$，则不等式 $f(x)-x>0$ 的解集为\mbox{（\quad）}

\keepwithprev
A. $(-\infty,-2028)\cup(2028,+\infty)$\qquad\qquad B. $\left(-\dfrac1{2028},\dfrac1{2028}\right)$

C. $(-2028,2028)$\qquad\qquad D. $(-2028,0)\cup(2028,+\infty)$

\ifanswer
\solbegin
\step{因为 $f(x)$ 是定义在 $\R$ 上的奇函数，所以 $f(-x)=-f(x)$，}
\step{令 $g(x)=\dfrac{f(x)}x$，则 $g(-x)=\dfrac{f(-x)}{-x}=\dfrac{-f(x)}{-x}=\dfrac{f(x)}x=g(x)$，}
\step{所以 $g(x)$ 为偶函数，}
\step{当 $x_1$、$x_2\in(0,+\infty)$ 且 $x_1\neq x_2$ 时，都有 $\dfrac{x_2f(x_1)-x_1f(x_2)}{x_1x_2(x_1-x_2)}>0$ 成立，}
\step{即 $\dfrac{\dfrac{f(x_1)}{x_1}-\dfrac{f(x_2)}{x_2}}{x_1-x_2}>0$，即 $g(x)$ 在 $(0,+\infty)$ 上严格增，}
\step{由偶函数对称性得 $g(x)$ 在 $(-\infty,0)$ 上严格减，}
\step{因为 $f(2028)=2028$，所以 $g(2028)=1$，}
\step{则不等式 $f(x)-x>0$ 可化为 $f(x)>x$，}
\step{当 $x=0$ 时，$f(0)=0$，显然不成立，}
\step{当 $x>0$ 时，不等式可化为 $\dfrac{f(x)}x>1$，即 $g(x)>1=g(2028)$，}
\step{所以 $\lvert x\rvert>2028$，解得 $x>2028$，}
\step{当 $x<0$ 时，不等式可化为 $\dfrac{f(x)}x<1$，即 $g(x)<1=g(2028)$，}
\step{所以 $\lvert x\rvert<2028$，解得 $-2028<x<0$，}
\step{故不等式的解集为 $(-2028,0)\cup(2028,+\infty)$. 故选 $D$.}
\solend
\fi

\end{enumerate}

\bigsec{三、解答题}

\begin{enumerate}
\setcounter{enumi}{16}

\item 已知函数 $f(x)=\dfrac{2x-3}{x+1}$.

(1) 试用函数单调性的定义证明：函数 $f(x)$ 在 $(-1,+\infty)$ 上是严格增函数；

(2) 求函数 $f(x)$ 在区间 $[1,4]$ 上的最大值和最小值.

\ifanswer
\solbegin
\step{(1) 设 $-1<x_1<x_2$，则 $f(x_1)-f(x_2)=\dfrac{2x_1-3}{x_1+1}-\dfrac{2x_2-3}{x_2+1}=\dfrac{5(x_1-x_2)}{(x_1+1)(x_2+1)}$}
\step{因为 $-1<x_1<x_2$，所以 $x_1-x_2<0$，$x_1+1>0$，$x_2+1>0$，}
\step{所以 $f(x_1)-f(x_2)<0$，所以 $f(x)$ 在 $(-1,+\infty)$ 上严格增.}
\step{(2) 由 (1) 得 $f(x)$ 在区间 $[1,4]$ 上严格增，}
\step{所以 $f(x)_{\min}=f(1)=-\dfrac12$，$f(x)_{\max}=f(4)=1$，}
\step{所以函数 $f(x)$ 在区间 $[1,4]$ 上的最大值为 1，最小值为 $-\dfrac12$.}
\solend
\fi

\item 设函数 $f(x)=x^2-mx+m$.

(1) 当 $m=2$ 时，求不等式 $f(x)<x$ 的解集；

(2) 若 $\exists x\in[1,2]$，使不等式 $f(x)<x+m-2$ 成立，求实数 $m$ 的取值范围.

\ifanswer
\solbegin
\step{(1) $m=2$ 时，函数 $f(x)=x^2-2x+2$，}
\step{不等式 $f(x)<x$ 可化为 $x^2-3x+2<0$，解得 $1<x<2$，}
\step{所以不等式的解集为 $\{x\mid1<x<2\}$；}
\step{(2) 若 $\exists x\in[1,2]$，使不等式 $f(x)<x+m-2$ 成立，即 $x^2-x+2<mx$，}
\step{所以 $m>x+\dfrac2x-1$，设 $g(x)=x+\dfrac2x-1$，$x\in[1,2]$，}
\step{则 $g(x)=x+\dfrac2x-1\geqslant2\sqrt{x\cdot\dfrac2x}-1=2\sqrt2-1$，}
\step{当且仅当 $x=\dfrac2x$，即 $x=\sqrt2$ 时取等号，}
\step{所以实数 $m$ 的取值范围是 $\{m\mid m>2\sqrt2-1\}$.}
\solend
\fi

\item 某企业原来生产某种产品 $x$ 万件）可获利 $w(x)$（万元），且满足\newline% 原卷“万件）”右括号无配对、疑漏印左括号“（”，照录
$w(x)=\begin{cases}20(x^2+17),0<x\leqslant2\\500-\dfrac{80}{x-1},2<x\leqslant5\end{cases}$.
现该企业计划引进新的生产设备和新的产品方案优化产品，优化后的产品的其他成本总投入为 $(20x+10)$ 万元. 由市场调研分析得知，当前产品供不应求. 记该企业优化后的产品的利润为 $f(x)$（单位：万元）.

(1) 求函数 $f(x)$ 的解析式；

(2) 当优化后的产品产量为多少万件时，该企业的利润 $f(x)$ 最大？最大利润是多少？请说明理由.

\ifanswer
\solbegin
\step{(1) 由题意得 $f(x)=\begin{cases}20x^2-20x+330,0<x\leqslant2\\-20x-\dfrac{80}{x-1}+490,2<x\leqslant5\end{cases}$；}
\step{(2) 当 $0<x\leqslant2$ 时，$f(x)=20(x^2-x)+330=20\left(x-\dfrac12\right)^2+325$，}
\step{故当 $x=2$ 时，$f(x)$ 取最大值，$f(2)=370$；}
\step{当 $2<x\leqslant5$ 时，$f(x)=-\left[20(x-1)+\dfrac{80}{x-1}\right]+470\leqslant-80+470=390$，}
\step{当且仅当 $20(x-1)=\dfrac{80}{x-1}$，即 $x=3$ 时，$f(3)=390$ 为最大值；}
\step{因为 $370<390$，}
\step{所以优化后的产品产量为 3 万件时，企业获最大利润 390 万元.}
\solend
\fi

\item 俄国数学家切比雪夫（1821—1894）是研究直线逼近函数的理论先驱. 对定义在非空集合 $I$ 上的函数 $f(x)$，以及函数 $g(x)=kx+b(k,b\in\R)$，切比雪夫将函数 $y=\lvert f(x)-g(x)\rvert,x\in I$ 的最大值称为 $f(x),g(x)$ 的“偏差”.

(1) 函数 $f(x)=x^2,x\in[0,1],g(x)=-x-2$，求 $f(x),g(x)$ 的“偏差”；

(2) 函数 $f(x)=\dfrac1x+1,x\in[1,2],g(x)=kx+1(k>0)$，若 $f(x),g(x)$ 的“偏差”为 2，求 $k$ 的值；

(3) 函数 $f(x)=x^2-x,x\in[0,3],g(x)=2x+b$，当 $f(x),g(x)$ 的“偏差”取最小值时，求 $b$ 的值，并求出“偏差”的最小值.

\ifanswer
% 注：原卷本题解析标签印作【解得】（全卷唯一，其余 18 题均印【解析】）；
%     本稿按全稿惯例由 \solbegin 统一印【解析】，正文内容照录。见文件头排版差异⑥。
\solbegin
\step{(1) $y=\lvert f(x)-g(x)\rvert=\lvert x^2+x+2\rvert=\left\lvert\left(x+\dfrac12\right)^2+\dfrac74\right\rvert,x\in[0,1]$，}
\step{因为 $x\in[0,1]$，所以 $\left(x+\dfrac12\right)^2+\dfrac74\in[2,4]$，}
\step{则 $y=\left\lvert\left(x+\dfrac12\right)^2+\dfrac74\right\rvert\in[2,4]$，$y_{\max}=4$，}
\step{所以函数 $f(x)$ 与 $g(x)$ 的“偏差”为 4.}
\step{(2) 令 $t=f(x)-g(x)=\dfrac1x+1-(kx+1)=\dfrac1x-kx,x\in[1,2]$，}
\step{因为 $k>0$，所以 $t=\dfrac1x-kx$ 是单调减函数，}
\step{所以 $t\in\left[\dfrac12-2k,1-k\right]$，由单调，$y=\lvert t\rvert,t\in\left[\dfrac12-2k,1-k\right]$，且 $y_{\max}=2$.}
\step{当 $\left\lvert\dfrac12-2k\right\rvert\leqslant\lvert1-k\rvert$，即 $0<k\leqslant\dfrac12$ 时，$\lvert1-k\rvert=2$，解得 $k=3$ 或 $k=-1$，}
\step{均不符合；}
\step{当 $\left\lvert\dfrac12-2k\right\rvert>\lvert1-k\rvert$，即 $k>\dfrac12$ 时，$\left\lvert\dfrac12-2k\right\rvert=2$，$\dfrac12-2k=2$ 或 $\dfrac12-2k=-2$，}
\step{解得 $k=\dfrac54$ 或 $k=-\dfrac34$（舍），所以 $k=\dfrac54$.}
\step{(3) $y=\lvert f(x)-g(x)\rvert=\lvert x^2-x-(2x+b)\rvert=\lvert x^2-3x-b\rvert$}
\step{$=\left\lvert\left(x-\dfrac32\right)^2-\dfrac94-b\right\rvert,x\in[0,3]$，}
\step{因为 $x\in[0,3]$，所以 $\left(x-\dfrac32\right)^2-\dfrac94-b\in\left[-\dfrac94-b,-b\right]$，}
\step{由 $y=\left\lvert\left(x-\dfrac32\right)^2-\dfrac94-b\right\rvert$，则 $y_{\max}=\max\left\{\lvert b\rvert,\left\lvert b+\dfrac94\right\rvert\right\}$，}
\step{令 $\lvert b\rvert\geqslant\left\lvert b+\dfrac94\right\rvert$，即 $b^2\geqslant\left(b+\dfrac94\right)^2$，解得 $b\leqslant-\dfrac98$，}
\step{$y_{\max}=\max\left\{\lvert b\rvert,\left\lvert b+\dfrac94\right\rvert\right\}=\begin{cases}-b,b\leqslant-\dfrac98\\b+\dfrac94,b>-\dfrac98\end{cases}$，}
\step{故当且仅当 $b=-\dfrac98$ 时，$(y_{\max})_{\min}=\dfrac98$，}
\step{故当 $b$ 的值为 $-\dfrac98$ 时，函数 $f(x)$ 与 $g(x)$ 的“偏差”取最小值 $\dfrac98$.}
\solend
\fi

\item 若函数 $f(x)$ 与 $g(x)$ 满足：对任意的 $x_1\in D$，总存在唯一的 $x_2\in D$，使 $f(x_1)g(x_2)=m$ 成立，则称 $f(x)$ 是 $g(x)$ 在区间 $D$ 上的“$m$ 阶伴随函数”；对任意的 $x_1\in D$，总存在唯一的 $x_2\in D$，使 $f(x_1)f(x_2)=m$ 成立，则称 $f(x)$ 是区间 $D$ 上的“$m$ 阶自伴函数”.

(1) 判断 $f(x)=x^2+1$ 是否为区间 $[0,3]$ 上的“2 阶自伴函数”？并说明理由：

(2) 若函数 $f(x)=3x-1$ 为区间 $\left[\dfrac12,b\right]$ 上的“1 阶自伴函数”，求 $b$ 的值；

(3) 若 $f(x)=\dfrac4{x+2}$ 是 $g(x)=x^2-2ax+a^2-1$ 在区间 $[0,2]$ 上的“2 阶伴随函数”，求实数 $a$ 的取值范围.

\ifanswer
\solbegin
\step{(1) 假定函数 $f(x)=x^2+1$ 是区间 $[0,3]$ 上的“2 阶自伴函数”，}
\step{取 $x_1=2,f(x_1)=f(2)=5$，由 $f(x_1)f(x_2)=2$，得 $5(x_2^2+1)=2$，}
\step{显然此方程无实数解，}
\step{所以函数 $f(x)=x^2+1$ 不是区间 $[0,3]$ 上的“2 阶自伴函数”.}
\step{(2) 函数 $f(x)=3x-1$ 为区间 $\left[\dfrac12,b\right]$ 上的“1 阶自伴函数”，}
\step{则对任意 $x_1\in\left[\dfrac12,b\right]$，总存在唯一的 $x_2\in\left[\dfrac12,b\right]$，使得 $f(x_1)f(x_2)=1$，}
\step{即 $3x_2-1=\dfrac1{3x_1-1}$，整理得 $x_2=\dfrac13+\dfrac1{3(3x_1-1)}$，}
\step{显然函数 $x_2=\dfrac13+\dfrac1{3(3x_1-1)}$ 在 $\left[\dfrac12,b\right]$ 上严格减，}
\step{且当 $x_1=\dfrac12$ 时，$x_2=1$，当 $x_1=b$ 时，$x_2=\dfrac13+\dfrac1{9b-3}$，}
\step{因此对 $\left[\dfrac12,b\right]$ 内的每一个 $x_1$，在 $\left[\dfrac13+\dfrac1{9b-3},1\right]$ 内有唯一 $x_2$ 值与之对应，}
\step{则 $x_2\in\left[\dfrac12,b\right]$，于是 $\left[\dfrac13+\dfrac1{9b-3},1\right]\subseteq\left[\dfrac12,b\right]$，应有 $\begin{cases}b\geqslant1\\\dfrac13+\dfrac1{9b-3}\geqslant\dfrac12\end{cases}$，}
\step{得 $\begin{cases}b\geqslant1\\b\leqslant1\end{cases}$，即 $b=1$，所以 $b$ 的值是 1.}
\step{(3) 由函数 $f(x)=\dfrac4{x+2}$ 在区间 $[0,2]$ 的值域为 $[1,2]$，}
\step{因为 $f(x)=\dfrac4{x+2}$ 是 $g(x)=x^2-2ax+a^2-1$ 在 $[0,2]$ 上的“2 阶伴随函数”，}
\step{则对任意的 $x_1\in[0,2]$，总存在唯一的 $x_2\in[0,2]$ 时，使得 $f(x_1)g(x_2)=2$，}% 原卷此处印“时”字（按上下文为衍字），照录
\step{所以 $g(x_2)=\dfrac2{f(x_1)}\in[1,2]$，}
\step{即 $g(x)=x^2-2ax+a^2-1$ 在 $[0,2]$ 区间上的值域必定包含区间 $[1,2]$，}
\step{且 $g(x)$ 值域在 $[1,2]$ 对应的自变量是唯一的，}
\step{又因为函数 $g(x)=x^2-2ax+a^2-1$ 开口向上，对称轴为 $x=a$，}
\step{1) 当 $a\leqslant0$ 时，$g(x)$ 在 $[0,2]$ 区间上严格增，则必有 $\begin{cases}g(x)_{\min}=g(0)\leqslant1\\g(x)_{\max}=g(2)\geqslant2\end{cases}$，}
\step{解得 $-\sqrt2\leqslant a\leqslant0$；}
\step{2) 当 $a\geqslant2$ 时，$g(x)$ 在 $[0,2]$ 区间上严格减，则必有 $\begin{cases}g(x)_{\min}=g(2)\leqslant1\\g(x)_{\max}=g(0)\geqslant2\end{cases}$，}
\step{解得 $2\leqslant a\leqslant2+\sqrt2$；}
\step{3) 当 $0<a\leqslant1$ 时，$g(x)$ 在 $[0,a]$ 上严格减，在 $[a,2]$ 上严格增，}
\step{则必有 $\begin{cases}g(0)<1\\g(x)_{\max}=g(2)\geqslant2\end{cases}$，解得 $0<a\leqslant2-\sqrt3$，}
\step{4) 当 $1<a<2$ 时，$g(x)$ 在 $[0,a]$ 上严格减，在 $[a,2]$ 上严格增，}
\step{则必有 $\begin{cases}g(2)<1\\g(x)_{\max}=g(0)\geqslant2\end{cases}$，解得 $\sqrt3\leqslant a<2$.}
\step{综上所述，$a$ 的范围 $[-\sqrt2,2-\sqrt3]\cup[\sqrt3,2+\sqrt2]$.}% 原卷“范围”后缺“是／为”字，照录
\solend
\fi

\end{enumerate}

\end{document}
"
%
% 原始材料：微信公众号“魔都数学加油站”文章，作者破晓，2026-11-09 17:52 发布；
% 连载“27学年高一122”，原文标题“27学年高一122——2026-2027年上海市华师大二附中高一上周末作业6”；
% 原文链接：https://mp.weixin.qq.com/s/pC5066JFKNkQQwjMqTK8Ow。
% 正文为 14 张 PNG 页（第 14 页 2560×3620，其余 13 页 4762×6735），
% 题目下印有【解析】，为讲评版。
% 题目、解析均据原图转录；卷面红色水印及公众号页脚未录入。全卷无插图（已逐页检索，无“如图”）。
%
% 卷首标题：原卷印作“2026-2027 年华师大二附中高一上周末作业 6”（公众号标题中的“上海市”
% 卷面标题行没有，本稿照卷面），年份区间按全稿惯例排 en dash；原卷未印副标题，本稿按题目内容
% 标为“函数”。副标题理由（以全部 14 页实际内容为准）：逐页复核 19 题，考的都是函数性质——
% 奇偶性与充要条件（Q3、Q13、Q15、Q16）、单调性与单调区间（Q4、Q7、Q14、Q17）、定义域／
% 值域／最值与复合函数（Q5、Q6、Q8、Q10、Q12、Q19、Q20）、分段函数与取整函数（Q7、Q8、
% Q11、Q15、Q19）、抽象函数与新定义（Q12、Q20、Q21）；不等式仅作求参数范围／解集的工具
% （Q9、Q14、Q16、Q18），集合与常用逻辑用语仅零星出现（Q3 的充要条件、Q11③ 的集合、
% Q20 的“非空集合 I”），均非独立章节。故主体为单一模块“函数”，与同校 高一102 周末作业 5
% 同（该函数模块卷亦按 高一83／105 成例只用模块名、不叠“集合与常用逻辑用语、不等式”）。
%
% 与原卷的排版差异：
%   ① 原文从第 3 题起，填空题编号为 3—12、选择题为 13—16、解答题为 17—21，本稿保留原题号；
%   ② 原卷每题下直接印【解析】，本稿按全稿惯例置于答案版；四道选择题括号内均未印答案、
%      答案在解析末句“故选 …”，括号内留空（全卷无只印【答案】的题、无 \ans{}）；
%   ③ 原卷未印副标题，本稿按“知识模块”口径补出；
%   ④ 原卷填空横线按全稿惯例排为 \blankline[4.4em]；选择题括号排 \mbox{（\quad）}、
%      选项段前按惯例加 \keepwithprev；解答题不留作答区（原卷为讲评版）；
%   ⑤ 原卷数集排斜体 $R$、$N_+$（下标加号），按全稿惯例改 \R、\N_+
%      （$N_+$ 下标加号与 $\N^*$ 是两种记号，照录下标、只改数集字母为正体，同 高一05 第 20 题成例）；
%      绝对值排 \lvert…\rvert、集合描述竖线排 \mid；
%      $f(x)_{min}$／$f(x)_{max}$／$y_{max}$／$(y_{max})_{min}$ 一类斜体下标按 高一96 成例
%      改正体 \min／\max；不等式一律排 \geqslant／\leqslant；
%   ⑥ 第 20 题解析标签原卷印作【解得】（全卷唯一，其余 18 题均印【解析】），本稿按全稿惯例
%      由 \solbegin 统一印【解析】，正文内容照录（行内加 % 注、并入文件头存疑清单）；
%   ⑦ 题干与解析的标点（半角／全角逗号、冒号、括号、句点、子题括号）均照原卷，
%      题目正文引号用弯引号“ ”。
%
% 子集符号：第 21 题 (2) 的 $\left[\frac13+\frac1{9b-3},1\right]\subseteq\left[\frac12,b\right]$
%   印有下划线，照录 \subseteq。全卷未出现判别式 \triangle／\Delta。
%
% 原卷存疑／笔误（均照录加注）：
%   第 11 题解析对② 的“$y=y+\Delta y$”（按上下文应为 $y=[y]+\Delta y$，原卷漏方括号）；
%   第 19 题题干“某种产品 $x$ 万件）可获利”右括号无配对（疑漏印左括号“（”）；
%   第 20 题解析标签【解得】（见排版差异⑥）；
%   第 21 题 (3) 解析“总存在唯一的 $x_2\in[0,2]$ 时，使得……”的“时”字（按上下文衍字）；
%   第 21 题 (3) 末句“综上所述，$a$ 的范围 $[-\sqrt2,2-\sqrt3]\cup[\sqrt3,2+\sqrt2]$”
%      “范围”后缺“是／为”字，照录。
%
% 行数核对：逐题数 \step 条数，合计 162 条（Q3：5；Q4：3；Q5：1；Q6：3；Q7：2；
% Q8：5；Q9：7；Q10：10；Q11：8；Q12：4；Q13：12；Q14：6；Q15：14；Q16：14；
% Q17：6；Q18：8；Q19：7；Q20：19；Q21：28），与原卷印刷行逐题相等
% （cases 块作一个印刷行计）。全卷无 \ans{}、无只印【答案】的题。
%
\documentclass[a4paper,11pt]{article}
\usepackage{common}

\begin{document}

\examtitlest[3]{2026--2027 年华师大二附中高一上周末作业 6}{函数}

\bigsec{一、填空题}

\begin{enumerate}
\setcounter{enumi}{2}

\item 已知 $p:a=0$，$q$：函数 $f(x)=x^3+x+a$ 是奇函数，则 $p$ 是 $q$ 的\blankline[4.4em]条件.

\ifanswer
\solbegin
\step{若 $a=0$，则 $f(x)=x^3+x$，}
\step{易得 $f(x)$ 的定义域为 $\R$，且 $f(-x)=-(x^3+x)=-f(x)$，则 $f(x)$ 为奇函数，}
\step{反之，若 $f(x)=x^3+x+a$ 是奇函数，}
\step{则有 $f(-x)+f(x)=-x^3-x+a+x^3+x+a=2a=0$，必有 $a=0$，}
\step{故 $p$ 是 $q$ 的充要条件.}
\solend
\fi

\item 函数 $y=\sqrt{-x^2+2x}$ 的单调递增区间是\blankline[4.4em].

\ifanswer
\solbegin
\step{由函数定义域得 $-x^2+2x\geqslant0$，解得 $0\leqslant x\leqslant2$，}
\step{二次函数 $y=-x^2+2x$ 开口向下，对称轴为 $x=1$，}
\step{结合复合函数的单调性得函数 $y=\sqrt{-x^2+2x}$ 的单调递增区间是 $[0,1]$.}
\solend
\fi

\item 设 $f(x)$ 定义域为 $[1,4]$，且 $f(x)=2x+1$，则 $f(x^2)$ 的最大值为\blankline[4.4em].

\ifanswer
\solbegin
\step{由 $x^2\in[1,4]$，得 $f(x^2)=2x^2+1\leqslant9$，当 $x=\pm2$ 时取等号，最大值为 9.}
\solend
\fi

\item 函数 $f(x)=\sqrt{x-2}-x$，$x\in(2,6]$ 的值域为\blankline[4.4em].

\ifanswer
\solbegin
\step{令 $t=\sqrt{x-2}$，则 $x=2+t^2$，由 $2<x\leqslant6$ 得 $0<t\leqslant2$，}
\step{原函数可化为 $g(t)=-t^2+t-2$，$g(0)=-2$，$g(2)=-4$，$g\left(\dfrac12\right)=-\dfrac74$，}
\step{由二次函数性质得 $g(t)\in\left[-4,-\dfrac74\right]$.}
\solend
\fi

\item 若函数 $f(x)=\lvert x-a+1\rvert$ 在区间 $[1,+\infty)$ 上是严格增函数，则实数 $a$ 的取值范围是\blankline[4.4em].

\ifanswer
\solbegin
\step{$f(x)=\begin{cases}x-a+1,x\geqslant a-1\\a-1-x,x<a-1\end{cases}$，显然 $f(x)$ 在 $[a-1,+\infty)$ 上严格增，}
\step{要使函数在区间 $[1,+\infty)$ 上是严格增函数，则 $a-1\leqslant1$，即 $a\leqslant2$.}
\solend
\fi

\item 若函数 $f(x)=\begin{cases}x^2-5,x\leqslant2\\2ax+1,x>2\end{cases}$ 的值域是 $\R$，则 $a$ 的取值范围是\blankline[4.4em].

\ifanswer
\solbegin
\step{当 $x\leqslant2$ 时，$f(x)=x^2-5$，其开口向上，对称轴为 $x=0$，值域为 $[-5,+\infty)$，}
\step{由函数 $f(x)=\begin{cases}x^2-5,x\leqslant2\\2ax+1,x>2\end{cases}$ 的值域是 $\R$，}
\step{则当 $x>2$ 时，$f(x)=2ax+1$ 的值域应包含 $(-\infty,-5)$，}
\step{所以 $f(x)$ 为严格减函数，所以 $\begin{cases}2\times2a+1\geqslant-5\\2a<0\end{cases}$，即 $\begin{cases}a\geqslant-\dfrac32\\a<0\end{cases}$，}
\step{解得 $-\dfrac32\leqslant a<0$，故 $a$ 的取值范围是 $\left[-\dfrac32,0\right)$.}
\solend
\fi

\item 已知函数 $f(x)=ax^2+\dfrac bx$ 的图象经过点 $(-1,0)$，且满足 $f(1)=2$. 若 $\forall x\in(0,+\infty)$，都有 $f(x)+f\left(\dfrac1x\right)-2m\geqslant0$ 恒成立，则实数 $m$ 的取值范围为\blankline[4.4em].

\ifanswer
\solbegin
\step{因为 $f(x)=ax^2+\dfrac bx$ 的图象经过点 $(-1,0)$，且满足 $f(1)=2$，}
\step{所以 $\begin{cases}a-b=0\\a+b=2\end{cases}$，所以 $a=1$，$b=1$，即 $f(x)=x^2+\dfrac1x$，}
\step{$\forall x\in(0,+\infty)$，$f(x)+f\left(\dfrac1x\right)-2m\geqslant0\Leftrightarrow2m\leqslant x^2+\dfrac1x+\dfrac1{x^2}+x$ 恒成立，}
\step{令函数 $g(x)=x^2+\dfrac1x+\dfrac1{x^2}+x=\left(x+\dfrac1x\right)^2+x+\dfrac1x-2$，}
\step{而当 $x>0$ 时，$x+\dfrac1x\geqslant2$，当且仅当 $x=1$ 时取等号，}
\step{因此 $g(x)\geqslant2^2+2-2=4$，当且仅当 $x=1$ 时取等号，则 $2m\leqslant4$，解得 $m\leqslant2$，}
\step{所以实数 $m\in(-\infty,2]$.}
\solend
\fi

\item 若定义在区间 $[a,b](a<b)$ 上的函数 $f(x)=k-\sqrt{x+1}$ 值域也为 $[a,b]$，则实数 $k$ 的取值范围是\blankline[4.4em].

\ifanswer
\solbegin
\step{函数 $f(x)=k-\sqrt{x+1}$ 在定义域 $[-1,+\infty)$ 严格减，}
\step{当 $f(x)$ 的定义域为 $[a,b]$ 时，$f(x)$ 的值域也为 $[a,b]$，$a<b$，}
\step{故 $f(a)=k-\sqrt{a+1}=b$，$f(b)=k-\sqrt{b+1}=a$，}
\step{两式相减得 $\sqrt{a+1}-\sqrt{b+1}=a-b=(a+1)-(b+1)$}
\step{$=(\sqrt{a+1}-\sqrt{b+1})(\sqrt{a+1}+\sqrt{b+1})$，}
\step{所以 $\sqrt{a+1}-\sqrt{b+1}<0$，即 $\sqrt{a+1}+\sqrt{b+1}=1$，}
\step{则 $k=\sqrt{b+1}+a=a+1-\sqrt{a+1}$，}
\step{令 $\lambda=\sqrt{a+1}\geqslant0$，得 $k=\lambda^2-\lambda=\left(\lambda-\dfrac12\right)^2-\dfrac14$，}
\step{又 $\sqrt{a+1}<\sqrt{b+1}$，因为 $\sqrt{a+1}+\sqrt{b+1}=1$，所以 $\sqrt{a+1}\in\left[0,\dfrac12\right)$，}
\step{所以 $0\leqslant\lambda<\dfrac12$，故实数 $k$ 的取值范围为 $\left(-\dfrac14,0\right]$.}
\solend
\fi

\item 函数 $f(x)=[x]$ 的函数值表示不超过 $x$ 的最大整数，例如：$[-3.5]=-4$，$[2.1]=2$，称 $f(x)=[x]$ 为取整函数，则下列关于取整函数的说法：①对任意的 $x\in\R$，$[-x]=-[x]$ ②对任意的 $x$、$y\in\R$，$[x]+[y]\leqslant[x+y]$ ③集合 $\{y\mid y=[nx]-n[x],x\in\R,n\in\N_+\}$ 共有 $n$ 个元素 ④$x>0$ 时，关于 $x$ 的方程 $[2x]=\left[\dfrac1x\right]$ 有无数个解 其中正确的有\blankline[4.4em].

\ifanswer
\solbegin
\step{对①，因为 $[-2.1]=-3\neq[2.1]=2$，所以①错误；}
\step{对②，设 $x=[x]+\Delta x$，$\Delta x\in[0,1)$，$y=y+\Delta y$，$\Delta y\in[0,1)$，}% 原卷印作 $y=y+\Delta y$（按上下文应为 $y=[y]+\Delta y$），照录
\step{$[x+y]=[[x]+\Delta x+[y]+\Delta y]=[x]+[y]+[\Delta x+\Delta y]$，又 $\Delta x+\Delta y\in[0,\ 2)$，}
\step{故 $\Delta x+\Delta y=0$ 或 1，故 $[x]+[y]\leqslant[x+y]$，所以②正确；}
\step{对③，$[nx]-n[x]=[n[x]+n\cdot\Delta x]-n[x]=n[x]+[n\cdot\Delta x]-n[x]=[n\cdot\Delta x]$，}
\step{因为 $\Delta x\in[0,1)$，所以 $0\leqslant n\cdot\Delta x<n$，$n\in\N_+$，所以③正确；}
\step{对④，当 $\dfrac12<x<1$ 时，$[2x]=\left[\dfrac1x\right]=1$，故方程存在无数个解，所以④正确.}
\step{故正确的是②③④.}
\solend
\fi

\item 已知函数 $f(x)$ 定义域为 $[0,1]$，且满足 $f(x)+f(1-x)=1$，$f(x)=2f\left(\dfrac x5\right)$，对任意 $x_1\leqslant x_2$，恒有 $f(x_1)\leqslant f(x_2)$，则 $f\left(\dfrac1{2026}\right)$ 的值为\blankline[4.4em].

\ifanswer
\solbegin
\step{$f(0)=2f(0)\Rightarrow f(0)=0$，从而 $f(1)=1$，$f\left(\dfrac15\right)=\dfrac12$，$f\left(\dfrac45\right)=1-f\left(\dfrac15\right)=\dfrac12$，}
\step{由单调性，$x\in\left[\dfrac15,\dfrac45\right]$ 时，$f(x)=\dfrac12$.}
\step{因为 $\dfrac15<\dfrac{625}{2026}<\dfrac45$，反复使用 $f\left(\dfrac x5\right)=\dfrac12 f(x)$，}
\step{得 $f\left(\dfrac1{2026}\right)=\dfrac1{2^4}f\left(\dfrac{625}{2026}\right)=\dfrac1{32}$.}
\solend
\fi

\end{enumerate}

\bigsec{二、选择题}

\begin{enumerate}
\setcounter{enumi}{12}

\item 已知 $f(x)$ 为 $\R$ 上的奇函数，$g(x)$ 为 $\R$ 上的偶函数，且 $g(x)\neq0$，则下列说法正确的是\mbox{（\quad）}

\keepwithprev
A. $f(x)+g(x)$ 为 $\R$ 上的奇函数\qquad\qquad B. $f(x)-g(x)$ 为 $\R$ 上的奇函数

C. $\dfrac{f(x)}{g(x)}$ 为 $\R$ 上的偶函数\qquad\qquad D. $\lvert f(x)g(x)\rvert$ 为 $\R$ 上的偶函数

\ifanswer
\solbegin
\step{因为 $f(x)$ 为 $\R$ 上的奇函数，$g(x)$ 为 $\R$ 上的偶函数，且 $g(x)\neq0$，}
\step{所以 $f(-x)=-f(x)$，$g(-x)=g(x)$，}
\step{所以 $f(-x)+g(-x)=-f(x)+g(x)\neq-[f(x)+g(x)]$，}
\step{故 $f(x)+g(x)$ 为非奇非偶函数，$A$ 错误；}
\step{同理，$f(x)-g(x)$ 为非奇非偶函数，$B$ 错误；}
\step{设 $F(x)=\dfrac{f(x)}{g(x)}$，则 $F(-x)=\dfrac{f(-x)}{g(-x)}=\dfrac{-f(x)}{g(x)}=-F(x)$，}
\step{所以 $F(x)$ 为奇函数，$C$ 错误；}
\step{设函数 $H(x)=\lvert f(x)g(x)\rvert$，}
\step{因为 $f(x)$ 为 $\R$ 上的奇函数，$g(x)$ 为 $\R$ 上的偶函数，且 $g(x)\neq0$，}
\step{所以 $f(-x)=-f(x)$，$g(-x)=g(x)$，则由函数奇偶性的定义}
\step{得 $H(-x)=\lvert f(-x)g(-x)\rvert=\lvert -f(x)g(x)\rvert=\lvert f(x)g(x)\rvert=H(x)$，$D$ 正确.}
\step{故选 $D$.}
\solend
\fi

\item 若函数 $f(x)$ 是定义域为 $(-\infty,0)\cup(0,+\infty)$ 的奇函数，在区间 $(0,+\infty)$ 上严格增，且 $f(-1)=0$，则不等式 $\dfrac{f(x-1)}{x-2}<0$ 的解集为\mbox{（\quad）}

\keepwithprev
A. $(-\infty,-2)$\qquad\qquad B. $(-2,0)$\qquad\qquad C. $(0,1)$\qquad\qquad D. $(1,2)$

\ifanswer
\solbegin
\step{若函数 $f(x)$ 是定义域为 $(-\infty,0)\cup(0,+\infty)$ 的奇函数，在区间 $(0,+\infty)$ 上严格增，}
\step{由奇函数对称性得，$f(x)$ 在 $(-\infty,0)$ 上严格增，}
\step{因为 $f(-1)=0$，所以 $f(1)=0$，}
\step{则不等式 $\dfrac{f(x-1)}{x-2}<0$ 可转化为 $\begin{cases}x>2\\f(x-1)<0\end{cases}$ 或 $\begin{cases}x<2\\f(x-1)>0\end{cases}$，}
\step{即 $\begin{cases}x>2\\x-1<-1\text{或}0<x-1<1\end{cases}$ 或 $\begin{cases}x<2\\x-1>1\text{或}-1<x-1<0\end{cases}$，解得 $0<x<1$.}
\step{故选 $C$.}
\solend
\fi

\item 若函数 $f(x)$ 同时满足：(1) 对于定义域上的任意 $x$，恒有 $f(x)+f(-x)=0$；(2) 对于定义域上的任意 $x_1$、$x_2$，当 $x_1\neq x_2$ 时，恒有 $x_1f(x_1)-x_1f(x_2)>x_2f(x_1)-x_2f(x_2)$，则称函数 $f(x)$ 为“理想函数”. 给出下列四个函数：① $f(x)=x^2$；② $f(x)=x^3$；③ $f(x)=\dfrac{2x-1}{2x+1}$；④ $f(x)=\begin{cases}x^2+4x,x\geqslant0\\-x^2+4x,x<0\end{cases}$，其中被称为“理想函数”的有\mbox{（\quad）}

\keepwithprev
A. 1 个\qquad\qquad B. 2 个\qquad\qquad C. 3 个\qquad\qquad D. 4 个

\ifanswer
\solbegin
\step{函数 $f(x)$ 同时满足：(1) 对于定义域上的任意 $x$，恒有 $f(x)+f(-x)=0$；}
\step{(2) 对于定义域上的任意 $x_1$、$x_2$，当 $x_1\neq x_2$ 时，}
\step{恒有 $x_1f(x_1)-x_1f(x_2)>x_2f(x_1)-x_2f(x_2)$，}
\step{整理得 $(x_1-x_2)[f(x_1)-f(x_2)]>0$，则称函数 $f(x)$ 为“理想函数”.}
\step{即函数满足 (1) 奇函数；(2) 函数在定义域内为严格增函数；}
\step{① $f(x)=x^2$；不满足函数为奇函数在定义域内也不严格增；}
\step{② $f(x)=x^3$；满足函数既为奇函数，也在定义域内严格增，}
\step{故该函数为理想函数；}
\step{③ $f(x)=\dfrac{2x-1}{2x+1}=\dfrac{2x+1-2}{2x+1}=-\dfrac2{2x+1}+1$，该函数不为奇函数，}
\step{但是函数在定义域内严格增，故该函数不为理想函数；}
\step{④ $f(x)=\begin{cases}x^2+4x,x\geqslant0\\-x^2+4x,x<0\end{cases}$，函数的关系式满足 $f(-x)=-f(x)$，}
\step{且在 $x\geqslant0$ 时，函数严格增，由于奇函数在对称区间上单调性相同，}
\step{故该函数为理想函数.}
\step{故选 $B$.}
\solend
\fi

\item 已知 $f(x)$ 是定义在 $\R$ 上的奇函数，当 $x_1$、$x_2\in(0,+\infty)$ 且 $x_1\neq x_2$ 时，都有\newline $\dfrac{x_2f(x_1)-x_1f(x_2)}{x_1x_2(x_1-x_2)}>0$ 成立，$f(2028)=2028$，则不等式 $f(x)-x>0$ 的解集为\mbox{（\quad）}

\keepwithprev
A. $(-\infty,-2028)\cup(2028,+\infty)$\qquad\qquad B. $\left(-\dfrac1{2028},\dfrac1{2028}\right)$

C. $(-2028,2028)$\qquad\qquad D. $(-2028,0)\cup(2028,+\infty)$

\ifanswer
\solbegin
\step{因为 $f(x)$ 是定义在 $\R$ 上的奇函数，所以 $f(-x)=-f(x)$，}
\step{令 $g(x)=\dfrac{f(x)}x$，则 $g(-x)=\dfrac{f(-x)}{-x}=\dfrac{-f(x)}{-x}=\dfrac{f(x)}x=g(x)$，}
\step{所以 $g(x)$ 为偶函数，}
\step{当 $x_1$、$x_2\in(0,+\infty)$ 且 $x_1\neq x_2$ 时，都有 $\dfrac{x_2f(x_1)-x_1f(x_2)}{x_1x_2(x_1-x_2)}>0$ 成立，}
\step{即 $\dfrac{\dfrac{f(x_1)}{x_1}-\dfrac{f(x_2)}{x_2}}{x_1-x_2}>0$，即 $g(x)$ 在 $(0,+\infty)$ 上严格增，}
\step{由偶函数对称性得 $g(x)$ 在 $(-\infty,0)$ 上严格减，}
\step{因为 $f(2028)=2028$，所以 $g(2028)=1$，}
\step{则不等式 $f(x)-x>0$ 可化为 $f(x)>x$，}
\step{当 $x=0$ 时，$f(0)=0$，显然不成立，}
\step{当 $x>0$ 时，不等式可化为 $\dfrac{f(x)}x>1$，即 $g(x)>1=g(2028)$，}
\step{所以 $\lvert x\rvert>2028$，解得 $x>2028$，}
\step{当 $x<0$ 时，不等式可化为 $\dfrac{f(x)}x<1$，即 $g(x)<1=g(2028)$，}
\step{所以 $\lvert x\rvert<2028$，解得 $-2028<x<0$，}
\step{故不等式的解集为 $(-2028,0)\cup(2028,+\infty)$. 故选 $D$.}
\solend
\fi

\end{enumerate}

\bigsec{三、解答题}

\begin{enumerate}
\setcounter{enumi}{16}

\item 已知函数 $f(x)=\dfrac{2x-3}{x+1}$.

(1) 试用函数单调性的定义证明：函数 $f(x)$ 在 $(-1,+\infty)$ 上是严格增函数；

(2) 求函数 $f(x)$ 在区间 $[1,4]$ 上的最大值和最小值.

\ifanswer
\solbegin
\step{(1) 设 $-1<x_1<x_2$，则 $f(x_1)-f(x_2)=\dfrac{2x_1-3}{x_1+1}-\dfrac{2x_2-3}{x_2+1}=\dfrac{5(x_1-x_2)}{(x_1+1)(x_2+1)}$}
\step{因为 $-1<x_1<x_2$，所以 $x_1-x_2<0$，$x_1+1>0$，$x_2+1>0$，}
\step{所以 $f(x_1)-f(x_2)<0$，所以 $f(x)$ 在 $(-1,+\infty)$ 上严格增.}
\step{(2) 由 (1) 得 $f(x)$ 在区间 $[1,4]$ 上严格增，}
\step{所以 $f(x)_{\min}=f(1)=-\dfrac12$，$f(x)_{\max}=f(4)=1$，}
\step{所以函数 $f(x)$ 在区间 $[1,4]$ 上的最大值为 1，最小值为 $-\dfrac12$.}
\solend
\fi

\item 设函数 $f(x)=x^2-mx+m$.

(1) 当 $m=2$ 时，求不等式 $f(x)<x$ 的解集；

(2) 若 $\exists x\in[1,2]$，使不等式 $f(x)<x+m-2$ 成立，求实数 $m$ 的取值范围.

\ifanswer
\solbegin
\step{(1) $m=2$ 时，函数 $f(x)=x^2-2x+2$，}
\step{不等式 $f(x)<x$ 可化为 $x^2-3x+2<0$，解得 $1<x<2$，}
\step{所以不等式的解集为 $\{x\mid1<x<2\}$；}
\step{(2) 若 $\exists x\in[1,2]$，使不等式 $f(x)<x+m-2$ 成立，即 $x^2-x+2<mx$，}
\step{所以 $m>x+\dfrac2x-1$，设 $g(x)=x+\dfrac2x-1$，$x\in[1,2]$，}
\step{则 $g(x)=x+\dfrac2x-1\geqslant2\sqrt{x\cdot\dfrac2x}-1=2\sqrt2-1$，}
\step{当且仅当 $x=\dfrac2x$，即 $x=\sqrt2$ 时取等号，}
\step{所以实数 $m$ 的取值范围是 $\{m\mid m>2\sqrt2-1\}$.}
\solend
\fi

\item 某企业原来生产某种产品 $x$ 万件）可获利 $w(x)$（万元），且满足\newline% 原卷“万件）”右括号无配对、疑漏印左括号“（”，照录
$w(x)=\begin{cases}20(x^2+17),0<x\leqslant2\\500-\dfrac{80}{x-1},2<x\leqslant5\end{cases}$.
现该企业计划引进新的生产设备和新的产品方案优化产品，优化后的产品的其他成本总投入为 $(20x+10)$ 万元. 由市场调研分析得知，当前产品供不应求. 记该企业优化后的产品的利润为 $f(x)$（单位：万元）.

(1) 求函数 $f(x)$ 的解析式；

(2) 当优化后的产品产量为多少万件时，该企业的利润 $f(x)$ 最大？最大利润是多少？请说明理由.

\ifanswer
\solbegin
\step{(1) 由题意得 $f(x)=\begin{cases}20x^2-20x+330,0<x\leqslant2\\-20x-\dfrac{80}{x-1}+490,2<x\leqslant5\end{cases}$；}
\step{(2) 当 $0<x\leqslant2$ 时，$f(x)=20(x^2-x)+330=20\left(x-\dfrac12\right)^2+325$，}
\step{故当 $x=2$ 时，$f(x)$ 取最大值，$f(2)=370$；}
\step{当 $2<x\leqslant5$ 时，$f(x)=-\left[20(x-1)+\dfrac{80}{x-1}\right]+470\leqslant-80+470=390$，}
\step{当且仅当 $20(x-1)=\dfrac{80}{x-1}$，即 $x=3$ 时，$f(3)=390$ 为最大值；}
\step{因为 $370<390$，}
\step{所以优化后的产品产量为 3 万件时，企业获最大利润 390 万元.}
\solend
\fi

\item 俄国数学家切比雪夫（1821—1894）是研究直线逼近函数的理论先驱. 对定义在非空集合 $I$ 上的函数 $f(x)$，以及函数 $g(x)=kx+b(k,b\in\R)$，切比雪夫将函数 $y=\lvert f(x)-g(x)\rvert,x\in I$ 的最大值称为 $f(x),g(x)$ 的“偏差”.

(1) 函数 $f(x)=x^2,x\in[0,1],g(x)=-x-2$，求 $f(x),g(x)$ 的“偏差”；

(2) 函数 $f(x)=\dfrac1x+1,x\in[1,2],g(x)=kx+1(k>0)$，若 $f(x),g(x)$ 的“偏差”为 2，求 $k$ 的值；

(3) 函数 $f(x)=x^2-x,x\in[0,3],g(x)=2x+b$，当 $f(x),g(x)$ 的“偏差”取最小值时，求 $b$ 的值，并求出“偏差”的最小值.

\ifanswer
% 注：原卷本题解析标签印作【解得】（全卷唯一，其余 18 题均印【解析】）；
%     本稿按全稿惯例由 \solbegin 统一印【解析】，正文内容照录。见文件头排版差异⑥。
\solbegin
\step{(1) $y=\lvert f(x)-g(x)\rvert=\lvert x^2+x+2\rvert=\left\lvert\left(x+\dfrac12\right)^2+\dfrac74\right\rvert,x\in[0,1]$，}
\step{因为 $x\in[0,1]$，所以 $\left(x+\dfrac12\right)^2+\dfrac74\in[2,4]$，}
\step{则 $y=\left\lvert\left(x+\dfrac12\right)^2+\dfrac74\right\rvert\in[2,4]$，$y_{\max}=4$，}
\step{所以函数 $f(x)$ 与 $g(x)$ 的“偏差”为 4.}
\step{(2) 令 $t=f(x)-g(x)=\dfrac1x+1-(kx+1)=\dfrac1x-kx,x\in[1,2]$，}
\step{因为 $k>0$，所以 $t=\dfrac1x-kx$ 是单调减函数，}
\step{所以 $t\in\left[\dfrac12-2k,1-k\right]$，由单调，$y=\lvert t\rvert,t\in\left[\dfrac12-2k,1-k\right]$，且 $y_{\max}=2$.}
\step{当 $\left\lvert\dfrac12-2k\right\rvert\leqslant\lvert1-k\rvert$，即 $0<k\leqslant\dfrac12$ 时，$\lvert1-k\rvert=2$，解得 $k=3$ 或 $k=-1$，}
\step{均不符合；}
\step{当 $\left\lvert\dfrac12-2k\right\rvert>\lvert1-k\rvert$，即 $k>\dfrac12$ 时，$\left\lvert\dfrac12-2k\right\rvert=2$，$\dfrac12-2k=2$ 或 $\dfrac12-2k=-2$，}
\step{解得 $k=\dfrac54$ 或 $k=-\dfrac34$（舍），所以 $k=\dfrac54$.}
\step{(3) $y=\lvert f(x)-g(x)\rvert=\lvert x^2-x-(2x+b)\rvert=\lvert x^2-3x-b\rvert$}
\step{$=\left\lvert\left(x-\dfrac32\right)^2-\dfrac94-b\right\rvert,x\in[0,3]$，}
\step{因为 $x\in[0,3]$，所以 $\left(x-\dfrac32\right)^2-\dfrac94-b\in\left[-\dfrac94-b,-b\right]$，}
\step{由 $y=\left\lvert\left(x-\dfrac32\right)^2-\dfrac94-b\right\rvert$，则 $y_{\max}=\max\left\{\lvert b\rvert,\left\lvert b+\dfrac94\right\rvert\right\}$，}
\step{令 $\lvert b\rvert\geqslant\left\lvert b+\dfrac94\right\rvert$，即 $b^2\geqslant\left(b+\dfrac94\right)^2$，解得 $b\leqslant-\dfrac98$，}
\step{$y_{\max}=\max\left\{\lvert b\rvert,\left\lvert b+\dfrac94\right\rvert\right\}=\begin{cases}-b,b\leqslant-\dfrac98\\b+\dfrac94,b>-\dfrac98\end{cases}$，}
\step{故当且仅当 $b=-\dfrac98$ 时，$(y_{\max})_{\min}=\dfrac98$，}
\step{故当 $b$ 的值为 $-\dfrac98$ 时，函数 $f(x)$ 与 $g(x)$ 的“偏差”取最小值 $\dfrac98$.}
\solend
\fi

\item 若函数 $f(x)$ 与 $g(x)$ 满足：对任意的 $x_1\in D$，总存在唯一的 $x_2\in D$，使 $f(x_1)g(x_2)=m$ 成立，则称 $f(x)$ 是 $g(x)$ 在区间 $D$ 上的“$m$ 阶伴随函数”；对任意的 $x_1\in D$，总存在唯一的 $x_2\in D$，使 $f(x_1)f(x_2)=m$ 成立，则称 $f(x)$ 是区间 $D$ 上的“$m$ 阶自伴函数”.

(1) 判断 $f(x)=x^2+1$ 是否为区间 $[0,3]$ 上的“2 阶自伴函数”？并说明理由：

(2) 若函数 $f(x)=3x-1$ 为区间 $\left[\dfrac12,b\right]$ 上的“1 阶自伴函数”，求 $b$ 的值；

(3) 若 $f(x)=\dfrac4{x+2}$ 是 $g(x)=x^2-2ax+a^2-1$ 在区间 $[0,2]$ 上的“2 阶伴随函数”，求实数 $a$ 的取值范围.

\ifanswer
\solbegin
\step{(1) 假定函数 $f(x)=x^2+1$ 是区间 $[0,3]$ 上的“2 阶自伴函数”，}
\step{取 $x_1=2,f(x_1)=f(2)=5$，由 $f(x_1)f(x_2)=2$，得 $5(x_2^2+1)=2$，}
\step{显然此方程无实数解，}
\step{所以函数 $f(x)=x^2+1$ 不是区间 $[0,3]$ 上的“2 阶自伴函数”.}
\step{(2) 函数 $f(x)=3x-1$ 为区间 $\left[\dfrac12,b\right]$ 上的“1 阶自伴函数”，}
\step{则对任意 $x_1\in\left[\dfrac12,b\right]$，总存在唯一的 $x_2\in\left[\dfrac12,b\right]$，使得 $f(x_1)f(x_2)=1$，}
\step{即 $3x_2-1=\dfrac1{3x_1-1}$，整理得 $x_2=\dfrac13+\dfrac1{3(3x_1-1)}$，}
\step{显然函数 $x_2=\dfrac13+\dfrac1{3(3x_1-1)}$ 在 $\left[\dfrac12,b\right]$ 上严格减，}
\step{且当 $x_1=\dfrac12$ 时，$x_2=1$，当 $x_1=b$ 时，$x_2=\dfrac13+\dfrac1{9b-3}$，}
\step{因此对 $\left[\dfrac12,b\right]$ 内的每一个 $x_1$，在 $\left[\dfrac13+\dfrac1{9b-3},1\right]$ 内有唯一 $x_2$ 值与之对应，}
\step{则 $x_2\in\left[\dfrac12,b\right]$，于是 $\left[\dfrac13+\dfrac1{9b-3},1\right]\subseteq\left[\dfrac12,b\right]$，应有 $\begin{cases}b\geqslant1\\\dfrac13+\dfrac1{9b-3}\geqslant\dfrac12\end{cases}$，}
\step{得 $\begin{cases}b\geqslant1\\b\leqslant1\end{cases}$，即 $b=1$，所以 $b$ 的值是 1.}
\step{(3) 由函数 $f(x)=\dfrac4{x+2}$ 在区间 $[0,2]$ 的值域为 $[1,2]$，}
\step{因为 $f(x)=\dfrac4{x+2}$ 是 $g(x)=x^2-2ax+a^2-1$ 在 $[0,2]$ 上的“2 阶伴随函数”，}
\step{则对任意的 $x_1\in[0,2]$，总存在唯一的 $x_2\in[0,2]$ 时，使得 $f(x_1)g(x_2)=2$，}% 原卷此处印“时”字（按上下文为衍字），照录
\step{所以 $g(x_2)=\dfrac2{f(x_1)}\in[1,2]$，}
\step{即 $g(x)=x^2-2ax+a^2-1$ 在 $[0,2]$ 区间上的值域必定包含区间 $[1,2]$，}
\step{且 $g(x)$ 值域在 $[1,2]$ 对应的自变量是唯一的，}
\step{又因为函数 $g(x)=x^2-2ax+a^2-1$ 开口向上，对称轴为 $x=a$，}
\step{1) 当 $a\leqslant0$ 时，$g(x)$ 在 $[0,2]$ 区间上严格增，则必有 $\begin{cases}g(x)_{\min}=g(0)\leqslant1\\g(x)_{\max}=g(2)\geqslant2\end{cases}$，}
\step{解得 $-\sqrt2\leqslant a\leqslant0$；}
\step{2) 当 $a\geqslant2$ 时，$g(x)$ 在 $[0,2]$ 区间上严格减，则必有 $\begin{cases}g(x)_{\min}=g(2)\leqslant1\\g(x)_{\max}=g(0)\geqslant2\end{cases}$，}
\step{解得 $2\leqslant a\leqslant2+\sqrt2$；}
\step{3) 当 $0<a\leqslant1$ 时，$g(x)$ 在 $[0,a]$ 上严格减，在 $[a,2]$ 上严格增，}
\step{则必有 $\begin{cases}g(0)<1\\g(x)_{\max}=g(2)\geqslant2\end{cases}$，解得 $0<a\leqslant2-\sqrt3$，}
\step{4) 当 $1<a<2$ 时，$g(x)$ 在 $[0,a]$ 上严格减，在 $[a,2]$ 上严格增，}
\step{则必有 $\begin{cases}g(2)<1\\g(x)_{\max}=g(0)\geqslant2\end{cases}$，解得 $\sqrt3\leqslant a<2$.}
\step{综上所述，$a$ 的范围 $[-\sqrt2,2-\sqrt3]\cup[\sqrt3,2+\sqrt2]$.}% 原卷“范围”后缺“是／为”字，照录
\solend
\fi

\end{enumerate}

\end{document}
"
%
% 原始材料：微信公众号“魔都数学加油站”文章，作者破晓，2026-11-09 17:52 发布；
% 连载“27学年高一122”，原文标题“27学年高一122——2026-2027年上海市华师大二附中高一上周末作业6”；
% 原文链接：https://mp.weixin.qq.com/s/pC5066JFKNkQQwjMqTK8Ow。
% 正文为 14 张 PNG 页（第 14 页 2560×3620，其余 13 页 4762×6735），
% 题目下印有【解析】，为讲评版。
% 题目、解析均据原图转录；卷面红色水印及公众号页脚未录入。全卷无插图（已逐页检索，无“如图”）。
%
% 卷首标题：原卷印作“2026-2027 年华师大二附中高一上周末作业 6”（公众号标题中的“上海市”
% 卷面标题行没有，本稿照卷面），年份区间按全稿惯例排 en dash；原卷未印副标题，本稿按题目内容
% 标为“函数”。副标题理由（以全部 14 页实际内容为准）：逐页复核 19 题，考的都是函数性质——
% 奇偶性与充要条件（Q3、Q13、Q15、Q16）、单调性与单调区间（Q4、Q7、Q14、Q17）、定义域／
% 值域／最值与复合函数（Q5、Q6、Q8、Q10、Q12、Q19、Q20）、分段函数与取整函数（Q7、Q8、
% Q11、Q15、Q19）、抽象函数与新定义（Q12、Q20、Q21）；不等式仅作求参数范围／解集的工具
% （Q9、Q14、Q16、Q18），集合与常用逻辑用语仅零星出现（Q3 的充要条件、Q11③ 的集合、
% Q20 的“非空集合 I”），均非独立章节。故主体为单一模块“函数”，与同校 高一102 周末作业 5
% 同（该函数模块卷亦按 高一83／105 成例只用模块名、不叠“集合与常用逻辑用语、不等式”）。
%
% 与原卷的排版差异：
%   ① 原文从第 3 题起，填空题编号为 3—12、选择题为 13—16、解答题为 17—21，本稿保留原题号；
%   ② 原卷每题下直接印【解析】，本稿按全稿惯例置于答案版；四道选择题括号内均未印答案、
%      答案在解析末句“故选 …”，括号内留空（全卷无只印【答案】的题、无 \ans{}）；
%   ③ 原卷未印副标题，本稿按“知识模块”口径补出；
%   ④ 原卷填空横线按全稿惯例排为 \blankline[4.4em]；选择题括号排 \mbox{（\quad）}、
%      选项段前按惯例加 \keepwithprev；解答题不留作答区（原卷为讲评版）；
%   ⑤ 原卷数集排斜体 $R$、$N_+$（下标加号），按全稿惯例改 \R、\N_+
%      （$N_+$ 下标加号与 $\N^*$ 是两种记号，照录下标、只改数集字母为正体，同 高一05 第 20 题成例）；
%      绝对值排 \lvert…\rvert、集合描述竖线排 \mid；
%      $f(x)_{min}$／$f(x)_{max}$／$y_{max}$／$(y_{max})_{min}$ 一类斜体下标按 高一96 成例
%      改正体 \min／\max；不等式一律排 \geqslant／\leqslant；
%   ⑥ 第 20 题解析标签原卷印作【解得】（全卷唯一，其余 18 题均印【解析】），本稿按全稿惯例
%      由 \solbegin 统一印【解析】，正文内容照录（行内加 % 注、并入文件头存疑清单）；
%   ⑦ 题干与解析的标点（半角／全角逗号、冒号、括号、句点、子题括号）均照原卷，
%      题目正文引号用弯引号“ ”。
%
% 子集符号：第 21 题 (2) 的 $\left[\frac13+\frac1{9b-3},1\right]\subseteq\left[\frac12,b\right]$
%   印有下划线，照录 \subseteq。全卷未出现判别式 \triangle／\Delta。
%
% 原卷存疑／笔误（均照录加注）：
%   第 11 题解析对② 的“$y=y+\Delta y$”（按上下文应为 $y=[y]+\Delta y$，原卷漏方括号）；
%   第 19 题题干“某种产品 $x$ 万件）可获利”右括号无配对（疑漏印左括号“（”）；
%   第 20 题解析标签【解得】（见排版差异⑥）；
%   第 21 题 (3) 解析“总存在唯一的 $x_2\in[0,2]$ 时，使得……”的“时”字（按上下文衍字）；
%   第 21 题 (3) 末句“综上所述，$a$ 的范围 $[-\sqrt2,2-\sqrt3]\cup[\sqrt3,2+\sqrt2]$”
%      “范围”后缺“是／为”字，照录。
%
% 行数核对：逐题数 \step 条数，合计 162 条（Q3：5；Q4：3；Q5：1；Q6：3；Q7：2；
% Q8：5；Q9：7；Q10：10；Q11：8；Q12：4；Q13：12；Q14：6；Q15：14；Q16：14；
% Q17：6；Q18：8；Q19：7；Q20：19；Q21：28），与原卷印刷行逐题相等
% （cases 块作一个印刷行计）。全卷无 \ans{}、无只印【答案】的题。
%
\documentclass[a4paper,11pt]{article}
\usepackage{common}

\begin{document}

\examtitlest[3]{2026--2027 年华师大二附中高一上周末作业 6}{函数}

\bigsec{一、填空题}

\begin{enumerate}
\setcounter{enumi}{2}

\item 已知 $p:a=0$，$q$：函数 $f(x)=x^3+x+a$ 是奇函数，则 $p$ 是 $q$ 的\blankline[4.4em]条件.

\ifanswer
\solbegin
\step{若 $a=0$，则 $f(x)=x^3+x$，}
\step{易得 $f(x)$ 的定义域为 $\R$，且 $f(-x)=-(x^3+x)=-f(x)$，则 $f(x)$ 为奇函数，}
\step{反之，若 $f(x)=x^3+x+a$ 是奇函数，}
\step{则有 $f(-x)+f(x)=-x^3-x+a+x^3+x+a=2a=0$，必有 $a=0$，}
\step{故 $p$ 是 $q$ 的充要条件.}
\solend
\fi

\item 函数 $y=\sqrt{-x^2+2x}$ 的单调递增区间是\blankline[4.4em].

\ifanswer
\solbegin
\step{由函数定义域得 $-x^2+2x\geqslant0$，解得 $0\leqslant x\leqslant2$，}
\step{二次函数 $y=-x^2+2x$ 开口向下，对称轴为 $x=1$，}
\step{结合复合函数的单调性得函数 $y=\sqrt{-x^2+2x}$ 的单调递增区间是 $[0,1]$.}
\solend
\fi

\item 设 $f(x)$ 定义域为 $[1,4]$，且 $f(x)=2x+1$，则 $f(x^2)$ 的最大值为\blankline[4.4em].

\ifanswer
\solbegin
\step{由 $x^2\in[1,4]$，得 $f(x^2)=2x^2+1\leqslant9$，当 $x=\pm2$ 时取等号，最大值为 9.}
\solend
\fi

\item 函数 $f(x)=\sqrt{x-2}-x$，$x\in(2,6]$ 的值域为\blankline[4.4em].

\ifanswer
\solbegin
\step{令 $t=\sqrt{x-2}$，则 $x=2+t^2$，由 $2<x\leqslant6$ 得 $0<t\leqslant2$，}
\step{原函数可化为 $g(t)=-t^2+t-2$，$g(0)=-2$，$g(2)=-4$，$g\left(\dfrac12\right)=-\dfrac74$，}
\step{由二次函数性质得 $g(t)\in\left[-4,-\dfrac74\right]$.}
\solend
\fi

\item 若函数 $f(x)=\lvert x-a+1\rvert$ 在区间 $[1,+\infty)$ 上是严格增函数，则实数 $a$ 的取值范围是\blankline[4.4em].

\ifanswer
\solbegin
\step{$f(x)=\begin{cases}x-a+1,x\geqslant a-1\\a-1-x,x<a-1\end{cases}$，显然 $f(x)$ 在 $[a-1,+\infty)$ 上严格增，}
\step{要使函数在区间 $[1,+\infty)$ 上是严格增函数，则 $a-1\leqslant1$，即 $a\leqslant2$.}
\solend
\fi

\item 若函数 $f(x)=\begin{cases}x^2-5,x\leqslant2\\2ax+1,x>2\end{cases}$ 的值域是 $\R$，则 $a$ 的取值范围是\blankline[4.4em].

\ifanswer
\solbegin
\step{当 $x\leqslant2$ 时，$f(x)=x^2-5$，其开口向上，对称轴为 $x=0$，值域为 $[-5,+\infty)$，}
\step{由函数 $f(x)=\begin{cases}x^2-5,x\leqslant2\\2ax+1,x>2\end{cases}$ 的值域是 $\R$，}
\step{则当 $x>2$ 时，$f(x)=2ax+1$ 的值域应包含 $(-\infty,-5)$，}
\step{所以 $f(x)$ 为严格减函数，所以 $\begin{cases}2\times2a+1\geqslant-5\\2a<0\end{cases}$，即 $\begin{cases}a\geqslant-\dfrac32\\a<0\end{cases}$，}
\step{解得 $-\dfrac32\leqslant a<0$，故 $a$ 的取值范围是 $\left[-\dfrac32,0\right)$.}
\solend
\fi

\item 已知函数 $f(x)=ax^2+\dfrac bx$ 的图象经过点 $(-1,0)$，且满足 $f(1)=2$. 若 $\forall x\in(0,+\infty)$，都有 $f(x)+f\left(\dfrac1x\right)-2m\geqslant0$ 恒成立，则实数 $m$ 的取值范围为\blankline[4.4em].

\ifanswer
\solbegin
\step{因为 $f(x)=ax^2+\dfrac bx$ 的图象经过点 $(-1,0)$，且满足 $f(1)=2$，}
\step{所以 $\begin{cases}a-b=0\\a+b=2\end{cases}$，所以 $a=1$，$b=1$，即 $f(x)=x^2+\dfrac1x$，}
\step{$\forall x\in(0,+\infty)$，$f(x)+f\left(\dfrac1x\right)-2m\geqslant0\Leftrightarrow2m\leqslant x^2+\dfrac1x+\dfrac1{x^2}+x$ 恒成立，}
\step{令函数 $g(x)=x^2+\dfrac1x+\dfrac1{x^2}+x=\left(x+\dfrac1x\right)^2+x+\dfrac1x-2$，}
\step{而当 $x>0$ 时，$x+\dfrac1x\geqslant2$，当且仅当 $x=1$ 时取等号，}
\step{因此 $g(x)\geqslant2^2+2-2=4$，当且仅当 $x=1$ 时取等号，则 $2m\leqslant4$，解得 $m\leqslant2$，}
\step{所以实数 $m\in(-\infty,2]$.}
\solend
\fi

\item 若定义在区间 $[a,b](a<b)$ 上的函数 $f(x)=k-\sqrt{x+1}$ 值域也为 $[a,b]$，则实数 $k$ 的取值范围是\blankline[4.4em].

\ifanswer
\solbegin
\step{函数 $f(x)=k-\sqrt{x+1}$ 在定义域 $[-1,+\infty)$ 严格减，}
\step{当 $f(x)$ 的定义域为 $[a,b]$ 时，$f(x)$ 的值域也为 $[a,b]$，$a<b$，}
\step{故 $f(a)=k-\sqrt{a+1}=b$，$f(b)=k-\sqrt{b+1}=a$，}
\step{两式相减得 $\sqrt{a+1}-\sqrt{b+1}=a-b=(a+1)-(b+1)$}
\step{$=(\sqrt{a+1}-\sqrt{b+1})(\sqrt{a+1}+\sqrt{b+1})$，}
\step{所以 $\sqrt{a+1}-\sqrt{b+1}<0$，即 $\sqrt{a+1}+\sqrt{b+1}=1$，}
\step{则 $k=\sqrt{b+1}+a=a+1-\sqrt{a+1}$，}
\step{令 $\lambda=\sqrt{a+1}\geqslant0$，得 $k=\lambda^2-\lambda=\left(\lambda-\dfrac12\right)^2-\dfrac14$，}
\step{又 $\sqrt{a+1}<\sqrt{b+1}$，因为 $\sqrt{a+1}+\sqrt{b+1}=1$，所以 $\sqrt{a+1}\in\left[0,\dfrac12\right)$，}
\step{所以 $0\leqslant\lambda<\dfrac12$，故实数 $k$ 的取值范围为 $\left(-\dfrac14,0\right]$.}
\solend
\fi

\item 函数 $f(x)=[x]$ 的函数值表示不超过 $x$ 的最大整数，例如：$[-3.5]=-4$，$[2.1]=2$，称 $f(x)=[x]$ 为取整函数，则下列关于取整函数的说法：①对任意的 $x\in\R$，$[-x]=-[x]$ ②对任意的 $x$、$y\in\R$，$[x]+[y]\leqslant[x+y]$ ③集合 $\{y\mid y=[nx]-n[x],x\in\R,n\in\N_+\}$ 共有 $n$ 个元素 ④$x>0$ 时，关于 $x$ 的方程 $[2x]=\left[\dfrac1x\right]$ 有无数个解 其中正确的有\blankline[4.4em].

\ifanswer
\solbegin
\step{对①，因为 $[-2.1]=-3\neq[2.1]=2$，所以①错误；}
\step{对②，设 $x=[x]+\Delta x$，$\Delta x\in[0,1)$，$y=y+\Delta y$，$\Delta y\in[0,1)$，}% 原卷印作 $y=y+\Delta y$（按上下文应为 $y=[y]+\Delta y$），照录
\step{$[x+y]=[[x]+\Delta x+[y]+\Delta y]=[x]+[y]+[\Delta x+\Delta y]$，又 $\Delta x+\Delta y\in[0,\ 2)$，}
\step{故 $\Delta x+\Delta y=0$ 或 1，故 $[x]+[y]\leqslant[x+y]$，所以②正确；}
\step{对③，$[nx]-n[x]=[n[x]+n\cdot\Delta x]-n[x]=n[x]+[n\cdot\Delta x]-n[x]=[n\cdot\Delta x]$，}
\step{因为 $\Delta x\in[0,1)$，所以 $0\leqslant n\cdot\Delta x<n$，$n\in\N_+$，所以③正确；}
\step{对④，当 $\dfrac12<x<1$ 时，$[2x]=\left[\dfrac1x\right]=1$，故方程存在无数个解，所以④正确.}
\step{故正确的是②③④.}
\solend
\fi

\item 已知函数 $f(x)$ 定义域为 $[0,1]$，且满足 $f(x)+f(1-x)=1$，$f(x)=2f\left(\dfrac x5\right)$，对任意 $x_1\leqslant x_2$，恒有 $f(x_1)\leqslant f(x_2)$，则 $f\left(\dfrac1{2026}\right)$ 的值为\blankline[4.4em].

\ifanswer
\solbegin
\step{$f(0)=2f(0)\Rightarrow f(0)=0$，从而 $f(1)=1$，$f\left(\dfrac15\right)=\dfrac12$，$f\left(\dfrac45\right)=1-f\left(\dfrac15\right)=\dfrac12$，}
\step{由单调性，$x\in\left[\dfrac15,\dfrac45\right]$ 时，$f(x)=\dfrac12$.}
\step{因为 $\dfrac15<\dfrac{625}{2026}<\dfrac45$，反复使用 $f\left(\dfrac x5\right)=\dfrac12 f(x)$，}
\step{得 $f\left(\dfrac1{2026}\right)=\dfrac1{2^4}f\left(\dfrac{625}{2026}\right)=\dfrac1{32}$.}
\solend
\fi

\end{enumerate}

\bigsec{二、选择题}

\begin{enumerate}
\setcounter{enumi}{12}

\item 已知 $f(x)$ 为 $\R$ 上的奇函数，$g(x)$ 为 $\R$ 上的偶函数，且 $g(x)\neq0$，则下列说法正确的是\mbox{（\quad）}

\keepwithprev
A. $f(x)+g(x)$ 为 $\R$ 上的奇函数\qquad\qquad B. $f(x)-g(x)$ 为 $\R$ 上的奇函数

C. $\dfrac{f(x)}{g(x)}$ 为 $\R$ 上的偶函数\qquad\qquad D. $\lvert f(x)g(x)\rvert$ 为 $\R$ 上的偶函数

\ifanswer
\solbegin
\step{因为 $f(x)$ 为 $\R$ 上的奇函数，$g(x)$ 为 $\R$ 上的偶函数，且 $g(x)\neq0$，}
\step{所以 $f(-x)=-f(x)$，$g(-x)=g(x)$，}
\step{所以 $f(-x)+g(-x)=-f(x)+g(x)\neq-[f(x)+g(x)]$，}
\step{故 $f(x)+g(x)$ 为非奇非偶函数，$A$ 错误；}
\step{同理，$f(x)-g(x)$ 为非奇非偶函数，$B$ 错误；}
\step{设 $F(x)=\dfrac{f(x)}{g(x)}$，则 $F(-x)=\dfrac{f(-x)}{g(-x)}=\dfrac{-f(x)}{g(x)}=-F(x)$，}
\step{所以 $F(x)$ 为奇函数，$C$ 错误；}
\step{设函数 $H(x)=\lvert f(x)g(x)\rvert$，}
\step{因为 $f(x)$ 为 $\R$ 上的奇函数，$g(x)$ 为 $\R$ 上的偶函数，且 $g(x)\neq0$，}
\step{所以 $f(-x)=-f(x)$，$g(-x)=g(x)$，则由函数奇偶性的定义}
\step{得 $H(-x)=\lvert f(-x)g(-x)\rvert=\lvert -f(x)g(x)\rvert=\lvert f(x)g(x)\rvert=H(x)$，$D$ 正确.}
\step{故选 $D$.}
\solend
\fi

\item 若函数 $f(x)$ 是定义域为 $(-\infty,0)\cup(0,+\infty)$ 的奇函数，在区间 $(0,+\infty)$ 上严格增，且 $f(-1)=0$，则不等式 $\dfrac{f(x-1)}{x-2}<0$ 的解集为\mbox{（\quad）}

\keepwithprev
A. $(-\infty,-2)$\qquad\qquad B. $(-2,0)$\qquad\qquad C. $(0,1)$\qquad\qquad D. $(1,2)$

\ifanswer
\solbegin
\step{若函数 $f(x)$ 是定义域为 $(-\infty,0)\cup(0,+\infty)$ 的奇函数，在区间 $(0,+\infty)$ 上严格增，}
\step{由奇函数对称性得，$f(x)$ 在 $(-\infty,0)$ 上严格增，}
\step{因为 $f(-1)=0$，所以 $f(1)=0$，}
\step{则不等式 $\dfrac{f(x-1)}{x-2}<0$ 可转化为 $\begin{cases}x>2\\f(x-1)<0\end{cases}$ 或 $\begin{cases}x<2\\f(x-1)>0\end{cases}$，}
\step{即 $\begin{cases}x>2\\x-1<-1\text{或}0<x-1<1\end{cases}$ 或 $\begin{cases}x<2\\x-1>1\text{或}-1<x-1<0\end{cases}$，解得 $0<x<1$.}
\step{故选 $C$.}
\solend
\fi

\item 若函数 $f(x)$ 同时满足：(1) 对于定义域上的任意 $x$，恒有 $f(x)+f(-x)=0$；(2) 对于定义域上的任意 $x_1$、$x_2$，当 $x_1\neq x_2$ 时，恒有 $x_1f(x_1)-x_1f(x_2)>x_2f(x_1)-x_2f(x_2)$，则称函数 $f(x)$ 为“理想函数”. 给出下列四个函数：① $f(x)=x^2$；② $f(x)=x^3$；③ $f(x)=\dfrac{2x-1}{2x+1}$；④ $f(x)=\begin{cases}x^2+4x,x\geqslant0\\-x^2+4x,x<0\end{cases}$，其中被称为“理想函数”的有\mbox{（\quad）}

\keepwithprev
A. 1 个\qquad\qquad B. 2 个\qquad\qquad C. 3 个\qquad\qquad D. 4 个

\ifanswer
\solbegin
\step{函数 $f(x)$ 同时满足：(1) 对于定义域上的任意 $x$，恒有 $f(x)+f(-x)=0$；}
\step{(2) 对于定义域上的任意 $x_1$、$x_2$，当 $x_1\neq x_2$ 时，}
\step{恒有 $x_1f(x_1)-x_1f(x_2)>x_2f(x_1)-x_2f(x_2)$，}
\step{整理得 $(x_1-x_2)[f(x_1)-f(x_2)]>0$，则称函数 $f(x)$ 为“理想函数”.}
\step{即函数满足 (1) 奇函数；(2) 函数在定义域内为严格增函数；}
\step{① $f(x)=x^2$；不满足函数为奇函数在定义域内也不严格增；}
\step{② $f(x)=x^3$；满足函数既为奇函数，也在定义域内严格增，}
\step{故该函数为理想函数；}
\step{③ $f(x)=\dfrac{2x-1}{2x+1}=\dfrac{2x+1-2}{2x+1}=-\dfrac2{2x+1}+1$，该函数不为奇函数，}
\step{但是函数在定义域内严格增，故该函数不为理想函数；}
\step{④ $f(x)=\begin{cases}x^2+4x,x\geqslant0\\-x^2+4x,x<0\end{cases}$，函数的关系式满足 $f(-x)=-f(x)$，}
\step{且在 $x\geqslant0$ 时，函数严格增，由于奇函数在对称区间上单调性相同，}
\step{故该函数为理想函数.}
\step{故选 $B$.}
\solend
\fi

\item 已知 $f(x)$ 是定义在 $\R$ 上的奇函数，当 $x_1$、$x_2\in(0,+\infty)$ 且 $x_1\neq x_2$ 时，都有\newline $\dfrac{x_2f(x_1)-x_1f(x_2)}{x_1x_2(x_1-x_2)}>0$ 成立，$f(2028)=2028$，则不等式 $f(x)-x>0$ 的解集为\mbox{（\quad）}

\keepwithprev
A. $(-\infty,-2028)\cup(2028,+\infty)$\qquad\qquad B. $\left(-\dfrac1{2028},\dfrac1{2028}\right)$

C. $(-2028,2028)$\qquad\qquad D. $(-2028,0)\cup(2028,+\infty)$

\ifanswer
\solbegin
\step{因为 $f(x)$ 是定义在 $\R$ 上的奇函数，所以 $f(-x)=-f(x)$，}
\step{令 $g(x)=\dfrac{f(x)}x$，则 $g(-x)=\dfrac{f(-x)}{-x}=\dfrac{-f(x)}{-x}=\dfrac{f(x)}x=g(x)$，}
\step{所以 $g(x)$ 为偶函数，}
\step{当 $x_1$、$x_2\in(0,+\infty)$ 且 $x_1\neq x_2$ 时，都有 $\dfrac{x_2f(x_1)-x_1f(x_2)}{x_1x_2(x_1-x_2)}>0$ 成立，}
\step{即 $\dfrac{\dfrac{f(x_1)}{x_1}-\dfrac{f(x_2)}{x_2}}{x_1-x_2}>0$，即 $g(x)$ 在 $(0,+\infty)$ 上严格增，}
\step{由偶函数对称性得 $g(x)$ 在 $(-\infty,0)$ 上严格减，}
\step{因为 $f(2028)=2028$，所以 $g(2028)=1$，}
\step{则不等式 $f(x)-x>0$ 可化为 $f(x)>x$，}
\step{当 $x=0$ 时，$f(0)=0$，显然不成立，}
\step{当 $x>0$ 时，不等式可化为 $\dfrac{f(x)}x>1$，即 $g(x)>1=g(2028)$，}
\step{所以 $\lvert x\rvert>2028$，解得 $x>2028$，}
\step{当 $x<0$ 时，不等式可化为 $\dfrac{f(x)}x<1$，即 $g(x)<1=g(2028)$，}
\step{所以 $\lvert x\rvert<2028$，解得 $-2028<x<0$，}
\step{故不等式的解集为 $(-2028,0)\cup(2028,+\infty)$. 故选 $D$.}
\solend
\fi

\end{enumerate}

\bigsec{三、解答题}

\begin{enumerate}
\setcounter{enumi}{16}

\item 已知函数 $f(x)=\dfrac{2x-3}{x+1}$.

(1) 试用函数单调性的定义证明：函数 $f(x)$ 在 $(-1,+\infty)$ 上是严格增函数；

(2) 求函数 $f(x)$ 在区间 $[1,4]$ 上的最大值和最小值.

\ifanswer
\solbegin
\step{(1) 设 $-1<x_1<x_2$，则 $f(x_1)-f(x_2)=\dfrac{2x_1-3}{x_1+1}-\dfrac{2x_2-3}{x_2+1}=\dfrac{5(x_1-x_2)}{(x_1+1)(x_2+1)}$}
\step{因为 $-1<x_1<x_2$，所以 $x_1-x_2<0$，$x_1+1>0$，$x_2+1>0$，}
\step{所以 $f(x_1)-f(x_2)<0$，所以 $f(x)$ 在 $(-1,+\infty)$ 上严格增.}
\step{(2) 由 (1) 得 $f(x)$ 在区间 $[1,4]$ 上严格增，}
\step{所以 $f(x)_{\min}=f(1)=-\dfrac12$，$f(x)_{\max}=f(4)=1$，}
\step{所以函数 $f(x)$ 在区间 $[1,4]$ 上的最大值为 1，最小值为 $-\dfrac12$.}
\solend
\fi

\item 设函数 $f(x)=x^2-mx+m$.

(1) 当 $m=2$ 时，求不等式 $f(x)<x$ 的解集；

(2) 若 $\exists x\in[1,2]$，使不等式 $f(x)<x+m-2$ 成立，求实数 $m$ 的取值范围.

\ifanswer
\solbegin
\step{(1) $m=2$ 时，函数 $f(x)=x^2-2x+2$，}
\step{不等式 $f(x)<x$ 可化为 $x^2-3x+2<0$，解得 $1<x<2$，}
\step{所以不等式的解集为 $\{x\mid1<x<2\}$；}
\step{(2) 若 $\exists x\in[1,2]$，使不等式 $f(x)<x+m-2$ 成立，即 $x^2-x+2<mx$，}
\step{所以 $m>x+\dfrac2x-1$，设 $g(x)=x+\dfrac2x-1$，$x\in[1,2]$，}
\step{则 $g(x)=x+\dfrac2x-1\geqslant2\sqrt{x\cdot\dfrac2x}-1=2\sqrt2-1$，}
\step{当且仅当 $x=\dfrac2x$，即 $x=\sqrt2$ 时取等号，}
\step{所以实数 $m$ 的取值范围是 $\{m\mid m>2\sqrt2-1\}$.}
\solend
\fi

\item 某企业原来生产某种产品 $x$ 万件）可获利 $w(x)$（万元），且满足\newline% 原卷“万件）”右括号无配对、疑漏印左括号“（”，照录
$w(x)=\begin{cases}20(x^2+17),0<x\leqslant2\\500-\dfrac{80}{x-1},2<x\leqslant5\end{cases}$.
现该企业计划引进新的生产设备和新的产品方案优化产品，优化后的产品的其他成本总投入为 $(20x+10)$ 万元. 由市场调研分析得知，当前产品供不应求. 记该企业优化后的产品的利润为 $f(x)$（单位：万元）.

(1) 求函数 $f(x)$ 的解析式；

(2) 当优化后的产品产量为多少万件时，该企业的利润 $f(x)$ 最大？最大利润是多少？请说明理由.

\ifanswer
\solbegin
\step{(1) 由题意得 $f(x)=\begin{cases}20x^2-20x+330,0<x\leqslant2\\-20x-\dfrac{80}{x-1}+490,2<x\leqslant5\end{cases}$；}
\step{(2) 当 $0<x\leqslant2$ 时，$f(x)=20(x^2-x)+330=20\left(x-\dfrac12\right)^2+325$，}
\step{故当 $x=2$ 时，$f(x)$ 取最大值，$f(2)=370$；}
\step{当 $2<x\leqslant5$ 时，$f(x)=-\left[20(x-1)+\dfrac{80}{x-1}\right]+470\leqslant-80+470=390$，}
\step{当且仅当 $20(x-1)=\dfrac{80}{x-1}$，即 $x=3$ 时，$f(3)=390$ 为最大值；}
\step{因为 $370<390$，}
\step{所以优化后的产品产量为 3 万件时，企业获最大利润 390 万元.}
\solend
\fi

\item 俄国数学家切比雪夫（1821—1894）是研究直线逼近函数的理论先驱. 对定义在非空集合 $I$ 上的函数 $f(x)$，以及函数 $g(x)=kx+b(k,b\in\R)$，切比雪夫将函数 $y=\lvert f(x)-g(x)\rvert,x\in I$ 的最大值称为 $f(x),g(x)$ 的“偏差”.

(1) 函数 $f(x)=x^2,x\in[0,1],g(x)=-x-2$，求 $f(x),g(x)$ 的“偏差”；

(2) 函数 $f(x)=\dfrac1x+1,x\in[1,2],g(x)=kx+1(k>0)$，若 $f(x),g(x)$ 的“偏差”为 2，求 $k$ 的值；

(3) 函数 $f(x)=x^2-x,x\in[0,3],g(x)=2x+b$，当 $f(x),g(x)$ 的“偏差”取最小值时，求 $b$ 的值，并求出“偏差”的最小值.

\ifanswer
% 注：原卷本题解析标签印作【解得】（全卷唯一，其余 18 题均印【解析】）；
%     本稿按全稿惯例由 \solbegin 统一印【解析】，正文内容照录。见文件头排版差异⑥。
\solbegin
\step{(1) $y=\lvert f(x)-g(x)\rvert=\lvert x^2+x+2\rvert=\left\lvert\left(x+\dfrac12\right)^2+\dfrac74\right\rvert,x\in[0,1]$，}
\step{因为 $x\in[0,1]$，所以 $\left(x+\dfrac12\right)^2+\dfrac74\in[2,4]$，}
\step{则 $y=\left\lvert\left(x+\dfrac12\right)^2+\dfrac74\right\rvert\in[2,4]$，$y_{\max}=4$，}
\step{所以函数 $f(x)$ 与 $g(x)$ 的“偏差”为 4.}
\step{(2) 令 $t=f(x)-g(x)=\dfrac1x+1-(kx+1)=\dfrac1x-kx,x\in[1,2]$，}
\step{因为 $k>0$，所以 $t=\dfrac1x-kx$ 是单调减函数，}
\step{所以 $t\in\left[\dfrac12-2k,1-k\right]$，由单调，$y=\lvert t\rvert,t\in\left[\dfrac12-2k,1-k\right]$，且 $y_{\max}=2$.}
\step{当 $\left\lvert\dfrac12-2k\right\rvert\leqslant\lvert1-k\rvert$，即 $0<k\leqslant\dfrac12$ 时，$\lvert1-k\rvert=2$，解得 $k=3$ 或 $k=-1$，}
\step{均不符合；}
\step{当 $\left\lvert\dfrac12-2k\right\rvert>\lvert1-k\rvert$，即 $k>\dfrac12$ 时，$\left\lvert\dfrac12-2k\right\rvert=2$，$\dfrac12-2k=2$ 或 $\dfrac12-2k=-2$，}
\step{解得 $k=\dfrac54$ 或 $k=-\dfrac34$（舍），所以 $k=\dfrac54$.}
\step{(3) $y=\lvert f(x)-g(x)\rvert=\lvert x^2-x-(2x+b)\rvert=\lvert x^2-3x-b\rvert$}
\step{$=\left\lvert\left(x-\dfrac32\right)^2-\dfrac94-b\right\rvert,x\in[0,3]$，}
\step{因为 $x\in[0,3]$，所以 $\left(x-\dfrac32\right)^2-\dfrac94-b\in\left[-\dfrac94-b,-b\right]$，}
\step{由 $y=\left\lvert\left(x-\dfrac32\right)^2-\dfrac94-b\right\rvert$，则 $y_{\max}=\max\left\{\lvert b\rvert,\left\lvert b+\dfrac94\right\rvert\right\}$，}
\step{令 $\lvert b\rvert\geqslant\left\lvert b+\dfrac94\right\rvert$，即 $b^2\geqslant\left(b+\dfrac94\right)^2$，解得 $b\leqslant-\dfrac98$，}
\step{$y_{\max}=\max\left\{\lvert b\rvert,\left\lvert b+\dfrac94\right\rvert\right\}=\begin{cases}-b,b\leqslant-\dfrac98\\b+\dfrac94,b>-\dfrac98\end{cases}$，}
\step{故当且仅当 $b=-\dfrac98$ 时，$(y_{\max})_{\min}=\dfrac98$，}
\step{故当 $b$ 的值为 $-\dfrac98$ 时，函数 $f(x)$ 与 $g(x)$ 的“偏差”取最小值 $\dfrac98$.}
\solend
\fi

\item 若函数 $f(x)$ 与 $g(x)$ 满足：对任意的 $x_1\in D$，总存在唯一的 $x_2\in D$，使 $f(x_1)g(x_2)=m$ 成立，则称 $f(x)$ 是 $g(x)$ 在区间 $D$ 上的“$m$ 阶伴随函数”；对任意的 $x_1\in D$，总存在唯一的 $x_2\in D$，使 $f(x_1)f(x_2)=m$ 成立，则称 $f(x)$ 是区间 $D$ 上的“$m$ 阶自伴函数”.

(1) 判断 $f(x)=x^2+1$ 是否为区间 $[0,3]$ 上的“2 阶自伴函数”？并说明理由：

(2) 若函数 $f(x)=3x-1$ 为区间 $\left[\dfrac12,b\right]$ 上的“1 阶自伴函数”，求 $b$ 的值；

(3) 若 $f(x)=\dfrac4{x+2}$ 是 $g(x)=x^2-2ax+a^2-1$ 在区间 $[0,2]$ 上的“2 阶伴随函数”，求实数 $a$ 的取值范围.

\ifanswer
\solbegin
\step{(1) 假定函数 $f(x)=x^2+1$ 是区间 $[0,3]$ 上的“2 阶自伴函数”，}
\step{取 $x_1=2,f(x_1)=f(2)=5$，由 $f(x_1)f(x_2)=2$，得 $5(x_2^2+1)=2$，}
\step{显然此方程无实数解，}
\step{所以函数 $f(x)=x^2+1$ 不是区间 $[0,3]$ 上的“2 阶自伴函数”.}
\step{(2) 函数 $f(x)=3x-1$ 为区间 $\left[\dfrac12,b\right]$ 上的“1 阶自伴函数”，}
\step{则对任意 $x_1\in\left[\dfrac12,b\right]$，总存在唯一的 $x_2\in\left[\dfrac12,b\right]$，使得 $f(x_1)f(x_2)=1$，}
\step{即 $3x_2-1=\dfrac1{3x_1-1}$，整理得 $x_2=\dfrac13+\dfrac1{3(3x_1-1)}$，}
\step{显然函数 $x_2=\dfrac13+\dfrac1{3(3x_1-1)}$ 在 $\left[\dfrac12,b\right]$ 上严格减，}
\step{且当 $x_1=\dfrac12$ 时，$x_2=1$，当 $x_1=b$ 时，$x_2=\dfrac13+\dfrac1{9b-3}$，}
\step{因此对 $\left[\dfrac12,b\right]$ 内的每一个 $x_1$，在 $\left[\dfrac13+\dfrac1{9b-3},1\right]$ 内有唯一 $x_2$ 值与之对应，}
\step{则 $x_2\in\left[\dfrac12,b\right]$，于是 $\left[\dfrac13+\dfrac1{9b-3},1\right]\subseteq\left[\dfrac12,b\right]$，应有 $\begin{cases}b\geqslant1\\\dfrac13+\dfrac1{9b-3}\geqslant\dfrac12\end{cases}$，}
\step{得 $\begin{cases}b\geqslant1\\b\leqslant1\end{cases}$，即 $b=1$，所以 $b$ 的值是 1.}
\step{(3) 由函数 $f(x)=\dfrac4{x+2}$ 在区间 $[0,2]$ 的值域为 $[1,2]$，}
\step{因为 $f(x)=\dfrac4{x+2}$ 是 $g(x)=x^2-2ax+a^2-1$ 在 $[0,2]$ 上的“2 阶伴随函数”，}
\step{则对任意的 $x_1\in[0,2]$，总存在唯一的 $x_2\in[0,2]$ 时，使得 $f(x_1)g(x_2)=2$，}% 原卷此处印“时”字（按上下文为衍字），照录
\step{所以 $g(x_2)=\dfrac2{f(x_1)}\in[1,2]$，}
\step{即 $g(x)=x^2-2ax+a^2-1$ 在 $[0,2]$ 区间上的值域必定包含区间 $[1,2]$，}
\step{且 $g(x)$ 值域在 $[1,2]$ 对应的自变量是唯一的，}
\step{又因为函数 $g(x)=x^2-2ax+a^2-1$ 开口向上，对称轴为 $x=a$，}
\step{1) 当 $a\leqslant0$ 时，$g(x)$ 在 $[0,2]$ 区间上严格增，则必有 $\begin{cases}g(x)_{\min}=g(0)\leqslant1\\g(x)_{\max}=g(2)\geqslant2\end{cases}$，}
\step{解得 $-\sqrt2\leqslant a\leqslant0$；}
\step{2) 当 $a\geqslant2$ 时，$g(x)$ 在 $[0,2]$ 区间上严格减，则必有 $\begin{cases}g(x)_{\min}=g(2)\leqslant1\\g(x)_{\max}=g(0)\geqslant2\end{cases}$，}
\step{解得 $2\leqslant a\leqslant2+\sqrt2$；}
\step{3) 当 $0<a\leqslant1$ 时，$g(x)$ 在 $[0,a]$ 上严格减，在 $[a,2]$ 上严格增，}
\step{则必有 $\begin{cases}g(0)<1\\g(x)_{\max}=g(2)\geqslant2\end{cases}$，解得 $0<a\leqslant2-\sqrt3$，}
\step{4) 当 $1<a<2$ 时，$g(x)$ 在 $[0,a]$ 上严格减，在 $[a,2]$ 上严格增，}
\step{则必有 $\begin{cases}g(2)<1\\g(x)_{\max}=g(0)\geqslant2\end{cases}$，解得 $\sqrt3\leqslant a<2$.}
\step{综上所述，$a$ 的范围 $[-\sqrt2,2-\sqrt3]\cup[\sqrt3,2+\sqrt2]$.}% 原卷“范围”后缺“是／为”字，照录
\solend
\fi

\end{enumerate}

\end{document}
"
% 紧凑版：xelatex "\def\iscompact{1}% 高一122_华师大二附中_周末作业6.tex
%   —— 高一上数学“2026-2027 年华师大二附中高一上周末作业 6”（试卷版/答案版）
% 试卷版：xelatex 高一122_华师大二附中_周末作业6.tex
% 答案版：xelatex "\def\isanswer{1}% 高一122_华师大二附中_周末作业6.tex
%   —— 高一上数学“2026-2027 年华师大二附中高一上周末作业 6”（试卷版/答案版）
% 试卷版：xelatex 高一122_华师大二附中_周末作业6.tex
% 答案版：xelatex "\def\isanswer{1}% 高一122_华师大二附中_周末作业6.tex
%   —— 高一上数学“2026-2027 年华师大二附中高一上周末作业 6”（试卷版/答案版）
% 试卷版：xelatex 高一122_华师大二附中_周末作业6.tex
% 答案版：xelatex "\def\isanswer{1}\input{高一122_华师大二附中_周末作业6.tex}"
% 紧凑版：xelatex "\def\iscompact{1}\input{高一122_华师大二附中_周末作业6.tex}"
%
% 原始材料：微信公众号“魔都数学加油站”文章，作者破晓，2026-11-09 17:52 发布；
% 连载“27学年高一122”，原文标题“27学年高一122——2026-2027年上海市华师大二附中高一上周末作业6”；
% 原文链接：https://mp.weixin.qq.com/s/pC5066JFKNkQQwjMqTK8Ow。
% 正文为 14 张 PNG 页（第 14 页 2560×3620，其余 13 页 4762×6735），
% 题目下印有【解析】，为讲评版。
% 题目、解析均据原图转录；卷面红色水印及公众号页脚未录入。全卷无插图（已逐页检索，无“如图”）。
%
% 卷首标题：原卷印作“2026-2027 年华师大二附中高一上周末作业 6”（公众号标题中的“上海市”
% 卷面标题行没有，本稿照卷面），年份区间按全稿惯例排 en dash；原卷未印副标题，本稿按题目内容
% 标为“函数”。副标题理由（以全部 14 页实际内容为准）：逐页复核 19 题，考的都是函数性质——
% 奇偶性与充要条件（Q3、Q13、Q15、Q16）、单调性与单调区间（Q4、Q7、Q14、Q17）、定义域／
% 值域／最值与复合函数（Q5、Q6、Q8、Q10、Q12、Q19、Q20）、分段函数与取整函数（Q7、Q8、
% Q11、Q15、Q19）、抽象函数与新定义（Q12、Q20、Q21）；不等式仅作求参数范围／解集的工具
% （Q9、Q14、Q16、Q18），集合与常用逻辑用语仅零星出现（Q3 的充要条件、Q11③ 的集合、
% Q20 的“非空集合 I”），均非独立章节。故主体为单一模块“函数”，与同校 高一102 周末作业 5
% 同（该函数模块卷亦按 高一83／105 成例只用模块名、不叠“集合与常用逻辑用语、不等式”）。
%
% 与原卷的排版差异：
%   ① 原文从第 3 题起，填空题编号为 3—12、选择题为 13—16、解答题为 17—21，本稿保留原题号；
%   ② 原卷每题下直接印【解析】，本稿按全稿惯例置于答案版；四道选择题括号内均未印答案、
%      答案在解析末句“故选 …”，括号内留空（全卷无只印【答案】的题、无 \ans{}）；
%   ③ 原卷未印副标题，本稿按“知识模块”口径补出；
%   ④ 原卷填空横线按全稿惯例排为 \blankline[4.4em]；选择题括号排 \mbox{（\quad）}、
%      选项段前按惯例加 \keepwithprev；解答题不留作答区（原卷为讲评版）；
%   ⑤ 原卷数集排斜体 $R$、$N_+$（下标加号），按全稿惯例改 \R、\N_+
%      （$N_+$ 下标加号与 $\N^*$ 是两种记号，照录下标、只改数集字母为正体，同 高一05 第 20 题成例）；
%      绝对值排 \lvert…\rvert、集合描述竖线排 \mid；
%      $f(x)_{min}$／$f(x)_{max}$／$y_{max}$／$(y_{max})_{min}$ 一类斜体下标按 高一96 成例
%      改正体 \min／\max；不等式一律排 \geqslant／\leqslant；
%   ⑥ 第 20 题解析标签原卷印作【解得】（全卷唯一，其余 18 题均印【解析】），本稿按全稿惯例
%      由 \solbegin 统一印【解析】，正文内容照录（行内加 % 注、并入文件头存疑清单）；
%   ⑦ 题干与解析的标点（半角／全角逗号、冒号、括号、句点、子题括号）均照原卷，
%      题目正文引号用弯引号“ ”。
%
% 子集符号：第 21 题 (2) 的 $\left[\frac13+\frac1{9b-3},1\right]\subseteq\left[\frac12,b\right]$
%   印有下划线，照录 \subseteq。全卷未出现判别式 \triangle／\Delta。
%
% 原卷存疑／笔误（均照录加注）：
%   第 11 题解析对② 的“$y=y+\Delta y$”（按上下文应为 $y=[y]+\Delta y$，原卷漏方括号）；
%   第 19 题题干“某种产品 $x$ 万件）可获利”右括号无配对（疑漏印左括号“（”）；
%   第 20 题解析标签【解得】（见排版差异⑥）；
%   第 21 题 (3) 解析“总存在唯一的 $x_2\in[0,2]$ 时，使得……”的“时”字（按上下文衍字）；
%   第 21 题 (3) 末句“综上所述，$a$ 的范围 $[-\sqrt2,2-\sqrt3]\cup[\sqrt3,2+\sqrt2]$”
%      “范围”后缺“是／为”字，照录。
%
% 行数核对：逐题数 \step 条数，合计 162 条（Q3：5；Q4：3；Q5：1；Q6：3；Q7：2；
% Q8：5；Q9：7；Q10：10；Q11：8；Q12：4；Q13：12；Q14：6；Q15：14；Q16：14；
% Q17：6；Q18：8；Q19：7；Q20：19；Q21：28），与原卷印刷行逐题相等
% （cases 块作一个印刷行计）。全卷无 \ans{}、无只印【答案】的题。
%
\documentclass[a4paper,11pt]{article}
\usepackage{common}

\begin{document}

\examtitlest[3]{2026--2027 年华师大二附中高一上周末作业 6}{函数}

\bigsec{一、填空题}

\begin{enumerate}
\setcounter{enumi}{2}

\item 已知 $p:a=0$，$q$：函数 $f(x)=x^3+x+a$ 是奇函数，则 $p$ 是 $q$ 的\blankline[4.4em]条件.

\ifanswer
\solbegin
\step{若 $a=0$，则 $f(x)=x^3+x$，}
\step{易得 $f(x)$ 的定义域为 $\R$，且 $f(-x)=-(x^3+x)=-f(x)$，则 $f(x)$ 为奇函数，}
\step{反之，若 $f(x)=x^3+x+a$ 是奇函数，}
\step{则有 $f(-x)+f(x)=-x^3-x+a+x^3+x+a=2a=0$，必有 $a=0$，}
\step{故 $p$ 是 $q$ 的充要条件.}
\solend
\fi

\item 函数 $y=\sqrt{-x^2+2x}$ 的单调递增区间是\blankline[4.4em].

\ifanswer
\solbegin
\step{由函数定义域得 $-x^2+2x\geqslant0$，解得 $0\leqslant x\leqslant2$，}
\step{二次函数 $y=-x^2+2x$ 开口向下，对称轴为 $x=1$，}
\step{结合复合函数的单调性得函数 $y=\sqrt{-x^2+2x}$ 的单调递增区间是 $[0,1]$.}
\solend
\fi

\item 设 $f(x)$ 定义域为 $[1,4]$，且 $f(x)=2x+1$，则 $f(x^2)$ 的最大值为\blankline[4.4em].

\ifanswer
\solbegin
\step{由 $x^2\in[1,4]$，得 $f(x^2)=2x^2+1\leqslant9$，当 $x=\pm2$ 时取等号，最大值为 9.}
\solend
\fi

\item 函数 $f(x)=\sqrt{x-2}-x$，$x\in(2,6]$ 的值域为\blankline[4.4em].

\ifanswer
\solbegin
\step{令 $t=\sqrt{x-2}$，则 $x=2+t^2$，由 $2<x\leqslant6$ 得 $0<t\leqslant2$，}
\step{原函数可化为 $g(t)=-t^2+t-2$，$g(0)=-2$，$g(2)=-4$，$g\left(\dfrac12\right)=-\dfrac74$，}
\step{由二次函数性质得 $g(t)\in\left[-4,-\dfrac74\right]$.}
\solend
\fi

\item 若函数 $f(x)=\lvert x-a+1\rvert$ 在区间 $[1,+\infty)$ 上是严格增函数，则实数 $a$ 的取值范围是\blankline[4.4em].

\ifanswer
\solbegin
\step{$f(x)=\begin{cases}x-a+1,x\geqslant a-1\\a-1-x,x<a-1\end{cases}$，显然 $f(x)$ 在 $[a-1,+\infty)$ 上严格增，}
\step{要使函数在区间 $[1,+\infty)$ 上是严格增函数，则 $a-1\leqslant1$，即 $a\leqslant2$.}
\solend
\fi

\item 若函数 $f(x)=\begin{cases}x^2-5,x\leqslant2\\2ax+1,x>2\end{cases}$ 的值域是 $\R$，则 $a$ 的取值范围是\blankline[4.4em].

\ifanswer
\solbegin
\step{当 $x\leqslant2$ 时，$f(x)=x^2-5$，其开口向上，对称轴为 $x=0$，值域为 $[-5,+\infty)$，}
\step{由函数 $f(x)=\begin{cases}x^2-5,x\leqslant2\\2ax+1,x>2\end{cases}$ 的值域是 $\R$，}
\step{则当 $x>2$ 时，$f(x)=2ax+1$ 的值域应包含 $(-\infty,-5)$，}
\step{所以 $f(x)$ 为严格减函数，所以 $\begin{cases}2\times2a+1\geqslant-5\\2a<0\end{cases}$，即 $\begin{cases}a\geqslant-\dfrac32\\a<0\end{cases}$，}
\step{解得 $-\dfrac32\leqslant a<0$，故 $a$ 的取值范围是 $\left[-\dfrac32,0\right)$.}
\solend
\fi

\item 已知函数 $f(x)=ax^2+\dfrac bx$ 的图象经过点 $(-1,0)$，且满足 $f(1)=2$. 若 $\forall x\in(0,+\infty)$，都有 $f(x)+f\left(\dfrac1x\right)-2m\geqslant0$ 恒成立，则实数 $m$ 的取值范围为\blankline[4.4em].

\ifanswer
\solbegin
\step{因为 $f(x)=ax^2+\dfrac bx$ 的图象经过点 $(-1,0)$，且满足 $f(1)=2$，}
\step{所以 $\begin{cases}a-b=0\\a+b=2\end{cases}$，所以 $a=1$，$b=1$，即 $f(x)=x^2+\dfrac1x$，}
\step{$\forall x\in(0,+\infty)$，$f(x)+f\left(\dfrac1x\right)-2m\geqslant0\Leftrightarrow2m\leqslant x^2+\dfrac1x+\dfrac1{x^2}+x$ 恒成立，}
\step{令函数 $g(x)=x^2+\dfrac1x+\dfrac1{x^2}+x=\left(x+\dfrac1x\right)^2+x+\dfrac1x-2$，}
\step{而当 $x>0$ 时，$x+\dfrac1x\geqslant2$，当且仅当 $x=1$ 时取等号，}
\step{因此 $g(x)\geqslant2^2+2-2=4$，当且仅当 $x=1$ 时取等号，则 $2m\leqslant4$，解得 $m\leqslant2$，}
\step{所以实数 $m\in(-\infty,2]$.}
\solend
\fi

\item 若定义在区间 $[a,b](a<b)$ 上的函数 $f(x)=k-\sqrt{x+1}$ 值域也为 $[a,b]$，则实数 $k$ 的取值范围是\blankline[4.4em].

\ifanswer
\solbegin
\step{函数 $f(x)=k-\sqrt{x+1}$ 在定义域 $[-1,+\infty)$ 严格减，}
\step{当 $f(x)$ 的定义域为 $[a,b]$ 时，$f(x)$ 的值域也为 $[a,b]$，$a<b$，}
\step{故 $f(a)=k-\sqrt{a+1}=b$，$f(b)=k-\sqrt{b+1}=a$，}
\step{两式相减得 $\sqrt{a+1}-\sqrt{b+1}=a-b=(a+1)-(b+1)$}
\step{$=(\sqrt{a+1}-\sqrt{b+1})(\sqrt{a+1}+\sqrt{b+1})$，}
\step{所以 $\sqrt{a+1}-\sqrt{b+1}<0$，即 $\sqrt{a+1}+\sqrt{b+1}=1$，}
\step{则 $k=\sqrt{b+1}+a=a+1-\sqrt{a+1}$，}
\step{令 $\lambda=\sqrt{a+1}\geqslant0$，得 $k=\lambda^2-\lambda=\left(\lambda-\dfrac12\right)^2-\dfrac14$，}
\step{又 $\sqrt{a+1}<\sqrt{b+1}$，因为 $\sqrt{a+1}+\sqrt{b+1}=1$，所以 $\sqrt{a+1}\in\left[0,\dfrac12\right)$，}
\step{所以 $0\leqslant\lambda<\dfrac12$，故实数 $k$ 的取值范围为 $\left(-\dfrac14,0\right]$.}
\solend
\fi

\item 函数 $f(x)=[x]$ 的函数值表示不超过 $x$ 的最大整数，例如：$[-3.5]=-4$，$[2.1]=2$，称 $f(x)=[x]$ 为取整函数，则下列关于取整函数的说法：①对任意的 $x\in\R$，$[-x]=-[x]$ ②对任意的 $x$、$y\in\R$，$[x]+[y]\leqslant[x+y]$ ③集合 $\{y\mid y=[nx]-n[x],x\in\R,n\in\N_+\}$ 共有 $n$ 个元素 ④$x>0$ 时，关于 $x$ 的方程 $[2x]=\left[\dfrac1x\right]$ 有无数个解 其中正确的有\blankline[4.4em].

\ifanswer
\solbegin
\step{对①，因为 $[-2.1]=-3\neq[2.1]=2$，所以①错误；}
\step{对②，设 $x=[x]+\Delta x$，$\Delta x\in[0,1)$，$y=y+\Delta y$，$\Delta y\in[0,1)$，}% 原卷印作 $y=y+\Delta y$（按上下文应为 $y=[y]+\Delta y$），照录
\step{$[x+y]=[[x]+\Delta x+[y]+\Delta y]=[x]+[y]+[\Delta x+\Delta y]$，又 $\Delta x+\Delta y\in[0,\ 2)$，}
\step{故 $\Delta x+\Delta y=0$ 或 1，故 $[x]+[y]\leqslant[x+y]$，所以②正确；}
\step{对③，$[nx]-n[x]=[n[x]+n\cdot\Delta x]-n[x]=n[x]+[n\cdot\Delta x]-n[x]=[n\cdot\Delta x]$，}
\step{因为 $\Delta x\in[0,1)$，所以 $0\leqslant n\cdot\Delta x<n$，$n\in\N_+$，所以③正确；}
\step{对④，当 $\dfrac12<x<1$ 时，$[2x]=\left[\dfrac1x\right]=1$，故方程存在无数个解，所以④正确.}
\step{故正确的是②③④.}
\solend
\fi

\item 已知函数 $f(x)$ 定义域为 $[0,1]$，且满足 $f(x)+f(1-x)=1$，$f(x)=2f\left(\dfrac x5\right)$，对任意 $x_1\leqslant x_2$，恒有 $f(x_1)\leqslant f(x_2)$，则 $f\left(\dfrac1{2026}\right)$ 的值为\blankline[4.4em].

\ifanswer
\solbegin
\step{$f(0)=2f(0)\Rightarrow f(0)=0$，从而 $f(1)=1$，$f\left(\dfrac15\right)=\dfrac12$，$f\left(\dfrac45\right)=1-f\left(\dfrac15\right)=\dfrac12$，}
\step{由单调性，$x\in\left[\dfrac15,\dfrac45\right]$ 时，$f(x)=\dfrac12$.}
\step{因为 $\dfrac15<\dfrac{625}{2026}<\dfrac45$，反复使用 $f\left(\dfrac x5\right)=\dfrac12 f(x)$，}
\step{得 $f\left(\dfrac1{2026}\right)=\dfrac1{2^4}f\left(\dfrac{625}{2026}\right)=\dfrac1{32}$.}
\solend
\fi

\end{enumerate}

\bigsec{二、选择题}

\begin{enumerate}
\setcounter{enumi}{12}

\item 已知 $f(x)$ 为 $\R$ 上的奇函数，$g(x)$ 为 $\R$ 上的偶函数，且 $g(x)\neq0$，则下列说法正确的是\mbox{（\quad）}

\keepwithprev
A. $f(x)+g(x)$ 为 $\R$ 上的奇函数\qquad\qquad B. $f(x)-g(x)$ 为 $\R$ 上的奇函数

C. $\dfrac{f(x)}{g(x)}$ 为 $\R$ 上的偶函数\qquad\qquad D. $\lvert f(x)g(x)\rvert$ 为 $\R$ 上的偶函数

\ifanswer
\solbegin
\step{因为 $f(x)$ 为 $\R$ 上的奇函数，$g(x)$ 为 $\R$ 上的偶函数，且 $g(x)\neq0$，}
\step{所以 $f(-x)=-f(x)$，$g(-x)=g(x)$，}
\step{所以 $f(-x)+g(-x)=-f(x)+g(x)\neq-[f(x)+g(x)]$，}
\step{故 $f(x)+g(x)$ 为非奇非偶函数，$A$ 错误；}
\step{同理，$f(x)-g(x)$ 为非奇非偶函数，$B$ 错误；}
\step{设 $F(x)=\dfrac{f(x)}{g(x)}$，则 $F(-x)=\dfrac{f(-x)}{g(-x)}=\dfrac{-f(x)}{g(x)}=-F(x)$，}
\step{所以 $F(x)$ 为奇函数，$C$ 错误；}
\step{设函数 $H(x)=\lvert f(x)g(x)\rvert$，}
\step{因为 $f(x)$ 为 $\R$ 上的奇函数，$g(x)$ 为 $\R$ 上的偶函数，且 $g(x)\neq0$，}
\step{所以 $f(-x)=-f(x)$，$g(-x)=g(x)$，则由函数奇偶性的定义}
\step{得 $H(-x)=\lvert f(-x)g(-x)\rvert=\lvert -f(x)g(x)\rvert=\lvert f(x)g(x)\rvert=H(x)$，$D$ 正确.}
\step{故选 $D$.}
\solend
\fi

\item 若函数 $f(x)$ 是定义域为 $(-\infty,0)\cup(0,+\infty)$ 的奇函数，在区间 $(0,+\infty)$ 上严格增，且 $f(-1)=0$，则不等式 $\dfrac{f(x-1)}{x-2}<0$ 的解集为\mbox{（\quad）}

\keepwithprev
A. $(-\infty,-2)$\qquad\qquad B. $(-2,0)$\qquad\qquad C. $(0,1)$\qquad\qquad D. $(1,2)$

\ifanswer
\solbegin
\step{若函数 $f(x)$ 是定义域为 $(-\infty,0)\cup(0,+\infty)$ 的奇函数，在区间 $(0,+\infty)$ 上严格增，}
\step{由奇函数对称性得，$f(x)$ 在 $(-\infty,0)$ 上严格增，}
\step{因为 $f(-1)=0$，所以 $f(1)=0$，}
\step{则不等式 $\dfrac{f(x-1)}{x-2}<0$ 可转化为 $\begin{cases}x>2\\f(x-1)<0\end{cases}$ 或 $\begin{cases}x<2\\f(x-1)>0\end{cases}$，}
\step{即 $\begin{cases}x>2\\x-1<-1\text{或}0<x-1<1\end{cases}$ 或 $\begin{cases}x<2\\x-1>1\text{或}-1<x-1<0\end{cases}$，解得 $0<x<1$.}
\step{故选 $C$.}
\solend
\fi

\item 若函数 $f(x)$ 同时满足：(1) 对于定义域上的任意 $x$，恒有 $f(x)+f(-x)=0$；(2) 对于定义域上的任意 $x_1$、$x_2$，当 $x_1\neq x_2$ 时，恒有 $x_1f(x_1)-x_1f(x_2)>x_2f(x_1)-x_2f(x_2)$，则称函数 $f(x)$ 为“理想函数”. 给出下列四个函数：① $f(x)=x^2$；② $f(x)=x^3$；③ $f(x)=\dfrac{2x-1}{2x+1}$；④ $f(x)=\begin{cases}x^2+4x,x\geqslant0\\-x^2+4x,x<0\end{cases}$，其中被称为“理想函数”的有\mbox{（\quad）}

\keepwithprev
A. 1 个\qquad\qquad B. 2 个\qquad\qquad C. 3 个\qquad\qquad D. 4 个

\ifanswer
\solbegin
\step{函数 $f(x)$ 同时满足：(1) 对于定义域上的任意 $x$，恒有 $f(x)+f(-x)=0$；}
\step{(2) 对于定义域上的任意 $x_1$、$x_2$，当 $x_1\neq x_2$ 时，}
\step{恒有 $x_1f(x_1)-x_1f(x_2)>x_2f(x_1)-x_2f(x_2)$，}
\step{整理得 $(x_1-x_2)[f(x_1)-f(x_2)]>0$，则称函数 $f(x)$ 为“理想函数”.}
\step{即函数满足 (1) 奇函数；(2) 函数在定义域内为严格增函数；}
\step{① $f(x)=x^2$；不满足函数为奇函数在定义域内也不严格增；}
\step{② $f(x)=x^3$；满足函数既为奇函数，也在定义域内严格增，}
\step{故该函数为理想函数；}
\step{③ $f(x)=\dfrac{2x-1}{2x+1}=\dfrac{2x+1-2}{2x+1}=-\dfrac2{2x+1}+1$，该函数不为奇函数，}
\step{但是函数在定义域内严格增，故该函数不为理想函数；}
\step{④ $f(x)=\begin{cases}x^2+4x,x\geqslant0\\-x^2+4x,x<0\end{cases}$，函数的关系式满足 $f(-x)=-f(x)$，}
\step{且在 $x\geqslant0$ 时，函数严格增，由于奇函数在对称区间上单调性相同，}
\step{故该函数为理想函数.}
\step{故选 $B$.}
\solend
\fi

\item 已知 $f(x)$ 是定义在 $\R$ 上的奇函数，当 $x_1$、$x_2\in(0,+\infty)$ 且 $x_1\neq x_2$ 时，都有\newline $\dfrac{x_2f(x_1)-x_1f(x_2)}{x_1x_2(x_1-x_2)}>0$ 成立，$f(2028)=2028$，则不等式 $f(x)-x>0$ 的解集为\mbox{（\quad）}

\keepwithprev
A. $(-\infty,-2028)\cup(2028,+\infty)$\qquad\qquad B. $\left(-\dfrac1{2028},\dfrac1{2028}\right)$

C. $(-2028,2028)$\qquad\qquad D. $(-2028,0)\cup(2028,+\infty)$

\ifanswer
\solbegin
\step{因为 $f(x)$ 是定义在 $\R$ 上的奇函数，所以 $f(-x)=-f(x)$，}
\step{令 $g(x)=\dfrac{f(x)}x$，则 $g(-x)=\dfrac{f(-x)}{-x}=\dfrac{-f(x)}{-x}=\dfrac{f(x)}x=g(x)$，}
\step{所以 $g(x)$ 为偶函数，}
\step{当 $x_1$、$x_2\in(0,+\infty)$ 且 $x_1\neq x_2$ 时，都有 $\dfrac{x_2f(x_1)-x_1f(x_2)}{x_1x_2(x_1-x_2)}>0$ 成立，}
\step{即 $\dfrac{\dfrac{f(x_1)}{x_1}-\dfrac{f(x_2)}{x_2}}{x_1-x_2}>0$，即 $g(x)$ 在 $(0,+\infty)$ 上严格增，}
\step{由偶函数对称性得 $g(x)$ 在 $(-\infty,0)$ 上严格减，}
\step{因为 $f(2028)=2028$，所以 $g(2028)=1$，}
\step{则不等式 $f(x)-x>0$ 可化为 $f(x)>x$，}
\step{当 $x=0$ 时，$f(0)=0$，显然不成立，}
\step{当 $x>0$ 时，不等式可化为 $\dfrac{f(x)}x>1$，即 $g(x)>1=g(2028)$，}
\step{所以 $\lvert x\rvert>2028$，解得 $x>2028$，}
\step{当 $x<0$ 时，不等式可化为 $\dfrac{f(x)}x<1$，即 $g(x)<1=g(2028)$，}
\step{所以 $\lvert x\rvert<2028$，解得 $-2028<x<0$，}
\step{故不等式的解集为 $(-2028,0)\cup(2028,+\infty)$. 故选 $D$.}
\solend
\fi

\end{enumerate}

\bigsec{三、解答题}

\begin{enumerate}
\setcounter{enumi}{16}

\item 已知函数 $f(x)=\dfrac{2x-3}{x+1}$.

(1) 试用函数单调性的定义证明：函数 $f(x)$ 在 $(-1,+\infty)$ 上是严格增函数；

(2) 求函数 $f(x)$ 在区间 $[1,4]$ 上的最大值和最小值.

\ifanswer
\solbegin
\step{(1) 设 $-1<x_1<x_2$，则 $f(x_1)-f(x_2)=\dfrac{2x_1-3}{x_1+1}-\dfrac{2x_2-3}{x_2+1}=\dfrac{5(x_1-x_2)}{(x_1+1)(x_2+1)}$}
\step{因为 $-1<x_1<x_2$，所以 $x_1-x_2<0$，$x_1+1>0$，$x_2+1>0$，}
\step{所以 $f(x_1)-f(x_2)<0$，所以 $f(x)$ 在 $(-1,+\infty)$ 上严格增.}
\step{(2) 由 (1) 得 $f(x)$ 在区间 $[1,4]$ 上严格增，}
\step{所以 $f(x)_{\min}=f(1)=-\dfrac12$，$f(x)_{\max}=f(4)=1$，}
\step{所以函数 $f(x)$ 在区间 $[1,4]$ 上的最大值为 1，最小值为 $-\dfrac12$.}
\solend
\fi

\item 设函数 $f(x)=x^2-mx+m$.

(1) 当 $m=2$ 时，求不等式 $f(x)<x$ 的解集；

(2) 若 $\exists x\in[1,2]$，使不等式 $f(x)<x+m-2$ 成立，求实数 $m$ 的取值范围.

\ifanswer
\solbegin
\step{(1) $m=2$ 时，函数 $f(x)=x^2-2x+2$，}
\step{不等式 $f(x)<x$ 可化为 $x^2-3x+2<0$，解得 $1<x<2$，}
\step{所以不等式的解集为 $\{x\mid1<x<2\}$；}
\step{(2) 若 $\exists x\in[1,2]$，使不等式 $f(x)<x+m-2$ 成立，即 $x^2-x+2<mx$，}
\step{所以 $m>x+\dfrac2x-1$，设 $g(x)=x+\dfrac2x-1$，$x\in[1,2]$，}
\step{则 $g(x)=x+\dfrac2x-1\geqslant2\sqrt{x\cdot\dfrac2x}-1=2\sqrt2-1$，}
\step{当且仅当 $x=\dfrac2x$，即 $x=\sqrt2$ 时取等号，}
\step{所以实数 $m$ 的取值范围是 $\{m\mid m>2\sqrt2-1\}$.}
\solend
\fi

\item 某企业原来生产某种产品 $x$ 万件）可获利 $w(x)$（万元），且满足\newline% 原卷“万件）”右括号无配对、疑漏印左括号“（”，照录
$w(x)=\begin{cases}20(x^2+17),0<x\leqslant2\\500-\dfrac{80}{x-1},2<x\leqslant5\end{cases}$.
现该企业计划引进新的生产设备和新的产品方案优化产品，优化后的产品的其他成本总投入为 $(20x+10)$ 万元. 由市场调研分析得知，当前产品供不应求. 记该企业优化后的产品的利润为 $f(x)$（单位：万元）.

(1) 求函数 $f(x)$ 的解析式；

(2) 当优化后的产品产量为多少万件时，该企业的利润 $f(x)$ 最大？最大利润是多少？请说明理由.

\ifanswer
\solbegin
\step{(1) 由题意得 $f(x)=\begin{cases}20x^2-20x+330,0<x\leqslant2\\-20x-\dfrac{80}{x-1}+490,2<x\leqslant5\end{cases}$；}
\step{(2) 当 $0<x\leqslant2$ 时，$f(x)=20(x^2-x)+330=20\left(x-\dfrac12\right)^2+325$，}
\step{故当 $x=2$ 时，$f(x)$ 取最大值，$f(2)=370$；}
\step{当 $2<x\leqslant5$ 时，$f(x)=-\left[20(x-1)+\dfrac{80}{x-1}\right]+470\leqslant-80+470=390$，}
\step{当且仅当 $20(x-1)=\dfrac{80}{x-1}$，即 $x=3$ 时，$f(3)=390$ 为最大值；}
\step{因为 $370<390$，}
\step{所以优化后的产品产量为 3 万件时，企业获最大利润 390 万元.}
\solend
\fi

\item 俄国数学家切比雪夫（1821—1894）是研究直线逼近函数的理论先驱. 对定义在非空集合 $I$ 上的函数 $f(x)$，以及函数 $g(x)=kx+b(k,b\in\R)$，切比雪夫将函数 $y=\lvert f(x)-g(x)\rvert,x\in I$ 的最大值称为 $f(x),g(x)$ 的“偏差”.

(1) 函数 $f(x)=x^2,x\in[0,1],g(x)=-x-2$，求 $f(x),g(x)$ 的“偏差”；

(2) 函数 $f(x)=\dfrac1x+1,x\in[1,2],g(x)=kx+1(k>0)$，若 $f(x),g(x)$ 的“偏差”为 2，求 $k$ 的值；

(3) 函数 $f(x)=x^2-x,x\in[0,3],g(x)=2x+b$，当 $f(x),g(x)$ 的“偏差”取最小值时，求 $b$ 的值，并求出“偏差”的最小值.

\ifanswer
% 注：原卷本题解析标签印作【解得】（全卷唯一，其余 18 题均印【解析】）；
%     本稿按全稿惯例由 \solbegin 统一印【解析】，正文内容照录。见文件头排版差异⑥。
\solbegin
\step{(1) $y=\lvert f(x)-g(x)\rvert=\lvert x^2+x+2\rvert=\left\lvert\left(x+\dfrac12\right)^2+\dfrac74\right\rvert,x\in[0,1]$，}
\step{因为 $x\in[0,1]$，所以 $\left(x+\dfrac12\right)^2+\dfrac74\in[2,4]$，}
\step{则 $y=\left\lvert\left(x+\dfrac12\right)^2+\dfrac74\right\rvert\in[2,4]$，$y_{\max}=4$，}
\step{所以函数 $f(x)$ 与 $g(x)$ 的“偏差”为 4.}
\step{(2) 令 $t=f(x)-g(x)=\dfrac1x+1-(kx+1)=\dfrac1x-kx,x\in[1,2]$，}
\step{因为 $k>0$，所以 $t=\dfrac1x-kx$ 是单调减函数，}
\step{所以 $t\in\left[\dfrac12-2k,1-k\right]$，由单调，$y=\lvert t\rvert,t\in\left[\dfrac12-2k,1-k\right]$，且 $y_{\max}=2$.}
\step{当 $\left\lvert\dfrac12-2k\right\rvert\leqslant\lvert1-k\rvert$，即 $0<k\leqslant\dfrac12$ 时，$\lvert1-k\rvert=2$，解得 $k=3$ 或 $k=-1$，}
\step{均不符合；}
\step{当 $\left\lvert\dfrac12-2k\right\rvert>\lvert1-k\rvert$，即 $k>\dfrac12$ 时，$\left\lvert\dfrac12-2k\right\rvert=2$，$\dfrac12-2k=2$ 或 $\dfrac12-2k=-2$，}
\step{解得 $k=\dfrac54$ 或 $k=-\dfrac34$（舍），所以 $k=\dfrac54$.}
\step{(3) $y=\lvert f(x)-g(x)\rvert=\lvert x^2-x-(2x+b)\rvert=\lvert x^2-3x-b\rvert$}
\step{$=\left\lvert\left(x-\dfrac32\right)^2-\dfrac94-b\right\rvert,x\in[0,3]$，}
\step{因为 $x\in[0,3]$，所以 $\left(x-\dfrac32\right)^2-\dfrac94-b\in\left[-\dfrac94-b,-b\right]$，}
\step{由 $y=\left\lvert\left(x-\dfrac32\right)^2-\dfrac94-b\right\rvert$，则 $y_{\max}=\max\left\{\lvert b\rvert,\left\lvert b+\dfrac94\right\rvert\right\}$，}
\step{令 $\lvert b\rvert\geqslant\left\lvert b+\dfrac94\right\rvert$，即 $b^2\geqslant\left(b+\dfrac94\right)^2$，解得 $b\leqslant-\dfrac98$，}
\step{$y_{\max}=\max\left\{\lvert b\rvert,\left\lvert b+\dfrac94\right\rvert\right\}=\begin{cases}-b,b\leqslant-\dfrac98\\b+\dfrac94,b>-\dfrac98\end{cases}$，}
\step{故当且仅当 $b=-\dfrac98$ 时，$(y_{\max})_{\min}=\dfrac98$，}
\step{故当 $b$ 的值为 $-\dfrac98$ 时，函数 $f(x)$ 与 $g(x)$ 的“偏差”取最小值 $\dfrac98$.}
\solend
\fi

\item 若函数 $f(x)$ 与 $g(x)$ 满足：对任意的 $x_1\in D$，总存在唯一的 $x_2\in D$，使 $f(x_1)g(x_2)=m$ 成立，则称 $f(x)$ 是 $g(x)$ 在区间 $D$ 上的“$m$ 阶伴随函数”；对任意的 $x_1\in D$，总存在唯一的 $x_2\in D$，使 $f(x_1)f(x_2)=m$ 成立，则称 $f(x)$ 是区间 $D$ 上的“$m$ 阶自伴函数”.

(1) 判断 $f(x)=x^2+1$ 是否为区间 $[0,3]$ 上的“2 阶自伴函数”？并说明理由：

(2) 若函数 $f(x)=3x-1$ 为区间 $\left[\dfrac12,b\right]$ 上的“1 阶自伴函数”，求 $b$ 的值；

(3) 若 $f(x)=\dfrac4{x+2}$ 是 $g(x)=x^2-2ax+a^2-1$ 在区间 $[0,2]$ 上的“2 阶伴随函数”，求实数 $a$ 的取值范围.

\ifanswer
\solbegin
\step{(1) 假定函数 $f(x)=x^2+1$ 是区间 $[0,3]$ 上的“2 阶自伴函数”，}
\step{取 $x_1=2,f(x_1)=f(2)=5$，由 $f(x_1)f(x_2)=2$，得 $5(x_2^2+1)=2$，}
\step{显然此方程无实数解，}
\step{所以函数 $f(x)=x^2+1$ 不是区间 $[0,3]$ 上的“2 阶自伴函数”.}
\step{(2) 函数 $f(x)=3x-1$ 为区间 $\left[\dfrac12,b\right]$ 上的“1 阶自伴函数”，}
\step{则对任意 $x_1\in\left[\dfrac12,b\right]$，总存在唯一的 $x_2\in\left[\dfrac12,b\right]$，使得 $f(x_1)f(x_2)=1$，}
\step{即 $3x_2-1=\dfrac1{3x_1-1}$，整理得 $x_2=\dfrac13+\dfrac1{3(3x_1-1)}$，}
\step{显然函数 $x_2=\dfrac13+\dfrac1{3(3x_1-1)}$ 在 $\left[\dfrac12,b\right]$ 上严格减，}
\step{且当 $x_1=\dfrac12$ 时，$x_2=1$，当 $x_1=b$ 时，$x_2=\dfrac13+\dfrac1{9b-3}$，}
\step{因此对 $\left[\dfrac12,b\right]$ 内的每一个 $x_1$，在 $\left[\dfrac13+\dfrac1{9b-3},1\right]$ 内有唯一 $x_2$ 值与之对应，}
\step{则 $x_2\in\left[\dfrac12,b\right]$，于是 $\left[\dfrac13+\dfrac1{9b-3},1\right]\subseteq\left[\dfrac12,b\right]$，应有 $\begin{cases}b\geqslant1\\\dfrac13+\dfrac1{9b-3}\geqslant\dfrac12\end{cases}$，}
\step{得 $\begin{cases}b\geqslant1\\b\leqslant1\end{cases}$，即 $b=1$，所以 $b$ 的值是 1.}
\step{(3) 由函数 $f(x)=\dfrac4{x+2}$ 在区间 $[0,2]$ 的值域为 $[1,2]$，}
\step{因为 $f(x)=\dfrac4{x+2}$ 是 $g(x)=x^2-2ax+a^2-1$ 在 $[0,2]$ 上的“2 阶伴随函数”，}
\step{则对任意的 $x_1\in[0,2]$，总存在唯一的 $x_2\in[0,2]$ 时，使得 $f(x_1)g(x_2)=2$，}% 原卷此处印“时”字（按上下文为衍字），照录
\step{所以 $g(x_2)=\dfrac2{f(x_1)}\in[1,2]$，}
\step{即 $g(x)=x^2-2ax+a^2-1$ 在 $[0,2]$ 区间上的值域必定包含区间 $[1,2]$，}
\step{且 $g(x)$ 值域在 $[1,2]$ 对应的自变量是唯一的，}
\step{又因为函数 $g(x)=x^2-2ax+a^2-1$ 开口向上，对称轴为 $x=a$，}
\step{1) 当 $a\leqslant0$ 时，$g(x)$ 在 $[0,2]$ 区间上严格增，则必有 $\begin{cases}g(x)_{\min}=g(0)\leqslant1\\g(x)_{\max}=g(2)\geqslant2\end{cases}$，}
\step{解得 $-\sqrt2\leqslant a\leqslant0$；}
\step{2) 当 $a\geqslant2$ 时，$g(x)$ 在 $[0,2]$ 区间上严格减，则必有 $\begin{cases}g(x)_{\min}=g(2)\leqslant1\\g(x)_{\max}=g(0)\geqslant2\end{cases}$，}
\step{解得 $2\leqslant a\leqslant2+\sqrt2$；}
\step{3) 当 $0<a\leqslant1$ 时，$g(x)$ 在 $[0,a]$ 上严格减，在 $[a,2]$ 上严格增，}
\step{则必有 $\begin{cases}g(0)<1\\g(x)_{\max}=g(2)\geqslant2\end{cases}$，解得 $0<a\leqslant2-\sqrt3$，}
\step{4) 当 $1<a<2$ 时，$g(x)$ 在 $[0,a]$ 上严格减，在 $[a,2]$ 上严格增，}
\step{则必有 $\begin{cases}g(2)<1\\g(x)_{\max}=g(0)\geqslant2\end{cases}$，解得 $\sqrt3\leqslant a<2$.}
\step{综上所述，$a$ 的范围 $[-\sqrt2,2-\sqrt3]\cup[\sqrt3,2+\sqrt2]$.}% 原卷“范围”后缺“是／为”字，照录
\solend
\fi

\end{enumerate}

\end{document}
"
% 紧凑版：xelatex "\def\iscompact{1}% 高一122_华师大二附中_周末作业6.tex
%   —— 高一上数学“2026-2027 年华师大二附中高一上周末作业 6”（试卷版/答案版）
% 试卷版：xelatex 高一122_华师大二附中_周末作业6.tex
% 答案版：xelatex "\def\isanswer{1}\input{高一122_华师大二附中_周末作业6.tex}"
% 紧凑版：xelatex "\def\iscompact{1}\input{高一122_华师大二附中_周末作业6.tex}"
%
% 原始材料：微信公众号“魔都数学加油站”文章，作者破晓，2026-11-09 17:52 发布；
% 连载“27学年高一122”，原文标题“27学年高一122——2026-2027年上海市华师大二附中高一上周末作业6”；
% 原文链接：https://mp.weixin.qq.com/s/pC5066JFKNkQQwjMqTK8Ow。
% 正文为 14 张 PNG 页（第 14 页 2560×3620，其余 13 页 4762×6735），
% 题目下印有【解析】，为讲评版。
% 题目、解析均据原图转录；卷面红色水印及公众号页脚未录入。全卷无插图（已逐页检索，无“如图”）。
%
% 卷首标题：原卷印作“2026-2027 年华师大二附中高一上周末作业 6”（公众号标题中的“上海市”
% 卷面标题行没有，本稿照卷面），年份区间按全稿惯例排 en dash；原卷未印副标题，本稿按题目内容
% 标为“函数”。副标题理由（以全部 14 页实际内容为准）：逐页复核 19 题，考的都是函数性质——
% 奇偶性与充要条件（Q3、Q13、Q15、Q16）、单调性与单调区间（Q4、Q7、Q14、Q17）、定义域／
% 值域／最值与复合函数（Q5、Q6、Q8、Q10、Q12、Q19、Q20）、分段函数与取整函数（Q7、Q8、
% Q11、Q15、Q19）、抽象函数与新定义（Q12、Q20、Q21）；不等式仅作求参数范围／解集的工具
% （Q9、Q14、Q16、Q18），集合与常用逻辑用语仅零星出现（Q3 的充要条件、Q11③ 的集合、
% Q20 的“非空集合 I”），均非独立章节。故主体为单一模块“函数”，与同校 高一102 周末作业 5
% 同（该函数模块卷亦按 高一83／105 成例只用模块名、不叠“集合与常用逻辑用语、不等式”）。
%
% 与原卷的排版差异：
%   ① 原文从第 3 题起，填空题编号为 3—12、选择题为 13—16、解答题为 17—21，本稿保留原题号；
%   ② 原卷每题下直接印【解析】，本稿按全稿惯例置于答案版；四道选择题括号内均未印答案、
%      答案在解析末句“故选 …”，括号内留空（全卷无只印【答案】的题、无 \ans{}）；
%   ③ 原卷未印副标题，本稿按“知识模块”口径补出；
%   ④ 原卷填空横线按全稿惯例排为 \blankline[4.4em]；选择题括号排 \mbox{（\quad）}、
%      选项段前按惯例加 \keepwithprev；解答题不留作答区（原卷为讲评版）；
%   ⑤ 原卷数集排斜体 $R$、$N_+$（下标加号），按全稿惯例改 \R、\N_+
%      （$N_+$ 下标加号与 $\N^*$ 是两种记号，照录下标、只改数集字母为正体，同 高一05 第 20 题成例）；
%      绝对值排 \lvert…\rvert、集合描述竖线排 \mid；
%      $f(x)_{min}$／$f(x)_{max}$／$y_{max}$／$(y_{max})_{min}$ 一类斜体下标按 高一96 成例
%      改正体 \min／\max；不等式一律排 \geqslant／\leqslant；
%   ⑥ 第 20 题解析标签原卷印作【解得】（全卷唯一，其余 18 题均印【解析】），本稿按全稿惯例
%      由 \solbegin 统一印【解析】，正文内容照录（行内加 % 注、并入文件头存疑清单）；
%   ⑦ 题干与解析的标点（半角／全角逗号、冒号、括号、句点、子题括号）均照原卷，
%      题目正文引号用弯引号“ ”。
%
% 子集符号：第 21 题 (2) 的 $\left[\frac13+\frac1{9b-3},1\right]\subseteq\left[\frac12,b\right]$
%   印有下划线，照录 \subseteq。全卷未出现判别式 \triangle／\Delta。
%
% 原卷存疑／笔误（均照录加注）：
%   第 11 题解析对② 的“$y=y+\Delta y$”（按上下文应为 $y=[y]+\Delta y$，原卷漏方括号）；
%   第 19 题题干“某种产品 $x$ 万件）可获利”右括号无配对（疑漏印左括号“（”）；
%   第 20 题解析标签【解得】（见排版差异⑥）；
%   第 21 题 (3) 解析“总存在唯一的 $x_2\in[0,2]$ 时，使得……”的“时”字（按上下文衍字）；
%   第 21 题 (3) 末句“综上所述，$a$ 的范围 $[-\sqrt2,2-\sqrt3]\cup[\sqrt3,2+\sqrt2]$”
%      “范围”后缺“是／为”字，照录。
%
% 行数核对：逐题数 \step 条数，合计 162 条（Q3：5；Q4：3；Q5：1；Q6：3；Q7：2；
% Q8：5；Q9：7；Q10：10；Q11：8；Q12：4；Q13：12；Q14：6；Q15：14；Q16：14；
% Q17：6；Q18：8；Q19：7；Q20：19；Q21：28），与原卷印刷行逐题相等
% （cases 块作一个印刷行计）。全卷无 \ans{}、无只印【答案】的题。
%
\documentclass[a4paper,11pt]{article}
\usepackage{common}

\begin{document}

\examtitlest[3]{2026--2027 年华师大二附中高一上周末作业 6}{函数}

\bigsec{一、填空题}

\begin{enumerate}
\setcounter{enumi}{2}

\item 已知 $p:a=0$，$q$：函数 $f(x)=x^3+x+a$ 是奇函数，则 $p$ 是 $q$ 的\blankline[4.4em]条件.

\ifanswer
\solbegin
\step{若 $a=0$，则 $f(x)=x^3+x$，}
\step{易得 $f(x)$ 的定义域为 $\R$，且 $f(-x)=-(x^3+x)=-f(x)$，则 $f(x)$ 为奇函数，}
\step{反之，若 $f(x)=x^3+x+a$ 是奇函数，}
\step{则有 $f(-x)+f(x)=-x^3-x+a+x^3+x+a=2a=0$，必有 $a=0$，}
\step{故 $p$ 是 $q$ 的充要条件.}
\solend
\fi

\item 函数 $y=\sqrt{-x^2+2x}$ 的单调递增区间是\blankline[4.4em].

\ifanswer
\solbegin
\step{由函数定义域得 $-x^2+2x\geqslant0$，解得 $0\leqslant x\leqslant2$，}
\step{二次函数 $y=-x^2+2x$ 开口向下，对称轴为 $x=1$，}
\step{结合复合函数的单调性得函数 $y=\sqrt{-x^2+2x}$ 的单调递增区间是 $[0,1]$.}
\solend
\fi

\item 设 $f(x)$ 定义域为 $[1,4]$，且 $f(x)=2x+1$，则 $f(x^2)$ 的最大值为\blankline[4.4em].

\ifanswer
\solbegin
\step{由 $x^2\in[1,4]$，得 $f(x^2)=2x^2+1\leqslant9$，当 $x=\pm2$ 时取等号，最大值为 9.}
\solend
\fi

\item 函数 $f(x)=\sqrt{x-2}-x$，$x\in(2,6]$ 的值域为\blankline[4.4em].

\ifanswer
\solbegin
\step{令 $t=\sqrt{x-2}$，则 $x=2+t^2$，由 $2<x\leqslant6$ 得 $0<t\leqslant2$，}
\step{原函数可化为 $g(t)=-t^2+t-2$，$g(0)=-2$，$g(2)=-4$，$g\left(\dfrac12\right)=-\dfrac74$，}
\step{由二次函数性质得 $g(t)\in\left[-4,-\dfrac74\right]$.}
\solend
\fi

\item 若函数 $f(x)=\lvert x-a+1\rvert$ 在区间 $[1,+\infty)$ 上是严格增函数，则实数 $a$ 的取值范围是\blankline[4.4em].

\ifanswer
\solbegin
\step{$f(x)=\begin{cases}x-a+1,x\geqslant a-1\\a-1-x,x<a-1\end{cases}$，显然 $f(x)$ 在 $[a-1,+\infty)$ 上严格增，}
\step{要使函数在区间 $[1,+\infty)$ 上是严格增函数，则 $a-1\leqslant1$，即 $a\leqslant2$.}
\solend
\fi

\item 若函数 $f(x)=\begin{cases}x^2-5,x\leqslant2\\2ax+1,x>2\end{cases}$ 的值域是 $\R$，则 $a$ 的取值范围是\blankline[4.4em].

\ifanswer
\solbegin
\step{当 $x\leqslant2$ 时，$f(x)=x^2-5$，其开口向上，对称轴为 $x=0$，值域为 $[-5,+\infty)$，}
\step{由函数 $f(x)=\begin{cases}x^2-5,x\leqslant2\\2ax+1,x>2\end{cases}$ 的值域是 $\R$，}
\step{则当 $x>2$ 时，$f(x)=2ax+1$ 的值域应包含 $(-\infty,-5)$，}
\step{所以 $f(x)$ 为严格减函数，所以 $\begin{cases}2\times2a+1\geqslant-5\\2a<0\end{cases}$，即 $\begin{cases}a\geqslant-\dfrac32\\a<0\end{cases}$，}
\step{解得 $-\dfrac32\leqslant a<0$，故 $a$ 的取值范围是 $\left[-\dfrac32,0\right)$.}
\solend
\fi

\item 已知函数 $f(x)=ax^2+\dfrac bx$ 的图象经过点 $(-1,0)$，且满足 $f(1)=2$. 若 $\forall x\in(0,+\infty)$，都有 $f(x)+f\left(\dfrac1x\right)-2m\geqslant0$ 恒成立，则实数 $m$ 的取值范围为\blankline[4.4em].

\ifanswer
\solbegin
\step{因为 $f(x)=ax^2+\dfrac bx$ 的图象经过点 $(-1,0)$，且满足 $f(1)=2$，}
\step{所以 $\begin{cases}a-b=0\\a+b=2\end{cases}$，所以 $a=1$，$b=1$，即 $f(x)=x^2+\dfrac1x$，}
\step{$\forall x\in(0,+\infty)$，$f(x)+f\left(\dfrac1x\right)-2m\geqslant0\Leftrightarrow2m\leqslant x^2+\dfrac1x+\dfrac1{x^2}+x$ 恒成立，}
\step{令函数 $g(x)=x^2+\dfrac1x+\dfrac1{x^2}+x=\left(x+\dfrac1x\right)^2+x+\dfrac1x-2$，}
\step{而当 $x>0$ 时，$x+\dfrac1x\geqslant2$，当且仅当 $x=1$ 时取等号，}
\step{因此 $g(x)\geqslant2^2+2-2=4$，当且仅当 $x=1$ 时取等号，则 $2m\leqslant4$，解得 $m\leqslant2$，}
\step{所以实数 $m\in(-\infty,2]$.}
\solend
\fi

\item 若定义在区间 $[a,b](a<b)$ 上的函数 $f(x)=k-\sqrt{x+1}$ 值域也为 $[a,b]$，则实数 $k$ 的取值范围是\blankline[4.4em].

\ifanswer
\solbegin
\step{函数 $f(x)=k-\sqrt{x+1}$ 在定义域 $[-1,+\infty)$ 严格减，}
\step{当 $f(x)$ 的定义域为 $[a,b]$ 时，$f(x)$ 的值域也为 $[a,b]$，$a<b$，}
\step{故 $f(a)=k-\sqrt{a+1}=b$，$f(b)=k-\sqrt{b+1}=a$，}
\step{两式相减得 $\sqrt{a+1}-\sqrt{b+1}=a-b=(a+1)-(b+1)$}
\step{$=(\sqrt{a+1}-\sqrt{b+1})(\sqrt{a+1}+\sqrt{b+1})$，}
\step{所以 $\sqrt{a+1}-\sqrt{b+1}<0$，即 $\sqrt{a+1}+\sqrt{b+1}=1$，}
\step{则 $k=\sqrt{b+1}+a=a+1-\sqrt{a+1}$，}
\step{令 $\lambda=\sqrt{a+1}\geqslant0$，得 $k=\lambda^2-\lambda=\left(\lambda-\dfrac12\right)^2-\dfrac14$，}
\step{又 $\sqrt{a+1}<\sqrt{b+1}$，因为 $\sqrt{a+1}+\sqrt{b+1}=1$，所以 $\sqrt{a+1}\in\left[0,\dfrac12\right)$，}
\step{所以 $0\leqslant\lambda<\dfrac12$，故实数 $k$ 的取值范围为 $\left(-\dfrac14,0\right]$.}
\solend
\fi

\item 函数 $f(x)=[x]$ 的函数值表示不超过 $x$ 的最大整数，例如：$[-3.5]=-4$，$[2.1]=2$，称 $f(x)=[x]$ 为取整函数，则下列关于取整函数的说法：①对任意的 $x\in\R$，$[-x]=-[x]$ ②对任意的 $x$、$y\in\R$，$[x]+[y]\leqslant[x+y]$ ③集合 $\{y\mid y=[nx]-n[x],x\in\R,n\in\N_+\}$ 共有 $n$ 个元素 ④$x>0$ 时，关于 $x$ 的方程 $[2x]=\left[\dfrac1x\right]$ 有无数个解 其中正确的有\blankline[4.4em].

\ifanswer
\solbegin
\step{对①，因为 $[-2.1]=-3\neq[2.1]=2$，所以①错误；}
\step{对②，设 $x=[x]+\Delta x$，$\Delta x\in[0,1)$，$y=y+\Delta y$，$\Delta y\in[0,1)$，}% 原卷印作 $y=y+\Delta y$（按上下文应为 $y=[y]+\Delta y$），照录
\step{$[x+y]=[[x]+\Delta x+[y]+\Delta y]=[x]+[y]+[\Delta x+\Delta y]$，又 $\Delta x+\Delta y\in[0,\ 2)$，}
\step{故 $\Delta x+\Delta y=0$ 或 1，故 $[x]+[y]\leqslant[x+y]$，所以②正确；}
\step{对③，$[nx]-n[x]=[n[x]+n\cdot\Delta x]-n[x]=n[x]+[n\cdot\Delta x]-n[x]=[n\cdot\Delta x]$，}
\step{因为 $\Delta x\in[0,1)$，所以 $0\leqslant n\cdot\Delta x<n$，$n\in\N_+$，所以③正确；}
\step{对④，当 $\dfrac12<x<1$ 时，$[2x]=\left[\dfrac1x\right]=1$，故方程存在无数个解，所以④正确.}
\step{故正确的是②③④.}
\solend
\fi

\item 已知函数 $f(x)$ 定义域为 $[0,1]$，且满足 $f(x)+f(1-x)=1$，$f(x)=2f\left(\dfrac x5\right)$，对任意 $x_1\leqslant x_2$，恒有 $f(x_1)\leqslant f(x_2)$，则 $f\left(\dfrac1{2026}\right)$ 的值为\blankline[4.4em].

\ifanswer
\solbegin
\step{$f(0)=2f(0)\Rightarrow f(0)=0$，从而 $f(1)=1$，$f\left(\dfrac15\right)=\dfrac12$，$f\left(\dfrac45\right)=1-f\left(\dfrac15\right)=\dfrac12$，}
\step{由单调性，$x\in\left[\dfrac15,\dfrac45\right]$ 时，$f(x)=\dfrac12$.}
\step{因为 $\dfrac15<\dfrac{625}{2026}<\dfrac45$，反复使用 $f\left(\dfrac x5\right)=\dfrac12 f(x)$，}
\step{得 $f\left(\dfrac1{2026}\right)=\dfrac1{2^4}f\left(\dfrac{625}{2026}\right)=\dfrac1{32}$.}
\solend
\fi

\end{enumerate}

\bigsec{二、选择题}

\begin{enumerate}
\setcounter{enumi}{12}

\item 已知 $f(x)$ 为 $\R$ 上的奇函数，$g(x)$ 为 $\R$ 上的偶函数，且 $g(x)\neq0$，则下列说法正确的是\mbox{（\quad）}

\keepwithprev
A. $f(x)+g(x)$ 为 $\R$ 上的奇函数\qquad\qquad B. $f(x)-g(x)$ 为 $\R$ 上的奇函数

C. $\dfrac{f(x)}{g(x)}$ 为 $\R$ 上的偶函数\qquad\qquad D. $\lvert f(x)g(x)\rvert$ 为 $\R$ 上的偶函数

\ifanswer
\solbegin
\step{因为 $f(x)$ 为 $\R$ 上的奇函数，$g(x)$ 为 $\R$ 上的偶函数，且 $g(x)\neq0$，}
\step{所以 $f(-x)=-f(x)$，$g(-x)=g(x)$，}
\step{所以 $f(-x)+g(-x)=-f(x)+g(x)\neq-[f(x)+g(x)]$，}
\step{故 $f(x)+g(x)$ 为非奇非偶函数，$A$ 错误；}
\step{同理，$f(x)-g(x)$ 为非奇非偶函数，$B$ 错误；}
\step{设 $F(x)=\dfrac{f(x)}{g(x)}$，则 $F(-x)=\dfrac{f(-x)}{g(-x)}=\dfrac{-f(x)}{g(x)}=-F(x)$，}
\step{所以 $F(x)$ 为奇函数，$C$ 错误；}
\step{设函数 $H(x)=\lvert f(x)g(x)\rvert$，}
\step{因为 $f(x)$ 为 $\R$ 上的奇函数，$g(x)$ 为 $\R$ 上的偶函数，且 $g(x)\neq0$，}
\step{所以 $f(-x)=-f(x)$，$g(-x)=g(x)$，则由函数奇偶性的定义}
\step{得 $H(-x)=\lvert f(-x)g(-x)\rvert=\lvert -f(x)g(x)\rvert=\lvert f(x)g(x)\rvert=H(x)$，$D$ 正确.}
\step{故选 $D$.}
\solend
\fi

\item 若函数 $f(x)$ 是定义域为 $(-\infty,0)\cup(0,+\infty)$ 的奇函数，在区间 $(0,+\infty)$ 上严格增，且 $f(-1)=0$，则不等式 $\dfrac{f(x-1)}{x-2}<0$ 的解集为\mbox{（\quad）}

\keepwithprev
A. $(-\infty,-2)$\qquad\qquad B. $(-2,0)$\qquad\qquad C. $(0,1)$\qquad\qquad D. $(1,2)$

\ifanswer
\solbegin
\step{若函数 $f(x)$ 是定义域为 $(-\infty,0)\cup(0,+\infty)$ 的奇函数，在区间 $(0,+\infty)$ 上严格增，}
\step{由奇函数对称性得，$f(x)$ 在 $(-\infty,0)$ 上严格增，}
\step{因为 $f(-1)=0$，所以 $f(1)=0$，}
\step{则不等式 $\dfrac{f(x-1)}{x-2}<0$ 可转化为 $\begin{cases}x>2\\f(x-1)<0\end{cases}$ 或 $\begin{cases}x<2\\f(x-1)>0\end{cases}$，}
\step{即 $\begin{cases}x>2\\x-1<-1\text{或}0<x-1<1\end{cases}$ 或 $\begin{cases}x<2\\x-1>1\text{或}-1<x-1<0\end{cases}$，解得 $0<x<1$.}
\step{故选 $C$.}
\solend
\fi

\item 若函数 $f(x)$ 同时满足：(1) 对于定义域上的任意 $x$，恒有 $f(x)+f(-x)=0$；(2) 对于定义域上的任意 $x_1$、$x_2$，当 $x_1\neq x_2$ 时，恒有 $x_1f(x_1)-x_1f(x_2)>x_2f(x_1)-x_2f(x_2)$，则称函数 $f(x)$ 为“理想函数”. 给出下列四个函数：① $f(x)=x^2$；② $f(x)=x^3$；③ $f(x)=\dfrac{2x-1}{2x+1}$；④ $f(x)=\begin{cases}x^2+4x,x\geqslant0\\-x^2+4x,x<0\end{cases}$，其中被称为“理想函数”的有\mbox{（\quad）}

\keepwithprev
A. 1 个\qquad\qquad B. 2 个\qquad\qquad C. 3 个\qquad\qquad D. 4 个

\ifanswer
\solbegin
\step{函数 $f(x)$ 同时满足：(1) 对于定义域上的任意 $x$，恒有 $f(x)+f(-x)=0$；}
\step{(2) 对于定义域上的任意 $x_1$、$x_2$，当 $x_1\neq x_2$ 时，}
\step{恒有 $x_1f(x_1)-x_1f(x_2)>x_2f(x_1)-x_2f(x_2)$，}
\step{整理得 $(x_1-x_2)[f(x_1)-f(x_2)]>0$，则称函数 $f(x)$ 为“理想函数”.}
\step{即函数满足 (1) 奇函数；(2) 函数在定义域内为严格增函数；}
\step{① $f(x)=x^2$；不满足函数为奇函数在定义域内也不严格增；}
\step{② $f(x)=x^3$；满足函数既为奇函数，也在定义域内严格增，}
\step{故该函数为理想函数；}
\step{③ $f(x)=\dfrac{2x-1}{2x+1}=\dfrac{2x+1-2}{2x+1}=-\dfrac2{2x+1}+1$，该函数不为奇函数，}
\step{但是函数在定义域内严格增，故该函数不为理想函数；}
\step{④ $f(x)=\begin{cases}x^2+4x,x\geqslant0\\-x^2+4x,x<0\end{cases}$，函数的关系式满足 $f(-x)=-f(x)$，}
\step{且在 $x\geqslant0$ 时，函数严格增，由于奇函数在对称区间上单调性相同，}
\step{故该函数为理想函数.}
\step{故选 $B$.}
\solend
\fi

\item 已知 $f(x)$ 是定义在 $\R$ 上的奇函数，当 $x_1$、$x_2\in(0,+\infty)$ 且 $x_1\neq x_2$ 时，都有\newline $\dfrac{x_2f(x_1)-x_1f(x_2)}{x_1x_2(x_1-x_2)}>0$ 成立，$f(2028)=2028$，则不等式 $f(x)-x>0$ 的解集为\mbox{（\quad）}

\keepwithprev
A. $(-\infty,-2028)\cup(2028,+\infty)$\qquad\qquad B. $\left(-\dfrac1{2028},\dfrac1{2028}\right)$

C. $(-2028,2028)$\qquad\qquad D. $(-2028,0)\cup(2028,+\infty)$

\ifanswer
\solbegin
\step{因为 $f(x)$ 是定义在 $\R$ 上的奇函数，所以 $f(-x)=-f(x)$，}
\step{令 $g(x)=\dfrac{f(x)}x$，则 $g(-x)=\dfrac{f(-x)}{-x}=\dfrac{-f(x)}{-x}=\dfrac{f(x)}x=g(x)$，}
\step{所以 $g(x)$ 为偶函数，}
\step{当 $x_1$、$x_2\in(0,+\infty)$ 且 $x_1\neq x_2$ 时，都有 $\dfrac{x_2f(x_1)-x_1f(x_2)}{x_1x_2(x_1-x_2)}>0$ 成立，}
\step{即 $\dfrac{\dfrac{f(x_1)}{x_1}-\dfrac{f(x_2)}{x_2}}{x_1-x_2}>0$，即 $g(x)$ 在 $(0,+\infty)$ 上严格增，}
\step{由偶函数对称性得 $g(x)$ 在 $(-\infty,0)$ 上严格减，}
\step{因为 $f(2028)=2028$，所以 $g(2028)=1$，}
\step{则不等式 $f(x)-x>0$ 可化为 $f(x)>x$，}
\step{当 $x=0$ 时，$f(0)=0$，显然不成立，}
\step{当 $x>0$ 时，不等式可化为 $\dfrac{f(x)}x>1$，即 $g(x)>1=g(2028)$，}
\step{所以 $\lvert x\rvert>2028$，解得 $x>2028$，}
\step{当 $x<0$ 时，不等式可化为 $\dfrac{f(x)}x<1$，即 $g(x)<1=g(2028)$，}
\step{所以 $\lvert x\rvert<2028$，解得 $-2028<x<0$，}
\step{故不等式的解集为 $(-2028,0)\cup(2028,+\infty)$. 故选 $D$.}
\solend
\fi

\end{enumerate}

\bigsec{三、解答题}

\begin{enumerate}
\setcounter{enumi}{16}

\item 已知函数 $f(x)=\dfrac{2x-3}{x+1}$.

(1) 试用函数单调性的定义证明：函数 $f(x)$ 在 $(-1,+\infty)$ 上是严格增函数；

(2) 求函数 $f(x)$ 在区间 $[1,4]$ 上的最大值和最小值.

\ifanswer
\solbegin
\step{(1) 设 $-1<x_1<x_2$，则 $f(x_1)-f(x_2)=\dfrac{2x_1-3}{x_1+1}-\dfrac{2x_2-3}{x_2+1}=\dfrac{5(x_1-x_2)}{(x_1+1)(x_2+1)}$}
\step{因为 $-1<x_1<x_2$，所以 $x_1-x_2<0$，$x_1+1>0$，$x_2+1>0$，}
\step{所以 $f(x_1)-f(x_2)<0$，所以 $f(x)$ 在 $(-1,+\infty)$ 上严格增.}
\step{(2) 由 (1) 得 $f(x)$ 在区间 $[1,4]$ 上严格增，}
\step{所以 $f(x)_{\min}=f(1)=-\dfrac12$，$f(x)_{\max}=f(4)=1$，}
\step{所以函数 $f(x)$ 在区间 $[1,4]$ 上的最大值为 1，最小值为 $-\dfrac12$.}
\solend
\fi

\item 设函数 $f(x)=x^2-mx+m$.

(1) 当 $m=2$ 时，求不等式 $f(x)<x$ 的解集；

(2) 若 $\exists x\in[1,2]$，使不等式 $f(x)<x+m-2$ 成立，求实数 $m$ 的取值范围.

\ifanswer
\solbegin
\step{(1) $m=2$ 时，函数 $f(x)=x^2-2x+2$，}
\step{不等式 $f(x)<x$ 可化为 $x^2-3x+2<0$，解得 $1<x<2$，}
\step{所以不等式的解集为 $\{x\mid1<x<2\}$；}
\step{(2) 若 $\exists x\in[1,2]$，使不等式 $f(x)<x+m-2$ 成立，即 $x^2-x+2<mx$，}
\step{所以 $m>x+\dfrac2x-1$，设 $g(x)=x+\dfrac2x-1$，$x\in[1,2]$，}
\step{则 $g(x)=x+\dfrac2x-1\geqslant2\sqrt{x\cdot\dfrac2x}-1=2\sqrt2-1$，}
\step{当且仅当 $x=\dfrac2x$，即 $x=\sqrt2$ 时取等号，}
\step{所以实数 $m$ 的取值范围是 $\{m\mid m>2\sqrt2-1\}$.}
\solend
\fi

\item 某企业原来生产某种产品 $x$ 万件）可获利 $w(x)$（万元），且满足\newline% 原卷“万件）”右括号无配对、疑漏印左括号“（”，照录
$w(x)=\begin{cases}20(x^2+17),0<x\leqslant2\\500-\dfrac{80}{x-1},2<x\leqslant5\end{cases}$.
现该企业计划引进新的生产设备和新的产品方案优化产品，优化后的产品的其他成本总投入为 $(20x+10)$ 万元. 由市场调研分析得知，当前产品供不应求. 记该企业优化后的产品的利润为 $f(x)$（单位：万元）.

(1) 求函数 $f(x)$ 的解析式；

(2) 当优化后的产品产量为多少万件时，该企业的利润 $f(x)$ 最大？最大利润是多少？请说明理由.

\ifanswer
\solbegin
\step{(1) 由题意得 $f(x)=\begin{cases}20x^2-20x+330,0<x\leqslant2\\-20x-\dfrac{80}{x-1}+490,2<x\leqslant5\end{cases}$；}
\step{(2) 当 $0<x\leqslant2$ 时，$f(x)=20(x^2-x)+330=20\left(x-\dfrac12\right)^2+325$，}
\step{故当 $x=2$ 时，$f(x)$ 取最大值，$f(2)=370$；}
\step{当 $2<x\leqslant5$ 时，$f(x)=-\left[20(x-1)+\dfrac{80}{x-1}\right]+470\leqslant-80+470=390$，}
\step{当且仅当 $20(x-1)=\dfrac{80}{x-1}$，即 $x=3$ 时，$f(3)=390$ 为最大值；}
\step{因为 $370<390$，}
\step{所以优化后的产品产量为 3 万件时，企业获最大利润 390 万元.}
\solend
\fi

\item 俄国数学家切比雪夫（1821—1894）是研究直线逼近函数的理论先驱. 对定义在非空集合 $I$ 上的函数 $f(x)$，以及函数 $g(x)=kx+b(k,b\in\R)$，切比雪夫将函数 $y=\lvert f(x)-g(x)\rvert,x\in I$ 的最大值称为 $f(x),g(x)$ 的“偏差”.

(1) 函数 $f(x)=x^2,x\in[0,1],g(x)=-x-2$，求 $f(x),g(x)$ 的“偏差”；

(2) 函数 $f(x)=\dfrac1x+1,x\in[1,2],g(x)=kx+1(k>0)$，若 $f(x),g(x)$ 的“偏差”为 2，求 $k$ 的值；

(3) 函数 $f(x)=x^2-x,x\in[0,3],g(x)=2x+b$，当 $f(x),g(x)$ 的“偏差”取最小值时，求 $b$ 的值，并求出“偏差”的最小值.

\ifanswer
% 注：原卷本题解析标签印作【解得】（全卷唯一，其余 18 题均印【解析】）；
%     本稿按全稿惯例由 \solbegin 统一印【解析】，正文内容照录。见文件头排版差异⑥。
\solbegin
\step{(1) $y=\lvert f(x)-g(x)\rvert=\lvert x^2+x+2\rvert=\left\lvert\left(x+\dfrac12\right)^2+\dfrac74\right\rvert,x\in[0,1]$，}
\step{因为 $x\in[0,1]$，所以 $\left(x+\dfrac12\right)^2+\dfrac74\in[2,4]$，}
\step{则 $y=\left\lvert\left(x+\dfrac12\right)^2+\dfrac74\right\rvert\in[2,4]$，$y_{\max}=4$，}
\step{所以函数 $f(x)$ 与 $g(x)$ 的“偏差”为 4.}
\step{(2) 令 $t=f(x)-g(x)=\dfrac1x+1-(kx+1)=\dfrac1x-kx,x\in[1,2]$，}
\step{因为 $k>0$，所以 $t=\dfrac1x-kx$ 是单调减函数，}
\step{所以 $t\in\left[\dfrac12-2k,1-k\right]$，由单调，$y=\lvert t\rvert,t\in\left[\dfrac12-2k,1-k\right]$，且 $y_{\max}=2$.}
\step{当 $\left\lvert\dfrac12-2k\right\rvert\leqslant\lvert1-k\rvert$，即 $0<k\leqslant\dfrac12$ 时，$\lvert1-k\rvert=2$，解得 $k=3$ 或 $k=-1$，}
\step{均不符合；}
\step{当 $\left\lvert\dfrac12-2k\right\rvert>\lvert1-k\rvert$，即 $k>\dfrac12$ 时，$\left\lvert\dfrac12-2k\right\rvert=2$，$\dfrac12-2k=2$ 或 $\dfrac12-2k=-2$，}
\step{解得 $k=\dfrac54$ 或 $k=-\dfrac34$（舍），所以 $k=\dfrac54$.}
\step{(3) $y=\lvert f(x)-g(x)\rvert=\lvert x^2-x-(2x+b)\rvert=\lvert x^2-3x-b\rvert$}
\step{$=\left\lvert\left(x-\dfrac32\right)^2-\dfrac94-b\right\rvert,x\in[0,3]$，}
\step{因为 $x\in[0,3]$，所以 $\left(x-\dfrac32\right)^2-\dfrac94-b\in\left[-\dfrac94-b,-b\right]$，}
\step{由 $y=\left\lvert\left(x-\dfrac32\right)^2-\dfrac94-b\right\rvert$，则 $y_{\max}=\max\left\{\lvert b\rvert,\left\lvert b+\dfrac94\right\rvert\right\}$，}
\step{令 $\lvert b\rvert\geqslant\left\lvert b+\dfrac94\right\rvert$，即 $b^2\geqslant\left(b+\dfrac94\right)^2$，解得 $b\leqslant-\dfrac98$，}
\step{$y_{\max}=\max\left\{\lvert b\rvert,\left\lvert b+\dfrac94\right\rvert\right\}=\begin{cases}-b,b\leqslant-\dfrac98\\b+\dfrac94,b>-\dfrac98\end{cases}$，}
\step{故当且仅当 $b=-\dfrac98$ 时，$(y_{\max})_{\min}=\dfrac98$，}
\step{故当 $b$ 的值为 $-\dfrac98$ 时，函数 $f(x)$ 与 $g(x)$ 的“偏差”取最小值 $\dfrac98$.}
\solend
\fi

\item 若函数 $f(x)$ 与 $g(x)$ 满足：对任意的 $x_1\in D$，总存在唯一的 $x_2\in D$，使 $f(x_1)g(x_2)=m$ 成立，则称 $f(x)$ 是 $g(x)$ 在区间 $D$ 上的“$m$ 阶伴随函数”；对任意的 $x_1\in D$，总存在唯一的 $x_2\in D$，使 $f(x_1)f(x_2)=m$ 成立，则称 $f(x)$ 是区间 $D$ 上的“$m$ 阶自伴函数”.

(1) 判断 $f(x)=x^2+1$ 是否为区间 $[0,3]$ 上的“2 阶自伴函数”？并说明理由：

(2) 若函数 $f(x)=3x-1$ 为区间 $\left[\dfrac12,b\right]$ 上的“1 阶自伴函数”，求 $b$ 的值；

(3) 若 $f(x)=\dfrac4{x+2}$ 是 $g(x)=x^2-2ax+a^2-1$ 在区间 $[0,2]$ 上的“2 阶伴随函数”，求实数 $a$ 的取值范围.

\ifanswer
\solbegin
\step{(1) 假定函数 $f(x)=x^2+1$ 是区间 $[0,3]$ 上的“2 阶自伴函数”，}
\step{取 $x_1=2,f(x_1)=f(2)=5$，由 $f(x_1)f(x_2)=2$，得 $5(x_2^2+1)=2$，}
\step{显然此方程无实数解，}
\step{所以函数 $f(x)=x^2+1$ 不是区间 $[0,3]$ 上的“2 阶自伴函数”.}
\step{(2) 函数 $f(x)=3x-1$ 为区间 $\left[\dfrac12,b\right]$ 上的“1 阶自伴函数”，}
\step{则对任意 $x_1\in\left[\dfrac12,b\right]$，总存在唯一的 $x_2\in\left[\dfrac12,b\right]$，使得 $f(x_1)f(x_2)=1$，}
\step{即 $3x_2-1=\dfrac1{3x_1-1}$，整理得 $x_2=\dfrac13+\dfrac1{3(3x_1-1)}$，}
\step{显然函数 $x_2=\dfrac13+\dfrac1{3(3x_1-1)}$ 在 $\left[\dfrac12,b\right]$ 上严格减，}
\step{且当 $x_1=\dfrac12$ 时，$x_2=1$，当 $x_1=b$ 时，$x_2=\dfrac13+\dfrac1{9b-3}$，}
\step{因此对 $\left[\dfrac12,b\right]$ 内的每一个 $x_1$，在 $\left[\dfrac13+\dfrac1{9b-3},1\right]$ 内有唯一 $x_2$ 值与之对应，}
\step{则 $x_2\in\left[\dfrac12,b\right]$，于是 $\left[\dfrac13+\dfrac1{9b-3},1\right]\subseteq\left[\dfrac12,b\right]$，应有 $\begin{cases}b\geqslant1\\\dfrac13+\dfrac1{9b-3}\geqslant\dfrac12\end{cases}$，}
\step{得 $\begin{cases}b\geqslant1\\b\leqslant1\end{cases}$，即 $b=1$，所以 $b$ 的值是 1.}
\step{(3) 由函数 $f(x)=\dfrac4{x+2}$ 在区间 $[0,2]$ 的值域为 $[1,2]$，}
\step{因为 $f(x)=\dfrac4{x+2}$ 是 $g(x)=x^2-2ax+a^2-1$ 在 $[0,2]$ 上的“2 阶伴随函数”，}
\step{则对任意的 $x_1\in[0,2]$，总存在唯一的 $x_2\in[0,2]$ 时，使得 $f(x_1)g(x_2)=2$，}% 原卷此处印“时”字（按上下文为衍字），照录
\step{所以 $g(x_2)=\dfrac2{f(x_1)}\in[1,2]$，}
\step{即 $g(x)=x^2-2ax+a^2-1$ 在 $[0,2]$ 区间上的值域必定包含区间 $[1,2]$，}
\step{且 $g(x)$ 值域在 $[1,2]$ 对应的自变量是唯一的，}
\step{又因为函数 $g(x)=x^2-2ax+a^2-1$ 开口向上，对称轴为 $x=a$，}
\step{1) 当 $a\leqslant0$ 时，$g(x)$ 在 $[0,2]$ 区间上严格增，则必有 $\begin{cases}g(x)_{\min}=g(0)\leqslant1\\g(x)_{\max}=g(2)\geqslant2\end{cases}$，}
\step{解得 $-\sqrt2\leqslant a\leqslant0$；}
\step{2) 当 $a\geqslant2$ 时，$g(x)$ 在 $[0,2]$ 区间上严格减，则必有 $\begin{cases}g(x)_{\min}=g(2)\leqslant1\\g(x)_{\max}=g(0)\geqslant2\end{cases}$，}
\step{解得 $2\leqslant a\leqslant2+\sqrt2$；}
\step{3) 当 $0<a\leqslant1$ 时，$g(x)$ 在 $[0,a]$ 上严格减，在 $[a,2]$ 上严格增，}
\step{则必有 $\begin{cases}g(0)<1\\g(x)_{\max}=g(2)\geqslant2\end{cases}$，解得 $0<a\leqslant2-\sqrt3$，}
\step{4) 当 $1<a<2$ 时，$g(x)$ 在 $[0,a]$ 上严格减，在 $[a,2]$ 上严格增，}
\step{则必有 $\begin{cases}g(2)<1\\g(x)_{\max}=g(0)\geqslant2\end{cases}$，解得 $\sqrt3\leqslant a<2$.}
\step{综上所述，$a$ 的范围 $[-\sqrt2,2-\sqrt3]\cup[\sqrt3,2+\sqrt2]$.}% 原卷“范围”后缺“是／为”字，照录
\solend
\fi

\end{enumerate}

\end{document}
"
%
% 原始材料：微信公众号“魔都数学加油站”文章，作者破晓，2026-11-09 17:52 发布；
% 连载“27学年高一122”，原文标题“27学年高一122——2026-2027年上海市华师大二附中高一上周末作业6”；
% 原文链接：https://mp.weixin.qq.com/s/pC5066JFKNkQQwjMqTK8Ow。
% 正文为 14 张 PNG 页（第 14 页 2560×3620，其余 13 页 4762×6735），
% 题目下印有【解析】，为讲评版。
% 题目、解析均据原图转录；卷面红色水印及公众号页脚未录入。全卷无插图（已逐页检索，无“如图”）。
%
% 卷首标题：原卷印作“2026-2027 年华师大二附中高一上周末作业 6”（公众号标题中的“上海市”
% 卷面标题行没有，本稿照卷面），年份区间按全稿惯例排 en dash；原卷未印副标题，本稿按题目内容
% 标为“函数”。副标题理由（以全部 14 页实际内容为准）：逐页复核 19 题，考的都是函数性质——
% 奇偶性与充要条件（Q3、Q13、Q15、Q16）、单调性与单调区间（Q4、Q7、Q14、Q17）、定义域／
% 值域／最值与复合函数（Q5、Q6、Q8、Q10、Q12、Q19、Q20）、分段函数与取整函数（Q7、Q8、
% Q11、Q15、Q19）、抽象函数与新定义（Q12、Q20、Q21）；不等式仅作求参数范围／解集的工具
% （Q9、Q14、Q16、Q18），集合与常用逻辑用语仅零星出现（Q3 的充要条件、Q11③ 的集合、
% Q20 的“非空集合 I”），均非独立章节。故主体为单一模块“函数”，与同校 高一102 周末作业 5
% 同（该函数模块卷亦按 高一83／105 成例只用模块名、不叠“集合与常用逻辑用语、不等式”）。
%
% 与原卷的排版差异：
%   ① 原文从第 3 题起，填空题编号为 3—12、选择题为 13—16、解答题为 17—21，本稿保留原题号；
%   ② 原卷每题下直接印【解析】，本稿按全稿惯例置于答案版；四道选择题括号内均未印答案、
%      答案在解析末句“故选 …”，括号内留空（全卷无只印【答案】的题、无 \ans{}）；
%   ③ 原卷未印副标题，本稿按“知识模块”口径补出；
%   ④ 原卷填空横线按全稿惯例排为 \blankline[4.4em]；选择题括号排 \mbox{（\quad）}、
%      选项段前按惯例加 \keepwithprev；解答题不留作答区（原卷为讲评版）；
%   ⑤ 原卷数集排斜体 $R$、$N_+$（下标加号），按全稿惯例改 \R、\N_+
%      （$N_+$ 下标加号与 $\N^*$ 是两种记号，照录下标、只改数集字母为正体，同 高一05 第 20 题成例）；
%      绝对值排 \lvert…\rvert、集合描述竖线排 \mid；
%      $f(x)_{min}$／$f(x)_{max}$／$y_{max}$／$(y_{max})_{min}$ 一类斜体下标按 高一96 成例
%      改正体 \min／\max；不等式一律排 \geqslant／\leqslant；
%   ⑥ 第 20 题解析标签原卷印作【解得】（全卷唯一，其余 18 题均印【解析】），本稿按全稿惯例
%      由 \solbegin 统一印【解析】，正文内容照录（行内加 % 注、并入文件头存疑清单）；
%   ⑦ 题干与解析的标点（半角／全角逗号、冒号、括号、句点、子题括号）均照原卷，
%      题目正文引号用弯引号“ ”。
%
% 子集符号：第 21 题 (2) 的 $\left[\frac13+\frac1{9b-3},1\right]\subseteq\left[\frac12,b\right]$
%   印有下划线，照录 \subseteq。全卷未出现判别式 \triangle／\Delta。
%
% 原卷存疑／笔误（均照录加注）：
%   第 11 题解析对② 的“$y=y+\Delta y$”（按上下文应为 $y=[y]+\Delta y$，原卷漏方括号）；
%   第 19 题题干“某种产品 $x$ 万件）可获利”右括号无配对（疑漏印左括号“（”）；
%   第 20 题解析标签【解得】（见排版差异⑥）；
%   第 21 题 (3) 解析“总存在唯一的 $x_2\in[0,2]$ 时，使得……”的“时”字（按上下文衍字）；
%   第 21 题 (3) 末句“综上所述，$a$ 的范围 $[-\sqrt2,2-\sqrt3]\cup[\sqrt3,2+\sqrt2]$”
%      “范围”后缺“是／为”字，照录。
%
% 行数核对：逐题数 \step 条数，合计 162 条（Q3：5；Q4：3；Q5：1；Q6：3；Q7：2；
% Q8：5；Q9：7；Q10：10；Q11：8；Q12：4；Q13：12；Q14：6；Q15：14；Q16：14；
% Q17：6；Q18：8；Q19：7；Q20：19；Q21：28），与原卷印刷行逐题相等
% （cases 块作一个印刷行计）。全卷无 \ans{}、无只印【答案】的题。
%
\documentclass[a4paper,11pt]{article}
\usepackage{common}

\begin{document}

\examtitlest[3]{2026--2027 年华师大二附中高一上周末作业 6}{函数}

\bigsec{一、填空题}

\begin{enumerate}
\setcounter{enumi}{2}

\item 已知 $p:a=0$，$q$：函数 $f(x)=x^3+x+a$ 是奇函数，则 $p$ 是 $q$ 的\blankline[4.4em]条件.

\ifanswer
\solbegin
\step{若 $a=0$，则 $f(x)=x^3+x$，}
\step{易得 $f(x)$ 的定义域为 $\R$，且 $f(-x)=-(x^3+x)=-f(x)$，则 $f(x)$ 为奇函数，}
\step{反之，若 $f(x)=x^3+x+a$ 是奇函数，}
\step{则有 $f(-x)+f(x)=-x^3-x+a+x^3+x+a=2a=0$，必有 $a=0$，}
\step{故 $p$ 是 $q$ 的充要条件.}
\solend
\fi

\item 函数 $y=\sqrt{-x^2+2x}$ 的单调递增区间是\blankline[4.4em].

\ifanswer
\solbegin
\step{由函数定义域得 $-x^2+2x\geqslant0$，解得 $0\leqslant x\leqslant2$，}
\step{二次函数 $y=-x^2+2x$ 开口向下，对称轴为 $x=1$，}
\step{结合复合函数的单调性得函数 $y=\sqrt{-x^2+2x}$ 的单调递增区间是 $[0,1]$.}
\solend
\fi

\item 设 $f(x)$ 定义域为 $[1,4]$，且 $f(x)=2x+1$，则 $f(x^2)$ 的最大值为\blankline[4.4em].

\ifanswer
\solbegin
\step{由 $x^2\in[1,4]$，得 $f(x^2)=2x^2+1\leqslant9$，当 $x=\pm2$ 时取等号，最大值为 9.}
\solend
\fi

\item 函数 $f(x)=\sqrt{x-2}-x$，$x\in(2,6]$ 的值域为\blankline[4.4em].

\ifanswer
\solbegin
\step{令 $t=\sqrt{x-2}$，则 $x=2+t^2$，由 $2<x\leqslant6$ 得 $0<t\leqslant2$，}
\step{原函数可化为 $g(t)=-t^2+t-2$，$g(0)=-2$，$g(2)=-4$，$g\left(\dfrac12\right)=-\dfrac74$，}
\step{由二次函数性质得 $g(t)\in\left[-4,-\dfrac74\right]$.}
\solend
\fi

\item 若函数 $f(x)=\lvert x-a+1\rvert$ 在区间 $[1,+\infty)$ 上是严格增函数，则实数 $a$ 的取值范围是\blankline[4.4em].

\ifanswer
\solbegin
\step{$f(x)=\begin{cases}x-a+1,x\geqslant a-1\\a-1-x,x<a-1\end{cases}$，显然 $f(x)$ 在 $[a-1,+\infty)$ 上严格增，}
\step{要使函数在区间 $[1,+\infty)$ 上是严格增函数，则 $a-1\leqslant1$，即 $a\leqslant2$.}
\solend
\fi

\item 若函数 $f(x)=\begin{cases}x^2-5,x\leqslant2\\2ax+1,x>2\end{cases}$ 的值域是 $\R$，则 $a$ 的取值范围是\blankline[4.4em].

\ifanswer
\solbegin
\step{当 $x\leqslant2$ 时，$f(x)=x^2-5$，其开口向上，对称轴为 $x=0$，值域为 $[-5,+\infty)$，}
\step{由函数 $f(x)=\begin{cases}x^2-5,x\leqslant2\\2ax+1,x>2\end{cases}$ 的值域是 $\R$，}
\step{则当 $x>2$ 时，$f(x)=2ax+1$ 的值域应包含 $(-\infty,-5)$，}
\step{所以 $f(x)$ 为严格减函数，所以 $\begin{cases}2\times2a+1\geqslant-5\\2a<0\end{cases}$，即 $\begin{cases}a\geqslant-\dfrac32\\a<0\end{cases}$，}
\step{解得 $-\dfrac32\leqslant a<0$，故 $a$ 的取值范围是 $\left[-\dfrac32,0\right)$.}
\solend
\fi

\item 已知函数 $f(x)=ax^2+\dfrac bx$ 的图象经过点 $(-1,0)$，且满足 $f(1)=2$. 若 $\forall x\in(0,+\infty)$，都有 $f(x)+f\left(\dfrac1x\right)-2m\geqslant0$ 恒成立，则实数 $m$ 的取值范围为\blankline[4.4em].

\ifanswer
\solbegin
\step{因为 $f(x)=ax^2+\dfrac bx$ 的图象经过点 $(-1,0)$，且满足 $f(1)=2$，}
\step{所以 $\begin{cases}a-b=0\\a+b=2\end{cases}$，所以 $a=1$，$b=1$，即 $f(x)=x^2+\dfrac1x$，}
\step{$\forall x\in(0,+\infty)$，$f(x)+f\left(\dfrac1x\right)-2m\geqslant0\Leftrightarrow2m\leqslant x^2+\dfrac1x+\dfrac1{x^2}+x$ 恒成立，}
\step{令函数 $g(x)=x^2+\dfrac1x+\dfrac1{x^2}+x=\left(x+\dfrac1x\right)^2+x+\dfrac1x-2$，}
\step{而当 $x>0$ 时，$x+\dfrac1x\geqslant2$，当且仅当 $x=1$ 时取等号，}
\step{因此 $g(x)\geqslant2^2+2-2=4$，当且仅当 $x=1$ 时取等号，则 $2m\leqslant4$，解得 $m\leqslant2$，}
\step{所以实数 $m\in(-\infty,2]$.}
\solend
\fi

\item 若定义在区间 $[a,b](a<b)$ 上的函数 $f(x)=k-\sqrt{x+1}$ 值域也为 $[a,b]$，则实数 $k$ 的取值范围是\blankline[4.4em].

\ifanswer
\solbegin
\step{函数 $f(x)=k-\sqrt{x+1}$ 在定义域 $[-1,+\infty)$ 严格减，}
\step{当 $f(x)$ 的定义域为 $[a,b]$ 时，$f(x)$ 的值域也为 $[a,b]$，$a<b$，}
\step{故 $f(a)=k-\sqrt{a+1}=b$，$f(b)=k-\sqrt{b+1}=a$，}
\step{两式相减得 $\sqrt{a+1}-\sqrt{b+1}=a-b=(a+1)-(b+1)$}
\step{$=(\sqrt{a+1}-\sqrt{b+1})(\sqrt{a+1}+\sqrt{b+1})$，}
\step{所以 $\sqrt{a+1}-\sqrt{b+1}<0$，即 $\sqrt{a+1}+\sqrt{b+1}=1$，}
\step{则 $k=\sqrt{b+1}+a=a+1-\sqrt{a+1}$，}
\step{令 $\lambda=\sqrt{a+1}\geqslant0$，得 $k=\lambda^2-\lambda=\left(\lambda-\dfrac12\right)^2-\dfrac14$，}
\step{又 $\sqrt{a+1}<\sqrt{b+1}$，因为 $\sqrt{a+1}+\sqrt{b+1}=1$，所以 $\sqrt{a+1}\in\left[0,\dfrac12\right)$，}
\step{所以 $0\leqslant\lambda<\dfrac12$，故实数 $k$ 的取值范围为 $\left(-\dfrac14,0\right]$.}
\solend
\fi

\item 函数 $f(x)=[x]$ 的函数值表示不超过 $x$ 的最大整数，例如：$[-3.5]=-4$，$[2.1]=2$，称 $f(x)=[x]$ 为取整函数，则下列关于取整函数的说法：①对任意的 $x\in\R$，$[-x]=-[x]$ ②对任意的 $x$、$y\in\R$，$[x]+[y]\leqslant[x+y]$ ③集合 $\{y\mid y=[nx]-n[x],x\in\R,n\in\N_+\}$ 共有 $n$ 个元素 ④$x>0$ 时，关于 $x$ 的方程 $[2x]=\left[\dfrac1x\right]$ 有无数个解 其中正确的有\blankline[4.4em].

\ifanswer
\solbegin
\step{对①，因为 $[-2.1]=-3\neq[2.1]=2$，所以①错误；}
\step{对②，设 $x=[x]+\Delta x$，$\Delta x\in[0,1)$，$y=y+\Delta y$，$\Delta y\in[0,1)$，}% 原卷印作 $y=y+\Delta y$（按上下文应为 $y=[y]+\Delta y$），照录
\step{$[x+y]=[[x]+\Delta x+[y]+\Delta y]=[x]+[y]+[\Delta x+\Delta y]$，又 $\Delta x+\Delta y\in[0,\ 2)$，}
\step{故 $\Delta x+\Delta y=0$ 或 1，故 $[x]+[y]\leqslant[x+y]$，所以②正确；}
\step{对③，$[nx]-n[x]=[n[x]+n\cdot\Delta x]-n[x]=n[x]+[n\cdot\Delta x]-n[x]=[n\cdot\Delta x]$，}
\step{因为 $\Delta x\in[0,1)$，所以 $0\leqslant n\cdot\Delta x<n$，$n\in\N_+$，所以③正确；}
\step{对④，当 $\dfrac12<x<1$ 时，$[2x]=\left[\dfrac1x\right]=1$，故方程存在无数个解，所以④正确.}
\step{故正确的是②③④.}
\solend
\fi

\item 已知函数 $f(x)$ 定义域为 $[0,1]$，且满足 $f(x)+f(1-x)=1$，$f(x)=2f\left(\dfrac x5\right)$，对任意 $x_1\leqslant x_2$，恒有 $f(x_1)\leqslant f(x_2)$，则 $f\left(\dfrac1{2026}\right)$ 的值为\blankline[4.4em].

\ifanswer
\solbegin
\step{$f(0)=2f(0)\Rightarrow f(0)=0$，从而 $f(1)=1$，$f\left(\dfrac15\right)=\dfrac12$，$f\left(\dfrac45\right)=1-f\left(\dfrac15\right)=\dfrac12$，}
\step{由单调性，$x\in\left[\dfrac15,\dfrac45\right]$ 时，$f(x)=\dfrac12$.}
\step{因为 $\dfrac15<\dfrac{625}{2026}<\dfrac45$，反复使用 $f\left(\dfrac x5\right)=\dfrac12 f(x)$，}
\step{得 $f\left(\dfrac1{2026}\right)=\dfrac1{2^4}f\left(\dfrac{625}{2026}\right)=\dfrac1{32}$.}
\solend
\fi

\end{enumerate}

\bigsec{二、选择题}

\begin{enumerate}
\setcounter{enumi}{12}

\item 已知 $f(x)$ 为 $\R$ 上的奇函数，$g(x)$ 为 $\R$ 上的偶函数，且 $g(x)\neq0$，则下列说法正确的是\mbox{（\quad）}

\keepwithprev
A. $f(x)+g(x)$ 为 $\R$ 上的奇函数\qquad\qquad B. $f(x)-g(x)$ 为 $\R$ 上的奇函数

C. $\dfrac{f(x)}{g(x)}$ 为 $\R$ 上的偶函数\qquad\qquad D. $\lvert f(x)g(x)\rvert$ 为 $\R$ 上的偶函数

\ifanswer
\solbegin
\step{因为 $f(x)$ 为 $\R$ 上的奇函数，$g(x)$ 为 $\R$ 上的偶函数，且 $g(x)\neq0$，}
\step{所以 $f(-x)=-f(x)$，$g(-x)=g(x)$，}
\step{所以 $f(-x)+g(-x)=-f(x)+g(x)\neq-[f(x)+g(x)]$，}
\step{故 $f(x)+g(x)$ 为非奇非偶函数，$A$ 错误；}
\step{同理，$f(x)-g(x)$ 为非奇非偶函数，$B$ 错误；}
\step{设 $F(x)=\dfrac{f(x)}{g(x)}$，则 $F(-x)=\dfrac{f(-x)}{g(-x)}=\dfrac{-f(x)}{g(x)}=-F(x)$，}
\step{所以 $F(x)$ 为奇函数，$C$ 错误；}
\step{设函数 $H(x)=\lvert f(x)g(x)\rvert$，}
\step{因为 $f(x)$ 为 $\R$ 上的奇函数，$g(x)$ 为 $\R$ 上的偶函数，且 $g(x)\neq0$，}
\step{所以 $f(-x)=-f(x)$，$g(-x)=g(x)$，则由函数奇偶性的定义}
\step{得 $H(-x)=\lvert f(-x)g(-x)\rvert=\lvert -f(x)g(x)\rvert=\lvert f(x)g(x)\rvert=H(x)$，$D$ 正确.}
\step{故选 $D$.}
\solend
\fi

\item 若函数 $f(x)$ 是定义域为 $(-\infty,0)\cup(0,+\infty)$ 的奇函数，在区间 $(0,+\infty)$ 上严格增，且 $f(-1)=0$，则不等式 $\dfrac{f(x-1)}{x-2}<0$ 的解集为\mbox{（\quad）}

\keepwithprev
A. $(-\infty,-2)$\qquad\qquad B. $(-2,0)$\qquad\qquad C. $(0,1)$\qquad\qquad D. $(1,2)$

\ifanswer
\solbegin
\step{若函数 $f(x)$ 是定义域为 $(-\infty,0)\cup(0,+\infty)$ 的奇函数，在区间 $(0,+\infty)$ 上严格增，}
\step{由奇函数对称性得，$f(x)$ 在 $(-\infty,0)$ 上严格增，}
\step{因为 $f(-1)=0$，所以 $f(1)=0$，}
\step{则不等式 $\dfrac{f(x-1)}{x-2}<0$ 可转化为 $\begin{cases}x>2\\f(x-1)<0\end{cases}$ 或 $\begin{cases}x<2\\f(x-1)>0\end{cases}$，}
\step{即 $\begin{cases}x>2\\x-1<-1\text{或}0<x-1<1\end{cases}$ 或 $\begin{cases}x<2\\x-1>1\text{或}-1<x-1<0\end{cases}$，解得 $0<x<1$.}
\step{故选 $C$.}
\solend
\fi

\item 若函数 $f(x)$ 同时满足：(1) 对于定义域上的任意 $x$，恒有 $f(x)+f(-x)=0$；(2) 对于定义域上的任意 $x_1$、$x_2$，当 $x_1\neq x_2$ 时，恒有 $x_1f(x_1)-x_1f(x_2)>x_2f(x_1)-x_2f(x_2)$，则称函数 $f(x)$ 为“理想函数”. 给出下列四个函数：① $f(x)=x^2$；② $f(x)=x^3$；③ $f(x)=\dfrac{2x-1}{2x+1}$；④ $f(x)=\begin{cases}x^2+4x,x\geqslant0\\-x^2+4x,x<0\end{cases}$，其中被称为“理想函数”的有\mbox{（\quad）}

\keepwithprev
A. 1 个\qquad\qquad B. 2 个\qquad\qquad C. 3 个\qquad\qquad D. 4 个

\ifanswer
\solbegin
\step{函数 $f(x)$ 同时满足：(1) 对于定义域上的任意 $x$，恒有 $f(x)+f(-x)=0$；}
\step{(2) 对于定义域上的任意 $x_1$、$x_2$，当 $x_1\neq x_2$ 时，}
\step{恒有 $x_1f(x_1)-x_1f(x_2)>x_2f(x_1)-x_2f(x_2)$，}
\step{整理得 $(x_1-x_2)[f(x_1)-f(x_2)]>0$，则称函数 $f(x)$ 为“理想函数”.}
\step{即函数满足 (1) 奇函数；(2) 函数在定义域内为严格增函数；}
\step{① $f(x)=x^2$；不满足函数为奇函数在定义域内也不严格增；}
\step{② $f(x)=x^3$；满足函数既为奇函数，也在定义域内严格增，}
\step{故该函数为理想函数；}
\step{③ $f(x)=\dfrac{2x-1}{2x+1}=\dfrac{2x+1-2}{2x+1}=-\dfrac2{2x+1}+1$，该函数不为奇函数，}
\step{但是函数在定义域内严格增，故该函数不为理想函数；}
\step{④ $f(x)=\begin{cases}x^2+4x,x\geqslant0\\-x^2+4x,x<0\end{cases}$，函数的关系式满足 $f(-x)=-f(x)$，}
\step{且在 $x\geqslant0$ 时，函数严格增，由于奇函数在对称区间上单调性相同，}
\step{故该函数为理想函数.}
\step{故选 $B$.}
\solend
\fi

\item 已知 $f(x)$ 是定义在 $\R$ 上的奇函数，当 $x_1$、$x_2\in(0,+\infty)$ 且 $x_1\neq x_2$ 时，都有\newline $\dfrac{x_2f(x_1)-x_1f(x_2)}{x_1x_2(x_1-x_2)}>0$ 成立，$f(2028)=2028$，则不等式 $f(x)-x>0$ 的解集为\mbox{（\quad）}

\keepwithprev
A. $(-\infty,-2028)\cup(2028,+\infty)$\qquad\qquad B. $\left(-\dfrac1{2028},\dfrac1{2028}\right)$

C. $(-2028,2028)$\qquad\qquad D. $(-2028,0)\cup(2028,+\infty)$

\ifanswer
\solbegin
\step{因为 $f(x)$ 是定义在 $\R$ 上的奇函数，所以 $f(-x)=-f(x)$，}
\step{令 $g(x)=\dfrac{f(x)}x$，则 $g(-x)=\dfrac{f(-x)}{-x}=\dfrac{-f(x)}{-x}=\dfrac{f(x)}x=g(x)$，}
\step{所以 $g(x)$ 为偶函数，}
\step{当 $x_1$、$x_2\in(0,+\infty)$ 且 $x_1\neq x_2$ 时，都有 $\dfrac{x_2f(x_1)-x_1f(x_2)}{x_1x_2(x_1-x_2)}>0$ 成立，}
\step{即 $\dfrac{\dfrac{f(x_1)}{x_1}-\dfrac{f(x_2)}{x_2}}{x_1-x_2}>0$，即 $g(x)$ 在 $(0,+\infty)$ 上严格增，}
\step{由偶函数对称性得 $g(x)$ 在 $(-\infty,0)$ 上严格减，}
\step{因为 $f(2028)=2028$，所以 $g(2028)=1$，}
\step{则不等式 $f(x)-x>0$ 可化为 $f(x)>x$，}
\step{当 $x=0$ 时，$f(0)=0$，显然不成立，}
\step{当 $x>0$ 时，不等式可化为 $\dfrac{f(x)}x>1$，即 $g(x)>1=g(2028)$，}
\step{所以 $\lvert x\rvert>2028$，解得 $x>2028$，}
\step{当 $x<0$ 时，不等式可化为 $\dfrac{f(x)}x<1$，即 $g(x)<1=g(2028)$，}
\step{所以 $\lvert x\rvert<2028$，解得 $-2028<x<0$，}
\step{故不等式的解集为 $(-2028,0)\cup(2028,+\infty)$. 故选 $D$.}
\solend
\fi

\end{enumerate}

\bigsec{三、解答题}

\begin{enumerate}
\setcounter{enumi}{16}

\item 已知函数 $f(x)=\dfrac{2x-3}{x+1}$.

(1) 试用函数单调性的定义证明：函数 $f(x)$ 在 $(-1,+\infty)$ 上是严格增函数；

(2) 求函数 $f(x)$ 在区间 $[1,4]$ 上的最大值和最小值.

\ifanswer
\solbegin
\step{(1) 设 $-1<x_1<x_2$，则 $f(x_1)-f(x_2)=\dfrac{2x_1-3}{x_1+1}-\dfrac{2x_2-3}{x_2+1}=\dfrac{5(x_1-x_2)}{(x_1+1)(x_2+1)}$}
\step{因为 $-1<x_1<x_2$，所以 $x_1-x_2<0$，$x_1+1>0$，$x_2+1>0$，}
\step{所以 $f(x_1)-f(x_2)<0$，所以 $f(x)$ 在 $(-1,+\infty)$ 上严格增.}
\step{(2) 由 (1) 得 $f(x)$ 在区间 $[1,4]$ 上严格增，}
\step{所以 $f(x)_{\min}=f(1)=-\dfrac12$，$f(x)_{\max}=f(4)=1$，}
\step{所以函数 $f(x)$ 在区间 $[1,4]$ 上的最大值为 1，最小值为 $-\dfrac12$.}
\solend
\fi

\item 设函数 $f(x)=x^2-mx+m$.

(1) 当 $m=2$ 时，求不等式 $f(x)<x$ 的解集；

(2) 若 $\exists x\in[1,2]$，使不等式 $f(x)<x+m-2$ 成立，求实数 $m$ 的取值范围.

\ifanswer
\solbegin
\step{(1) $m=2$ 时，函数 $f(x)=x^2-2x+2$，}
\step{不等式 $f(x)<x$ 可化为 $x^2-3x+2<0$，解得 $1<x<2$，}
\step{所以不等式的解集为 $\{x\mid1<x<2\}$；}
\step{(2) 若 $\exists x\in[1,2]$，使不等式 $f(x)<x+m-2$ 成立，即 $x^2-x+2<mx$，}
\step{所以 $m>x+\dfrac2x-1$，设 $g(x)=x+\dfrac2x-1$，$x\in[1,2]$，}
\step{则 $g(x)=x+\dfrac2x-1\geqslant2\sqrt{x\cdot\dfrac2x}-1=2\sqrt2-1$，}
\step{当且仅当 $x=\dfrac2x$，即 $x=\sqrt2$ 时取等号，}
\step{所以实数 $m$ 的取值范围是 $\{m\mid m>2\sqrt2-1\}$.}
\solend
\fi

\item 某企业原来生产某种产品 $x$ 万件）可获利 $w(x)$（万元），且满足\newline% 原卷“万件）”右括号无配对、疑漏印左括号“（”，照录
$w(x)=\begin{cases}20(x^2+17),0<x\leqslant2\\500-\dfrac{80}{x-1},2<x\leqslant5\end{cases}$.
现该企业计划引进新的生产设备和新的产品方案优化产品，优化后的产品的其他成本总投入为 $(20x+10)$ 万元. 由市场调研分析得知，当前产品供不应求. 记该企业优化后的产品的利润为 $f(x)$（单位：万元）.

(1) 求函数 $f(x)$ 的解析式；

(2) 当优化后的产品产量为多少万件时，该企业的利润 $f(x)$ 最大？最大利润是多少？请说明理由.

\ifanswer
\solbegin
\step{(1) 由题意得 $f(x)=\begin{cases}20x^2-20x+330,0<x\leqslant2\\-20x-\dfrac{80}{x-1}+490,2<x\leqslant5\end{cases}$；}
\step{(2) 当 $0<x\leqslant2$ 时，$f(x)=20(x^2-x)+330=20\left(x-\dfrac12\right)^2+325$，}
\step{故当 $x=2$ 时，$f(x)$ 取最大值，$f(2)=370$；}
\step{当 $2<x\leqslant5$ 时，$f(x)=-\left[20(x-1)+\dfrac{80}{x-1}\right]+470\leqslant-80+470=390$，}
\step{当且仅当 $20(x-1)=\dfrac{80}{x-1}$，即 $x=3$ 时，$f(3)=390$ 为最大值；}
\step{因为 $370<390$，}
\step{所以优化后的产品产量为 3 万件时，企业获最大利润 390 万元.}
\solend
\fi

\item 俄国数学家切比雪夫（1821—1894）是研究直线逼近函数的理论先驱. 对定义在非空集合 $I$ 上的函数 $f(x)$，以及函数 $g(x)=kx+b(k,b\in\R)$，切比雪夫将函数 $y=\lvert f(x)-g(x)\rvert,x\in I$ 的最大值称为 $f(x),g(x)$ 的“偏差”.

(1) 函数 $f(x)=x^2,x\in[0,1],g(x)=-x-2$，求 $f(x),g(x)$ 的“偏差”；

(2) 函数 $f(x)=\dfrac1x+1,x\in[1,2],g(x)=kx+1(k>0)$，若 $f(x),g(x)$ 的“偏差”为 2，求 $k$ 的值；

(3) 函数 $f(x)=x^2-x,x\in[0,3],g(x)=2x+b$，当 $f(x),g(x)$ 的“偏差”取最小值时，求 $b$ 的值，并求出“偏差”的最小值.

\ifanswer
% 注：原卷本题解析标签印作【解得】（全卷唯一，其余 18 题均印【解析】）；
%     本稿按全稿惯例由 \solbegin 统一印【解析】，正文内容照录。见文件头排版差异⑥。
\solbegin
\step{(1) $y=\lvert f(x)-g(x)\rvert=\lvert x^2+x+2\rvert=\left\lvert\left(x+\dfrac12\right)^2+\dfrac74\right\rvert,x\in[0,1]$，}
\step{因为 $x\in[0,1]$，所以 $\left(x+\dfrac12\right)^2+\dfrac74\in[2,4]$，}
\step{则 $y=\left\lvert\left(x+\dfrac12\right)^2+\dfrac74\right\rvert\in[2,4]$，$y_{\max}=4$，}
\step{所以函数 $f(x)$ 与 $g(x)$ 的“偏差”为 4.}
\step{(2) 令 $t=f(x)-g(x)=\dfrac1x+1-(kx+1)=\dfrac1x-kx,x\in[1,2]$，}
\step{因为 $k>0$，所以 $t=\dfrac1x-kx$ 是单调减函数，}
\step{所以 $t\in\left[\dfrac12-2k,1-k\right]$，由单调，$y=\lvert t\rvert,t\in\left[\dfrac12-2k,1-k\right]$，且 $y_{\max}=2$.}
\step{当 $\left\lvert\dfrac12-2k\right\rvert\leqslant\lvert1-k\rvert$，即 $0<k\leqslant\dfrac12$ 时，$\lvert1-k\rvert=2$，解得 $k=3$ 或 $k=-1$，}
\step{均不符合；}
\step{当 $\left\lvert\dfrac12-2k\right\rvert>\lvert1-k\rvert$，即 $k>\dfrac12$ 时，$\left\lvert\dfrac12-2k\right\rvert=2$，$\dfrac12-2k=2$ 或 $\dfrac12-2k=-2$，}
\step{解得 $k=\dfrac54$ 或 $k=-\dfrac34$（舍），所以 $k=\dfrac54$.}
\step{(3) $y=\lvert f(x)-g(x)\rvert=\lvert x^2-x-(2x+b)\rvert=\lvert x^2-3x-b\rvert$}
\step{$=\left\lvert\left(x-\dfrac32\right)^2-\dfrac94-b\right\rvert,x\in[0,3]$，}
\step{因为 $x\in[0,3]$，所以 $\left(x-\dfrac32\right)^2-\dfrac94-b\in\left[-\dfrac94-b,-b\right]$，}
\step{由 $y=\left\lvert\left(x-\dfrac32\right)^2-\dfrac94-b\right\rvert$，则 $y_{\max}=\max\left\{\lvert b\rvert,\left\lvert b+\dfrac94\right\rvert\right\}$，}
\step{令 $\lvert b\rvert\geqslant\left\lvert b+\dfrac94\right\rvert$，即 $b^2\geqslant\left(b+\dfrac94\right)^2$，解得 $b\leqslant-\dfrac98$，}
\step{$y_{\max}=\max\left\{\lvert b\rvert,\left\lvert b+\dfrac94\right\rvert\right\}=\begin{cases}-b,b\leqslant-\dfrac98\\b+\dfrac94,b>-\dfrac98\end{cases}$，}
\step{故当且仅当 $b=-\dfrac98$ 时，$(y_{\max})_{\min}=\dfrac98$，}
\step{故当 $b$ 的值为 $-\dfrac98$ 时，函数 $f(x)$ 与 $g(x)$ 的“偏差”取最小值 $\dfrac98$.}
\solend
\fi

\item 若函数 $f(x)$ 与 $g(x)$ 满足：对任意的 $x_1\in D$，总存在唯一的 $x_2\in D$，使 $f(x_1)g(x_2)=m$ 成立，则称 $f(x)$ 是 $g(x)$ 在区间 $D$ 上的“$m$ 阶伴随函数”；对任意的 $x_1\in D$，总存在唯一的 $x_2\in D$，使 $f(x_1)f(x_2)=m$ 成立，则称 $f(x)$ 是区间 $D$ 上的“$m$ 阶自伴函数”.

(1) 判断 $f(x)=x^2+1$ 是否为区间 $[0,3]$ 上的“2 阶自伴函数”？并说明理由：

(2) 若函数 $f(x)=3x-1$ 为区间 $\left[\dfrac12,b\right]$ 上的“1 阶自伴函数”，求 $b$ 的值；

(3) 若 $f(x)=\dfrac4{x+2}$ 是 $g(x)=x^2-2ax+a^2-1$ 在区间 $[0,2]$ 上的“2 阶伴随函数”，求实数 $a$ 的取值范围.

\ifanswer
\solbegin
\step{(1) 假定函数 $f(x)=x^2+1$ 是区间 $[0,3]$ 上的“2 阶自伴函数”，}
\step{取 $x_1=2,f(x_1)=f(2)=5$，由 $f(x_1)f(x_2)=2$，得 $5(x_2^2+1)=2$，}
\step{显然此方程无实数解，}
\step{所以函数 $f(x)=x^2+1$ 不是区间 $[0,3]$ 上的“2 阶自伴函数”.}
\step{(2) 函数 $f(x)=3x-1$ 为区间 $\left[\dfrac12,b\right]$ 上的“1 阶自伴函数”，}
\step{则对任意 $x_1\in\left[\dfrac12,b\right]$，总存在唯一的 $x_2\in\left[\dfrac12,b\right]$，使得 $f(x_1)f(x_2)=1$，}
\step{即 $3x_2-1=\dfrac1{3x_1-1}$，整理得 $x_2=\dfrac13+\dfrac1{3(3x_1-1)}$，}
\step{显然函数 $x_2=\dfrac13+\dfrac1{3(3x_1-1)}$ 在 $\left[\dfrac12,b\right]$ 上严格减，}
\step{且当 $x_1=\dfrac12$ 时，$x_2=1$，当 $x_1=b$ 时，$x_2=\dfrac13+\dfrac1{9b-3}$，}
\step{因此对 $\left[\dfrac12,b\right]$ 内的每一个 $x_1$，在 $\left[\dfrac13+\dfrac1{9b-3},1\right]$ 内有唯一 $x_2$ 值与之对应，}
\step{则 $x_2\in\left[\dfrac12,b\right]$，于是 $\left[\dfrac13+\dfrac1{9b-3},1\right]\subseteq\left[\dfrac12,b\right]$，应有 $\begin{cases}b\geqslant1\\\dfrac13+\dfrac1{9b-3}\geqslant\dfrac12\end{cases}$，}
\step{得 $\begin{cases}b\geqslant1\\b\leqslant1\end{cases}$，即 $b=1$，所以 $b$ 的值是 1.}
\step{(3) 由函数 $f(x)=\dfrac4{x+2}$ 在区间 $[0,2]$ 的值域为 $[1,2]$，}
\step{因为 $f(x)=\dfrac4{x+2}$ 是 $g(x)=x^2-2ax+a^2-1$ 在 $[0,2]$ 上的“2 阶伴随函数”，}
\step{则对任意的 $x_1\in[0,2]$，总存在唯一的 $x_2\in[0,2]$ 时，使得 $f(x_1)g(x_2)=2$，}% 原卷此处印“时”字（按上下文为衍字），照录
\step{所以 $g(x_2)=\dfrac2{f(x_1)}\in[1,2]$，}
\step{即 $g(x)=x^2-2ax+a^2-1$ 在 $[0,2]$ 区间上的值域必定包含区间 $[1,2]$，}
\step{且 $g(x)$ 值域在 $[1,2]$ 对应的自变量是唯一的，}
\step{又因为函数 $g(x)=x^2-2ax+a^2-1$ 开口向上，对称轴为 $x=a$，}
\step{1) 当 $a\leqslant0$ 时，$g(x)$ 在 $[0,2]$ 区间上严格增，则必有 $\begin{cases}g(x)_{\min}=g(0)\leqslant1\\g(x)_{\max}=g(2)\geqslant2\end{cases}$，}
\step{解得 $-\sqrt2\leqslant a\leqslant0$；}
\step{2) 当 $a\geqslant2$ 时，$g(x)$ 在 $[0,2]$ 区间上严格减，则必有 $\begin{cases}g(x)_{\min}=g(2)\leqslant1\\g(x)_{\max}=g(0)\geqslant2\end{cases}$，}
\step{解得 $2\leqslant a\leqslant2+\sqrt2$；}
\step{3) 当 $0<a\leqslant1$ 时，$g(x)$ 在 $[0,a]$ 上严格减，在 $[a,2]$ 上严格增，}
\step{则必有 $\begin{cases}g(0)<1\\g(x)_{\max}=g(2)\geqslant2\end{cases}$，解得 $0<a\leqslant2-\sqrt3$，}
\step{4) 当 $1<a<2$ 时，$g(x)$ 在 $[0,a]$ 上严格减，在 $[a,2]$ 上严格增，}
\step{则必有 $\begin{cases}g(2)<1\\g(x)_{\max}=g(0)\geqslant2\end{cases}$，解得 $\sqrt3\leqslant a<2$.}
\step{综上所述，$a$ 的范围 $[-\sqrt2,2-\sqrt3]\cup[\sqrt3,2+\sqrt2]$.}% 原卷“范围”后缺“是／为”字，照录
\solend
\fi

\end{enumerate}

\end{document}
"
% 紧凑版：xelatex "\def\iscompact{1}% 高一122_华师大二附中_周末作业6.tex
%   —— 高一上数学“2026-2027 年华师大二附中高一上周末作业 6”（试卷版/答案版）
% 试卷版：xelatex 高一122_华师大二附中_周末作业6.tex
% 答案版：xelatex "\def\isanswer{1}% 高一122_华师大二附中_周末作业6.tex
%   —— 高一上数学“2026-2027 年华师大二附中高一上周末作业 6”（试卷版/答案版）
% 试卷版：xelatex 高一122_华师大二附中_周末作业6.tex
% 答案版：xelatex "\def\isanswer{1}\input{高一122_华师大二附中_周末作业6.tex}"
% 紧凑版：xelatex "\def\iscompact{1}\input{高一122_华师大二附中_周末作业6.tex}"
%
% 原始材料：微信公众号“魔都数学加油站”文章，作者破晓，2026-11-09 17:52 发布；
% 连载“27学年高一122”，原文标题“27学年高一122——2026-2027年上海市华师大二附中高一上周末作业6”；
% 原文链接：https://mp.weixin.qq.com/s/pC5066JFKNkQQwjMqTK8Ow。
% 正文为 14 张 PNG 页（第 14 页 2560×3620，其余 13 页 4762×6735），
% 题目下印有【解析】，为讲评版。
% 题目、解析均据原图转录；卷面红色水印及公众号页脚未录入。全卷无插图（已逐页检索，无“如图”）。
%
% 卷首标题：原卷印作“2026-2027 年华师大二附中高一上周末作业 6”（公众号标题中的“上海市”
% 卷面标题行没有，本稿照卷面），年份区间按全稿惯例排 en dash；原卷未印副标题，本稿按题目内容
% 标为“函数”。副标题理由（以全部 14 页实际内容为准）：逐页复核 19 题，考的都是函数性质——
% 奇偶性与充要条件（Q3、Q13、Q15、Q16）、单调性与单调区间（Q4、Q7、Q14、Q17）、定义域／
% 值域／最值与复合函数（Q5、Q6、Q8、Q10、Q12、Q19、Q20）、分段函数与取整函数（Q7、Q8、
% Q11、Q15、Q19）、抽象函数与新定义（Q12、Q20、Q21）；不等式仅作求参数范围／解集的工具
% （Q9、Q14、Q16、Q18），集合与常用逻辑用语仅零星出现（Q3 的充要条件、Q11③ 的集合、
% Q20 的“非空集合 I”），均非独立章节。故主体为单一模块“函数”，与同校 高一102 周末作业 5
% 同（该函数模块卷亦按 高一83／105 成例只用模块名、不叠“集合与常用逻辑用语、不等式”）。
%
% 与原卷的排版差异：
%   ① 原文从第 3 题起，填空题编号为 3—12、选择题为 13—16、解答题为 17—21，本稿保留原题号；
%   ② 原卷每题下直接印【解析】，本稿按全稿惯例置于答案版；四道选择题括号内均未印答案、
%      答案在解析末句“故选 …”，括号内留空（全卷无只印【答案】的题、无 \ans{}）；
%   ③ 原卷未印副标题，本稿按“知识模块”口径补出；
%   ④ 原卷填空横线按全稿惯例排为 \blankline[4.4em]；选择题括号排 \mbox{（\quad）}、
%      选项段前按惯例加 \keepwithprev；解答题不留作答区（原卷为讲评版）；
%   ⑤ 原卷数集排斜体 $R$、$N_+$（下标加号），按全稿惯例改 \R、\N_+
%      （$N_+$ 下标加号与 $\N^*$ 是两种记号，照录下标、只改数集字母为正体，同 高一05 第 20 题成例）；
%      绝对值排 \lvert…\rvert、集合描述竖线排 \mid；
%      $f(x)_{min}$／$f(x)_{max}$／$y_{max}$／$(y_{max})_{min}$ 一类斜体下标按 高一96 成例
%      改正体 \min／\max；不等式一律排 \geqslant／\leqslant；
%   ⑥ 第 20 题解析标签原卷印作【解得】（全卷唯一，其余 18 题均印【解析】），本稿按全稿惯例
%      由 \solbegin 统一印【解析】，正文内容照录（行内加 % 注、并入文件头存疑清单）；
%   ⑦ 题干与解析的标点（半角／全角逗号、冒号、括号、句点、子题括号）均照原卷，
%      题目正文引号用弯引号“ ”。
%
% 子集符号：第 21 题 (2) 的 $\left[\frac13+\frac1{9b-3},1\right]\subseteq\left[\frac12,b\right]$
%   印有下划线，照录 \subseteq。全卷未出现判别式 \triangle／\Delta。
%
% 原卷存疑／笔误（均照录加注）：
%   第 11 题解析对② 的“$y=y+\Delta y$”（按上下文应为 $y=[y]+\Delta y$，原卷漏方括号）；
%   第 19 题题干“某种产品 $x$ 万件）可获利”右括号无配对（疑漏印左括号“（”）；
%   第 20 题解析标签【解得】（见排版差异⑥）；
%   第 21 题 (3) 解析“总存在唯一的 $x_2\in[0,2]$ 时，使得……”的“时”字（按上下文衍字）；
%   第 21 题 (3) 末句“综上所述，$a$ 的范围 $[-\sqrt2,2-\sqrt3]\cup[\sqrt3,2+\sqrt2]$”
%      “范围”后缺“是／为”字，照录。
%
% 行数核对：逐题数 \step 条数，合计 162 条（Q3：5；Q4：3；Q5：1；Q6：3；Q7：2；
% Q8：5；Q9：7；Q10：10；Q11：8；Q12：4；Q13：12；Q14：6；Q15：14；Q16：14；
% Q17：6；Q18：8；Q19：7；Q20：19；Q21：28），与原卷印刷行逐题相等
% （cases 块作一个印刷行计）。全卷无 \ans{}、无只印【答案】的题。
%
\documentclass[a4paper,11pt]{article}
\usepackage{common}

\begin{document}

\examtitlest[3]{2026--2027 年华师大二附中高一上周末作业 6}{函数}

\bigsec{一、填空题}

\begin{enumerate}
\setcounter{enumi}{2}

\item 已知 $p:a=0$，$q$：函数 $f(x)=x^3+x+a$ 是奇函数，则 $p$ 是 $q$ 的\blankline[4.4em]条件.

\ifanswer
\solbegin
\step{若 $a=0$，则 $f(x)=x^3+x$，}
\step{易得 $f(x)$ 的定义域为 $\R$，且 $f(-x)=-(x^3+x)=-f(x)$，则 $f(x)$ 为奇函数，}
\step{反之，若 $f(x)=x^3+x+a$ 是奇函数，}
\step{则有 $f(-x)+f(x)=-x^3-x+a+x^3+x+a=2a=0$，必有 $a=0$，}
\step{故 $p$ 是 $q$ 的充要条件.}
\solend
\fi

\item 函数 $y=\sqrt{-x^2+2x}$ 的单调递增区间是\blankline[4.4em].

\ifanswer
\solbegin
\step{由函数定义域得 $-x^2+2x\geqslant0$，解得 $0\leqslant x\leqslant2$，}
\step{二次函数 $y=-x^2+2x$ 开口向下，对称轴为 $x=1$，}
\step{结合复合函数的单调性得函数 $y=\sqrt{-x^2+2x}$ 的单调递增区间是 $[0,1]$.}
\solend
\fi

\item 设 $f(x)$ 定义域为 $[1,4]$，且 $f(x)=2x+1$，则 $f(x^2)$ 的最大值为\blankline[4.4em].

\ifanswer
\solbegin
\step{由 $x^2\in[1,4]$，得 $f(x^2)=2x^2+1\leqslant9$，当 $x=\pm2$ 时取等号，最大值为 9.}
\solend
\fi

\item 函数 $f(x)=\sqrt{x-2}-x$，$x\in(2,6]$ 的值域为\blankline[4.4em].

\ifanswer
\solbegin
\step{令 $t=\sqrt{x-2}$，则 $x=2+t^2$，由 $2<x\leqslant6$ 得 $0<t\leqslant2$，}
\step{原函数可化为 $g(t)=-t^2+t-2$，$g(0)=-2$，$g(2)=-4$，$g\left(\dfrac12\right)=-\dfrac74$，}
\step{由二次函数性质得 $g(t)\in\left[-4,-\dfrac74\right]$.}
\solend
\fi

\item 若函数 $f(x)=\lvert x-a+1\rvert$ 在区间 $[1,+\infty)$ 上是严格增函数，则实数 $a$ 的取值范围是\blankline[4.4em].

\ifanswer
\solbegin
\step{$f(x)=\begin{cases}x-a+1,x\geqslant a-1\\a-1-x,x<a-1\end{cases}$，显然 $f(x)$ 在 $[a-1,+\infty)$ 上严格增，}
\step{要使函数在区间 $[1,+\infty)$ 上是严格增函数，则 $a-1\leqslant1$，即 $a\leqslant2$.}
\solend
\fi

\item 若函数 $f(x)=\begin{cases}x^2-5,x\leqslant2\\2ax+1,x>2\end{cases}$ 的值域是 $\R$，则 $a$ 的取值范围是\blankline[4.4em].

\ifanswer
\solbegin
\step{当 $x\leqslant2$ 时，$f(x)=x^2-5$，其开口向上，对称轴为 $x=0$，值域为 $[-5,+\infty)$，}
\step{由函数 $f(x)=\begin{cases}x^2-5,x\leqslant2\\2ax+1,x>2\end{cases}$ 的值域是 $\R$，}
\step{则当 $x>2$ 时，$f(x)=2ax+1$ 的值域应包含 $(-\infty,-5)$，}
\step{所以 $f(x)$ 为严格减函数，所以 $\begin{cases}2\times2a+1\geqslant-5\\2a<0\end{cases}$，即 $\begin{cases}a\geqslant-\dfrac32\\a<0\end{cases}$，}
\step{解得 $-\dfrac32\leqslant a<0$，故 $a$ 的取值范围是 $\left[-\dfrac32,0\right)$.}
\solend
\fi

\item 已知函数 $f(x)=ax^2+\dfrac bx$ 的图象经过点 $(-1,0)$，且满足 $f(1)=2$. 若 $\forall x\in(0,+\infty)$，都有 $f(x)+f\left(\dfrac1x\right)-2m\geqslant0$ 恒成立，则实数 $m$ 的取值范围为\blankline[4.4em].

\ifanswer
\solbegin
\step{因为 $f(x)=ax^2+\dfrac bx$ 的图象经过点 $(-1,0)$，且满足 $f(1)=2$，}
\step{所以 $\begin{cases}a-b=0\\a+b=2\end{cases}$，所以 $a=1$，$b=1$，即 $f(x)=x^2+\dfrac1x$，}
\step{$\forall x\in(0,+\infty)$，$f(x)+f\left(\dfrac1x\right)-2m\geqslant0\Leftrightarrow2m\leqslant x^2+\dfrac1x+\dfrac1{x^2}+x$ 恒成立，}
\step{令函数 $g(x)=x^2+\dfrac1x+\dfrac1{x^2}+x=\left(x+\dfrac1x\right)^2+x+\dfrac1x-2$，}
\step{而当 $x>0$ 时，$x+\dfrac1x\geqslant2$，当且仅当 $x=1$ 时取等号，}
\step{因此 $g(x)\geqslant2^2+2-2=4$，当且仅当 $x=1$ 时取等号，则 $2m\leqslant4$，解得 $m\leqslant2$，}
\step{所以实数 $m\in(-\infty,2]$.}
\solend
\fi

\item 若定义在区间 $[a,b](a<b)$ 上的函数 $f(x)=k-\sqrt{x+1}$ 值域也为 $[a,b]$，则实数 $k$ 的取值范围是\blankline[4.4em].

\ifanswer
\solbegin
\step{函数 $f(x)=k-\sqrt{x+1}$ 在定义域 $[-1,+\infty)$ 严格减，}
\step{当 $f(x)$ 的定义域为 $[a,b]$ 时，$f(x)$ 的值域也为 $[a,b]$，$a<b$，}
\step{故 $f(a)=k-\sqrt{a+1}=b$，$f(b)=k-\sqrt{b+1}=a$，}
\step{两式相减得 $\sqrt{a+1}-\sqrt{b+1}=a-b=(a+1)-(b+1)$}
\step{$=(\sqrt{a+1}-\sqrt{b+1})(\sqrt{a+1}+\sqrt{b+1})$，}
\step{所以 $\sqrt{a+1}-\sqrt{b+1}<0$，即 $\sqrt{a+1}+\sqrt{b+1}=1$，}
\step{则 $k=\sqrt{b+1}+a=a+1-\sqrt{a+1}$，}
\step{令 $\lambda=\sqrt{a+1}\geqslant0$，得 $k=\lambda^2-\lambda=\left(\lambda-\dfrac12\right)^2-\dfrac14$，}
\step{又 $\sqrt{a+1}<\sqrt{b+1}$，因为 $\sqrt{a+1}+\sqrt{b+1}=1$，所以 $\sqrt{a+1}\in\left[0,\dfrac12\right)$，}
\step{所以 $0\leqslant\lambda<\dfrac12$，故实数 $k$ 的取值范围为 $\left(-\dfrac14,0\right]$.}
\solend
\fi

\item 函数 $f(x)=[x]$ 的函数值表示不超过 $x$ 的最大整数，例如：$[-3.5]=-4$，$[2.1]=2$，称 $f(x)=[x]$ 为取整函数，则下列关于取整函数的说法：①对任意的 $x\in\R$，$[-x]=-[x]$ ②对任意的 $x$、$y\in\R$，$[x]+[y]\leqslant[x+y]$ ③集合 $\{y\mid y=[nx]-n[x],x\in\R,n\in\N_+\}$ 共有 $n$ 个元素 ④$x>0$ 时，关于 $x$ 的方程 $[2x]=\left[\dfrac1x\right]$ 有无数个解 其中正确的有\blankline[4.4em].

\ifanswer
\solbegin
\step{对①，因为 $[-2.1]=-3\neq[2.1]=2$，所以①错误；}
\step{对②，设 $x=[x]+\Delta x$，$\Delta x\in[0,1)$，$y=y+\Delta y$，$\Delta y\in[0,1)$，}% 原卷印作 $y=y+\Delta y$（按上下文应为 $y=[y]+\Delta y$），照录
\step{$[x+y]=[[x]+\Delta x+[y]+\Delta y]=[x]+[y]+[\Delta x+\Delta y]$，又 $\Delta x+\Delta y\in[0,\ 2)$，}
\step{故 $\Delta x+\Delta y=0$ 或 1，故 $[x]+[y]\leqslant[x+y]$，所以②正确；}
\step{对③，$[nx]-n[x]=[n[x]+n\cdot\Delta x]-n[x]=n[x]+[n\cdot\Delta x]-n[x]=[n\cdot\Delta x]$，}
\step{因为 $\Delta x\in[0,1)$，所以 $0\leqslant n\cdot\Delta x<n$，$n\in\N_+$，所以③正确；}
\step{对④，当 $\dfrac12<x<1$ 时，$[2x]=\left[\dfrac1x\right]=1$，故方程存在无数个解，所以④正确.}
\step{故正确的是②③④.}
\solend
\fi

\item 已知函数 $f(x)$ 定义域为 $[0,1]$，且满足 $f(x)+f(1-x)=1$，$f(x)=2f\left(\dfrac x5\right)$，对任意 $x_1\leqslant x_2$，恒有 $f(x_1)\leqslant f(x_2)$，则 $f\left(\dfrac1{2026}\right)$ 的值为\blankline[4.4em].

\ifanswer
\solbegin
\step{$f(0)=2f(0)\Rightarrow f(0)=0$，从而 $f(1)=1$，$f\left(\dfrac15\right)=\dfrac12$，$f\left(\dfrac45\right)=1-f\left(\dfrac15\right)=\dfrac12$，}
\step{由单调性，$x\in\left[\dfrac15,\dfrac45\right]$ 时，$f(x)=\dfrac12$.}
\step{因为 $\dfrac15<\dfrac{625}{2026}<\dfrac45$，反复使用 $f\left(\dfrac x5\right)=\dfrac12 f(x)$，}
\step{得 $f\left(\dfrac1{2026}\right)=\dfrac1{2^4}f\left(\dfrac{625}{2026}\right)=\dfrac1{32}$.}
\solend
\fi

\end{enumerate}

\bigsec{二、选择题}

\begin{enumerate}
\setcounter{enumi}{12}

\item 已知 $f(x)$ 为 $\R$ 上的奇函数，$g(x)$ 为 $\R$ 上的偶函数，且 $g(x)\neq0$，则下列说法正确的是\mbox{（\quad）}

\keepwithprev
A. $f(x)+g(x)$ 为 $\R$ 上的奇函数\qquad\qquad B. $f(x)-g(x)$ 为 $\R$ 上的奇函数

C. $\dfrac{f(x)}{g(x)}$ 为 $\R$ 上的偶函数\qquad\qquad D. $\lvert f(x)g(x)\rvert$ 为 $\R$ 上的偶函数

\ifanswer
\solbegin
\step{因为 $f(x)$ 为 $\R$ 上的奇函数，$g(x)$ 为 $\R$ 上的偶函数，且 $g(x)\neq0$，}
\step{所以 $f(-x)=-f(x)$，$g(-x)=g(x)$，}
\step{所以 $f(-x)+g(-x)=-f(x)+g(x)\neq-[f(x)+g(x)]$，}
\step{故 $f(x)+g(x)$ 为非奇非偶函数，$A$ 错误；}
\step{同理，$f(x)-g(x)$ 为非奇非偶函数，$B$ 错误；}
\step{设 $F(x)=\dfrac{f(x)}{g(x)}$，则 $F(-x)=\dfrac{f(-x)}{g(-x)}=\dfrac{-f(x)}{g(x)}=-F(x)$，}
\step{所以 $F(x)$ 为奇函数，$C$ 错误；}
\step{设函数 $H(x)=\lvert f(x)g(x)\rvert$，}
\step{因为 $f(x)$ 为 $\R$ 上的奇函数，$g(x)$ 为 $\R$ 上的偶函数，且 $g(x)\neq0$，}
\step{所以 $f(-x)=-f(x)$，$g(-x)=g(x)$，则由函数奇偶性的定义}
\step{得 $H(-x)=\lvert f(-x)g(-x)\rvert=\lvert -f(x)g(x)\rvert=\lvert f(x)g(x)\rvert=H(x)$，$D$ 正确.}
\step{故选 $D$.}
\solend
\fi

\item 若函数 $f(x)$ 是定义域为 $(-\infty,0)\cup(0,+\infty)$ 的奇函数，在区间 $(0,+\infty)$ 上严格增，且 $f(-1)=0$，则不等式 $\dfrac{f(x-1)}{x-2}<0$ 的解集为\mbox{（\quad）}

\keepwithprev
A. $(-\infty,-2)$\qquad\qquad B. $(-2,0)$\qquad\qquad C. $(0,1)$\qquad\qquad D. $(1,2)$

\ifanswer
\solbegin
\step{若函数 $f(x)$ 是定义域为 $(-\infty,0)\cup(0,+\infty)$ 的奇函数，在区间 $(0,+\infty)$ 上严格增，}
\step{由奇函数对称性得，$f(x)$ 在 $(-\infty,0)$ 上严格增，}
\step{因为 $f(-1)=0$，所以 $f(1)=0$，}
\step{则不等式 $\dfrac{f(x-1)}{x-2}<0$ 可转化为 $\begin{cases}x>2\\f(x-1)<0\end{cases}$ 或 $\begin{cases}x<2\\f(x-1)>0\end{cases}$，}
\step{即 $\begin{cases}x>2\\x-1<-1\text{或}0<x-1<1\end{cases}$ 或 $\begin{cases}x<2\\x-1>1\text{或}-1<x-1<0\end{cases}$，解得 $0<x<1$.}
\step{故选 $C$.}
\solend
\fi

\item 若函数 $f(x)$ 同时满足：(1) 对于定义域上的任意 $x$，恒有 $f(x)+f(-x)=0$；(2) 对于定义域上的任意 $x_1$、$x_2$，当 $x_1\neq x_2$ 时，恒有 $x_1f(x_1)-x_1f(x_2)>x_2f(x_1)-x_2f(x_2)$，则称函数 $f(x)$ 为“理想函数”. 给出下列四个函数：① $f(x)=x^2$；② $f(x)=x^3$；③ $f(x)=\dfrac{2x-1}{2x+1}$；④ $f(x)=\begin{cases}x^2+4x,x\geqslant0\\-x^2+4x,x<0\end{cases}$，其中被称为“理想函数”的有\mbox{（\quad）}

\keepwithprev
A. 1 个\qquad\qquad B. 2 个\qquad\qquad C. 3 个\qquad\qquad D. 4 个

\ifanswer
\solbegin
\step{函数 $f(x)$ 同时满足：(1) 对于定义域上的任意 $x$，恒有 $f(x)+f(-x)=0$；}
\step{(2) 对于定义域上的任意 $x_1$、$x_2$，当 $x_1\neq x_2$ 时，}
\step{恒有 $x_1f(x_1)-x_1f(x_2)>x_2f(x_1)-x_2f(x_2)$，}
\step{整理得 $(x_1-x_2)[f(x_1)-f(x_2)]>0$，则称函数 $f(x)$ 为“理想函数”.}
\step{即函数满足 (1) 奇函数；(2) 函数在定义域内为严格增函数；}
\step{① $f(x)=x^2$；不满足函数为奇函数在定义域内也不严格增；}
\step{② $f(x)=x^3$；满足函数既为奇函数，也在定义域内严格增，}
\step{故该函数为理想函数；}
\step{③ $f(x)=\dfrac{2x-1}{2x+1}=\dfrac{2x+1-2}{2x+1}=-\dfrac2{2x+1}+1$，该函数不为奇函数，}
\step{但是函数在定义域内严格增，故该函数不为理想函数；}
\step{④ $f(x)=\begin{cases}x^2+4x,x\geqslant0\\-x^2+4x,x<0\end{cases}$，函数的关系式满足 $f(-x)=-f(x)$，}
\step{且在 $x\geqslant0$ 时，函数严格增，由于奇函数在对称区间上单调性相同，}
\step{故该函数为理想函数.}
\step{故选 $B$.}
\solend
\fi

\item 已知 $f(x)$ 是定义在 $\R$ 上的奇函数，当 $x_1$、$x_2\in(0,+\infty)$ 且 $x_1\neq x_2$ 时，都有\newline $\dfrac{x_2f(x_1)-x_1f(x_2)}{x_1x_2(x_1-x_2)}>0$ 成立，$f(2028)=2028$，则不等式 $f(x)-x>0$ 的解集为\mbox{（\quad）}

\keepwithprev
A. $(-\infty,-2028)\cup(2028,+\infty)$\qquad\qquad B. $\left(-\dfrac1{2028},\dfrac1{2028}\right)$

C. $(-2028,2028)$\qquad\qquad D. $(-2028,0)\cup(2028,+\infty)$

\ifanswer
\solbegin
\step{因为 $f(x)$ 是定义在 $\R$ 上的奇函数，所以 $f(-x)=-f(x)$，}
\step{令 $g(x)=\dfrac{f(x)}x$，则 $g(-x)=\dfrac{f(-x)}{-x}=\dfrac{-f(x)}{-x}=\dfrac{f(x)}x=g(x)$，}
\step{所以 $g(x)$ 为偶函数，}
\step{当 $x_1$、$x_2\in(0,+\infty)$ 且 $x_1\neq x_2$ 时，都有 $\dfrac{x_2f(x_1)-x_1f(x_2)}{x_1x_2(x_1-x_2)}>0$ 成立，}
\step{即 $\dfrac{\dfrac{f(x_1)}{x_1}-\dfrac{f(x_2)}{x_2}}{x_1-x_2}>0$，即 $g(x)$ 在 $(0,+\infty)$ 上严格增，}
\step{由偶函数对称性得 $g(x)$ 在 $(-\infty,0)$ 上严格减，}
\step{因为 $f(2028)=2028$，所以 $g(2028)=1$，}
\step{则不等式 $f(x)-x>0$ 可化为 $f(x)>x$，}
\step{当 $x=0$ 时，$f(0)=0$，显然不成立，}
\step{当 $x>0$ 时，不等式可化为 $\dfrac{f(x)}x>1$，即 $g(x)>1=g(2028)$，}
\step{所以 $\lvert x\rvert>2028$，解得 $x>2028$，}
\step{当 $x<0$ 时，不等式可化为 $\dfrac{f(x)}x<1$，即 $g(x)<1=g(2028)$，}
\step{所以 $\lvert x\rvert<2028$，解得 $-2028<x<0$，}
\step{故不等式的解集为 $(-2028,0)\cup(2028,+\infty)$. 故选 $D$.}
\solend
\fi

\end{enumerate}

\bigsec{三、解答题}

\begin{enumerate}
\setcounter{enumi}{16}

\item 已知函数 $f(x)=\dfrac{2x-3}{x+1}$.

(1) 试用函数单调性的定义证明：函数 $f(x)$ 在 $(-1,+\infty)$ 上是严格增函数；

(2) 求函数 $f(x)$ 在区间 $[1,4]$ 上的最大值和最小值.

\ifanswer
\solbegin
\step{(1) 设 $-1<x_1<x_2$，则 $f(x_1)-f(x_2)=\dfrac{2x_1-3}{x_1+1}-\dfrac{2x_2-3}{x_2+1}=\dfrac{5(x_1-x_2)}{(x_1+1)(x_2+1)}$}
\step{因为 $-1<x_1<x_2$，所以 $x_1-x_2<0$，$x_1+1>0$，$x_2+1>0$，}
\step{所以 $f(x_1)-f(x_2)<0$，所以 $f(x)$ 在 $(-1,+\infty)$ 上严格增.}
\step{(2) 由 (1) 得 $f(x)$ 在区间 $[1,4]$ 上严格增，}
\step{所以 $f(x)_{\min}=f(1)=-\dfrac12$，$f(x)_{\max}=f(4)=1$，}
\step{所以函数 $f(x)$ 在区间 $[1,4]$ 上的最大值为 1，最小值为 $-\dfrac12$.}
\solend
\fi

\item 设函数 $f(x)=x^2-mx+m$.

(1) 当 $m=2$ 时，求不等式 $f(x)<x$ 的解集；

(2) 若 $\exists x\in[1,2]$，使不等式 $f(x)<x+m-2$ 成立，求实数 $m$ 的取值范围.

\ifanswer
\solbegin
\step{(1) $m=2$ 时，函数 $f(x)=x^2-2x+2$，}
\step{不等式 $f(x)<x$ 可化为 $x^2-3x+2<0$，解得 $1<x<2$，}
\step{所以不等式的解集为 $\{x\mid1<x<2\}$；}
\step{(2) 若 $\exists x\in[1,2]$，使不等式 $f(x)<x+m-2$ 成立，即 $x^2-x+2<mx$，}
\step{所以 $m>x+\dfrac2x-1$，设 $g(x)=x+\dfrac2x-1$，$x\in[1,2]$，}
\step{则 $g(x)=x+\dfrac2x-1\geqslant2\sqrt{x\cdot\dfrac2x}-1=2\sqrt2-1$，}
\step{当且仅当 $x=\dfrac2x$，即 $x=\sqrt2$ 时取等号，}
\step{所以实数 $m$ 的取值范围是 $\{m\mid m>2\sqrt2-1\}$.}
\solend
\fi

\item 某企业原来生产某种产品 $x$ 万件）可获利 $w(x)$（万元），且满足\newline% 原卷“万件）”右括号无配对、疑漏印左括号“（”，照录
$w(x)=\begin{cases}20(x^2+17),0<x\leqslant2\\500-\dfrac{80}{x-1},2<x\leqslant5\end{cases}$.
现该企业计划引进新的生产设备和新的产品方案优化产品，优化后的产品的其他成本总投入为 $(20x+10)$ 万元. 由市场调研分析得知，当前产品供不应求. 记该企业优化后的产品的利润为 $f(x)$（单位：万元）.

(1) 求函数 $f(x)$ 的解析式；

(2) 当优化后的产品产量为多少万件时，该企业的利润 $f(x)$ 最大？最大利润是多少？请说明理由.

\ifanswer
\solbegin
\step{(1) 由题意得 $f(x)=\begin{cases}20x^2-20x+330,0<x\leqslant2\\-20x-\dfrac{80}{x-1}+490,2<x\leqslant5\end{cases}$；}
\step{(2) 当 $0<x\leqslant2$ 时，$f(x)=20(x^2-x)+330=20\left(x-\dfrac12\right)^2+325$，}
\step{故当 $x=2$ 时，$f(x)$ 取最大值，$f(2)=370$；}
\step{当 $2<x\leqslant5$ 时，$f(x)=-\left[20(x-1)+\dfrac{80}{x-1}\right]+470\leqslant-80+470=390$，}
\step{当且仅当 $20(x-1)=\dfrac{80}{x-1}$，即 $x=3$ 时，$f(3)=390$ 为最大值；}
\step{因为 $370<390$，}
\step{所以优化后的产品产量为 3 万件时，企业获最大利润 390 万元.}
\solend
\fi

\item 俄国数学家切比雪夫（1821—1894）是研究直线逼近函数的理论先驱. 对定义在非空集合 $I$ 上的函数 $f(x)$，以及函数 $g(x)=kx+b(k,b\in\R)$，切比雪夫将函数 $y=\lvert f(x)-g(x)\rvert,x\in I$ 的最大值称为 $f(x),g(x)$ 的“偏差”.

(1) 函数 $f(x)=x^2,x\in[0,1],g(x)=-x-2$，求 $f(x),g(x)$ 的“偏差”；

(2) 函数 $f(x)=\dfrac1x+1,x\in[1,2],g(x)=kx+1(k>0)$，若 $f(x),g(x)$ 的“偏差”为 2，求 $k$ 的值；

(3) 函数 $f(x)=x^2-x,x\in[0,3],g(x)=2x+b$，当 $f(x),g(x)$ 的“偏差”取最小值时，求 $b$ 的值，并求出“偏差”的最小值.

\ifanswer
% 注：原卷本题解析标签印作【解得】（全卷唯一，其余 18 题均印【解析】）；
%     本稿按全稿惯例由 \solbegin 统一印【解析】，正文内容照录。见文件头排版差异⑥。
\solbegin
\step{(1) $y=\lvert f(x)-g(x)\rvert=\lvert x^2+x+2\rvert=\left\lvert\left(x+\dfrac12\right)^2+\dfrac74\right\rvert,x\in[0,1]$，}
\step{因为 $x\in[0,1]$，所以 $\left(x+\dfrac12\right)^2+\dfrac74\in[2,4]$，}
\step{则 $y=\left\lvert\left(x+\dfrac12\right)^2+\dfrac74\right\rvert\in[2,4]$，$y_{\max}=4$，}
\step{所以函数 $f(x)$ 与 $g(x)$ 的“偏差”为 4.}
\step{(2) 令 $t=f(x)-g(x)=\dfrac1x+1-(kx+1)=\dfrac1x-kx,x\in[1,2]$，}
\step{因为 $k>0$，所以 $t=\dfrac1x-kx$ 是单调减函数，}
\step{所以 $t\in\left[\dfrac12-2k,1-k\right]$，由单调，$y=\lvert t\rvert,t\in\left[\dfrac12-2k,1-k\right]$，且 $y_{\max}=2$.}
\step{当 $\left\lvert\dfrac12-2k\right\rvert\leqslant\lvert1-k\rvert$，即 $0<k\leqslant\dfrac12$ 时，$\lvert1-k\rvert=2$，解得 $k=3$ 或 $k=-1$，}
\step{均不符合；}
\step{当 $\left\lvert\dfrac12-2k\right\rvert>\lvert1-k\rvert$，即 $k>\dfrac12$ 时，$\left\lvert\dfrac12-2k\right\rvert=2$，$\dfrac12-2k=2$ 或 $\dfrac12-2k=-2$，}
\step{解得 $k=\dfrac54$ 或 $k=-\dfrac34$（舍），所以 $k=\dfrac54$.}
\step{(3) $y=\lvert f(x)-g(x)\rvert=\lvert x^2-x-(2x+b)\rvert=\lvert x^2-3x-b\rvert$}
\step{$=\left\lvert\left(x-\dfrac32\right)^2-\dfrac94-b\right\rvert,x\in[0,3]$，}
\step{因为 $x\in[0,3]$，所以 $\left(x-\dfrac32\right)^2-\dfrac94-b\in\left[-\dfrac94-b,-b\right]$，}
\step{由 $y=\left\lvert\left(x-\dfrac32\right)^2-\dfrac94-b\right\rvert$，则 $y_{\max}=\max\left\{\lvert b\rvert,\left\lvert b+\dfrac94\right\rvert\right\}$，}
\step{令 $\lvert b\rvert\geqslant\left\lvert b+\dfrac94\right\rvert$，即 $b^2\geqslant\left(b+\dfrac94\right)^2$，解得 $b\leqslant-\dfrac98$，}
\step{$y_{\max}=\max\left\{\lvert b\rvert,\left\lvert b+\dfrac94\right\rvert\right\}=\begin{cases}-b,b\leqslant-\dfrac98\\b+\dfrac94,b>-\dfrac98\end{cases}$，}
\step{故当且仅当 $b=-\dfrac98$ 时，$(y_{\max})_{\min}=\dfrac98$，}
\step{故当 $b$ 的值为 $-\dfrac98$ 时，函数 $f(x)$ 与 $g(x)$ 的“偏差”取最小值 $\dfrac98$.}
\solend
\fi

\item 若函数 $f(x)$ 与 $g(x)$ 满足：对任意的 $x_1\in D$，总存在唯一的 $x_2\in D$，使 $f(x_1)g(x_2)=m$ 成立，则称 $f(x)$ 是 $g(x)$ 在区间 $D$ 上的“$m$ 阶伴随函数”；对任意的 $x_1\in D$，总存在唯一的 $x_2\in D$，使 $f(x_1)f(x_2)=m$ 成立，则称 $f(x)$ 是区间 $D$ 上的“$m$ 阶自伴函数”.

(1) 判断 $f(x)=x^2+1$ 是否为区间 $[0,3]$ 上的“2 阶自伴函数”？并说明理由：

(2) 若函数 $f(x)=3x-1$ 为区间 $\left[\dfrac12,b\right]$ 上的“1 阶自伴函数”，求 $b$ 的值；

(3) 若 $f(x)=\dfrac4{x+2}$ 是 $g(x)=x^2-2ax+a^2-1$ 在区间 $[0,2]$ 上的“2 阶伴随函数”，求实数 $a$ 的取值范围.

\ifanswer
\solbegin
\step{(1) 假定函数 $f(x)=x^2+1$ 是区间 $[0,3]$ 上的“2 阶自伴函数”，}
\step{取 $x_1=2,f(x_1)=f(2)=5$，由 $f(x_1)f(x_2)=2$，得 $5(x_2^2+1)=2$，}
\step{显然此方程无实数解，}
\step{所以函数 $f(x)=x^2+1$ 不是区间 $[0,3]$ 上的“2 阶自伴函数”.}
\step{(2) 函数 $f(x)=3x-1$ 为区间 $\left[\dfrac12,b\right]$ 上的“1 阶自伴函数”，}
\step{则对任意 $x_1\in\left[\dfrac12,b\right]$，总存在唯一的 $x_2\in\left[\dfrac12,b\right]$，使得 $f(x_1)f(x_2)=1$，}
\step{即 $3x_2-1=\dfrac1{3x_1-1}$，整理得 $x_2=\dfrac13+\dfrac1{3(3x_1-1)}$，}
\step{显然函数 $x_2=\dfrac13+\dfrac1{3(3x_1-1)}$ 在 $\left[\dfrac12,b\right]$ 上严格减，}
\step{且当 $x_1=\dfrac12$ 时，$x_2=1$，当 $x_1=b$ 时，$x_2=\dfrac13+\dfrac1{9b-3}$，}
\step{因此对 $\left[\dfrac12,b\right]$ 内的每一个 $x_1$，在 $\left[\dfrac13+\dfrac1{9b-3},1\right]$ 内有唯一 $x_2$ 值与之对应，}
\step{则 $x_2\in\left[\dfrac12,b\right]$，于是 $\left[\dfrac13+\dfrac1{9b-3},1\right]\subseteq\left[\dfrac12,b\right]$，应有 $\begin{cases}b\geqslant1\\\dfrac13+\dfrac1{9b-3}\geqslant\dfrac12\end{cases}$，}
\step{得 $\begin{cases}b\geqslant1\\b\leqslant1\end{cases}$，即 $b=1$，所以 $b$ 的值是 1.}
\step{(3) 由函数 $f(x)=\dfrac4{x+2}$ 在区间 $[0,2]$ 的值域为 $[1,2]$，}
\step{因为 $f(x)=\dfrac4{x+2}$ 是 $g(x)=x^2-2ax+a^2-1$ 在 $[0,2]$ 上的“2 阶伴随函数”，}
\step{则对任意的 $x_1\in[0,2]$，总存在唯一的 $x_2\in[0,2]$ 时，使得 $f(x_1)g(x_2)=2$，}% 原卷此处印“时”字（按上下文为衍字），照录
\step{所以 $g(x_2)=\dfrac2{f(x_1)}\in[1,2]$，}
\step{即 $g(x)=x^2-2ax+a^2-1$ 在 $[0,2]$ 区间上的值域必定包含区间 $[1,2]$，}
\step{且 $g(x)$ 值域在 $[1,2]$ 对应的自变量是唯一的，}
\step{又因为函数 $g(x)=x^2-2ax+a^2-1$ 开口向上，对称轴为 $x=a$，}
\step{1) 当 $a\leqslant0$ 时，$g(x)$ 在 $[0,2]$ 区间上严格增，则必有 $\begin{cases}g(x)_{\min}=g(0)\leqslant1\\g(x)_{\max}=g(2)\geqslant2\end{cases}$，}
\step{解得 $-\sqrt2\leqslant a\leqslant0$；}
\step{2) 当 $a\geqslant2$ 时，$g(x)$ 在 $[0,2]$ 区间上严格减，则必有 $\begin{cases}g(x)_{\min}=g(2)\leqslant1\\g(x)_{\max}=g(0)\geqslant2\end{cases}$，}
\step{解得 $2\leqslant a\leqslant2+\sqrt2$；}
\step{3) 当 $0<a\leqslant1$ 时，$g(x)$ 在 $[0,a]$ 上严格减，在 $[a,2]$ 上严格增，}
\step{则必有 $\begin{cases}g(0)<1\\g(x)_{\max}=g(2)\geqslant2\end{cases}$，解得 $0<a\leqslant2-\sqrt3$，}
\step{4) 当 $1<a<2$ 时，$g(x)$ 在 $[0,a]$ 上严格减，在 $[a,2]$ 上严格增，}
\step{则必有 $\begin{cases}g(2)<1\\g(x)_{\max}=g(0)\geqslant2\end{cases}$，解得 $\sqrt3\leqslant a<2$.}
\step{综上所述，$a$ 的范围 $[-\sqrt2,2-\sqrt3]\cup[\sqrt3,2+\sqrt2]$.}% 原卷“范围”后缺“是／为”字，照录
\solend
\fi

\end{enumerate}

\end{document}
"
% 紧凑版：xelatex "\def\iscompact{1}% 高一122_华师大二附中_周末作业6.tex
%   —— 高一上数学“2026-2027 年华师大二附中高一上周末作业 6”（试卷版/答案版）
% 试卷版：xelatex 高一122_华师大二附中_周末作业6.tex
% 答案版：xelatex "\def\isanswer{1}\input{高一122_华师大二附中_周末作业6.tex}"
% 紧凑版：xelatex "\def\iscompact{1}\input{高一122_华师大二附中_周末作业6.tex}"
%
% 原始材料：微信公众号“魔都数学加油站”文章，作者破晓，2026-11-09 17:52 发布；
% 连载“27学年高一122”，原文标题“27学年高一122——2026-2027年上海市华师大二附中高一上周末作业6”；
% 原文链接：https://mp.weixin.qq.com/s/pC5066JFKNkQQwjMqTK8Ow。
% 正文为 14 张 PNG 页（第 14 页 2560×3620，其余 13 页 4762×6735），
% 题目下印有【解析】，为讲评版。
% 题目、解析均据原图转录；卷面红色水印及公众号页脚未录入。全卷无插图（已逐页检索，无“如图”）。
%
% 卷首标题：原卷印作“2026-2027 年华师大二附中高一上周末作业 6”（公众号标题中的“上海市”
% 卷面标题行没有，本稿照卷面），年份区间按全稿惯例排 en dash；原卷未印副标题，本稿按题目内容
% 标为“函数”。副标题理由（以全部 14 页实际内容为准）：逐页复核 19 题，考的都是函数性质——
% 奇偶性与充要条件（Q3、Q13、Q15、Q16）、单调性与单调区间（Q4、Q7、Q14、Q17）、定义域／
% 值域／最值与复合函数（Q5、Q6、Q8、Q10、Q12、Q19、Q20）、分段函数与取整函数（Q7、Q8、
% Q11、Q15、Q19）、抽象函数与新定义（Q12、Q20、Q21）；不等式仅作求参数范围／解集的工具
% （Q9、Q14、Q16、Q18），集合与常用逻辑用语仅零星出现（Q3 的充要条件、Q11③ 的集合、
% Q20 的“非空集合 I”），均非独立章节。故主体为单一模块“函数”，与同校 高一102 周末作业 5
% 同（该函数模块卷亦按 高一83／105 成例只用模块名、不叠“集合与常用逻辑用语、不等式”）。
%
% 与原卷的排版差异：
%   ① 原文从第 3 题起，填空题编号为 3—12、选择题为 13—16、解答题为 17—21，本稿保留原题号；
%   ② 原卷每题下直接印【解析】，本稿按全稿惯例置于答案版；四道选择题括号内均未印答案、
%      答案在解析末句“故选 …”，括号内留空（全卷无只印【答案】的题、无 \ans{}）；
%   ③ 原卷未印副标题，本稿按“知识模块”口径补出；
%   ④ 原卷填空横线按全稿惯例排为 \blankline[4.4em]；选择题括号排 \mbox{（\quad）}、
%      选项段前按惯例加 \keepwithprev；解答题不留作答区（原卷为讲评版）；
%   ⑤ 原卷数集排斜体 $R$、$N_+$（下标加号），按全稿惯例改 \R、\N_+
%      （$N_+$ 下标加号与 $\N^*$ 是两种记号，照录下标、只改数集字母为正体，同 高一05 第 20 题成例）；
%      绝对值排 \lvert…\rvert、集合描述竖线排 \mid；
%      $f(x)_{min}$／$f(x)_{max}$／$y_{max}$／$(y_{max})_{min}$ 一类斜体下标按 高一96 成例
%      改正体 \min／\max；不等式一律排 \geqslant／\leqslant；
%   ⑥ 第 20 题解析标签原卷印作【解得】（全卷唯一，其余 18 题均印【解析】），本稿按全稿惯例
%      由 \solbegin 统一印【解析】，正文内容照录（行内加 % 注、并入文件头存疑清单）；
%   ⑦ 题干与解析的标点（半角／全角逗号、冒号、括号、句点、子题括号）均照原卷，
%      题目正文引号用弯引号“ ”。
%
% 子集符号：第 21 题 (2) 的 $\left[\frac13+\frac1{9b-3},1\right]\subseteq\left[\frac12,b\right]$
%   印有下划线，照录 \subseteq。全卷未出现判别式 \triangle／\Delta。
%
% 原卷存疑／笔误（均照录加注）：
%   第 11 题解析对② 的“$y=y+\Delta y$”（按上下文应为 $y=[y]+\Delta y$，原卷漏方括号）；
%   第 19 题题干“某种产品 $x$ 万件）可获利”右括号无配对（疑漏印左括号“（”）；
%   第 20 题解析标签【解得】（见排版差异⑥）；
%   第 21 题 (3) 解析“总存在唯一的 $x_2\in[0,2]$ 时，使得……”的“时”字（按上下文衍字）；
%   第 21 题 (3) 末句“综上所述，$a$ 的范围 $[-\sqrt2,2-\sqrt3]\cup[\sqrt3,2+\sqrt2]$”
%      “范围”后缺“是／为”字，照录。
%
% 行数核对：逐题数 \step 条数，合计 162 条（Q3：5；Q4：3；Q5：1；Q6：3；Q7：2；
% Q8：5；Q9：7；Q10：10；Q11：8；Q12：4；Q13：12；Q14：6；Q15：14；Q16：14；
% Q17：6；Q18：8；Q19：7；Q20：19；Q21：28），与原卷印刷行逐题相等
% （cases 块作一个印刷行计）。全卷无 \ans{}、无只印【答案】的题。
%
\documentclass[a4paper,11pt]{article}
\usepackage{common}

\begin{document}

\examtitlest[3]{2026--2027 年华师大二附中高一上周末作业 6}{函数}

\bigsec{一、填空题}

\begin{enumerate}
\setcounter{enumi}{2}

\item 已知 $p:a=0$，$q$：函数 $f(x)=x^3+x+a$ 是奇函数，则 $p$ 是 $q$ 的\blankline[4.4em]条件.

\ifanswer
\solbegin
\step{若 $a=0$，则 $f(x)=x^3+x$，}
\step{易得 $f(x)$ 的定义域为 $\R$，且 $f(-x)=-(x^3+x)=-f(x)$，则 $f(x)$ 为奇函数，}
\step{反之，若 $f(x)=x^3+x+a$ 是奇函数，}
\step{则有 $f(-x)+f(x)=-x^3-x+a+x^3+x+a=2a=0$，必有 $a=0$，}
\step{故 $p$ 是 $q$ 的充要条件.}
\solend
\fi

\item 函数 $y=\sqrt{-x^2+2x}$ 的单调递增区间是\blankline[4.4em].

\ifanswer
\solbegin
\step{由函数定义域得 $-x^2+2x\geqslant0$，解得 $0\leqslant x\leqslant2$，}
\step{二次函数 $y=-x^2+2x$ 开口向下，对称轴为 $x=1$，}
\step{结合复合函数的单调性得函数 $y=\sqrt{-x^2+2x}$ 的单调递增区间是 $[0,1]$.}
\solend
\fi

\item 设 $f(x)$ 定义域为 $[1,4]$，且 $f(x)=2x+1$，则 $f(x^2)$ 的最大值为\blankline[4.4em].

\ifanswer
\solbegin
\step{由 $x^2\in[1,4]$，得 $f(x^2)=2x^2+1\leqslant9$，当 $x=\pm2$ 时取等号，最大值为 9.}
\solend
\fi

\item 函数 $f(x)=\sqrt{x-2}-x$，$x\in(2,6]$ 的值域为\blankline[4.4em].

\ifanswer
\solbegin
\step{令 $t=\sqrt{x-2}$，则 $x=2+t^2$，由 $2<x\leqslant6$ 得 $0<t\leqslant2$，}
\step{原函数可化为 $g(t)=-t^2+t-2$，$g(0)=-2$，$g(2)=-4$，$g\left(\dfrac12\right)=-\dfrac74$，}
\step{由二次函数性质得 $g(t)\in\left[-4,-\dfrac74\right]$.}
\solend
\fi

\item 若函数 $f(x)=\lvert x-a+1\rvert$ 在区间 $[1,+\infty)$ 上是严格增函数，则实数 $a$ 的取值范围是\blankline[4.4em].

\ifanswer
\solbegin
\step{$f(x)=\begin{cases}x-a+1,x\geqslant a-1\\a-1-x,x<a-1\end{cases}$，显然 $f(x)$ 在 $[a-1,+\infty)$ 上严格增，}
\step{要使函数在区间 $[1,+\infty)$ 上是严格增函数，则 $a-1\leqslant1$，即 $a\leqslant2$.}
\solend
\fi

\item 若函数 $f(x)=\begin{cases}x^2-5,x\leqslant2\\2ax+1,x>2\end{cases}$ 的值域是 $\R$，则 $a$ 的取值范围是\blankline[4.4em].

\ifanswer
\solbegin
\step{当 $x\leqslant2$ 时，$f(x)=x^2-5$，其开口向上，对称轴为 $x=0$，值域为 $[-5,+\infty)$，}
\step{由函数 $f(x)=\begin{cases}x^2-5,x\leqslant2\\2ax+1,x>2\end{cases}$ 的值域是 $\R$，}
\step{则当 $x>2$ 时，$f(x)=2ax+1$ 的值域应包含 $(-\infty,-5)$，}
\step{所以 $f(x)$ 为严格减函数，所以 $\begin{cases}2\times2a+1\geqslant-5\\2a<0\end{cases}$，即 $\begin{cases}a\geqslant-\dfrac32\\a<0\end{cases}$，}
\step{解得 $-\dfrac32\leqslant a<0$，故 $a$ 的取值范围是 $\left[-\dfrac32,0\right)$.}
\solend
\fi

\item 已知函数 $f(x)=ax^2+\dfrac bx$ 的图象经过点 $(-1,0)$，且满足 $f(1)=2$. 若 $\forall x\in(0,+\infty)$，都有 $f(x)+f\left(\dfrac1x\right)-2m\geqslant0$ 恒成立，则实数 $m$ 的取值范围为\blankline[4.4em].

\ifanswer
\solbegin
\step{因为 $f(x)=ax^2+\dfrac bx$ 的图象经过点 $(-1,0)$，且满足 $f(1)=2$，}
\step{所以 $\begin{cases}a-b=0\\a+b=2\end{cases}$，所以 $a=1$，$b=1$，即 $f(x)=x^2+\dfrac1x$，}
\step{$\forall x\in(0,+\infty)$，$f(x)+f\left(\dfrac1x\right)-2m\geqslant0\Leftrightarrow2m\leqslant x^2+\dfrac1x+\dfrac1{x^2}+x$ 恒成立，}
\step{令函数 $g(x)=x^2+\dfrac1x+\dfrac1{x^2}+x=\left(x+\dfrac1x\right)^2+x+\dfrac1x-2$，}
\step{而当 $x>0$ 时，$x+\dfrac1x\geqslant2$，当且仅当 $x=1$ 时取等号，}
\step{因此 $g(x)\geqslant2^2+2-2=4$，当且仅当 $x=1$ 时取等号，则 $2m\leqslant4$，解得 $m\leqslant2$，}
\step{所以实数 $m\in(-\infty,2]$.}
\solend
\fi

\item 若定义在区间 $[a,b](a<b)$ 上的函数 $f(x)=k-\sqrt{x+1}$ 值域也为 $[a,b]$，则实数 $k$ 的取值范围是\blankline[4.4em].

\ifanswer
\solbegin
\step{函数 $f(x)=k-\sqrt{x+1}$ 在定义域 $[-1,+\infty)$ 严格减，}
\step{当 $f(x)$ 的定义域为 $[a,b]$ 时，$f(x)$ 的值域也为 $[a,b]$，$a<b$，}
\step{故 $f(a)=k-\sqrt{a+1}=b$，$f(b)=k-\sqrt{b+1}=a$，}
\step{两式相减得 $\sqrt{a+1}-\sqrt{b+1}=a-b=(a+1)-(b+1)$}
\step{$=(\sqrt{a+1}-\sqrt{b+1})(\sqrt{a+1}+\sqrt{b+1})$，}
\step{所以 $\sqrt{a+1}-\sqrt{b+1}<0$，即 $\sqrt{a+1}+\sqrt{b+1}=1$，}
\step{则 $k=\sqrt{b+1}+a=a+1-\sqrt{a+1}$，}
\step{令 $\lambda=\sqrt{a+1}\geqslant0$，得 $k=\lambda^2-\lambda=\left(\lambda-\dfrac12\right)^2-\dfrac14$，}
\step{又 $\sqrt{a+1}<\sqrt{b+1}$，因为 $\sqrt{a+1}+\sqrt{b+1}=1$，所以 $\sqrt{a+1}\in\left[0,\dfrac12\right)$，}
\step{所以 $0\leqslant\lambda<\dfrac12$，故实数 $k$ 的取值范围为 $\left(-\dfrac14,0\right]$.}
\solend
\fi

\item 函数 $f(x)=[x]$ 的函数值表示不超过 $x$ 的最大整数，例如：$[-3.5]=-4$，$[2.1]=2$，称 $f(x)=[x]$ 为取整函数，则下列关于取整函数的说法：①对任意的 $x\in\R$，$[-x]=-[x]$ ②对任意的 $x$、$y\in\R$，$[x]+[y]\leqslant[x+y]$ ③集合 $\{y\mid y=[nx]-n[x],x\in\R,n\in\N_+\}$ 共有 $n$ 个元素 ④$x>0$ 时，关于 $x$ 的方程 $[2x]=\left[\dfrac1x\right]$ 有无数个解 其中正确的有\blankline[4.4em].

\ifanswer
\solbegin
\step{对①，因为 $[-2.1]=-3\neq[2.1]=2$，所以①错误；}
\step{对②，设 $x=[x]+\Delta x$，$\Delta x\in[0,1)$，$y=y+\Delta y$，$\Delta y\in[0,1)$，}% 原卷印作 $y=y+\Delta y$（按上下文应为 $y=[y]+\Delta y$），照录
\step{$[x+y]=[[x]+\Delta x+[y]+\Delta y]=[x]+[y]+[\Delta x+\Delta y]$，又 $\Delta x+\Delta y\in[0,\ 2)$，}
\step{故 $\Delta x+\Delta y=0$ 或 1，故 $[x]+[y]\leqslant[x+y]$，所以②正确；}
\step{对③，$[nx]-n[x]=[n[x]+n\cdot\Delta x]-n[x]=n[x]+[n\cdot\Delta x]-n[x]=[n\cdot\Delta x]$，}
\step{因为 $\Delta x\in[0,1)$，所以 $0\leqslant n\cdot\Delta x<n$，$n\in\N_+$，所以③正确；}
\step{对④，当 $\dfrac12<x<1$ 时，$[2x]=\left[\dfrac1x\right]=1$，故方程存在无数个解，所以④正确.}
\step{故正确的是②③④.}
\solend
\fi

\item 已知函数 $f(x)$ 定义域为 $[0,1]$，且满足 $f(x)+f(1-x)=1$，$f(x)=2f\left(\dfrac x5\right)$，对任意 $x_1\leqslant x_2$，恒有 $f(x_1)\leqslant f(x_2)$，则 $f\left(\dfrac1{2026}\right)$ 的值为\blankline[4.4em].

\ifanswer
\solbegin
\step{$f(0)=2f(0)\Rightarrow f(0)=0$，从而 $f(1)=1$，$f\left(\dfrac15\right)=\dfrac12$，$f\left(\dfrac45\right)=1-f\left(\dfrac15\right)=\dfrac12$，}
\step{由单调性，$x\in\left[\dfrac15,\dfrac45\right]$ 时，$f(x)=\dfrac12$.}
\step{因为 $\dfrac15<\dfrac{625}{2026}<\dfrac45$，反复使用 $f\left(\dfrac x5\right)=\dfrac12 f(x)$，}
\step{得 $f\left(\dfrac1{2026}\right)=\dfrac1{2^4}f\left(\dfrac{625}{2026}\right)=\dfrac1{32}$.}
\solend
\fi

\end{enumerate}

\bigsec{二、选择题}

\begin{enumerate}
\setcounter{enumi}{12}

\item 已知 $f(x)$ 为 $\R$ 上的奇函数，$g(x)$ 为 $\R$ 上的偶函数，且 $g(x)\neq0$，则下列说法正确的是\mbox{（\quad）}

\keepwithprev
A. $f(x)+g(x)$ 为 $\R$ 上的奇函数\qquad\qquad B. $f(x)-g(x)$ 为 $\R$ 上的奇函数

C. $\dfrac{f(x)}{g(x)}$ 为 $\R$ 上的偶函数\qquad\qquad D. $\lvert f(x)g(x)\rvert$ 为 $\R$ 上的偶函数

\ifanswer
\solbegin
\step{因为 $f(x)$ 为 $\R$ 上的奇函数，$g(x)$ 为 $\R$ 上的偶函数，且 $g(x)\neq0$，}
\step{所以 $f(-x)=-f(x)$，$g(-x)=g(x)$，}
\step{所以 $f(-x)+g(-x)=-f(x)+g(x)\neq-[f(x)+g(x)]$，}
\step{故 $f(x)+g(x)$ 为非奇非偶函数，$A$ 错误；}
\step{同理，$f(x)-g(x)$ 为非奇非偶函数，$B$ 错误；}
\step{设 $F(x)=\dfrac{f(x)}{g(x)}$，则 $F(-x)=\dfrac{f(-x)}{g(-x)}=\dfrac{-f(x)}{g(x)}=-F(x)$，}
\step{所以 $F(x)$ 为奇函数，$C$ 错误；}
\step{设函数 $H(x)=\lvert f(x)g(x)\rvert$，}
\step{因为 $f(x)$ 为 $\R$ 上的奇函数，$g(x)$ 为 $\R$ 上的偶函数，且 $g(x)\neq0$，}
\step{所以 $f(-x)=-f(x)$，$g(-x)=g(x)$，则由函数奇偶性的定义}
\step{得 $H(-x)=\lvert f(-x)g(-x)\rvert=\lvert -f(x)g(x)\rvert=\lvert f(x)g(x)\rvert=H(x)$，$D$ 正确.}
\step{故选 $D$.}
\solend
\fi

\item 若函数 $f(x)$ 是定义域为 $(-\infty,0)\cup(0,+\infty)$ 的奇函数，在区间 $(0,+\infty)$ 上严格增，且 $f(-1)=0$，则不等式 $\dfrac{f(x-1)}{x-2}<0$ 的解集为\mbox{（\quad）}

\keepwithprev
A. $(-\infty,-2)$\qquad\qquad B. $(-2,0)$\qquad\qquad C. $(0,1)$\qquad\qquad D. $(1,2)$

\ifanswer
\solbegin
\step{若函数 $f(x)$ 是定义域为 $(-\infty,0)\cup(0,+\infty)$ 的奇函数，在区间 $(0,+\infty)$ 上严格增，}
\step{由奇函数对称性得，$f(x)$ 在 $(-\infty,0)$ 上严格增，}
\step{因为 $f(-1)=0$，所以 $f(1)=0$，}
\step{则不等式 $\dfrac{f(x-1)}{x-2}<0$ 可转化为 $\begin{cases}x>2\\f(x-1)<0\end{cases}$ 或 $\begin{cases}x<2\\f(x-1)>0\end{cases}$，}
\step{即 $\begin{cases}x>2\\x-1<-1\text{或}0<x-1<1\end{cases}$ 或 $\begin{cases}x<2\\x-1>1\text{或}-1<x-1<0\end{cases}$，解得 $0<x<1$.}
\step{故选 $C$.}
\solend
\fi

\item 若函数 $f(x)$ 同时满足：(1) 对于定义域上的任意 $x$，恒有 $f(x)+f(-x)=0$；(2) 对于定义域上的任意 $x_1$、$x_2$，当 $x_1\neq x_2$ 时，恒有 $x_1f(x_1)-x_1f(x_2)>x_2f(x_1)-x_2f(x_2)$，则称函数 $f(x)$ 为“理想函数”. 给出下列四个函数：① $f(x)=x^2$；② $f(x)=x^3$；③ $f(x)=\dfrac{2x-1}{2x+1}$；④ $f(x)=\begin{cases}x^2+4x,x\geqslant0\\-x^2+4x,x<0\end{cases}$，其中被称为“理想函数”的有\mbox{（\quad）}

\keepwithprev
A. 1 个\qquad\qquad B. 2 个\qquad\qquad C. 3 个\qquad\qquad D. 4 个

\ifanswer
\solbegin
\step{函数 $f(x)$ 同时满足：(1) 对于定义域上的任意 $x$，恒有 $f(x)+f(-x)=0$；}
\step{(2) 对于定义域上的任意 $x_1$、$x_2$，当 $x_1\neq x_2$ 时，}
\step{恒有 $x_1f(x_1)-x_1f(x_2)>x_2f(x_1)-x_2f(x_2)$，}
\step{整理得 $(x_1-x_2)[f(x_1)-f(x_2)]>0$，则称函数 $f(x)$ 为“理想函数”.}
\step{即函数满足 (1) 奇函数；(2) 函数在定义域内为严格增函数；}
\step{① $f(x)=x^2$；不满足函数为奇函数在定义域内也不严格增；}
\step{② $f(x)=x^3$；满足函数既为奇函数，也在定义域内严格增，}
\step{故该函数为理想函数；}
\step{③ $f(x)=\dfrac{2x-1}{2x+1}=\dfrac{2x+1-2}{2x+1}=-\dfrac2{2x+1}+1$，该函数不为奇函数，}
\step{但是函数在定义域内严格增，故该函数不为理想函数；}
\step{④ $f(x)=\begin{cases}x^2+4x,x\geqslant0\\-x^2+4x,x<0\end{cases}$，函数的关系式满足 $f(-x)=-f(x)$，}
\step{且在 $x\geqslant0$ 时，函数严格增，由于奇函数在对称区间上单调性相同，}
\step{故该函数为理想函数.}
\step{故选 $B$.}
\solend
\fi

\item 已知 $f(x)$ 是定义在 $\R$ 上的奇函数，当 $x_1$、$x_2\in(0,+\infty)$ 且 $x_1\neq x_2$ 时，都有\newline $\dfrac{x_2f(x_1)-x_1f(x_2)}{x_1x_2(x_1-x_2)}>0$ 成立，$f(2028)=2028$，则不等式 $f(x)-x>0$ 的解集为\mbox{（\quad）}

\keepwithprev
A. $(-\infty,-2028)\cup(2028,+\infty)$\qquad\qquad B. $\left(-\dfrac1{2028},\dfrac1{2028}\right)$

C. $(-2028,2028)$\qquad\qquad D. $(-2028,0)\cup(2028,+\infty)$

\ifanswer
\solbegin
\step{因为 $f(x)$ 是定义在 $\R$ 上的奇函数，所以 $f(-x)=-f(x)$，}
\step{令 $g(x)=\dfrac{f(x)}x$，则 $g(-x)=\dfrac{f(-x)}{-x}=\dfrac{-f(x)}{-x}=\dfrac{f(x)}x=g(x)$，}
\step{所以 $g(x)$ 为偶函数，}
\step{当 $x_1$、$x_2\in(0,+\infty)$ 且 $x_1\neq x_2$ 时，都有 $\dfrac{x_2f(x_1)-x_1f(x_2)}{x_1x_2(x_1-x_2)}>0$ 成立，}
\step{即 $\dfrac{\dfrac{f(x_1)}{x_1}-\dfrac{f(x_2)}{x_2}}{x_1-x_2}>0$，即 $g(x)$ 在 $(0,+\infty)$ 上严格增，}
\step{由偶函数对称性得 $g(x)$ 在 $(-\infty,0)$ 上严格减，}
\step{因为 $f(2028)=2028$，所以 $g(2028)=1$，}
\step{则不等式 $f(x)-x>0$ 可化为 $f(x)>x$，}
\step{当 $x=0$ 时，$f(0)=0$，显然不成立，}
\step{当 $x>0$ 时，不等式可化为 $\dfrac{f(x)}x>1$，即 $g(x)>1=g(2028)$，}
\step{所以 $\lvert x\rvert>2028$，解得 $x>2028$，}
\step{当 $x<0$ 时，不等式可化为 $\dfrac{f(x)}x<1$，即 $g(x)<1=g(2028)$，}
\step{所以 $\lvert x\rvert<2028$，解得 $-2028<x<0$，}
\step{故不等式的解集为 $(-2028,0)\cup(2028,+\infty)$. 故选 $D$.}
\solend
\fi

\end{enumerate}

\bigsec{三、解答题}

\begin{enumerate}
\setcounter{enumi}{16}

\item 已知函数 $f(x)=\dfrac{2x-3}{x+1}$.

(1) 试用函数单调性的定义证明：函数 $f(x)$ 在 $(-1,+\infty)$ 上是严格增函数；

(2) 求函数 $f(x)$ 在区间 $[1,4]$ 上的最大值和最小值.

\ifanswer
\solbegin
\step{(1) 设 $-1<x_1<x_2$，则 $f(x_1)-f(x_2)=\dfrac{2x_1-3}{x_1+1}-\dfrac{2x_2-3}{x_2+1}=\dfrac{5(x_1-x_2)}{(x_1+1)(x_2+1)}$}
\step{因为 $-1<x_1<x_2$，所以 $x_1-x_2<0$，$x_1+1>0$，$x_2+1>0$，}
\step{所以 $f(x_1)-f(x_2)<0$，所以 $f(x)$ 在 $(-1,+\infty)$ 上严格增.}
\step{(2) 由 (1) 得 $f(x)$ 在区间 $[1,4]$ 上严格增，}
\step{所以 $f(x)_{\min}=f(1)=-\dfrac12$，$f(x)_{\max}=f(4)=1$，}
\step{所以函数 $f(x)$ 在区间 $[1,4]$ 上的最大值为 1，最小值为 $-\dfrac12$.}
\solend
\fi

\item 设函数 $f(x)=x^2-mx+m$.

(1) 当 $m=2$ 时，求不等式 $f(x)<x$ 的解集；

(2) 若 $\exists x\in[1,2]$，使不等式 $f(x)<x+m-2$ 成立，求实数 $m$ 的取值范围.

\ifanswer
\solbegin
\step{(1) $m=2$ 时，函数 $f(x)=x^2-2x+2$，}
\step{不等式 $f(x)<x$ 可化为 $x^2-3x+2<0$，解得 $1<x<2$，}
\step{所以不等式的解集为 $\{x\mid1<x<2\}$；}
\step{(2) 若 $\exists x\in[1,2]$，使不等式 $f(x)<x+m-2$ 成立，即 $x^2-x+2<mx$，}
\step{所以 $m>x+\dfrac2x-1$，设 $g(x)=x+\dfrac2x-1$，$x\in[1,2]$，}
\step{则 $g(x)=x+\dfrac2x-1\geqslant2\sqrt{x\cdot\dfrac2x}-1=2\sqrt2-1$，}
\step{当且仅当 $x=\dfrac2x$，即 $x=\sqrt2$ 时取等号，}
\step{所以实数 $m$ 的取值范围是 $\{m\mid m>2\sqrt2-1\}$.}
\solend
\fi

\item 某企业原来生产某种产品 $x$ 万件）可获利 $w(x)$（万元），且满足\newline% 原卷“万件）”右括号无配对、疑漏印左括号“（”，照录
$w(x)=\begin{cases}20(x^2+17),0<x\leqslant2\\500-\dfrac{80}{x-1},2<x\leqslant5\end{cases}$.
现该企业计划引进新的生产设备和新的产品方案优化产品，优化后的产品的其他成本总投入为 $(20x+10)$ 万元. 由市场调研分析得知，当前产品供不应求. 记该企业优化后的产品的利润为 $f(x)$（单位：万元）.

(1) 求函数 $f(x)$ 的解析式；

(2) 当优化后的产品产量为多少万件时，该企业的利润 $f(x)$ 最大？最大利润是多少？请说明理由.

\ifanswer
\solbegin
\step{(1) 由题意得 $f(x)=\begin{cases}20x^2-20x+330,0<x\leqslant2\\-20x-\dfrac{80}{x-1}+490,2<x\leqslant5\end{cases}$；}
\step{(2) 当 $0<x\leqslant2$ 时，$f(x)=20(x^2-x)+330=20\left(x-\dfrac12\right)^2+325$，}
\step{故当 $x=2$ 时，$f(x)$ 取最大值，$f(2)=370$；}
\step{当 $2<x\leqslant5$ 时，$f(x)=-\left[20(x-1)+\dfrac{80}{x-1}\right]+470\leqslant-80+470=390$，}
\step{当且仅当 $20(x-1)=\dfrac{80}{x-1}$，即 $x=3$ 时，$f(3)=390$ 为最大值；}
\step{因为 $370<390$，}
\step{所以优化后的产品产量为 3 万件时，企业获最大利润 390 万元.}
\solend
\fi

\item 俄国数学家切比雪夫（1821—1894）是研究直线逼近函数的理论先驱. 对定义在非空集合 $I$ 上的函数 $f(x)$，以及函数 $g(x)=kx+b(k,b\in\R)$，切比雪夫将函数 $y=\lvert f(x)-g(x)\rvert,x\in I$ 的最大值称为 $f(x),g(x)$ 的“偏差”.

(1) 函数 $f(x)=x^2,x\in[0,1],g(x)=-x-2$，求 $f(x),g(x)$ 的“偏差”；

(2) 函数 $f(x)=\dfrac1x+1,x\in[1,2],g(x)=kx+1(k>0)$，若 $f(x),g(x)$ 的“偏差”为 2，求 $k$ 的值；

(3) 函数 $f(x)=x^2-x,x\in[0,3],g(x)=2x+b$，当 $f(x),g(x)$ 的“偏差”取最小值时，求 $b$ 的值，并求出“偏差”的最小值.

\ifanswer
% 注：原卷本题解析标签印作【解得】（全卷唯一，其余 18 题均印【解析】）；
%     本稿按全稿惯例由 \solbegin 统一印【解析】，正文内容照录。见文件头排版差异⑥。
\solbegin
\step{(1) $y=\lvert f(x)-g(x)\rvert=\lvert x^2+x+2\rvert=\left\lvert\left(x+\dfrac12\right)^2+\dfrac74\right\rvert,x\in[0,1]$，}
\step{因为 $x\in[0,1]$，所以 $\left(x+\dfrac12\right)^2+\dfrac74\in[2,4]$，}
\step{则 $y=\left\lvert\left(x+\dfrac12\right)^2+\dfrac74\right\rvert\in[2,4]$，$y_{\max}=4$，}
\step{所以函数 $f(x)$ 与 $g(x)$ 的“偏差”为 4.}
\step{(2) 令 $t=f(x)-g(x)=\dfrac1x+1-(kx+1)=\dfrac1x-kx,x\in[1,2]$，}
\step{因为 $k>0$，所以 $t=\dfrac1x-kx$ 是单调减函数，}
\step{所以 $t\in\left[\dfrac12-2k,1-k\right]$，由单调，$y=\lvert t\rvert,t\in\left[\dfrac12-2k,1-k\right]$，且 $y_{\max}=2$.}
\step{当 $\left\lvert\dfrac12-2k\right\rvert\leqslant\lvert1-k\rvert$，即 $0<k\leqslant\dfrac12$ 时，$\lvert1-k\rvert=2$，解得 $k=3$ 或 $k=-1$，}
\step{均不符合；}
\step{当 $\left\lvert\dfrac12-2k\right\rvert>\lvert1-k\rvert$，即 $k>\dfrac12$ 时，$\left\lvert\dfrac12-2k\right\rvert=2$，$\dfrac12-2k=2$ 或 $\dfrac12-2k=-2$，}
\step{解得 $k=\dfrac54$ 或 $k=-\dfrac34$（舍），所以 $k=\dfrac54$.}
\step{(3) $y=\lvert f(x)-g(x)\rvert=\lvert x^2-x-(2x+b)\rvert=\lvert x^2-3x-b\rvert$}
\step{$=\left\lvert\left(x-\dfrac32\right)^2-\dfrac94-b\right\rvert,x\in[0,3]$，}
\step{因为 $x\in[0,3]$，所以 $\left(x-\dfrac32\right)^2-\dfrac94-b\in\left[-\dfrac94-b,-b\right]$，}
\step{由 $y=\left\lvert\left(x-\dfrac32\right)^2-\dfrac94-b\right\rvert$，则 $y_{\max}=\max\left\{\lvert b\rvert,\left\lvert b+\dfrac94\right\rvert\right\}$，}
\step{令 $\lvert b\rvert\geqslant\left\lvert b+\dfrac94\right\rvert$，即 $b^2\geqslant\left(b+\dfrac94\right)^2$，解得 $b\leqslant-\dfrac98$，}
\step{$y_{\max}=\max\left\{\lvert b\rvert,\left\lvert b+\dfrac94\right\rvert\right\}=\begin{cases}-b,b\leqslant-\dfrac98\\b+\dfrac94,b>-\dfrac98\end{cases}$，}
\step{故当且仅当 $b=-\dfrac98$ 时，$(y_{\max})_{\min}=\dfrac98$，}
\step{故当 $b$ 的值为 $-\dfrac98$ 时，函数 $f(x)$ 与 $g(x)$ 的“偏差”取最小值 $\dfrac98$.}
\solend
\fi

\item 若函数 $f(x)$ 与 $g(x)$ 满足：对任意的 $x_1\in D$，总存在唯一的 $x_2\in D$，使 $f(x_1)g(x_2)=m$ 成立，则称 $f(x)$ 是 $g(x)$ 在区间 $D$ 上的“$m$ 阶伴随函数”；对任意的 $x_1\in D$，总存在唯一的 $x_2\in D$，使 $f(x_1)f(x_2)=m$ 成立，则称 $f(x)$ 是区间 $D$ 上的“$m$ 阶自伴函数”.

(1) 判断 $f(x)=x^2+1$ 是否为区间 $[0,3]$ 上的“2 阶自伴函数”？并说明理由：

(2) 若函数 $f(x)=3x-1$ 为区间 $\left[\dfrac12,b\right]$ 上的“1 阶自伴函数”，求 $b$ 的值；

(3) 若 $f(x)=\dfrac4{x+2}$ 是 $g(x)=x^2-2ax+a^2-1$ 在区间 $[0,2]$ 上的“2 阶伴随函数”，求实数 $a$ 的取值范围.

\ifanswer
\solbegin
\step{(1) 假定函数 $f(x)=x^2+1$ 是区间 $[0,3]$ 上的“2 阶自伴函数”，}
\step{取 $x_1=2,f(x_1)=f(2)=5$，由 $f(x_1)f(x_2)=2$，得 $5(x_2^2+1)=2$，}
\step{显然此方程无实数解，}
\step{所以函数 $f(x)=x^2+1$ 不是区间 $[0,3]$ 上的“2 阶自伴函数”.}
\step{(2) 函数 $f(x)=3x-1$ 为区间 $\left[\dfrac12,b\right]$ 上的“1 阶自伴函数”，}
\step{则对任意 $x_1\in\left[\dfrac12,b\right]$，总存在唯一的 $x_2\in\left[\dfrac12,b\right]$，使得 $f(x_1)f(x_2)=1$，}
\step{即 $3x_2-1=\dfrac1{3x_1-1}$，整理得 $x_2=\dfrac13+\dfrac1{3(3x_1-1)}$，}
\step{显然函数 $x_2=\dfrac13+\dfrac1{3(3x_1-1)}$ 在 $\left[\dfrac12,b\right]$ 上严格减，}
\step{且当 $x_1=\dfrac12$ 时，$x_2=1$，当 $x_1=b$ 时，$x_2=\dfrac13+\dfrac1{9b-3}$，}
\step{因此对 $\left[\dfrac12,b\right]$ 内的每一个 $x_1$，在 $\left[\dfrac13+\dfrac1{9b-3},1\right]$ 内有唯一 $x_2$ 值与之对应，}
\step{则 $x_2\in\left[\dfrac12,b\right]$，于是 $\left[\dfrac13+\dfrac1{9b-3},1\right]\subseteq\left[\dfrac12,b\right]$，应有 $\begin{cases}b\geqslant1\\\dfrac13+\dfrac1{9b-3}\geqslant\dfrac12\end{cases}$，}
\step{得 $\begin{cases}b\geqslant1\\b\leqslant1\end{cases}$，即 $b=1$，所以 $b$ 的值是 1.}
\step{(3) 由函数 $f(x)=\dfrac4{x+2}$ 在区间 $[0,2]$ 的值域为 $[1,2]$，}
\step{因为 $f(x)=\dfrac4{x+2}$ 是 $g(x)=x^2-2ax+a^2-1$ 在 $[0,2]$ 上的“2 阶伴随函数”，}
\step{则对任意的 $x_1\in[0,2]$，总存在唯一的 $x_2\in[0,2]$ 时，使得 $f(x_1)g(x_2)=2$，}% 原卷此处印“时”字（按上下文为衍字），照录
\step{所以 $g(x_2)=\dfrac2{f(x_1)}\in[1,2]$，}
\step{即 $g(x)=x^2-2ax+a^2-1$ 在 $[0,2]$ 区间上的值域必定包含区间 $[1,2]$，}
\step{且 $g(x)$ 值域在 $[1,2]$ 对应的自变量是唯一的，}
\step{又因为函数 $g(x)=x^2-2ax+a^2-1$ 开口向上，对称轴为 $x=a$，}
\step{1) 当 $a\leqslant0$ 时，$g(x)$ 在 $[0,2]$ 区间上严格增，则必有 $\begin{cases}g(x)_{\min}=g(0)\leqslant1\\g(x)_{\max}=g(2)\geqslant2\end{cases}$，}
\step{解得 $-\sqrt2\leqslant a\leqslant0$；}
\step{2) 当 $a\geqslant2$ 时，$g(x)$ 在 $[0,2]$ 区间上严格减，则必有 $\begin{cases}g(x)_{\min}=g(2)\leqslant1\\g(x)_{\max}=g(0)\geqslant2\end{cases}$，}
\step{解得 $2\leqslant a\leqslant2+\sqrt2$；}
\step{3) 当 $0<a\leqslant1$ 时，$g(x)$ 在 $[0,a]$ 上严格减，在 $[a,2]$ 上严格增，}
\step{则必有 $\begin{cases}g(0)<1\\g(x)_{\max}=g(2)\geqslant2\end{cases}$，解得 $0<a\leqslant2-\sqrt3$，}
\step{4) 当 $1<a<2$ 时，$g(x)$ 在 $[0,a]$ 上严格减，在 $[a,2]$ 上严格增，}
\step{则必有 $\begin{cases}g(2)<1\\g(x)_{\max}=g(0)\geqslant2\end{cases}$，解得 $\sqrt3\leqslant a<2$.}
\step{综上所述，$a$ 的范围 $[-\sqrt2,2-\sqrt3]\cup[\sqrt3,2+\sqrt2]$.}% 原卷“范围”后缺“是／为”字，照录
\solend
\fi

\end{enumerate}

\end{document}
"
%
% 原始材料：微信公众号“魔都数学加油站”文章，作者破晓，2026-11-09 17:52 发布；
% 连载“27学年高一122”，原文标题“27学年高一122——2026-2027年上海市华师大二附中高一上周末作业6”；
% 原文链接：https://mp.weixin.qq.com/s/pC5066JFKNkQQwjMqTK8Ow。
% 正文为 14 张 PNG 页（第 14 页 2560×3620，其余 13 页 4762×6735），
% 题目下印有【解析】，为讲评版。
% 题目、解析均据原图转录；卷面红色水印及公众号页脚未录入。全卷无插图（已逐页检索，无“如图”）。
%
% 卷首标题：原卷印作“2026-2027 年华师大二附中高一上周末作业 6”（公众号标题中的“上海市”
% 卷面标题行没有，本稿照卷面），年份区间按全稿惯例排 en dash；原卷未印副标题，本稿按题目内容
% 标为“函数”。副标题理由（以全部 14 页实际内容为准）：逐页复核 19 题，考的都是函数性质——
% 奇偶性与充要条件（Q3、Q13、Q15、Q16）、单调性与单调区间（Q4、Q7、Q14、Q17）、定义域／
% 值域／最值与复合函数（Q5、Q6、Q8、Q10、Q12、Q19、Q20）、分段函数与取整函数（Q7、Q8、
% Q11、Q15、Q19）、抽象函数与新定义（Q12、Q20、Q21）；不等式仅作求参数范围／解集的工具
% （Q9、Q14、Q16、Q18），集合与常用逻辑用语仅零星出现（Q3 的充要条件、Q11③ 的集合、
% Q20 的“非空集合 I”），均非独立章节。故主体为单一模块“函数”，与同校 高一102 周末作业 5
% 同（该函数模块卷亦按 高一83／105 成例只用模块名、不叠“集合与常用逻辑用语、不等式”）。
%
% 与原卷的排版差异：
%   ① 原文从第 3 题起，填空题编号为 3—12、选择题为 13—16、解答题为 17—21，本稿保留原题号；
%   ② 原卷每题下直接印【解析】，本稿按全稿惯例置于答案版；四道选择题括号内均未印答案、
%      答案在解析末句“故选 …”，括号内留空（全卷无只印【答案】的题、无 \ans{}）；
%   ③ 原卷未印副标题，本稿按“知识模块”口径补出；
%   ④ 原卷填空横线按全稿惯例排为 \blankline[4.4em]；选择题括号排 \mbox{（\quad）}、
%      选项段前按惯例加 \keepwithprev；解答题不留作答区（原卷为讲评版）；
%   ⑤ 原卷数集排斜体 $R$、$N_+$（下标加号），按全稿惯例改 \R、\N_+
%      （$N_+$ 下标加号与 $\N^*$ 是两种记号，照录下标、只改数集字母为正体，同 高一05 第 20 题成例）；
%      绝对值排 \lvert…\rvert、集合描述竖线排 \mid；
%      $f(x)_{min}$／$f(x)_{max}$／$y_{max}$／$(y_{max})_{min}$ 一类斜体下标按 高一96 成例
%      改正体 \min／\max；不等式一律排 \geqslant／\leqslant；
%   ⑥ 第 20 题解析标签原卷印作【解得】（全卷唯一，其余 18 题均印【解析】），本稿按全稿惯例
%      由 \solbegin 统一印【解析】，正文内容照录（行内加 % 注、并入文件头存疑清单）；
%   ⑦ 题干与解析的标点（半角／全角逗号、冒号、括号、句点、子题括号）均照原卷，
%      题目正文引号用弯引号“ ”。
%
% 子集符号：第 21 题 (2) 的 $\left[\frac13+\frac1{9b-3},1\right]\subseteq\left[\frac12,b\right]$
%   印有下划线，照录 \subseteq。全卷未出现判别式 \triangle／\Delta。
%
% 原卷存疑／笔误（均照录加注）：
%   第 11 题解析对② 的“$y=y+\Delta y$”（按上下文应为 $y=[y]+\Delta y$，原卷漏方括号）；
%   第 19 题题干“某种产品 $x$ 万件）可获利”右括号无配对（疑漏印左括号“（”）；
%   第 20 题解析标签【解得】（见排版差异⑥）；
%   第 21 题 (3) 解析“总存在唯一的 $x_2\in[0,2]$ 时，使得……”的“时”字（按上下文衍字）；
%   第 21 题 (3) 末句“综上所述，$a$ 的范围 $[-\sqrt2,2-\sqrt3]\cup[\sqrt3,2+\sqrt2]$”
%      “范围”后缺“是／为”字，照录。
%
% 行数核对：逐题数 \step 条数，合计 162 条（Q3：5；Q4：3；Q5：1；Q6：3；Q7：2；
% Q8：5；Q9：7；Q10：10；Q11：8；Q12：4；Q13：12；Q14：6；Q15：14；Q16：14；
% Q17：6；Q18：8；Q19：7；Q20：19；Q21：28），与原卷印刷行逐题相等
% （cases 块作一个印刷行计）。全卷无 \ans{}、无只印【答案】的题。
%
\documentclass[a4paper,11pt]{article}
\usepackage{common}

\begin{document}

\examtitlest[3]{2026--2027 年华师大二附中高一上周末作业 6}{函数}

\bigsec{一、填空题}

\begin{enumerate}
\setcounter{enumi}{2}

\item 已知 $p:a=0$，$q$：函数 $f(x)=x^3+x+a$ 是奇函数，则 $p$ 是 $q$ 的\blankline[4.4em]条件.

\ifanswer
\solbegin
\step{若 $a=0$，则 $f(x)=x^3+x$，}
\step{易得 $f(x)$ 的定义域为 $\R$，且 $f(-x)=-(x^3+x)=-f(x)$，则 $f(x)$ 为奇函数，}
\step{反之，若 $f(x)=x^3+x+a$ 是奇函数，}
\step{则有 $f(-x)+f(x)=-x^3-x+a+x^3+x+a=2a=0$，必有 $a=0$，}
\step{故 $p$ 是 $q$ 的充要条件.}
\solend
\fi

\item 函数 $y=\sqrt{-x^2+2x}$ 的单调递增区间是\blankline[4.4em].

\ifanswer
\solbegin
\step{由函数定义域得 $-x^2+2x\geqslant0$，解得 $0\leqslant x\leqslant2$，}
\step{二次函数 $y=-x^2+2x$ 开口向下，对称轴为 $x=1$，}
\step{结合复合函数的单调性得函数 $y=\sqrt{-x^2+2x}$ 的单调递增区间是 $[0,1]$.}
\solend
\fi

\item 设 $f(x)$ 定义域为 $[1,4]$，且 $f(x)=2x+1$，则 $f(x^2)$ 的最大值为\blankline[4.4em].

\ifanswer
\solbegin
\step{由 $x^2\in[1,4]$，得 $f(x^2)=2x^2+1\leqslant9$，当 $x=\pm2$ 时取等号，最大值为 9.}
\solend
\fi

\item 函数 $f(x)=\sqrt{x-2}-x$，$x\in(2,6]$ 的值域为\blankline[4.4em].

\ifanswer
\solbegin
\step{令 $t=\sqrt{x-2}$，则 $x=2+t^2$，由 $2<x\leqslant6$ 得 $0<t\leqslant2$，}
\step{原函数可化为 $g(t)=-t^2+t-2$，$g(0)=-2$，$g(2)=-4$，$g\left(\dfrac12\right)=-\dfrac74$，}
\step{由二次函数性质得 $g(t)\in\left[-4,-\dfrac74\right]$.}
\solend
\fi

\item 若函数 $f(x)=\lvert x-a+1\rvert$ 在区间 $[1,+\infty)$ 上是严格增函数，则实数 $a$ 的取值范围是\blankline[4.4em].

\ifanswer
\solbegin
\step{$f(x)=\begin{cases}x-a+1,x\geqslant a-1\\a-1-x,x<a-1\end{cases}$，显然 $f(x)$ 在 $[a-1,+\infty)$ 上严格增，}
\step{要使函数在区间 $[1,+\infty)$ 上是严格增函数，则 $a-1\leqslant1$，即 $a\leqslant2$.}
\solend
\fi

\item 若函数 $f(x)=\begin{cases}x^2-5,x\leqslant2\\2ax+1,x>2\end{cases}$ 的值域是 $\R$，则 $a$ 的取值范围是\blankline[4.4em].

\ifanswer
\solbegin
\step{当 $x\leqslant2$ 时，$f(x)=x^2-5$，其开口向上，对称轴为 $x=0$，值域为 $[-5,+\infty)$，}
\step{由函数 $f(x)=\begin{cases}x^2-5,x\leqslant2\\2ax+1,x>2\end{cases}$ 的值域是 $\R$，}
\step{则当 $x>2$ 时，$f(x)=2ax+1$ 的值域应包含 $(-\infty,-5)$，}
\step{所以 $f(x)$ 为严格减函数，所以 $\begin{cases}2\times2a+1\geqslant-5\\2a<0\end{cases}$，即 $\begin{cases}a\geqslant-\dfrac32\\a<0\end{cases}$，}
\step{解得 $-\dfrac32\leqslant a<0$，故 $a$ 的取值范围是 $\left[-\dfrac32,0\right)$.}
\solend
\fi

\item 已知函数 $f(x)=ax^2+\dfrac bx$ 的图象经过点 $(-1,0)$，且满足 $f(1)=2$. 若 $\forall x\in(0,+\infty)$，都有 $f(x)+f\left(\dfrac1x\right)-2m\geqslant0$ 恒成立，则实数 $m$ 的取值范围为\blankline[4.4em].

\ifanswer
\solbegin
\step{因为 $f(x)=ax^2+\dfrac bx$ 的图象经过点 $(-1,0)$，且满足 $f(1)=2$，}
\step{所以 $\begin{cases}a-b=0\\a+b=2\end{cases}$，所以 $a=1$，$b=1$，即 $f(x)=x^2+\dfrac1x$，}
\step{$\forall x\in(0,+\infty)$，$f(x)+f\left(\dfrac1x\right)-2m\geqslant0\Leftrightarrow2m\leqslant x^2+\dfrac1x+\dfrac1{x^2}+x$ 恒成立，}
\step{令函数 $g(x)=x^2+\dfrac1x+\dfrac1{x^2}+x=\left(x+\dfrac1x\right)^2+x+\dfrac1x-2$，}
\step{而当 $x>0$ 时，$x+\dfrac1x\geqslant2$，当且仅当 $x=1$ 时取等号，}
\step{因此 $g(x)\geqslant2^2+2-2=4$，当且仅当 $x=1$ 时取等号，则 $2m\leqslant4$，解得 $m\leqslant2$，}
\step{所以实数 $m\in(-\infty,2]$.}
\solend
\fi

\item 若定义在区间 $[a,b](a<b)$ 上的函数 $f(x)=k-\sqrt{x+1}$ 值域也为 $[a,b]$，则实数 $k$ 的取值范围是\blankline[4.4em].

\ifanswer
\solbegin
\step{函数 $f(x)=k-\sqrt{x+1}$ 在定义域 $[-1,+\infty)$ 严格减，}
\step{当 $f(x)$ 的定义域为 $[a,b]$ 时，$f(x)$ 的值域也为 $[a,b]$，$a<b$，}
\step{故 $f(a)=k-\sqrt{a+1}=b$，$f(b)=k-\sqrt{b+1}=a$，}
\step{两式相减得 $\sqrt{a+1}-\sqrt{b+1}=a-b=(a+1)-(b+1)$}
\step{$=(\sqrt{a+1}-\sqrt{b+1})(\sqrt{a+1}+\sqrt{b+1})$，}
\step{所以 $\sqrt{a+1}-\sqrt{b+1}<0$，即 $\sqrt{a+1}+\sqrt{b+1}=1$，}
\step{则 $k=\sqrt{b+1}+a=a+1-\sqrt{a+1}$，}
\step{令 $\lambda=\sqrt{a+1}\geqslant0$，得 $k=\lambda^2-\lambda=\left(\lambda-\dfrac12\right)^2-\dfrac14$，}
\step{又 $\sqrt{a+1}<\sqrt{b+1}$，因为 $\sqrt{a+1}+\sqrt{b+1}=1$，所以 $\sqrt{a+1}\in\left[0,\dfrac12\right)$，}
\step{所以 $0\leqslant\lambda<\dfrac12$，故实数 $k$ 的取值范围为 $\left(-\dfrac14,0\right]$.}
\solend
\fi

\item 函数 $f(x)=[x]$ 的函数值表示不超过 $x$ 的最大整数，例如：$[-3.5]=-4$，$[2.1]=2$，称 $f(x)=[x]$ 为取整函数，则下列关于取整函数的说法：①对任意的 $x\in\R$，$[-x]=-[x]$ ②对任意的 $x$、$y\in\R$，$[x]+[y]\leqslant[x+y]$ ③集合 $\{y\mid y=[nx]-n[x],x\in\R,n\in\N_+\}$ 共有 $n$ 个元素 ④$x>0$ 时，关于 $x$ 的方程 $[2x]=\left[\dfrac1x\right]$ 有无数个解 其中正确的有\blankline[4.4em].

\ifanswer
\solbegin
\step{对①，因为 $[-2.1]=-3\neq[2.1]=2$，所以①错误；}
\step{对②，设 $x=[x]+\Delta x$，$\Delta x\in[0,1)$，$y=y+\Delta y$，$\Delta y\in[0,1)$，}% 原卷印作 $y=y+\Delta y$（按上下文应为 $y=[y]+\Delta y$），照录
\step{$[x+y]=[[x]+\Delta x+[y]+\Delta y]=[x]+[y]+[\Delta x+\Delta y]$，又 $\Delta x+\Delta y\in[0,\ 2)$，}
\step{故 $\Delta x+\Delta y=0$ 或 1，故 $[x]+[y]\leqslant[x+y]$，所以②正确；}
\step{对③，$[nx]-n[x]=[n[x]+n\cdot\Delta x]-n[x]=n[x]+[n\cdot\Delta x]-n[x]=[n\cdot\Delta x]$，}
\step{因为 $\Delta x\in[0,1)$，所以 $0\leqslant n\cdot\Delta x<n$，$n\in\N_+$，所以③正确；}
\step{对④，当 $\dfrac12<x<1$ 时，$[2x]=\left[\dfrac1x\right]=1$，故方程存在无数个解，所以④正确.}
\step{故正确的是②③④.}
\solend
\fi

\item 已知函数 $f(x)$ 定义域为 $[0,1]$，且满足 $f(x)+f(1-x)=1$，$f(x)=2f\left(\dfrac x5\right)$，对任意 $x_1\leqslant x_2$，恒有 $f(x_1)\leqslant f(x_2)$，则 $f\left(\dfrac1{2026}\right)$ 的值为\blankline[4.4em].

\ifanswer
\solbegin
\step{$f(0)=2f(0)\Rightarrow f(0)=0$，从而 $f(1)=1$，$f\left(\dfrac15\right)=\dfrac12$，$f\left(\dfrac45\right)=1-f\left(\dfrac15\right)=\dfrac12$，}
\step{由单调性，$x\in\left[\dfrac15,\dfrac45\right]$ 时，$f(x)=\dfrac12$.}
\step{因为 $\dfrac15<\dfrac{625}{2026}<\dfrac45$，反复使用 $f\left(\dfrac x5\right)=\dfrac12 f(x)$，}
\step{得 $f\left(\dfrac1{2026}\right)=\dfrac1{2^4}f\left(\dfrac{625}{2026}\right)=\dfrac1{32}$.}
\solend
\fi

\end{enumerate}

\bigsec{二、选择题}

\begin{enumerate}
\setcounter{enumi}{12}

\item 已知 $f(x)$ 为 $\R$ 上的奇函数，$g(x)$ 为 $\R$ 上的偶函数，且 $g(x)\neq0$，则下列说法正确的是\mbox{（\quad）}

\keepwithprev
A. $f(x)+g(x)$ 为 $\R$ 上的奇函数\qquad\qquad B. $f(x)-g(x)$ 为 $\R$ 上的奇函数

C. $\dfrac{f(x)}{g(x)}$ 为 $\R$ 上的偶函数\qquad\qquad D. $\lvert f(x)g(x)\rvert$ 为 $\R$ 上的偶函数

\ifanswer
\solbegin
\step{因为 $f(x)$ 为 $\R$ 上的奇函数，$g(x)$ 为 $\R$ 上的偶函数，且 $g(x)\neq0$，}
\step{所以 $f(-x)=-f(x)$，$g(-x)=g(x)$，}
\step{所以 $f(-x)+g(-x)=-f(x)+g(x)\neq-[f(x)+g(x)]$，}
\step{故 $f(x)+g(x)$ 为非奇非偶函数，$A$ 错误；}
\step{同理，$f(x)-g(x)$ 为非奇非偶函数，$B$ 错误；}
\step{设 $F(x)=\dfrac{f(x)}{g(x)}$，则 $F(-x)=\dfrac{f(-x)}{g(-x)}=\dfrac{-f(x)}{g(x)}=-F(x)$，}
\step{所以 $F(x)$ 为奇函数，$C$ 错误；}
\step{设函数 $H(x)=\lvert f(x)g(x)\rvert$，}
\step{因为 $f(x)$ 为 $\R$ 上的奇函数，$g(x)$ 为 $\R$ 上的偶函数，且 $g(x)\neq0$，}
\step{所以 $f(-x)=-f(x)$，$g(-x)=g(x)$，则由函数奇偶性的定义}
\step{得 $H(-x)=\lvert f(-x)g(-x)\rvert=\lvert -f(x)g(x)\rvert=\lvert f(x)g(x)\rvert=H(x)$，$D$ 正确.}
\step{故选 $D$.}
\solend
\fi

\item 若函数 $f(x)$ 是定义域为 $(-\infty,0)\cup(0,+\infty)$ 的奇函数，在区间 $(0,+\infty)$ 上严格增，且 $f(-1)=0$，则不等式 $\dfrac{f(x-1)}{x-2}<0$ 的解集为\mbox{（\quad）}

\keepwithprev
A. $(-\infty,-2)$\qquad\qquad B. $(-2,0)$\qquad\qquad C. $(0,1)$\qquad\qquad D. $(1,2)$

\ifanswer
\solbegin
\step{若函数 $f(x)$ 是定义域为 $(-\infty,0)\cup(0,+\infty)$ 的奇函数，在区间 $(0,+\infty)$ 上严格增，}
\step{由奇函数对称性得，$f(x)$ 在 $(-\infty,0)$ 上严格增，}
\step{因为 $f(-1)=0$，所以 $f(1)=0$，}
\step{则不等式 $\dfrac{f(x-1)}{x-2}<0$ 可转化为 $\begin{cases}x>2\\f(x-1)<0\end{cases}$ 或 $\begin{cases}x<2\\f(x-1)>0\end{cases}$，}
\step{即 $\begin{cases}x>2\\x-1<-1\text{或}0<x-1<1\end{cases}$ 或 $\begin{cases}x<2\\x-1>1\text{或}-1<x-1<0\end{cases}$，解得 $0<x<1$.}
\step{故选 $C$.}
\solend
\fi

\item 若函数 $f(x)$ 同时满足：(1) 对于定义域上的任意 $x$，恒有 $f(x)+f(-x)=0$；(2) 对于定义域上的任意 $x_1$、$x_2$，当 $x_1\neq x_2$ 时，恒有 $x_1f(x_1)-x_1f(x_2)>x_2f(x_1)-x_2f(x_2)$，则称函数 $f(x)$ 为“理想函数”. 给出下列四个函数：① $f(x)=x^2$；② $f(x)=x^3$；③ $f(x)=\dfrac{2x-1}{2x+1}$；④ $f(x)=\begin{cases}x^2+4x,x\geqslant0\\-x^2+4x,x<0\end{cases}$，其中被称为“理想函数”的有\mbox{（\quad）}

\keepwithprev
A. 1 个\qquad\qquad B. 2 个\qquad\qquad C. 3 个\qquad\qquad D. 4 个

\ifanswer
\solbegin
\step{函数 $f(x)$ 同时满足：(1) 对于定义域上的任意 $x$，恒有 $f(x)+f(-x)=0$；}
\step{(2) 对于定义域上的任意 $x_1$、$x_2$，当 $x_1\neq x_2$ 时，}
\step{恒有 $x_1f(x_1)-x_1f(x_2)>x_2f(x_1)-x_2f(x_2)$，}
\step{整理得 $(x_1-x_2)[f(x_1)-f(x_2)]>0$，则称函数 $f(x)$ 为“理想函数”.}
\step{即函数满足 (1) 奇函数；(2) 函数在定义域内为严格增函数；}
\step{① $f(x)=x^2$；不满足函数为奇函数在定义域内也不严格增；}
\step{② $f(x)=x^3$；满足函数既为奇函数，也在定义域内严格增，}
\step{故该函数为理想函数；}
\step{③ $f(x)=\dfrac{2x-1}{2x+1}=\dfrac{2x+1-2}{2x+1}=-\dfrac2{2x+1}+1$，该函数不为奇函数，}
\step{但是函数在定义域内严格增，故该函数不为理想函数；}
\step{④ $f(x)=\begin{cases}x^2+4x,x\geqslant0\\-x^2+4x,x<0\end{cases}$，函数的关系式满足 $f(-x)=-f(x)$，}
\step{且在 $x\geqslant0$ 时，函数严格增，由于奇函数在对称区间上单调性相同，}
\step{故该函数为理想函数.}
\step{故选 $B$.}
\solend
\fi

\item 已知 $f(x)$ 是定义在 $\R$ 上的奇函数，当 $x_1$、$x_2\in(0,+\infty)$ 且 $x_1\neq x_2$ 时，都有\newline $\dfrac{x_2f(x_1)-x_1f(x_2)}{x_1x_2(x_1-x_2)}>0$ 成立，$f(2028)=2028$，则不等式 $f(x)-x>0$ 的解集为\mbox{（\quad）}

\keepwithprev
A. $(-\infty,-2028)\cup(2028,+\infty)$\qquad\qquad B. $\left(-\dfrac1{2028},\dfrac1{2028}\right)$

C. $(-2028,2028)$\qquad\qquad D. $(-2028,0)\cup(2028,+\infty)$

\ifanswer
\solbegin
\step{因为 $f(x)$ 是定义在 $\R$ 上的奇函数，所以 $f(-x)=-f(x)$，}
\step{令 $g(x)=\dfrac{f(x)}x$，则 $g(-x)=\dfrac{f(-x)}{-x}=\dfrac{-f(x)}{-x}=\dfrac{f(x)}x=g(x)$，}
\step{所以 $g(x)$ 为偶函数，}
\step{当 $x_1$、$x_2\in(0,+\infty)$ 且 $x_1\neq x_2$ 时，都有 $\dfrac{x_2f(x_1)-x_1f(x_2)}{x_1x_2(x_1-x_2)}>0$ 成立，}
\step{即 $\dfrac{\dfrac{f(x_1)}{x_1}-\dfrac{f(x_2)}{x_2}}{x_1-x_2}>0$，即 $g(x)$ 在 $(0,+\infty)$ 上严格增，}
\step{由偶函数对称性得 $g(x)$ 在 $(-\infty,0)$ 上严格减，}
\step{因为 $f(2028)=2028$，所以 $g(2028)=1$，}
\step{则不等式 $f(x)-x>0$ 可化为 $f(x)>x$，}
\step{当 $x=0$ 时，$f(0)=0$，显然不成立，}
\step{当 $x>0$ 时，不等式可化为 $\dfrac{f(x)}x>1$，即 $g(x)>1=g(2028)$，}
\step{所以 $\lvert x\rvert>2028$，解得 $x>2028$，}
\step{当 $x<0$ 时，不等式可化为 $\dfrac{f(x)}x<1$，即 $g(x)<1=g(2028)$，}
\step{所以 $\lvert x\rvert<2028$，解得 $-2028<x<0$，}
\step{故不等式的解集为 $(-2028,0)\cup(2028,+\infty)$. 故选 $D$.}
\solend
\fi

\end{enumerate}

\bigsec{三、解答题}

\begin{enumerate}
\setcounter{enumi}{16}

\item 已知函数 $f(x)=\dfrac{2x-3}{x+1}$.

(1) 试用函数单调性的定义证明：函数 $f(x)$ 在 $(-1,+\infty)$ 上是严格增函数；

(2) 求函数 $f(x)$ 在区间 $[1,4]$ 上的最大值和最小值.

\ifanswer
\solbegin
\step{(1) 设 $-1<x_1<x_2$，则 $f(x_1)-f(x_2)=\dfrac{2x_1-3}{x_1+1}-\dfrac{2x_2-3}{x_2+1}=\dfrac{5(x_1-x_2)}{(x_1+1)(x_2+1)}$}
\step{因为 $-1<x_1<x_2$，所以 $x_1-x_2<0$，$x_1+1>0$，$x_2+1>0$，}
\step{所以 $f(x_1)-f(x_2)<0$，所以 $f(x)$ 在 $(-1,+\infty)$ 上严格增.}
\step{(2) 由 (1) 得 $f(x)$ 在区间 $[1,4]$ 上严格增，}
\step{所以 $f(x)_{\min}=f(1)=-\dfrac12$，$f(x)_{\max}=f(4)=1$，}
\step{所以函数 $f(x)$ 在区间 $[1,4]$ 上的最大值为 1，最小值为 $-\dfrac12$.}
\solend
\fi

\item 设函数 $f(x)=x^2-mx+m$.

(1) 当 $m=2$ 时，求不等式 $f(x)<x$ 的解集；

(2) 若 $\exists x\in[1,2]$，使不等式 $f(x)<x+m-2$ 成立，求实数 $m$ 的取值范围.

\ifanswer
\solbegin
\step{(1) $m=2$ 时，函数 $f(x)=x^2-2x+2$，}
\step{不等式 $f(x)<x$ 可化为 $x^2-3x+2<0$，解得 $1<x<2$，}
\step{所以不等式的解集为 $\{x\mid1<x<2\}$；}
\step{(2) 若 $\exists x\in[1,2]$，使不等式 $f(x)<x+m-2$ 成立，即 $x^2-x+2<mx$，}
\step{所以 $m>x+\dfrac2x-1$，设 $g(x)=x+\dfrac2x-1$，$x\in[1,2]$，}
\step{则 $g(x)=x+\dfrac2x-1\geqslant2\sqrt{x\cdot\dfrac2x}-1=2\sqrt2-1$，}
\step{当且仅当 $x=\dfrac2x$，即 $x=\sqrt2$ 时取等号，}
\step{所以实数 $m$ 的取值范围是 $\{m\mid m>2\sqrt2-1\}$.}
\solend
\fi

\item 某企业原来生产某种产品 $x$ 万件）可获利 $w(x)$（万元），且满足\newline% 原卷“万件）”右括号无配对、疑漏印左括号“（”，照录
$w(x)=\begin{cases}20(x^2+17),0<x\leqslant2\\500-\dfrac{80}{x-1},2<x\leqslant5\end{cases}$.
现该企业计划引进新的生产设备和新的产品方案优化产品，优化后的产品的其他成本总投入为 $(20x+10)$ 万元. 由市场调研分析得知，当前产品供不应求. 记该企业优化后的产品的利润为 $f(x)$（单位：万元）.

(1) 求函数 $f(x)$ 的解析式；

(2) 当优化后的产品产量为多少万件时，该企业的利润 $f(x)$ 最大？最大利润是多少？请说明理由.

\ifanswer
\solbegin
\step{(1) 由题意得 $f(x)=\begin{cases}20x^2-20x+330,0<x\leqslant2\\-20x-\dfrac{80}{x-1}+490,2<x\leqslant5\end{cases}$；}
\step{(2) 当 $0<x\leqslant2$ 时，$f(x)=20(x^2-x)+330=20\left(x-\dfrac12\right)^2+325$，}
\step{故当 $x=2$ 时，$f(x)$ 取最大值，$f(2)=370$；}
\step{当 $2<x\leqslant5$ 时，$f(x)=-\left[20(x-1)+\dfrac{80}{x-1}\right]+470\leqslant-80+470=390$，}
\step{当且仅当 $20(x-1)=\dfrac{80}{x-1}$，即 $x=3$ 时，$f(3)=390$ 为最大值；}
\step{因为 $370<390$，}
\step{所以优化后的产品产量为 3 万件时，企业获最大利润 390 万元.}
\solend
\fi

\item 俄国数学家切比雪夫（1821—1894）是研究直线逼近函数的理论先驱. 对定义在非空集合 $I$ 上的函数 $f(x)$，以及函数 $g(x)=kx+b(k,b\in\R)$，切比雪夫将函数 $y=\lvert f(x)-g(x)\rvert,x\in I$ 的最大值称为 $f(x),g(x)$ 的“偏差”.

(1) 函数 $f(x)=x^2,x\in[0,1],g(x)=-x-2$，求 $f(x),g(x)$ 的“偏差”；

(2) 函数 $f(x)=\dfrac1x+1,x\in[1,2],g(x)=kx+1(k>0)$，若 $f(x),g(x)$ 的“偏差”为 2，求 $k$ 的值；

(3) 函数 $f(x)=x^2-x,x\in[0,3],g(x)=2x+b$，当 $f(x),g(x)$ 的“偏差”取最小值时，求 $b$ 的值，并求出“偏差”的最小值.

\ifanswer
% 注：原卷本题解析标签印作【解得】（全卷唯一，其余 18 题均印【解析】）；
%     本稿按全稿惯例由 \solbegin 统一印【解析】，正文内容照录。见文件头排版差异⑥。
\solbegin
\step{(1) $y=\lvert f(x)-g(x)\rvert=\lvert x^2+x+2\rvert=\left\lvert\left(x+\dfrac12\right)^2+\dfrac74\right\rvert,x\in[0,1]$，}
\step{因为 $x\in[0,1]$，所以 $\left(x+\dfrac12\right)^2+\dfrac74\in[2,4]$，}
\step{则 $y=\left\lvert\left(x+\dfrac12\right)^2+\dfrac74\right\rvert\in[2,4]$，$y_{\max}=4$，}
\step{所以函数 $f(x)$ 与 $g(x)$ 的“偏差”为 4.}
\step{(2) 令 $t=f(x)-g(x)=\dfrac1x+1-(kx+1)=\dfrac1x-kx,x\in[1,2]$，}
\step{因为 $k>0$，所以 $t=\dfrac1x-kx$ 是单调减函数，}
\step{所以 $t\in\left[\dfrac12-2k,1-k\right]$，由单调，$y=\lvert t\rvert,t\in\left[\dfrac12-2k,1-k\right]$，且 $y_{\max}=2$.}
\step{当 $\left\lvert\dfrac12-2k\right\rvert\leqslant\lvert1-k\rvert$，即 $0<k\leqslant\dfrac12$ 时，$\lvert1-k\rvert=2$，解得 $k=3$ 或 $k=-1$，}
\step{均不符合；}
\step{当 $\left\lvert\dfrac12-2k\right\rvert>\lvert1-k\rvert$，即 $k>\dfrac12$ 时，$\left\lvert\dfrac12-2k\right\rvert=2$，$\dfrac12-2k=2$ 或 $\dfrac12-2k=-2$，}
\step{解得 $k=\dfrac54$ 或 $k=-\dfrac34$（舍），所以 $k=\dfrac54$.}
\step{(3) $y=\lvert f(x)-g(x)\rvert=\lvert x^2-x-(2x+b)\rvert=\lvert x^2-3x-b\rvert$}
\step{$=\left\lvert\left(x-\dfrac32\right)^2-\dfrac94-b\right\rvert,x\in[0,3]$，}
\step{因为 $x\in[0,3]$，所以 $\left(x-\dfrac32\right)^2-\dfrac94-b\in\left[-\dfrac94-b,-b\right]$，}
\step{由 $y=\left\lvert\left(x-\dfrac32\right)^2-\dfrac94-b\right\rvert$，则 $y_{\max}=\max\left\{\lvert b\rvert,\left\lvert b+\dfrac94\right\rvert\right\}$，}
\step{令 $\lvert b\rvert\geqslant\left\lvert b+\dfrac94\right\rvert$，即 $b^2\geqslant\left(b+\dfrac94\right)^2$，解得 $b\leqslant-\dfrac98$，}
\step{$y_{\max}=\max\left\{\lvert b\rvert,\left\lvert b+\dfrac94\right\rvert\right\}=\begin{cases}-b,b\leqslant-\dfrac98\\b+\dfrac94,b>-\dfrac98\end{cases}$，}
\step{故当且仅当 $b=-\dfrac98$ 时，$(y_{\max})_{\min}=\dfrac98$，}
\step{故当 $b$ 的值为 $-\dfrac98$ 时，函数 $f(x)$ 与 $g(x)$ 的“偏差”取最小值 $\dfrac98$.}
\solend
\fi

\item 若函数 $f(x)$ 与 $g(x)$ 满足：对任意的 $x_1\in D$，总存在唯一的 $x_2\in D$，使 $f(x_1)g(x_2)=m$ 成立，则称 $f(x)$ 是 $g(x)$ 在区间 $D$ 上的“$m$ 阶伴随函数”；对任意的 $x_1\in D$，总存在唯一的 $x_2\in D$，使 $f(x_1)f(x_2)=m$ 成立，则称 $f(x)$ 是区间 $D$ 上的“$m$ 阶自伴函数”.

(1) 判断 $f(x)=x^2+1$ 是否为区间 $[0,3]$ 上的“2 阶自伴函数”？并说明理由：

(2) 若函数 $f(x)=3x-1$ 为区间 $\left[\dfrac12,b\right]$ 上的“1 阶自伴函数”，求 $b$ 的值；

(3) 若 $f(x)=\dfrac4{x+2}$ 是 $g(x)=x^2-2ax+a^2-1$ 在区间 $[0,2]$ 上的“2 阶伴随函数”，求实数 $a$ 的取值范围.

\ifanswer
\solbegin
\step{(1) 假定函数 $f(x)=x^2+1$ 是区间 $[0,3]$ 上的“2 阶自伴函数”，}
\step{取 $x_1=2,f(x_1)=f(2)=5$，由 $f(x_1)f(x_2)=2$，得 $5(x_2^2+1)=2$，}
\step{显然此方程无实数解，}
\step{所以函数 $f(x)=x^2+1$ 不是区间 $[0,3]$ 上的“2 阶自伴函数”.}
\step{(2) 函数 $f(x)=3x-1$ 为区间 $\left[\dfrac12,b\right]$ 上的“1 阶自伴函数”，}
\step{则对任意 $x_1\in\left[\dfrac12,b\right]$，总存在唯一的 $x_2\in\left[\dfrac12,b\right]$，使得 $f(x_1)f(x_2)=1$，}
\step{即 $3x_2-1=\dfrac1{3x_1-1}$，整理得 $x_2=\dfrac13+\dfrac1{3(3x_1-1)}$，}
\step{显然函数 $x_2=\dfrac13+\dfrac1{3(3x_1-1)}$ 在 $\left[\dfrac12,b\right]$ 上严格减，}
\step{且当 $x_1=\dfrac12$ 时，$x_2=1$，当 $x_1=b$ 时，$x_2=\dfrac13+\dfrac1{9b-3}$，}
\step{因此对 $\left[\dfrac12,b\right]$ 内的每一个 $x_1$，在 $\left[\dfrac13+\dfrac1{9b-3},1\right]$ 内有唯一 $x_2$ 值与之对应，}
\step{则 $x_2\in\left[\dfrac12,b\right]$，于是 $\left[\dfrac13+\dfrac1{9b-3},1\right]\subseteq\left[\dfrac12,b\right]$，应有 $\begin{cases}b\geqslant1\\\dfrac13+\dfrac1{9b-3}\geqslant\dfrac12\end{cases}$，}
\step{得 $\begin{cases}b\geqslant1\\b\leqslant1\end{cases}$，即 $b=1$，所以 $b$ 的值是 1.}
\step{(3) 由函数 $f(x)=\dfrac4{x+2}$ 在区间 $[0,2]$ 的值域为 $[1,2]$，}
\step{因为 $f(x)=\dfrac4{x+2}$ 是 $g(x)=x^2-2ax+a^2-1$ 在 $[0,2]$ 上的“2 阶伴随函数”，}
\step{则对任意的 $x_1\in[0,2]$，总存在唯一的 $x_2\in[0,2]$ 时，使得 $f(x_1)g(x_2)=2$，}% 原卷此处印“时”字（按上下文为衍字），照录
\step{所以 $g(x_2)=\dfrac2{f(x_1)}\in[1,2]$，}
\step{即 $g(x)=x^2-2ax+a^2-1$ 在 $[0,2]$ 区间上的值域必定包含区间 $[1,2]$，}
\step{且 $g(x)$ 值域在 $[1,2]$ 对应的自变量是唯一的，}
\step{又因为函数 $g(x)=x^2-2ax+a^2-1$ 开口向上，对称轴为 $x=a$，}
\step{1) 当 $a\leqslant0$ 时，$g(x)$ 在 $[0,2]$ 区间上严格增，则必有 $\begin{cases}g(x)_{\min}=g(0)\leqslant1\\g(x)_{\max}=g(2)\geqslant2\end{cases}$，}
\step{解得 $-\sqrt2\leqslant a\leqslant0$；}
\step{2) 当 $a\geqslant2$ 时，$g(x)$ 在 $[0,2]$ 区间上严格减，则必有 $\begin{cases}g(x)_{\min}=g(2)\leqslant1\\g(x)_{\max}=g(0)\geqslant2\end{cases}$，}
\step{解得 $2\leqslant a\leqslant2+\sqrt2$；}
\step{3) 当 $0<a\leqslant1$ 时，$g(x)$ 在 $[0,a]$ 上严格减，在 $[a,2]$ 上严格增，}
\step{则必有 $\begin{cases}g(0)<1\\g(x)_{\max}=g(2)\geqslant2\end{cases}$，解得 $0<a\leqslant2-\sqrt3$，}
\step{4) 当 $1<a<2$ 时，$g(x)$ 在 $[0,a]$ 上严格减，在 $[a,2]$ 上严格增，}
\step{则必有 $\begin{cases}g(2)<1\\g(x)_{\max}=g(0)\geqslant2\end{cases}$，解得 $\sqrt3\leqslant a<2$.}
\step{综上所述，$a$ 的范围 $[-\sqrt2,2-\sqrt3]\cup[\sqrt3,2+\sqrt2]$.}% 原卷“范围”后缺“是／为”字，照录
\solend
\fi

\end{enumerate}

\end{document}
"
%
% 原始材料：微信公众号“魔都数学加油站”文章，作者破晓，2026-11-09 17:52 发布；
% 连载“27学年高一122”，原文标题“27学年高一122——2026-2027年上海市华师大二附中高一上周末作业6”；
% 原文链接：https://mp.weixin.qq.com/s/pC5066JFKNkQQwjMqTK8Ow。
% 正文为 14 张 PNG 页（第 14 页 2560×3620，其余 13 页 4762×6735），
% 题目下印有【解析】，为讲评版。
% 题目、解析均据原图转录；卷面红色水印及公众号页脚未录入。全卷无插图（已逐页检索，无“如图”）。
%
% 卷首标题：原卷印作“2026-2027 年华师大二附中高一上周末作业 6”（公众号标题中的“上海市”
% 卷面标题行没有，本稿照卷面），年份区间按全稿惯例排 en dash；原卷未印副标题，本稿按题目内容
% 标为“函数”。副标题理由（以全部 14 页实际内容为准）：逐页复核 19 题，考的都是函数性质——
% 奇偶性与充要条件（Q3、Q13、Q15、Q16）、单调性与单调区间（Q4、Q7、Q14、Q17）、定义域／
% 值域／最值与复合函数（Q5、Q6、Q8、Q10、Q12、Q19、Q20）、分段函数与取整函数（Q7、Q8、
% Q11、Q15、Q19）、抽象函数与新定义（Q12、Q20、Q21）；不等式仅作求参数范围／解集的工具
% （Q9、Q14、Q16、Q18），集合与常用逻辑用语仅零星出现（Q3 的充要条件、Q11③ 的集合、
% Q20 的“非空集合 I”），均非独立章节。故主体为单一模块“函数”，与同校 高一102 周末作业 5
% 同（该函数模块卷亦按 高一83／105 成例只用模块名、不叠“集合与常用逻辑用语、不等式”）。
%
% 与原卷的排版差异：
%   ① 原文从第 3 题起，填空题编号为 3—12、选择题为 13—16、解答题为 17—21，本稿保留原题号；
%   ② 原卷每题下直接印【解析】，本稿按全稿惯例置于答案版；四道选择题括号内均未印答案、
%      答案在解析末句“故选 …”，括号内留空（全卷无只印【答案】的题、无 \ans{}）；
%   ③ 原卷未印副标题，本稿按“知识模块”口径补出；
%   ④ 原卷填空横线按全稿惯例排为 \blankline[4.4em]；选择题括号排 \mbox{（\quad）}、
%      选项段前按惯例加 \keepwithprev；解答题不留作答区（原卷为讲评版）；
%   ⑤ 原卷数集排斜体 $R$、$N_+$（下标加号），按全稿惯例改 \R、\N_+
%      （$N_+$ 下标加号与 $\N^*$ 是两种记号，照录下标、只改数集字母为正体，同 高一05 第 20 题成例）；
%      绝对值排 \lvert…\rvert、集合描述竖线排 \mid；
%      $f(x)_{min}$／$f(x)_{max}$／$y_{max}$／$(y_{max})_{min}$ 一类斜体下标按 高一96 成例
%      改正体 \min／\max；不等式一律排 \geqslant／\leqslant；
%   ⑥ 第 20 题解析标签原卷印作【解得】（全卷唯一，其余 18 题均印【解析】），本稿按全稿惯例
%      由 \solbegin 统一印【解析】，正文内容照录（行内加 % 注、并入文件头存疑清单）；
%   ⑦ 题干与解析的标点（半角／全角逗号、冒号、括号、句点、子题括号）均照原卷，
%      题目正文引号用弯引号“ ”。
%
% 子集符号：第 21 题 (2) 的 $\left[\frac13+\frac1{9b-3},1\right]\subseteq\left[\frac12,b\right]$
%   印有下划线，照录 \subseteq。全卷未出现判别式 \triangle／\Delta。
%
% 原卷存疑／笔误（均照录加注）：
%   第 11 题解析对② 的“$y=y+\Delta y$”（按上下文应为 $y=[y]+\Delta y$，原卷漏方括号）；
%   第 19 题题干“某种产品 $x$ 万件）可获利”右括号无配对（疑漏印左括号“（”）；
%   第 20 题解析标签【解得】（见排版差异⑥）；
%   第 21 题 (3) 解析“总存在唯一的 $x_2\in[0,2]$ 时，使得……”的“时”字（按上下文衍字）；
%   第 21 题 (3) 末句“综上所述，$a$ 的范围 $[-\sqrt2,2-\sqrt3]\cup[\sqrt3,2+\sqrt2]$”
%      “范围”后缺“是／为”字，照录。
%
% 行数核对：逐题数 \step 条数，合计 162 条（Q3：5；Q4：3；Q5：1；Q6：3；Q7：2；
% Q8：5；Q9：7；Q10：10；Q11：8；Q12：4；Q13：12；Q14：6；Q15：14；Q16：14；
% Q17：6；Q18：8；Q19：7；Q20：19；Q21：28），与原卷印刷行逐题相等
% （cases 块作一个印刷行计）。全卷无 \ans{}、无只印【答案】的题。
%
\documentclass[a4paper,11pt]{article}
\usepackage{common}

\begin{document}

\examtitlest[3]{2026--2027 年华师大二附中高一上周末作业 6}{函数}

\bigsec{一、填空题}

\begin{enumerate}
\setcounter{enumi}{2}

\item 已知 $p:a=0$，$q$：函数 $f(x)=x^3+x+a$ 是奇函数，则 $p$ 是 $q$ 的\blankline[4.4em]条件.

\ifanswer
\solbegin
\step{若 $a=0$，则 $f(x)=x^3+x$，}
\step{易得 $f(x)$ 的定义域为 $\R$，且 $f(-x)=-(x^3+x)=-f(x)$，则 $f(x)$ 为奇函数，}
\step{反之，若 $f(x)=x^3+x+a$ 是奇函数，}
\step{则有 $f(-x)+f(x)=-x^3-x+a+x^3+x+a=2a=0$，必有 $a=0$，}
\step{故 $p$ 是 $q$ 的充要条件.}
\solend
\fi

\item 函数 $y=\sqrt{-x^2+2x}$ 的单调递增区间是\blankline[4.4em].

\ifanswer
\solbegin
\step{由函数定义域得 $-x^2+2x\geqslant0$，解得 $0\leqslant x\leqslant2$，}
\step{二次函数 $y=-x^2+2x$ 开口向下，对称轴为 $x=1$，}
\step{结合复合函数的单调性得函数 $y=\sqrt{-x^2+2x}$ 的单调递增区间是 $[0,1]$.}
\solend
\fi

\item 设 $f(x)$ 定义域为 $[1,4]$，且 $f(x)=2x+1$，则 $f(x^2)$ 的最大值为\blankline[4.4em].

\ifanswer
\solbegin
\step{由 $x^2\in[1,4]$，得 $f(x^2)=2x^2+1\leqslant9$，当 $x=\pm2$ 时取等号，最大值为 9.}
\solend
\fi

\item 函数 $f(x)=\sqrt{x-2}-x$，$x\in(2,6]$ 的值域为\blankline[4.4em].

\ifanswer
\solbegin
\step{令 $t=\sqrt{x-2}$，则 $x=2+t^2$，由 $2<x\leqslant6$ 得 $0<t\leqslant2$，}
\step{原函数可化为 $g(t)=-t^2+t-2$，$g(0)=-2$，$g(2)=-4$，$g\left(\dfrac12\right)=-\dfrac74$，}
\step{由二次函数性质得 $g(t)\in\left[-4,-\dfrac74\right]$.}
\solend
\fi

\item 若函数 $f(x)=\lvert x-a+1\rvert$ 在区间 $[1,+\infty)$ 上是严格增函数，则实数 $a$ 的取值范围是\blankline[4.4em].

\ifanswer
\solbegin
\step{$f(x)=\begin{cases}x-a+1,x\geqslant a-1\\a-1-x,x<a-1\end{cases}$，显然 $f(x)$ 在 $[a-1,+\infty)$ 上严格增，}
\step{要使函数在区间 $[1,+\infty)$ 上是严格增函数，则 $a-1\leqslant1$，即 $a\leqslant2$.}
\solend
\fi

\item 若函数 $f(x)=\begin{cases}x^2-5,x\leqslant2\\2ax+1,x>2\end{cases}$ 的值域是 $\R$，则 $a$ 的取值范围是\blankline[4.4em].

\ifanswer
\solbegin
\step{当 $x\leqslant2$ 时，$f(x)=x^2-5$，其开口向上，对称轴为 $x=0$，值域为 $[-5,+\infty)$，}
\step{由函数 $f(x)=\begin{cases}x^2-5,x\leqslant2\\2ax+1,x>2\end{cases}$ 的值域是 $\R$，}
\step{则当 $x>2$ 时，$f(x)=2ax+1$ 的值域应包含 $(-\infty,-5)$，}
\step{所以 $f(x)$ 为严格减函数，所以 $\begin{cases}2\times2a+1\geqslant-5\\2a<0\end{cases}$，即 $\begin{cases}a\geqslant-\dfrac32\\a<0\end{cases}$，}
\step{解得 $-\dfrac32\leqslant a<0$，故 $a$ 的取值范围是 $\left[-\dfrac32,0\right)$.}
\solend
\fi

\item 已知函数 $f(x)=ax^2+\dfrac bx$ 的图象经过点 $(-1,0)$，且满足 $f(1)=2$. 若 $\forall x\in(0,+\infty)$，都有 $f(x)+f\left(\dfrac1x\right)-2m\geqslant0$ 恒成立，则实数 $m$ 的取值范围为\blankline[4.4em].

\ifanswer
\solbegin
\step{因为 $f(x)=ax^2+\dfrac bx$ 的图象经过点 $(-1,0)$，且满足 $f(1)=2$，}
\step{所以 $\begin{cases}a-b=0\\a+b=2\end{cases}$，所以 $a=1$，$b=1$，即 $f(x)=x^2+\dfrac1x$，}
\step{$\forall x\in(0,+\infty)$，$f(x)+f\left(\dfrac1x\right)-2m\geqslant0\Leftrightarrow2m\leqslant x^2+\dfrac1x+\dfrac1{x^2}+x$ 恒成立，}
\step{令函数 $g(x)=x^2+\dfrac1x+\dfrac1{x^2}+x=\left(x+\dfrac1x\right)^2+x+\dfrac1x-2$，}
\step{而当 $x>0$ 时，$x+\dfrac1x\geqslant2$，当且仅当 $x=1$ 时取等号，}
\step{因此 $g(x)\geqslant2^2+2-2=4$，当且仅当 $x=1$ 时取等号，则 $2m\leqslant4$，解得 $m\leqslant2$，}
\step{所以实数 $m\in(-\infty,2]$.}
\solend
\fi

\item 若定义在区间 $[a,b](a<b)$ 上的函数 $f(x)=k-\sqrt{x+1}$ 值域也为 $[a,b]$，则实数 $k$ 的取值范围是\blankline[4.4em].

\ifanswer
\solbegin
\step{函数 $f(x)=k-\sqrt{x+1}$ 在定义域 $[-1,+\infty)$ 严格减，}
\step{当 $f(x)$ 的定义域为 $[a,b]$ 时，$f(x)$ 的值域也为 $[a,b]$，$a<b$，}
\step{故 $f(a)=k-\sqrt{a+1}=b$，$f(b)=k-\sqrt{b+1}=a$，}
\step{两式相减得 $\sqrt{a+1}-\sqrt{b+1}=a-b=(a+1)-(b+1)$}
\step{$=(\sqrt{a+1}-\sqrt{b+1})(\sqrt{a+1}+\sqrt{b+1})$，}
\step{所以 $\sqrt{a+1}-\sqrt{b+1}<0$，即 $\sqrt{a+1}+\sqrt{b+1}=1$，}
\step{则 $k=\sqrt{b+1}+a=a+1-\sqrt{a+1}$，}
\step{令 $\lambda=\sqrt{a+1}\geqslant0$，得 $k=\lambda^2-\lambda=\left(\lambda-\dfrac12\right)^2-\dfrac14$，}
\step{又 $\sqrt{a+1}<\sqrt{b+1}$，因为 $\sqrt{a+1}+\sqrt{b+1}=1$，所以 $\sqrt{a+1}\in\left[0,\dfrac12\right)$，}
\step{所以 $0\leqslant\lambda<\dfrac12$，故实数 $k$ 的取值范围为 $\left(-\dfrac14,0\right]$.}
\solend
\fi

\item 函数 $f(x)=[x]$ 的函数值表示不超过 $x$ 的最大整数，例如：$[-3.5]=-4$，$[2.1]=2$，称 $f(x)=[x]$ 为取整函数，则下列关于取整函数的说法：①对任意的 $x\in\R$，$[-x]=-[x]$ ②对任意的 $x$、$y\in\R$，$[x]+[y]\leqslant[x+y]$ ③集合 $\{y\mid y=[nx]-n[x],x\in\R,n\in\N_+\}$ 共有 $n$ 个元素 ④$x>0$ 时，关于 $x$ 的方程 $[2x]=\left[\dfrac1x\right]$ 有无数个解 其中正确的有\blankline[4.4em].

\ifanswer
\solbegin
\step{对①，因为 $[-2.1]=-3\neq[2.1]=2$，所以①错误；}
\step{对②，设 $x=[x]+\Delta x$，$\Delta x\in[0,1)$，$y=y+\Delta y$，$\Delta y\in[0,1)$，}% 原卷印作 $y=y+\Delta y$（按上下文应为 $y=[y]+\Delta y$），照录
\step{$[x+y]=[[x]+\Delta x+[y]+\Delta y]=[x]+[y]+[\Delta x+\Delta y]$，又 $\Delta x+\Delta y\in[0,\ 2)$，}
\step{故 $\Delta x+\Delta y=0$ 或 1，故 $[x]+[y]\leqslant[x+y]$，所以②正确；}
\step{对③，$[nx]-n[x]=[n[x]+n\cdot\Delta x]-n[x]=n[x]+[n\cdot\Delta x]-n[x]=[n\cdot\Delta x]$，}
\step{因为 $\Delta x\in[0,1)$，所以 $0\leqslant n\cdot\Delta x<n$，$n\in\N_+$，所以③正确；}
\step{对④，当 $\dfrac12<x<1$ 时，$[2x]=\left[\dfrac1x\right]=1$，故方程存在无数个解，所以④正确.}
\step{故正确的是②③④.}
\solend
\fi

\item 已知函数 $f(x)$ 定义域为 $[0,1]$，且满足 $f(x)+f(1-x)=1$，$f(x)=2f\left(\dfrac x5\right)$，对任意 $x_1\leqslant x_2$，恒有 $f(x_1)\leqslant f(x_2)$，则 $f\left(\dfrac1{2026}\right)$ 的值为\blankline[4.4em].

\ifanswer
\solbegin
\step{$f(0)=2f(0)\Rightarrow f(0)=0$，从而 $f(1)=1$，$f\left(\dfrac15\right)=\dfrac12$，$f\left(\dfrac45\right)=1-f\left(\dfrac15\right)=\dfrac12$，}
\step{由单调性，$x\in\left[\dfrac15,\dfrac45\right]$ 时，$f(x)=\dfrac12$.}
\step{因为 $\dfrac15<\dfrac{625}{2026}<\dfrac45$，反复使用 $f\left(\dfrac x5\right)=\dfrac12 f(x)$，}
\step{得 $f\left(\dfrac1{2026}\right)=\dfrac1{2^4}f\left(\dfrac{625}{2026}\right)=\dfrac1{32}$.}
\solend
\fi

\end{enumerate}

\bigsec{二、选择题}

\begin{enumerate}
\setcounter{enumi}{12}

\item 已知 $f(x)$ 为 $\R$ 上的奇函数，$g(x)$ 为 $\R$ 上的偶函数，且 $g(x)\neq0$，则下列说法正确的是\mbox{（\quad）}

\keepwithprev
A. $f(x)+g(x)$ 为 $\R$ 上的奇函数\qquad\qquad B. $f(x)-g(x)$ 为 $\R$ 上的奇函数

C. $\dfrac{f(x)}{g(x)}$ 为 $\R$ 上的偶函数\qquad\qquad D. $\lvert f(x)g(x)\rvert$ 为 $\R$ 上的偶函数

\ifanswer
\solbegin
\step{因为 $f(x)$ 为 $\R$ 上的奇函数，$g(x)$ 为 $\R$ 上的偶函数，且 $g(x)\neq0$，}
\step{所以 $f(-x)=-f(x)$，$g(-x)=g(x)$，}
\step{所以 $f(-x)+g(-x)=-f(x)+g(x)\neq-[f(x)+g(x)]$，}
\step{故 $f(x)+g(x)$ 为非奇非偶函数，$A$ 错误；}
\step{同理，$f(x)-g(x)$ 为非奇非偶函数，$B$ 错误；}
\step{设 $F(x)=\dfrac{f(x)}{g(x)}$，则 $F(-x)=\dfrac{f(-x)}{g(-x)}=\dfrac{-f(x)}{g(x)}=-F(x)$，}
\step{所以 $F(x)$ 为奇函数，$C$ 错误；}
\step{设函数 $H(x)=\lvert f(x)g(x)\rvert$，}
\step{因为 $f(x)$ 为 $\R$ 上的奇函数，$g(x)$ 为 $\R$ 上的偶函数，且 $g(x)\neq0$，}
\step{所以 $f(-x)=-f(x)$，$g(-x)=g(x)$，则由函数奇偶性的定义}
\step{得 $H(-x)=\lvert f(-x)g(-x)\rvert=\lvert -f(x)g(x)\rvert=\lvert f(x)g(x)\rvert=H(x)$，$D$ 正确.}
\step{故选 $D$.}
\solend
\fi

\item 若函数 $f(x)$ 是定义域为 $(-\infty,0)\cup(0,+\infty)$ 的奇函数，在区间 $(0,+\infty)$ 上严格增，且 $f(-1)=0$，则不等式 $\dfrac{f(x-1)}{x-2}<0$ 的解集为\mbox{（\quad）}

\keepwithprev
A. $(-\infty,-2)$\qquad\qquad B. $(-2,0)$\qquad\qquad C. $(0,1)$\qquad\qquad D. $(1,2)$

\ifanswer
\solbegin
\step{若函数 $f(x)$ 是定义域为 $(-\infty,0)\cup(0,+\infty)$ 的奇函数，在区间 $(0,+\infty)$ 上严格增，}
\step{由奇函数对称性得，$f(x)$ 在 $(-\infty,0)$ 上严格增，}
\step{因为 $f(-1)=0$，所以 $f(1)=0$，}
\step{则不等式 $\dfrac{f(x-1)}{x-2}<0$ 可转化为 $\begin{cases}x>2\\f(x-1)<0\end{cases}$ 或 $\begin{cases}x<2\\f(x-1)>0\end{cases}$，}
\step{即 $\begin{cases}x>2\\x-1<-1\text{或}0<x-1<1\end{cases}$ 或 $\begin{cases}x<2\\x-1>1\text{或}-1<x-1<0\end{cases}$，解得 $0<x<1$.}
\step{故选 $C$.}
\solend
\fi

\item 若函数 $f(x)$ 同时满足：(1) 对于定义域上的任意 $x$，恒有 $f(x)+f(-x)=0$；(2) 对于定义域上的任意 $x_1$、$x_2$，当 $x_1\neq x_2$ 时，恒有 $x_1f(x_1)-x_1f(x_2)>x_2f(x_1)-x_2f(x_2)$，则称函数 $f(x)$ 为“理想函数”. 给出下列四个函数：① $f(x)=x^2$；② $f(x)=x^3$；③ $f(x)=\dfrac{2x-1}{2x+1}$；④ $f(x)=\begin{cases}x^2+4x,x\geqslant0\\-x^2+4x,x<0\end{cases}$，其中被称为“理想函数”的有\mbox{（\quad）}

\keepwithprev
A. 1 个\qquad\qquad B. 2 个\qquad\qquad C. 3 个\qquad\qquad D. 4 个

\ifanswer
\solbegin
\step{函数 $f(x)$ 同时满足：(1) 对于定义域上的任意 $x$，恒有 $f(x)+f(-x)=0$；}
\step{(2) 对于定义域上的任意 $x_1$、$x_2$，当 $x_1\neq x_2$ 时，}
\step{恒有 $x_1f(x_1)-x_1f(x_2)>x_2f(x_1)-x_2f(x_2)$，}
\step{整理得 $(x_1-x_2)[f(x_1)-f(x_2)]>0$，则称函数 $f(x)$ 为“理想函数”.}
\step{即函数满足 (1) 奇函数；(2) 函数在定义域内为严格增函数；}
\step{① $f(x)=x^2$；不满足函数为奇函数在定义域内也不严格增；}
\step{② $f(x)=x^3$；满足函数既为奇函数，也在定义域内严格增，}
\step{故该函数为理想函数；}
\step{③ $f(x)=\dfrac{2x-1}{2x+1}=\dfrac{2x+1-2}{2x+1}=-\dfrac2{2x+1}+1$，该函数不为奇函数，}
\step{但是函数在定义域内严格增，故该函数不为理想函数；}
\step{④ $f(x)=\begin{cases}x^2+4x,x\geqslant0\\-x^2+4x,x<0\end{cases}$，函数的关系式满足 $f(-x)=-f(x)$，}
\step{且在 $x\geqslant0$ 时，函数严格增，由于奇函数在对称区间上单调性相同，}
\step{故该函数为理想函数.}
\step{故选 $B$.}
\solend
\fi

\item 已知 $f(x)$ 是定义在 $\R$ 上的奇函数，当 $x_1$、$x_2\in(0,+\infty)$ 且 $x_1\neq x_2$ 时，都有\newline $\dfrac{x_2f(x_1)-x_1f(x_2)}{x_1x_2(x_1-x_2)}>0$ 成立，$f(2028)=2028$，则不等式 $f(x)-x>0$ 的解集为\mbox{（\quad）}

\keepwithprev
A. $(-\infty,-2028)\cup(2028,+\infty)$\qquad\qquad B. $\left(-\dfrac1{2028},\dfrac1{2028}\right)$

C. $(-2028,2028)$\qquad\qquad D. $(-2028,0)\cup(2028,+\infty)$

\ifanswer
\solbegin
\step{因为 $f(x)$ 是定义在 $\R$ 上的奇函数，所以 $f(-x)=-f(x)$，}
\step{令 $g(x)=\dfrac{f(x)}x$，则 $g(-x)=\dfrac{f(-x)}{-x}=\dfrac{-f(x)}{-x}=\dfrac{f(x)}x=g(x)$，}
\step{所以 $g(x)$ 为偶函数，}
\step{当 $x_1$、$x_2\in(0,+\infty)$ 且 $x_1\neq x_2$ 时，都有 $\dfrac{x_2f(x_1)-x_1f(x_2)}{x_1x_2(x_1-x_2)}>0$ 成立，}
\step{即 $\dfrac{\dfrac{f(x_1)}{x_1}-\dfrac{f(x_2)}{x_2}}{x_1-x_2}>0$，即 $g(x)$ 在 $(0,+\infty)$ 上严格增，}
\step{由偶函数对称性得 $g(x)$ 在 $(-\infty,0)$ 上严格减，}
\step{因为 $f(2028)=2028$，所以 $g(2028)=1$，}
\step{则不等式 $f(x)-x>0$ 可化为 $f(x)>x$，}
\step{当 $x=0$ 时，$f(0)=0$，显然不成立，}
\step{当 $x>0$ 时，不等式可化为 $\dfrac{f(x)}x>1$，即 $g(x)>1=g(2028)$，}
\step{所以 $\lvert x\rvert>2028$，解得 $x>2028$，}
\step{当 $x<0$ 时，不等式可化为 $\dfrac{f(x)}x<1$，即 $g(x)<1=g(2028)$，}
\step{所以 $\lvert x\rvert<2028$，解得 $-2028<x<0$，}
\step{故不等式的解集为 $(-2028,0)\cup(2028,+\infty)$. 故选 $D$.}
\solend
\fi

\end{enumerate}

\bigsec{三、解答题}

\begin{enumerate}
\setcounter{enumi}{16}

\item 已知函数 $f(x)=\dfrac{2x-3}{x+1}$.

(1) 试用函数单调性的定义证明：函数 $f(x)$ 在 $(-1,+\infty)$ 上是严格增函数；

(2) 求函数 $f(x)$ 在区间 $[1,4]$ 上的最大值和最小值.

\ifanswer
\solbegin
\step{(1) 设 $-1<x_1<x_2$，则 $f(x_1)-f(x_2)=\dfrac{2x_1-3}{x_1+1}-\dfrac{2x_2-3}{x_2+1}=\dfrac{5(x_1-x_2)}{(x_1+1)(x_2+1)}$}
\step{因为 $-1<x_1<x_2$，所以 $x_1-x_2<0$，$x_1+1>0$，$x_2+1>0$，}
\step{所以 $f(x_1)-f(x_2)<0$，所以 $f(x)$ 在 $(-1,+\infty)$ 上严格增.}
\step{(2) 由 (1) 得 $f(x)$ 在区间 $[1,4]$ 上严格增，}
\step{所以 $f(x)_{\min}=f(1)=-\dfrac12$，$f(x)_{\max}=f(4)=1$，}
\step{所以函数 $f(x)$ 在区间 $[1,4]$ 上的最大值为 1，最小值为 $-\dfrac12$.}
\solend
\fi

\item 设函数 $f(x)=x^2-mx+m$.

(1) 当 $m=2$ 时，求不等式 $f(x)<x$ 的解集；

(2) 若 $\exists x\in[1,2]$，使不等式 $f(x)<x+m-2$ 成立，求实数 $m$ 的取值范围.

\ifanswer
\solbegin
\step{(1) $m=2$ 时，函数 $f(x)=x^2-2x+2$，}
\step{不等式 $f(x)<x$ 可化为 $x^2-3x+2<0$，解得 $1<x<2$，}
\step{所以不等式的解集为 $\{x\mid1<x<2\}$；}
\step{(2) 若 $\exists x\in[1,2]$，使不等式 $f(x)<x+m-2$ 成立，即 $x^2-x+2<mx$，}
\step{所以 $m>x+\dfrac2x-1$，设 $g(x)=x+\dfrac2x-1$，$x\in[1,2]$，}
\step{则 $g(x)=x+\dfrac2x-1\geqslant2\sqrt{x\cdot\dfrac2x}-1=2\sqrt2-1$，}
\step{当且仅当 $x=\dfrac2x$，即 $x=\sqrt2$ 时取等号，}
\step{所以实数 $m$ 的取值范围是 $\{m\mid m>2\sqrt2-1\}$.}
\solend
\fi

\item 某企业原来生产某种产品 $x$ 万件）可获利 $w(x)$（万元），且满足\newline% 原卷“万件）”右括号无配对、疑漏印左括号“（”，照录
$w(x)=\begin{cases}20(x^2+17),0<x\leqslant2\\500-\dfrac{80}{x-1},2<x\leqslant5\end{cases}$.
现该企业计划引进新的生产设备和新的产品方案优化产品，优化后的产品的其他成本总投入为 $(20x+10)$ 万元. 由市场调研分析得知，当前产品供不应求. 记该企业优化后的产品的利润为 $f(x)$（单位：万元）.

(1) 求函数 $f(x)$ 的解析式；

(2) 当优化后的产品产量为多少万件时，该企业的利润 $f(x)$ 最大？最大利润是多少？请说明理由.

\ifanswer
\solbegin
\step{(1) 由题意得 $f(x)=\begin{cases}20x^2-20x+330,0<x\leqslant2\\-20x-\dfrac{80}{x-1}+490,2<x\leqslant5\end{cases}$；}
\step{(2) 当 $0<x\leqslant2$ 时，$f(x)=20(x^2-x)+330=20\left(x-\dfrac12\right)^2+325$，}
\step{故当 $x=2$ 时，$f(x)$ 取最大值，$f(2)=370$；}
\step{当 $2<x\leqslant5$ 时，$f(x)=-\left[20(x-1)+\dfrac{80}{x-1}\right]+470\leqslant-80+470=390$，}
\step{当且仅当 $20(x-1)=\dfrac{80}{x-1}$，即 $x=3$ 时，$f(3)=390$ 为最大值；}
\step{因为 $370<390$，}
\step{所以优化后的产品产量为 3 万件时，企业获最大利润 390 万元.}
\solend
\fi

\item 俄国数学家切比雪夫（1821—1894）是研究直线逼近函数的理论先驱. 对定义在非空集合 $I$ 上的函数 $f(x)$，以及函数 $g(x)=kx+b(k,b\in\R)$，切比雪夫将函数 $y=\lvert f(x)-g(x)\rvert,x\in I$ 的最大值称为 $f(x),g(x)$ 的“偏差”.

(1) 函数 $f(x)=x^2,x\in[0,1],g(x)=-x-2$，求 $f(x),g(x)$ 的“偏差”；

(2) 函数 $f(x)=\dfrac1x+1,x\in[1,2],g(x)=kx+1(k>0)$，若 $f(x),g(x)$ 的“偏差”为 2，求 $k$ 的值；

(3) 函数 $f(x)=x^2-x,x\in[0,3],g(x)=2x+b$，当 $f(x),g(x)$ 的“偏差”取最小值时，求 $b$ 的值，并求出“偏差”的最小值.

\ifanswer
% 注：原卷本题解析标签印作【解得】（全卷唯一，其余 18 题均印【解析】）；
%     本稿按全稿惯例由 \solbegin 统一印【解析】，正文内容照录。见文件头排版差异⑥。
\solbegin
\step{(1) $y=\lvert f(x)-g(x)\rvert=\lvert x^2+x+2\rvert=\left\lvert\left(x+\dfrac12\right)^2+\dfrac74\right\rvert,x\in[0,1]$，}
\step{因为 $x\in[0,1]$，所以 $\left(x+\dfrac12\right)^2+\dfrac74\in[2,4]$，}
\step{则 $y=\left\lvert\left(x+\dfrac12\right)^2+\dfrac74\right\rvert\in[2,4]$，$y_{\max}=4$，}
\step{所以函数 $f(x)$ 与 $g(x)$ 的“偏差”为 4.}
\step{(2) 令 $t=f(x)-g(x)=\dfrac1x+1-(kx+1)=\dfrac1x-kx,x\in[1,2]$，}
\step{因为 $k>0$，所以 $t=\dfrac1x-kx$ 是单调减函数，}
\step{所以 $t\in\left[\dfrac12-2k,1-k\right]$，由单调，$y=\lvert t\rvert,t\in\left[\dfrac12-2k,1-k\right]$，且 $y_{\max}=2$.}
\step{当 $\left\lvert\dfrac12-2k\right\rvert\leqslant\lvert1-k\rvert$，即 $0<k\leqslant\dfrac12$ 时，$\lvert1-k\rvert=2$，解得 $k=3$ 或 $k=-1$，}
\step{均不符合；}
\step{当 $\left\lvert\dfrac12-2k\right\rvert>\lvert1-k\rvert$，即 $k>\dfrac12$ 时，$\left\lvert\dfrac12-2k\right\rvert=2$，$\dfrac12-2k=2$ 或 $\dfrac12-2k=-2$，}
\step{解得 $k=\dfrac54$ 或 $k=-\dfrac34$（舍），所以 $k=\dfrac54$.}
\step{(3) $y=\lvert f(x)-g(x)\rvert=\lvert x^2-x-(2x+b)\rvert=\lvert x^2-3x-b\rvert$}
\step{$=\left\lvert\left(x-\dfrac32\right)^2-\dfrac94-b\right\rvert,x\in[0,3]$，}
\step{因为 $x\in[0,3]$，所以 $\left(x-\dfrac32\right)^2-\dfrac94-b\in\left[-\dfrac94-b,-b\right]$，}
\step{由 $y=\left\lvert\left(x-\dfrac32\right)^2-\dfrac94-b\right\rvert$，则 $y_{\max}=\max\left\{\lvert b\rvert,\left\lvert b+\dfrac94\right\rvert\right\}$，}
\step{令 $\lvert b\rvert\geqslant\left\lvert b+\dfrac94\right\rvert$，即 $b^2\geqslant\left(b+\dfrac94\right)^2$，解得 $b\leqslant-\dfrac98$，}
\step{$y_{\max}=\max\left\{\lvert b\rvert,\left\lvert b+\dfrac94\right\rvert\right\}=\begin{cases}-b,b\leqslant-\dfrac98\\b+\dfrac94,b>-\dfrac98\end{cases}$，}
\step{故当且仅当 $b=-\dfrac98$ 时，$(y_{\max})_{\min}=\dfrac98$，}
\step{故当 $b$ 的值为 $-\dfrac98$ 时，函数 $f(x)$ 与 $g(x)$ 的“偏差”取最小值 $\dfrac98$.}
\solend
\fi

\item 若函数 $f(x)$ 与 $g(x)$ 满足：对任意的 $x_1\in D$，总存在唯一的 $x_2\in D$，使 $f(x_1)g(x_2)=m$ 成立，则称 $f(x)$ 是 $g(x)$ 在区间 $D$ 上的“$m$ 阶伴随函数”；对任意的 $x_1\in D$，总存在唯一的 $x_2\in D$，使 $f(x_1)f(x_2)=m$ 成立，则称 $f(x)$ 是区间 $D$ 上的“$m$ 阶自伴函数”.

(1) 判断 $f(x)=x^2+1$ 是否为区间 $[0,3]$ 上的“2 阶自伴函数”？并说明理由：

(2) 若函数 $f(x)=3x-1$ 为区间 $\left[\dfrac12,b\right]$ 上的“1 阶自伴函数”，求 $b$ 的值；

(3) 若 $f(x)=\dfrac4{x+2}$ 是 $g(x)=x^2-2ax+a^2-1$ 在区间 $[0,2]$ 上的“2 阶伴随函数”，求实数 $a$ 的取值范围.

\ifanswer
\solbegin
\step{(1) 假定函数 $f(x)=x^2+1$ 是区间 $[0,3]$ 上的“2 阶自伴函数”，}
\step{取 $x_1=2,f(x_1)=f(2)=5$，由 $f(x_1)f(x_2)=2$，得 $5(x_2^2+1)=2$，}
\step{显然此方程无实数解，}
\step{所以函数 $f(x)=x^2+1$ 不是区间 $[0,3]$ 上的“2 阶自伴函数”.}
\step{(2) 函数 $f(x)=3x-1$ 为区间 $\left[\dfrac12,b\right]$ 上的“1 阶自伴函数”，}
\step{则对任意 $x_1\in\left[\dfrac12,b\right]$，总存在唯一的 $x_2\in\left[\dfrac12,b\right]$，使得 $f(x_1)f(x_2)=1$，}
\step{即 $3x_2-1=\dfrac1{3x_1-1}$，整理得 $x_2=\dfrac13+\dfrac1{3(3x_1-1)}$，}
\step{显然函数 $x_2=\dfrac13+\dfrac1{3(3x_1-1)}$ 在 $\left[\dfrac12,b\right]$ 上严格减，}
\step{且当 $x_1=\dfrac12$ 时，$x_2=1$，当 $x_1=b$ 时，$x_2=\dfrac13+\dfrac1{9b-3}$，}
\step{因此对 $\left[\dfrac12,b\right]$ 内的每一个 $x_1$，在 $\left[\dfrac13+\dfrac1{9b-3},1\right]$ 内有唯一 $x_2$ 值与之对应，}
\step{则 $x_2\in\left[\dfrac12,b\right]$，于是 $\left[\dfrac13+\dfrac1{9b-3},1\right]\subseteq\left[\dfrac12,b\right]$，应有 $\begin{cases}b\geqslant1\\\dfrac13+\dfrac1{9b-3}\geqslant\dfrac12\end{cases}$，}
\step{得 $\begin{cases}b\geqslant1\\b\leqslant1\end{cases}$，即 $b=1$，所以 $b$ 的值是 1.}
\step{(3) 由函数 $f(x)=\dfrac4{x+2}$ 在区间 $[0,2]$ 的值域为 $[1,2]$，}
\step{因为 $f(x)=\dfrac4{x+2}$ 是 $g(x)=x^2-2ax+a^2-1$ 在 $[0,2]$ 上的“2 阶伴随函数”，}
\step{则对任意的 $x_1\in[0,2]$，总存在唯一的 $x_2\in[0,2]$ 时，使得 $f(x_1)g(x_2)=2$，}% 原卷此处印“时”字（按上下文为衍字），照录
\step{所以 $g(x_2)=\dfrac2{f(x_1)}\in[1,2]$，}
\step{即 $g(x)=x^2-2ax+a^2-1$ 在 $[0,2]$ 区间上的值域必定包含区间 $[1,2]$，}
\step{且 $g(x)$ 值域在 $[1,2]$ 对应的自变量是唯一的，}
\step{又因为函数 $g(x)=x^2-2ax+a^2-1$ 开口向上，对称轴为 $x=a$，}
\step{1) 当 $a\leqslant0$ 时，$g(x)$ 在 $[0,2]$ 区间上严格增，则必有 $\begin{cases}g(x)_{\min}=g(0)\leqslant1\\g(x)_{\max}=g(2)\geqslant2\end{cases}$，}
\step{解得 $-\sqrt2\leqslant a\leqslant0$；}
\step{2) 当 $a\geqslant2$ 时，$g(x)$ 在 $[0,2]$ 区间上严格减，则必有 $\begin{cases}g(x)_{\min}=g(2)\leqslant1\\g(x)_{\max}=g(0)\geqslant2\end{cases}$，}
\step{解得 $2\leqslant a\leqslant2+\sqrt2$；}
\step{3) 当 $0<a\leqslant1$ 时，$g(x)$ 在 $[0,a]$ 上严格减，在 $[a,2]$ 上严格增，}
\step{则必有 $\begin{cases}g(0)<1\\g(x)_{\max}=g(2)\geqslant2\end{cases}$，解得 $0<a\leqslant2-\sqrt3$，}
\step{4) 当 $1<a<2$ 时，$g(x)$ 在 $[0,a]$ 上严格减，在 $[a,2]$ 上严格增，}
\step{则必有 $\begin{cases}g(2)<1\\g(x)_{\max}=g(0)\geqslant2\end{cases}$，解得 $\sqrt3\leqslant a<2$.}
\step{综上所述，$a$ 的范围 $[-\sqrt2,2-\sqrt3]\cup[\sqrt3,2+\sqrt2]$.}% 原卷“范围”后缺“是／为”字，照录
\solend
\fi

\end{enumerate}

\end{document}
"
%
% 原始材料：微信公众号“魔都数学加油站”文章，作者破晓，2026-11-09 17:52 发布；
% 连载“27学年高一122”，原文标题“27学年高一122——2026-2027年上海市华师大二附中高一上周末作业6”；
% 原文链接：https://mp.weixin.qq.com/s/pC5066JFKNkQQwjMqTK8Ow。
% 正文为 14 张 PNG 页（第 14 页 2560×3620，其余 13 页 4762×6735），
% 题目下印有【解析】，为讲评版。
% 题目、解析均据原图转录；卷面红色水印及公众号页脚未录入。全卷无插图（已逐页检索，无“如图”）。
%
% 卷首标题：原卷印作“2026-2027 年华师大二附中高一上周末作业 6”（公众号标题中的“上海市”
% 卷面标题行没有，本稿照卷面），年份区间按全稿惯例排 en dash；原卷未印副标题，本稿按题目内容
% 标为“函数”。副标题理由（以全部 14 页实际内容为准）：逐页复核 19 题，考的都是函数性质——
% 奇偶性与充要条件（Q3、Q13、Q15、Q16）、单调性与单调区间（Q4、Q7、Q14、Q17）、定义域／
% 值域／最值与复合函数（Q5、Q6、Q8、Q10、Q12、Q19、Q20）、分段函数与取整函数（Q7、Q8、
% Q11、Q15、Q19）、抽象函数与新定义（Q12、Q20、Q21）；不等式仅作求参数范围／解集的工具
% （Q9、Q14、Q16、Q18），集合与常用逻辑用语仅零星出现（Q3 的充要条件、Q11③ 的集合、
% Q20 的“非空集合 I”），均非独立章节。故主体为单一模块“函数”，与同校 高一102 周末作业 5
% 同（该函数模块卷亦按 高一83／105 成例只用模块名、不叠“集合与常用逻辑用语、不等式”）。
%
% 与原卷的排版差异：
%   ① 原文从第 3 题起，填空题编号为 3—12、选择题为 13—16、解答题为 17—21，本稿保留原题号；
%   ② 原卷每题下直接印【解析】，本稿按全稿惯例置于答案版；四道选择题括号内均未印答案、
%      答案在解析末句“故选 …”，括号内留空（全卷无只印【答案】的题、无 \ans{}）；
%   ③ 原卷未印副标题，本稿按“知识模块”口径补出；
%   ④ 原卷填空横线按全稿惯例排为 \blankline[4.4em]；选择题括号排 \mbox{（\quad）}、
%      选项段前按惯例加 \keepwithprev；解答题不留作答区（原卷为讲评版）；
%   ⑤ 原卷数集排斜体 $R$、$N_+$（下标加号），按全稿惯例改 \R、\N_+
%      （$N_+$ 下标加号与 $\N^*$ 是两种记号，照录下标、只改数集字母为正体，同 高一05 第 20 题成例）；
%      绝对值排 \lvert…\rvert、集合描述竖线排 \mid；
%      $f(x)_{min}$／$f(x)_{max}$／$y_{max}$／$(y_{max})_{min}$ 一类斜体下标按 高一96 成例
%      改正体 \min／\max；不等式一律排 \geqslant／\leqslant；
%   ⑥ 第 20 题解析标签原卷印作【解得】（全卷唯一，其余 18 题均印【解析】），本稿按全稿惯例
%      由 \solbegin 统一印【解析】，正文内容照录（行内加 % 注、并入文件头存疑清单）；
%   ⑦ 题干与解析的标点（半角／全角逗号、冒号、括号、句点、子题括号）均照原卷，
%      题目正文引号用弯引号“ ”。
%
% 子集符号：第 21 题 (2) 的 $\left[\frac13+\frac1{9b-3},1\right]\subseteq\left[\frac12,b\right]$
%   印有下划线，照录 \subseteq。全卷未出现判别式 \triangle／\Delta。
%
% 原卷存疑／笔误（均照录加注）：
%   第 11 题解析对② 的“$y=y+\Delta y$”（按上下文应为 $y=[y]+\Delta y$，原卷漏方括号）；
%   第 19 题题干“某种产品 $x$ 万件）可获利”右括号无配对（疑漏印左括号“（”）；
%   第 20 题解析标签【解得】（见排版差异⑥）；
%   第 21 题 (3) 解析“总存在唯一的 $x_2\in[0,2]$ 时，使得……”的“时”字（按上下文衍字）；
%   第 21 题 (3) 末句“综上所述，$a$ 的范围 $[-\sqrt2,2-\sqrt3]\cup[\sqrt3,2+\sqrt2]$”
%      “范围”后缺“是／为”字，照录。
%
% 行数核对：逐题数 \step 条数，合计 162 条（Q3：5；Q4：3；Q5：1；Q6：3；Q7：2；
% Q8：5；Q9：7；Q10：10；Q11：8；Q12：4；Q13：12；Q14：6；Q15：14；Q16：14；
% Q17：6；Q18：8；Q19：7；Q20：19；Q21：28），与原卷印刷行逐题相等
% （cases 块作一个印刷行计）。全卷无 \ans{}、无只印【答案】的题。
%
\documentclass[a4paper,11pt]{article}
\usepackage{common}

\begin{document}

\examtitlest[3]{2026--2027 年华师大二附中高一上周末作业 6}{函数}

\bigsec{一、填空题}

\begin{enumerate}
\setcounter{enumi}{2}

\item 已知 $p:a=0$，$q$：函数 $f(x)=x^3+x+a$ 是奇函数，则 $p$ 是 $q$ 的\blankline[4.4em]条件.

\ifanswer
\solbegin
\step{若 $a=0$，则 $f(x)=x^3+x$，}
\step{易得 $f(x)$ 的定义域为 $\R$，且 $f(-x)=-(x^3+x)=-f(x)$，则 $f(x)$ 为奇函数，}
\step{反之，若 $f(x)=x^3+x+a$ 是奇函数，}
\step{则有 $f(-x)+f(x)=-x^3-x+a+x^3+x+a=2a=0$，必有 $a=0$，}
\step{故 $p$ 是 $q$ 的充要条件.}
\solend
\fi

\item 函数 $y=\sqrt{-x^2+2x}$ 的单调递增区间是\blankline[4.4em].

\ifanswer
\solbegin
\step{由函数定义域得 $-x^2+2x\geqslant0$，解得 $0\leqslant x\leqslant2$，}
\step{二次函数 $y=-x^2+2x$ 开口向下，对称轴为 $x=1$，}
\step{结合复合函数的单调性得函数 $y=\sqrt{-x^2+2x}$ 的单调递增区间是 $[0,1]$.}
\solend
\fi

\item 设 $f(x)$ 定义域为 $[1,4]$，且 $f(x)=2x+1$，则 $f(x^2)$ 的最大值为\blankline[4.4em].

\ifanswer
\solbegin
\step{由 $x^2\in[1,4]$，得 $f(x^2)=2x^2+1\leqslant9$，当 $x=\pm2$ 时取等号，最大值为 9.}
\solend
\fi

\item 函数 $f(x)=\sqrt{x-2}-x$，$x\in(2,6]$ 的值域为\blankline[4.4em].

\ifanswer
\solbegin
\step{令 $t=\sqrt{x-2}$，则 $x=2+t^2$，由 $2<x\leqslant6$ 得 $0<t\leqslant2$，}
\step{原函数可化为 $g(t)=-t^2+t-2$，$g(0)=-2$，$g(2)=-4$，$g\left(\dfrac12\right)=-\dfrac74$，}
\step{由二次函数性质得 $g(t)\in\left[-4,-\dfrac74\right]$.}
\solend
\fi

\item 若函数 $f(x)=\lvert x-a+1\rvert$ 在区间 $[1,+\infty)$ 上是严格增函数，则实数 $a$ 的取值范围是\blankline[4.4em].

\ifanswer
\solbegin
\step{$f(x)=\begin{cases}x-a+1,x\geqslant a-1\\a-1-x,x<a-1\end{cases}$，显然 $f(x)$ 在 $[a-1,+\infty)$ 上严格增，}
\step{要使函数在区间 $[1,+\infty)$ 上是严格增函数，则 $a-1\leqslant1$，即 $a\leqslant2$.}
\solend
\fi

\item 若函数 $f(x)=\begin{cases}x^2-5,x\leqslant2\\2ax+1,x>2\end{cases}$ 的值域是 $\R$，则 $a$ 的取值范围是\blankline[4.4em].

\ifanswer
\solbegin
\step{当 $x\leqslant2$ 时，$f(x)=x^2-5$，其开口向上，对称轴为 $x=0$，值域为 $[-5,+\infty)$，}
\step{由函数 $f(x)=\begin{cases}x^2-5,x\leqslant2\\2ax+1,x>2\end{cases}$ 的值域是 $\R$，}
\step{则当 $x>2$ 时，$f(x)=2ax+1$ 的值域应包含 $(-\infty,-5)$，}
\step{所以 $f(x)$ 为严格减函数，所以 $\begin{cases}2\times2a+1\geqslant-5\\2a<0\end{cases}$，即 $\begin{cases}a\geqslant-\dfrac32\\a<0\end{cases}$，}
\step{解得 $-\dfrac32\leqslant a<0$，故 $a$ 的取值范围是 $\left[-\dfrac32,0\right)$.}
\solend
\fi

\item 已知函数 $f(x)=ax^2+\dfrac bx$ 的图象经过点 $(-1,0)$，且满足 $f(1)=2$. 若 $\forall x\in(0,+\infty)$，都有 $f(x)+f\left(\dfrac1x\right)-2m\geqslant0$ 恒成立，则实数 $m$ 的取值范围为\blankline[4.4em].

\ifanswer
\solbegin
\step{因为 $f(x)=ax^2+\dfrac bx$ 的图象经过点 $(-1,0)$，且满足 $f(1)=2$，}
\step{所以 $\begin{cases}a-b=0\\a+b=2\end{cases}$，所以 $a=1$，$b=1$，即 $f(x)=x^2+\dfrac1x$，}
\step{$\forall x\in(0,+\infty)$，$f(x)+f\left(\dfrac1x\right)-2m\geqslant0\Leftrightarrow2m\leqslant x^2+\dfrac1x+\dfrac1{x^2}+x$ 恒成立，}
\step{令函数 $g(x)=x^2+\dfrac1x+\dfrac1{x^2}+x=\left(x+\dfrac1x\right)^2+x+\dfrac1x-2$，}
\step{而当 $x>0$ 时，$x+\dfrac1x\geqslant2$，当且仅当 $x=1$ 时取等号，}
\step{因此 $g(x)\geqslant2^2+2-2=4$，当且仅当 $x=1$ 时取等号，则 $2m\leqslant4$，解得 $m\leqslant2$，}
\step{所以实数 $m\in(-\infty,2]$.}
\solend
\fi

\item 若定义在区间 $[a,b](a<b)$ 上的函数 $f(x)=k-\sqrt{x+1}$ 值域也为 $[a,b]$，则实数 $k$ 的取值范围是\blankline[4.4em].

\ifanswer
\solbegin
\step{函数 $f(x)=k-\sqrt{x+1}$ 在定义域 $[-1,+\infty)$ 严格减，}
\step{当 $f(x)$ 的定义域为 $[a,b]$ 时，$f(x)$ 的值域也为 $[a,b]$，$a<b$，}
\step{故 $f(a)=k-\sqrt{a+1}=b$，$f(b)=k-\sqrt{b+1}=a$，}
\step{两式相减得 $\sqrt{a+1}-\sqrt{b+1}=a-b=(a+1)-(b+1)$}
\step{$=(\sqrt{a+1}-\sqrt{b+1})(\sqrt{a+1}+\sqrt{b+1})$，}
\step{所以 $\sqrt{a+1}-\sqrt{b+1}<0$，即 $\sqrt{a+1}+\sqrt{b+1}=1$，}
\step{则 $k=\sqrt{b+1}+a=a+1-\sqrt{a+1}$，}
\step{令 $\lambda=\sqrt{a+1}\geqslant0$，得 $k=\lambda^2-\lambda=\left(\lambda-\dfrac12\right)^2-\dfrac14$，}
\step{又 $\sqrt{a+1}<\sqrt{b+1}$，因为 $\sqrt{a+1}+\sqrt{b+1}=1$，所以 $\sqrt{a+1}\in\left[0,\dfrac12\right)$，}
\step{所以 $0\leqslant\lambda<\dfrac12$，故实数 $k$ 的取值范围为 $\left(-\dfrac14,0\right]$.}
\solend
\fi

\item 函数 $f(x)=[x]$ 的函数值表示不超过 $x$ 的最大整数，例如：$[-3.5]=-4$，$[2.1]=2$，称 $f(x)=[x]$ 为取整函数，则下列关于取整函数的说法：①对任意的 $x\in\R$，$[-x]=-[x]$ ②对任意的 $x$、$y\in\R$，$[x]+[y]\leqslant[x+y]$ ③集合 $\{y\mid y=[nx]-n[x],x\in\R,n\in\N_+\}$ 共有 $n$ 个元素 ④$x>0$ 时，关于 $x$ 的方程 $[2x]=\left[\dfrac1x\right]$ 有无数个解 其中正确的有\blankline[4.4em].

\ifanswer
\solbegin
\step{对①，因为 $[-2.1]=-3\neq[2.1]=2$，所以①错误；}
\step{对②，设 $x=[x]+\Delta x$，$\Delta x\in[0,1)$，$y=y+\Delta y$，$\Delta y\in[0,1)$，}% 原卷印作 $y=y+\Delta y$（按上下文应为 $y=[y]+\Delta y$），照录
\step{$[x+y]=[[x]+\Delta x+[y]+\Delta y]=[x]+[y]+[\Delta x+\Delta y]$，又 $\Delta x+\Delta y\in[0,\ 2)$，}
\step{故 $\Delta x+\Delta y=0$ 或 1，故 $[x]+[y]\leqslant[x+y]$，所以②正确；}
\step{对③，$[nx]-n[x]=[n[x]+n\cdot\Delta x]-n[x]=n[x]+[n\cdot\Delta x]-n[x]=[n\cdot\Delta x]$，}
\step{因为 $\Delta x\in[0,1)$，所以 $0\leqslant n\cdot\Delta x<n$，$n\in\N_+$，所以③正确；}
\step{对④，当 $\dfrac12<x<1$ 时，$[2x]=\left[\dfrac1x\right]=1$，故方程存在无数个解，所以④正确.}
\step{故正确的是②③④.}
\solend
\fi

\item 已知函数 $f(x)$ 定义域为 $[0,1]$，且满足 $f(x)+f(1-x)=1$，$f(x)=2f\left(\dfrac x5\right)$，对任意 $x_1\leqslant x_2$，恒有 $f(x_1)\leqslant f(x_2)$，则 $f\left(\dfrac1{2026}\right)$ 的值为\blankline[4.4em].

\ifanswer
\solbegin
\step{$f(0)=2f(0)\Rightarrow f(0)=0$，从而 $f(1)=1$，$f\left(\dfrac15\right)=\dfrac12$，$f\left(\dfrac45\right)=1-f\left(\dfrac15\right)=\dfrac12$，}
\step{由单调性，$x\in\left[\dfrac15,\dfrac45\right]$ 时，$f(x)=\dfrac12$.}
\step{因为 $\dfrac15<\dfrac{625}{2026}<\dfrac45$，反复使用 $f\left(\dfrac x5\right)=\dfrac12 f(x)$，}
\step{得 $f\left(\dfrac1{2026}\right)=\dfrac1{2^4}f\left(\dfrac{625}{2026}\right)=\dfrac1{32}$.}
\solend
\fi

\end{enumerate}

\bigsec{二、选择题}

\begin{enumerate}
\setcounter{enumi}{12}

\item 已知 $f(x)$ 为 $\R$ 上的奇函数，$g(x)$ 为 $\R$ 上的偶函数，且 $g(x)\neq0$，则下列说法正确的是\mbox{（\quad）}

\keepwithprev
A. $f(x)+g(x)$ 为 $\R$ 上的奇函数\qquad\qquad B. $f(x)-g(x)$ 为 $\R$ 上的奇函数

C. $\dfrac{f(x)}{g(x)}$ 为 $\R$ 上的偶函数\qquad\qquad D. $\lvert f(x)g(x)\rvert$ 为 $\R$ 上的偶函数

\ifanswer
\solbegin
\step{因为 $f(x)$ 为 $\R$ 上的奇函数，$g(x)$ 为 $\R$ 上的偶函数，且 $g(x)\neq0$，}
\step{所以 $f(-x)=-f(x)$，$g(-x)=g(x)$，}
\step{所以 $f(-x)+g(-x)=-f(x)+g(x)\neq-[f(x)+g(x)]$，}
\step{故 $f(x)+g(x)$ 为非奇非偶函数，$A$ 错误；}
\step{同理，$f(x)-g(x)$ 为非奇非偶函数，$B$ 错误；}
\step{设 $F(x)=\dfrac{f(x)}{g(x)}$，则 $F(-x)=\dfrac{f(-x)}{g(-x)}=\dfrac{-f(x)}{g(x)}=-F(x)$，}
\step{所以 $F(x)$ 为奇函数，$C$ 错误；}
\step{设函数 $H(x)=\lvert f(x)g(x)\rvert$，}
\step{因为 $f(x)$ 为 $\R$ 上的奇函数，$g(x)$ 为 $\R$ 上的偶函数，且 $g(x)\neq0$，}
\step{所以 $f(-x)=-f(x)$，$g(-x)=g(x)$，则由函数奇偶性的定义}
\step{得 $H(-x)=\lvert f(-x)g(-x)\rvert=\lvert -f(x)g(x)\rvert=\lvert f(x)g(x)\rvert=H(x)$，$D$ 正确.}
\step{故选 $D$.}
\solend
\fi

\item 若函数 $f(x)$ 是定义域为 $(-\infty,0)\cup(0,+\infty)$ 的奇函数，在区间 $(0,+\infty)$ 上严格增，且 $f(-1)=0$，则不等式 $\dfrac{f(x-1)}{x-2}<0$ 的解集为\mbox{（\quad）}

\keepwithprev
A. $(-\infty,-2)$\qquad\qquad B. $(-2,0)$\qquad\qquad C. $(0,1)$\qquad\qquad D. $(1,2)$

\ifanswer
\solbegin
\step{若函数 $f(x)$ 是定义域为 $(-\infty,0)\cup(0,+\infty)$ 的奇函数，在区间 $(0,+\infty)$ 上严格增，}
\step{由奇函数对称性得，$f(x)$ 在 $(-\infty,0)$ 上严格增，}
\step{因为 $f(-1)=0$，所以 $f(1)=0$，}
\step{则不等式 $\dfrac{f(x-1)}{x-2}<0$ 可转化为 $\begin{cases}x>2\\f(x-1)<0\end{cases}$ 或 $\begin{cases}x<2\\f(x-1)>0\end{cases}$，}
\step{即 $\begin{cases}x>2\\x-1<-1\text{或}0<x-1<1\end{cases}$ 或 $\begin{cases}x<2\\x-1>1\text{或}-1<x-1<0\end{cases}$，解得 $0<x<1$.}
\step{故选 $C$.}
\solend
\fi

\item 若函数 $f(x)$ 同时满足：(1) 对于定义域上的任意 $x$，恒有 $f(x)+f(-x)=0$；(2) 对于定义域上的任意 $x_1$、$x_2$，当 $x_1\neq x_2$ 时，恒有 $x_1f(x_1)-x_1f(x_2)>x_2f(x_1)-x_2f(x_2)$，则称函数 $f(x)$ 为“理想函数”. 给出下列四个函数：① $f(x)=x^2$；② $f(x)=x^3$；③ $f(x)=\dfrac{2x-1}{2x+1}$；④ $f(x)=\begin{cases}x^2+4x,x\geqslant0\\-x^2+4x,x<0\end{cases}$，其中被称为“理想函数”的有\mbox{（\quad）}

\keepwithprev
A. 1 个\qquad\qquad B. 2 个\qquad\qquad C. 3 个\qquad\qquad D. 4 个

\ifanswer
\solbegin
\step{函数 $f(x)$ 同时满足：(1) 对于定义域上的任意 $x$，恒有 $f(x)+f(-x)=0$；}
\step{(2) 对于定义域上的任意 $x_1$、$x_2$，当 $x_1\neq x_2$ 时，}
\step{恒有 $x_1f(x_1)-x_1f(x_2)>x_2f(x_1)-x_2f(x_2)$，}
\step{整理得 $(x_1-x_2)[f(x_1)-f(x_2)]>0$，则称函数 $f(x)$ 为“理想函数”.}
\step{即函数满足 (1) 奇函数；(2) 函数在定义域内为严格增函数；}
\step{① $f(x)=x^2$；不满足函数为奇函数在定义域内也不严格增；}
\step{② $f(x)=x^3$；满足函数既为奇函数，也在定义域内严格增，}
\step{故该函数为理想函数；}
\step{③ $f(x)=\dfrac{2x-1}{2x+1}=\dfrac{2x+1-2}{2x+1}=-\dfrac2{2x+1}+1$，该函数不为奇函数，}
\step{但是函数在定义域内严格增，故该函数不为理想函数；}
\step{④ $f(x)=\begin{cases}x^2+4x,x\geqslant0\\-x^2+4x,x<0\end{cases}$，函数的关系式满足 $f(-x)=-f(x)$，}
\step{且在 $x\geqslant0$ 时，函数严格增，由于奇函数在对称区间上单调性相同，}
\step{故该函数为理想函数.}
\step{故选 $B$.}
\solend
\fi

\item 已知 $f(x)$ 是定义在 $\R$ 上的奇函数，当 $x_1$、$x_2\in(0,+\infty)$ 且 $x_1\neq x_2$ 时，都有\newline $\dfrac{x_2f(x_1)-x_1f(x_2)}{x_1x_2(x_1-x_2)}>0$ 成立，$f(2028)=2028$，则不等式 $f(x)-x>0$ 的解集为\mbox{（\quad）}

\keepwithprev
A. $(-\infty,-2028)\cup(2028,+\infty)$\qquad\qquad B. $\left(-\dfrac1{2028},\dfrac1{2028}\right)$

C. $(-2028,2028)$\qquad\qquad D. $(-2028,0)\cup(2028,+\infty)$

\ifanswer
\solbegin
\step{因为 $f(x)$ 是定义在 $\R$ 上的奇函数，所以 $f(-x)=-f(x)$，}
\step{令 $g(x)=\dfrac{f(x)}x$，则 $g(-x)=\dfrac{f(-x)}{-x}=\dfrac{-f(x)}{-x}=\dfrac{f(x)}x=g(x)$，}
\step{所以 $g(x)$ 为偶函数，}
\step{当 $x_1$、$x_2\in(0,+\infty)$ 且 $x_1\neq x_2$ 时，都有 $\dfrac{x_2f(x_1)-x_1f(x_2)}{x_1x_2(x_1-x_2)}>0$ 成立，}
\step{即 $\dfrac{\dfrac{f(x_1)}{x_1}-\dfrac{f(x_2)}{x_2}}{x_1-x_2}>0$，即 $g(x)$ 在 $(0,+\infty)$ 上严格增，}
\step{由偶函数对称性得 $g(x)$ 在 $(-\infty,0)$ 上严格减，}
\step{因为 $f(2028)=2028$，所以 $g(2028)=1$，}
\step{则不等式 $f(x)-x>0$ 可化为 $f(x)>x$，}
\step{当 $x=0$ 时，$f(0)=0$，显然不成立，}
\step{当 $x>0$ 时，不等式可化为 $\dfrac{f(x)}x>1$，即 $g(x)>1=g(2028)$，}
\step{所以 $\lvert x\rvert>2028$，解得 $x>2028$，}
\step{当 $x<0$ 时，不等式可化为 $\dfrac{f(x)}x<1$，即 $g(x)<1=g(2028)$，}
\step{所以 $\lvert x\rvert<2028$，解得 $-2028<x<0$，}
\step{故不等式的解集为 $(-2028,0)\cup(2028,+\infty)$. 故选 $D$.}
\solend
\fi

\end{enumerate}

\bigsec{三、解答题}

\begin{enumerate}
\setcounter{enumi}{16}

\item 已知函数 $f(x)=\dfrac{2x-3}{x+1}$.

(1) 试用函数单调性的定义证明：函数 $f(x)$ 在 $(-1,+\infty)$ 上是严格增函数；

(2) 求函数 $f(x)$ 在区间 $[1,4]$ 上的最大值和最小值.

\ifanswer
\solbegin
\step{(1) 设 $-1<x_1<x_2$，则 $f(x_1)-f(x_2)=\dfrac{2x_1-3}{x_1+1}-\dfrac{2x_2-3}{x_2+1}=\dfrac{5(x_1-x_2)}{(x_1+1)(x_2+1)}$}
\step{因为 $-1<x_1<x_2$，所以 $x_1-x_2<0$，$x_1+1>0$，$x_2+1>0$，}
\step{所以 $f(x_1)-f(x_2)<0$，所以 $f(x)$ 在 $(-1,+\infty)$ 上严格增.}
\step{(2) 由 (1) 得 $f(x)$ 在区间 $[1,4]$ 上严格增，}
\step{所以 $f(x)_{\min}=f(1)=-\dfrac12$，$f(x)_{\max}=f(4)=1$，}
\step{所以函数 $f(x)$ 在区间 $[1,4]$ 上的最大值为 1，最小值为 $-\dfrac12$.}
\solend
\fi

\item 设函数 $f(x)=x^2-mx+m$.

(1) 当 $m=2$ 时，求不等式 $f(x)<x$ 的解集；

(2) 若 $\exists x\in[1,2]$，使不等式 $f(x)<x+m-2$ 成立，求实数 $m$ 的取值范围.

\ifanswer
\solbegin
\step{(1) $m=2$ 时，函数 $f(x)=x^2-2x+2$，}
\step{不等式 $f(x)<x$ 可化为 $x^2-3x+2<0$，解得 $1<x<2$，}
\step{所以不等式的解集为 $\{x\mid1<x<2\}$；}
\step{(2) 若 $\exists x\in[1,2]$，使不等式 $f(x)<x+m-2$ 成立，即 $x^2-x+2<mx$，}
\step{所以 $m>x+\dfrac2x-1$，设 $g(x)=x+\dfrac2x-1$，$x\in[1,2]$，}
\step{则 $g(x)=x+\dfrac2x-1\geqslant2\sqrt{x\cdot\dfrac2x}-1=2\sqrt2-1$，}
\step{当且仅当 $x=\dfrac2x$，即 $x=\sqrt2$ 时取等号，}
\step{所以实数 $m$ 的取值范围是 $\{m\mid m>2\sqrt2-1\}$.}
\solend
\fi

\item 某企业原来生产某种产品 $x$ 万件）可获利 $w(x)$（万元），且满足\newline% 原卷“万件）”右括号无配对、疑漏印左括号“（”，照录
$w(x)=\begin{cases}20(x^2+17),0<x\leqslant2\\500-\dfrac{80}{x-1},2<x\leqslant5\end{cases}$.
现该企业计划引进新的生产设备和新的产品方案优化产品，优化后的产品的其他成本总投入为 $(20x+10)$ 万元. 由市场调研分析得知，当前产品供不应求. 记该企业优化后的产品的利润为 $f(x)$（单位：万元）.

(1) 求函数 $f(x)$ 的解析式；

(2) 当优化后的产品产量为多少万件时，该企业的利润 $f(x)$ 最大？最大利润是多少？请说明理由.

\ifanswer
\solbegin
\step{(1) 由题意得 $f(x)=\begin{cases}20x^2-20x+330,0<x\leqslant2\\-20x-\dfrac{80}{x-1}+490,2<x\leqslant5\end{cases}$；}
\step{(2) 当 $0<x\leqslant2$ 时，$f(x)=20(x^2-x)+330=20\left(x-\dfrac12\right)^2+325$，}
\step{故当 $x=2$ 时，$f(x)$ 取最大值，$f(2)=370$；}
\step{当 $2<x\leqslant5$ 时，$f(x)=-\left[20(x-1)+\dfrac{80}{x-1}\right]+470\leqslant-80+470=390$，}
\step{当且仅当 $20(x-1)=\dfrac{80}{x-1}$，即 $x=3$ 时，$f(3)=390$ 为最大值；}
\step{因为 $370<390$，}
\step{所以优化后的产品产量为 3 万件时，企业获最大利润 390 万元.}
\solend
\fi

\item 俄国数学家切比雪夫（1821—1894）是研究直线逼近函数的理论先驱. 对定义在非空集合 $I$ 上的函数 $f(x)$，以及函数 $g(x)=kx+b(k,b\in\R)$，切比雪夫将函数 $y=\lvert f(x)-g(x)\rvert,x\in I$ 的最大值称为 $f(x),g(x)$ 的“偏差”.

(1) 函数 $f(x)=x^2,x\in[0,1],g(x)=-x-2$，求 $f(x),g(x)$ 的“偏差”；

(2) 函数 $f(x)=\dfrac1x+1,x\in[1,2],g(x)=kx+1(k>0)$，若 $f(x),g(x)$ 的“偏差”为 2，求 $k$ 的值；

(3) 函数 $f(x)=x^2-x,x\in[0,3],g(x)=2x+b$，当 $f(x),g(x)$ 的“偏差”取最小值时，求 $b$ 的值，并求出“偏差”的最小值.

\ifanswer
% 注：原卷本题解析标签印作【解得】（全卷唯一，其余 18 题均印【解析】）；
%     本稿按全稿惯例由 \solbegin 统一印【解析】，正文内容照录。见文件头排版差异⑥。
\solbegin
\step{(1) $y=\lvert f(x)-g(x)\rvert=\lvert x^2+x+2\rvert=\left\lvert\left(x+\dfrac12\right)^2+\dfrac74\right\rvert,x\in[0,1]$，}
\step{因为 $x\in[0,1]$，所以 $\left(x+\dfrac12\right)^2+\dfrac74\in[2,4]$，}
\step{则 $y=\left\lvert\left(x+\dfrac12\right)^2+\dfrac74\right\rvert\in[2,4]$，$y_{\max}=4$，}
\step{所以函数 $f(x)$ 与 $g(x)$ 的“偏差”为 4.}
\step{(2) 令 $t=f(x)-g(x)=\dfrac1x+1-(kx+1)=\dfrac1x-kx,x\in[1,2]$，}
\step{因为 $k>0$，所以 $t=\dfrac1x-kx$ 是单调减函数，}
\step{所以 $t\in\left[\dfrac12-2k,1-k\right]$，由单调，$y=\lvert t\rvert,t\in\left[\dfrac12-2k,1-k\right]$，且 $y_{\max}=2$.}
\step{当 $\left\lvert\dfrac12-2k\right\rvert\leqslant\lvert1-k\rvert$，即 $0<k\leqslant\dfrac12$ 时，$\lvert1-k\rvert=2$，解得 $k=3$ 或 $k=-1$，}
\step{均不符合；}
\step{当 $\left\lvert\dfrac12-2k\right\rvert>\lvert1-k\rvert$，即 $k>\dfrac12$ 时，$\left\lvert\dfrac12-2k\right\rvert=2$，$\dfrac12-2k=2$ 或 $\dfrac12-2k=-2$，}
\step{解得 $k=\dfrac54$ 或 $k=-\dfrac34$（舍），所以 $k=\dfrac54$.}
\step{(3) $y=\lvert f(x)-g(x)\rvert=\lvert x^2-x-(2x+b)\rvert=\lvert x^2-3x-b\rvert$}
\step{$=\left\lvert\left(x-\dfrac32\right)^2-\dfrac94-b\right\rvert,x\in[0,3]$，}
\step{因为 $x\in[0,3]$，所以 $\left(x-\dfrac32\right)^2-\dfrac94-b\in\left[-\dfrac94-b,-b\right]$，}
\step{由 $y=\left\lvert\left(x-\dfrac32\right)^2-\dfrac94-b\right\rvert$，则 $y_{\max}=\max\left\{\lvert b\rvert,\left\lvert b+\dfrac94\right\rvert\right\}$，}
\step{令 $\lvert b\rvert\geqslant\left\lvert b+\dfrac94\right\rvert$，即 $b^2\geqslant\left(b+\dfrac94\right)^2$，解得 $b\leqslant-\dfrac98$，}
\step{$y_{\max}=\max\left\{\lvert b\rvert,\left\lvert b+\dfrac94\right\rvert\right\}=\begin{cases}-b,b\leqslant-\dfrac98\\b+\dfrac94,b>-\dfrac98\end{cases}$，}
\step{故当且仅当 $b=-\dfrac98$ 时，$(y_{\max})_{\min}=\dfrac98$，}
\step{故当 $b$ 的值为 $-\dfrac98$ 时，函数 $f(x)$ 与 $g(x)$ 的“偏差”取最小值 $\dfrac98$.}
\solend
\fi

\item 若函数 $f(x)$ 与 $g(x)$ 满足：对任意的 $x_1\in D$，总存在唯一的 $x_2\in D$，使 $f(x_1)g(x_2)=m$ 成立，则称 $f(x)$ 是 $g(x)$ 在区间 $D$ 上的“$m$ 阶伴随函数”；对任意的 $x_1\in D$，总存在唯一的 $x_2\in D$，使 $f(x_1)f(x_2)=m$ 成立，则称 $f(x)$ 是区间 $D$ 上的“$m$ 阶自伴函数”.

(1) 判断 $f(x)=x^2+1$ 是否为区间 $[0,3]$ 上的“2 阶自伴函数”？并说明理由：

(2) 若函数 $f(x)=3x-1$ 为区间 $\left[\dfrac12,b\right]$ 上的“1 阶自伴函数”，求 $b$ 的值；

(3) 若 $f(x)=\dfrac4{x+2}$ 是 $g(x)=x^2-2ax+a^2-1$ 在区间 $[0,2]$ 上的“2 阶伴随函数”，求实数 $a$ 的取值范围.

\ifanswer
\solbegin
\step{(1) 假定函数 $f(x)=x^2+1$ 是区间 $[0,3]$ 上的“2 阶自伴函数”，}
\step{取 $x_1=2,f(x_1)=f(2)=5$，由 $f(x_1)f(x_2)=2$，得 $5(x_2^2+1)=2$，}
\step{显然此方程无实数解，}
\step{所以函数 $f(x)=x^2+1$ 不是区间 $[0,3]$ 上的“2 阶自伴函数”.}
\step{(2) 函数 $f(x)=3x-1$ 为区间 $\left[\dfrac12,b\right]$ 上的“1 阶自伴函数”，}
\step{则对任意 $x_1\in\left[\dfrac12,b\right]$，总存在唯一的 $x_2\in\left[\dfrac12,b\right]$，使得 $f(x_1)f(x_2)=1$，}
\step{即 $3x_2-1=\dfrac1{3x_1-1}$，整理得 $x_2=\dfrac13+\dfrac1{3(3x_1-1)}$，}
\step{显然函数 $x_2=\dfrac13+\dfrac1{3(3x_1-1)}$ 在 $\left[\dfrac12,b\right]$ 上严格减，}
\step{且当 $x_1=\dfrac12$ 时，$x_2=1$，当 $x_1=b$ 时，$x_2=\dfrac13+\dfrac1{9b-3}$，}
\step{因此对 $\left[\dfrac12,b\right]$ 内的每一个 $x_1$，在 $\left[\dfrac13+\dfrac1{9b-3},1\right]$ 内有唯一 $x_2$ 值与之对应，}
\step{则 $x_2\in\left[\dfrac12,b\right]$，于是 $\left[\dfrac13+\dfrac1{9b-3},1\right]\subseteq\left[\dfrac12,b\right]$，应有 $\begin{cases}b\geqslant1\\\dfrac13+\dfrac1{9b-3}\geqslant\dfrac12\end{cases}$，}
\step{得 $\begin{cases}b\geqslant1\\b\leqslant1\end{cases}$，即 $b=1$，所以 $b$ 的值是 1.}
\step{(3) 由函数 $f(x)=\dfrac4{x+2}$ 在区间 $[0,2]$ 的值域为 $[1,2]$，}
\step{因为 $f(x)=\dfrac4{x+2}$ 是 $g(x)=x^2-2ax+a^2-1$ 在 $[0,2]$ 上的“2 阶伴随函数”，}
\step{则对任意的 $x_1\in[0,2]$，总存在唯一的 $x_2\in[0,2]$ 时，使得 $f(x_1)g(x_2)=2$，}% 原卷此处印“时”字（按上下文为衍字），照录
\step{所以 $g(x_2)=\dfrac2{f(x_1)}\in[1,2]$，}
\step{即 $g(x)=x^2-2ax+a^2-1$ 在 $[0,2]$ 区间上的值域必定包含区间 $[1,2]$，}
\step{且 $g(x)$ 值域在 $[1,2]$ 对应的自变量是唯一的，}
\step{又因为函数 $g(x)=x^2-2ax+a^2-1$ 开口向上，对称轴为 $x=a$，}
\step{1) 当 $a\leqslant0$ 时，$g(x)$ 在 $[0,2]$ 区间上严格增，则必有 $\begin{cases}g(x)_{\min}=g(0)\leqslant1\\g(x)_{\max}=g(2)\geqslant2\end{cases}$，}
\step{解得 $-\sqrt2\leqslant a\leqslant0$；}
\step{2) 当 $a\geqslant2$ 时，$g(x)$ 在 $[0,2]$ 区间上严格减，则必有 $\begin{cases}g(x)_{\min}=g(2)\leqslant1\\g(x)_{\max}=g(0)\geqslant2\end{cases}$，}
\step{解得 $2\leqslant a\leqslant2+\sqrt2$；}
\step{3) 当 $0<a\leqslant1$ 时，$g(x)$ 在 $[0,a]$ 上严格减，在 $[a,2]$ 上严格增，}
\step{则必有 $\begin{cases}g(0)<1\\g(x)_{\max}=g(2)\geqslant2\end{cases}$，解得 $0<a\leqslant2-\sqrt3$，}
\step{4) 当 $1<a<2$ 时，$g(x)$ 在 $[0,a]$ 上严格减，在 $[a,2]$ 上严格增，}
\step{则必有 $\begin{cases}g(2)<1\\g(x)_{\max}=g(0)\geqslant2\end{cases}$，解得 $\sqrt3\leqslant a<2$.}
\step{综上所述，$a$ 的范围 $[-\sqrt2,2-\sqrt3]\cup[\sqrt3,2+\sqrt2]$.}% 原卷“范围”后缺“是／为”字，照录
\solend
\fi

\end{enumerate}

\end{document}
